% concatenated from main.tex, math_commands.tex, abstract.tex, intro.tex, overview.tex, refutation_overview.tex, Ainverse.tex, Qanalysis.tex, Qapprox.tex, vertical.tex, iterative_graph.tex, defer_charging.tex, direct_trace_method_Q.tex, truncation.tex, verification-param.tex, appendix.tex, defer_constructon.tex, defer_deviation.tex, scalar.tex, variance_deviation.tex, MP_appendix.tex, defer_Q.tex, defer_truncation.tex
%\documentclass[12pt,letterpaper]{article}
\documentclass[12pt]{article}
\usepackage[utf8]{inputenc}

%\usepackage{macros}
% author notesa
\usepackage{svgcolor}
\usepackage{csquotes}
\def\showauthornotes{0}


% \ifnum\showauthornotes=1

% \newcommand{\jnote}[1]{\textcolor{blue}{ {\textbf{(Jeff: #1)}}}}
% \newcommand{\anote}[1]{\textcolor{cyan}{ {\textbf{(Aaron: #1)}}}}
% \newcommand{\mnote}[1]{\textcolor{red}{ {\textbf{(Madhur: #1)}}}}
% \newcommand{\snote}[1]{\textcolor{orange}{ {\textbf{(Sofia: #1)}}}}

% \else

% \newcommand{\jnote}[1]{}
% \newcommand{\anote}[1]{}
% \newcommand{\mnote}[1]{}
% \newcommand{\snote}[1]{}
% \fi

%%%%% NEW MATH DEFINITIONS %%%%%

% Common
\def \N {\mathbb{N}}
\def \R {\mathbb{R}}
\def \C {\mathbb{C}}
\def \Q {\mathbb{Q}}
\def \Z {\mathbb{Z}}
\def \F {\mathbb{F}}
\def \E {\mathbb{E}}
\def \PE {\wt{\mathbb{E}}}
\def \eps {\epsilon}
\def \al {\alpha}
\renewcommand{\Pr}{\mathop{\bf Pr\/}}
\newcommand{\wt}[1]{\widetilde{#1}}
\newcommand{\wh}[1]{\widehat{#1}}
\newcommand{\ol}[1]{\overline{#1}}
\newcommand{\ul}[1]{\underline{#1}}
\newcommand{\ot}{\otimes}
\newcommand{\zo}{\{0,1\}}
\newcommand{\pmo}{\{\pm1\}}
\newcommand{\Var}{\mathrm{Var}}
\newcommand{\Cov}{\mathrm{Cov}}
\newcommand{\imp}{\Longrightarrow}
\newcommand{\brac}[1]{\left\langle #1 \right\rangle}
\newcommand{\KL}{d_{\mathrm{KL}}}
\newcommand{\TV}{d_{\mathrm{TV}}}
\def\ceil#1{\lceil #1 \rceil}
\def\floor#1{\lfloor #1 \rfloor}
 \def\1{\bm{1}}
\def \st {\text{s.t.}}
\newcommand{\TODO}[1]{\textcolor{red}{#1}}
\newcommand{\calF}{\mathcal{F}}
% mathsf 
\newcommand{\SA}{\mathsf{SA}}
\newcommand{\SOS}{\mathsf{SOS}}
\newcommand{\pE}{\widetilde{\E}}


\usepackage{amsmath,amssymb,amsthm,amsfonts,latexsym,bm,bbm,xspace,graphicx,float,mathtools,mathdots,physics}
\usepackage{braket,caption,subcaption,ellipsis,xcolor,textcomp,hhline,pifont,combelow,booktabs}
\usepackage[colorlinks=true, allcolors=blue]{hyperref}
\usepackage{color}
\usepackage{times}
\usepackage{fullpage}
\usepackage{tikz-cd}
\usepackage[shortlabels]{enumitem}
\usepackage{thm-restate}
\usepackage{cleveref}
\usepackage{mdframed}

\usepackage{mathrsfs}

% complexity classes
\newcommand{\PTIME}{\mathsf{P}}
\newcommand{\LOGSPACE}{\mathsf{L}}
\newcommand{\ZPP}{\mathsf{ZPP}}
\newcommand{\RP}{\mathsf{RP}}
\newcommand{\BPP}{\mathsf{BPP}}
\newcommand{\NP}{\mathsf{NP}}
\newcommand{\coNP}{\mathsf{coNP}}
\newcommand{\MA}{\mathsf{MA}}
\newcommand{\PH}{\mathsf{PH}}
\newcommand{\TC}{\mathsf{TC}}
\newcommand{\AC}{\mathsf{AC}}
\newcommand{\SC}{\mathsf{SC}}
\newcommand{\SZK}{\mathsf{SZK}}
\newcommand{\AM}{\mathsf{AM}}
\newcommand{\IP}{\mathsf{IP}}
\newcommand{\PSPACE}{\mathsf{PSPACE}}
\newcommand{\EXP}{\mathsf{EXP}}
\newcommand{\MIP}{\mathsf{MIP}}
\newcommand{\NEXP}{\mathsf{NEXP}}
\newcommand{\BQP}{\mathsf{BQP}}
\newcommand{\PP}{\mathsf{PP}}
\newcommand{\SP}{\mathsf{\#P}}

% Environments
% \renewenvironment{proof}{{\emph{Proof.}  }}{\hfill$\square$ \vspace{8pt}}
\newtheorem{theorem}{Theorem}[section]
\newtheorem{lemma}[theorem]{Lemma}
\newtheorem{claim}[theorem]{Claim}
\newtheorem{proposition}[theorem]{Proposition}
\newtheorem{fact}[theorem]{Fact}
\newtheorem{corollary}[theorem]{Corollary}
\newtheorem{conjecture}[theorem]{Conjecture}
\newtheorem{question}[theorem]{Question}
\newtheorem{answer}[theorem]{Answer}
\newtheorem{hypothesis}[theorem]{Hypothesis}
\newtheorem{exercise}[theorem]{Exercise}
% \theoremstyle{definition}
\newtheorem{definition}[theorem]{Definition}
\newtheorem{remark}[theorem]{Remark}
\newtheorem{notation}[theorem]{Notation}
\newtheorem{observation}[theorem]{Observation}
\newtheorem{example}[theorem]{Example}
\newtheorem{examples}[theorem]{Example}
\newtheorem{assumption}{Assumption}

% Math terms
\newcommand{\poly}{\operatorname{poly}}
\newcommand{\polylog}{\operatorname{polylog}}
\newcommand{\size}{\textnormal{size}}
\newcommand{\avg}{\mathop{\textnormal{avg}}}
\newcommand{\sgn}{\textnormal{sgn}}
\newcommand{\dist}{\textnormal{dist}}
\newcommand{\vol}{\textnormal{vol}}
\newcommand{\spn}{\textnormal{span}}
\newcommand{\supp}{\textnormal{supp}}
% \newcommand{\tr}{\operatorname{tr}}
\newcommand{\codim}{\operatorname{codim}}
\newcommand{\diag}{\operatorname{diag}}
\newcommand{\conv}{\operatorname{conv}}


% Mark sections of captions for referring to divisions of figures
\newcommand{\figleft}{{\em (Left)}}
\newcommand{\figcenter}{{\em (Center)}}
\newcommand{\figright}{{\em (Right)}}
\newcommand{\figtop}{{\em (Top)}}
\newcommand{\figbottom}{{\em (Bottom)}}
\newcommand{\captiona}{{\em (a)}}
\newcommand{\captionb}{{\em (b)}}
\newcommand{\captionc}{{\em (c)}}
\newcommand{\captiond}{{\em (d)}}

% Highlight a newly defined term
\newcommand{\newterm}[1]{{\bf #1}}


% Figure reference, lower-case.
\def\figref#1{figure~\ref{#1}}
% Figure reference, capital. For start of sentence
\def\Figref#1{Figure~\ref{#1}}
\def\twofigref#1#2{figures \ref{#1} and \ref{#2}}
\def\quadfigref#1#2#3#4{figures \ref{#1}, \ref{#2}, \ref{#3} and \ref{#4}}
% Section reference, lower-case.
\def\secref#1{section~\ref{#1}}
% Section reference, capital.
\def\Secref#1{Section~\ref{#1}}
% Reference to two sections.
\def\twosecrefs#1#2{sections \ref{#1} and \ref{#2}}
% Reference to three sections.
\def\secrefs#1#2#3{sections \ref{#1}, \ref{#2} and \ref{#3}}
% Reference to an equation, lower-case.
%\def\eqref#1{equation~\ref{#1}}
% Reference to an equation, upper case
\def\Eqref#1{Equation~\ref{#1}}
% A raw reference to an equation---avoid using if possible
\def\plaineqref#1{\ref{#1}}
% Reference to a chapter, lower-case.
\def\chapref#1{chapter~\ref{#1}}
% Reference to an equation, upper case.
\def\Chapref#1{Chapter~\ref{#1}}
% Reference to a range of chapters
\def\rangechapref#1#2{chapters\ref{#1}--\ref{#2}}
% Reference to an algorithm, lower-case.
\def\algref#1{algorithm~\ref{#1}}
% Reference to an algorithm, upper case.
\def\Algref#1{Algorithm~\ref{#1}}
\def\twoalgref#1#2{algorithms \ref{#1} and \ref{#2}}
\def\Twoalgref#1#2{Algorithms \ref{#1} and \ref{#2}}
% Reference to a part, lower case
\def\partref#1{part~\ref{#1}}
% Reference to a part, upper case
\def\Partref#1{Part~\ref{#1}}
\def\twopartref#1#2{parts \ref{#1} and \ref{#2}}

\newcommand{\train}{\mathcal{D}}
\newcommand{\valid}{\mathcal{D_{\mathrm{valid}}}}
\newcommand{\test}{\mathcal{D_{\mathrm{test}}}}


% Random variables
\def\reta{{\textnormal{$\eta$}}}
\def\ra{{\textnormal{a}}}
\def\rb{{\textnormal{b}}}
\def\rc{{\textnormal{c}}}
\def\rd{{\textnormal{d}}}
\def\re{{\textnormal{e}}}
\def\rf{{\textnormal{f}}}
\def\rg{{\textnormal{g}}}
\def\rh{{\textnormal{h}}}
\def\ri{{\textnormal{i}}}
\def\rj{{\textnormal{j}}}
\def\rk{{\textnormal{k}}}
\def\rl{{\textnormal{l}}}
% rm is already a command, just don't name any random variables m
\def\rn{{\textnormal{n}}}
\def\ro{{\textnormal{o}}}
\def\rp{{\textnormal{p}}}
\def\rq{{\textnormal{q}}}
\def\rr{{\textnormal{r}}}
\def\rs{{\textnormal{s}}}
\def\rt{{\textnormal{t}}}
\def\ru{{\textnormal{u}}}
\def\rv{{\textnormal{v}}}
\def\rw{{\textnormal{w}}}
\def\rx{{\textnormal{x}}}
\def\ry{{\textnormal{y}}}
\def\rz{{\textnormal{z}}}

% Random vectors
\def\rvepsilon{{\mathbf{\epsilon}}}
\def\rvtheta{{\mathbf{\theta}}}
\def\rva{{\mathbf{a}}}
\def\rvb{{\mathbf{b}}}
\def\rvc{{\mathbf{c}}}
\def\rvd{{\mathbf{d}}}
\def\rve{{\mathbf{e}}}
\def\rvf{{\mathbf{f}}}
\def\rvg{{\mathbf{g}}}
\def\rvh{{\mathbf{h}}}
\def\rvu{{\mathbf{i}}}
\def\rvj{{\mathbf{j}}}
\def\rvk{{\mathbf{k}}}
\def\rvl{{\mathbf{l}}}
\def\rvm{{\mathbf{m}}}
\def\rvn{{\mathbf{n}}}
\def\rvo{{\mathbf{o}}}
\def\rvp{{\mathbf{p}}}
\def\rvq{{\mathbf{q}}}
\def\rvr{{\mathbf{r}}}
\def\rvs{{\mathbf{s}}}
\def\rvt{{\mathbf{t}}}
\def\rvu{{\mathbf{u}}}
\def\rvv{{\mathbf{v}}}
\def\rvw{{\mathbf{w}}}
\def\rvx{{\mathbf{x}}}
\def\rvy{{\mathbf{y}}}
\def\rvz{{\mathbf{z}}}

% Elements of random vectors
\def\erva{{\textnormal{a}}}
\def\ervb{{\textnormal{b}}}
\def\ervc{{\textnormal{c}}}
\def\ervd{{\textnormal{d}}}
\def\erve{{\textnormal{e}}}
\def\ervf{{\textnormal{f}}}
\def\ervg{{\textnormal{g}}}
\def\ervh{{\textnormal{h}}}
\def\ervi{{\textnormal{i}}}
\def\ervj{{\textnormal{j}}}
\def\ervk{{\textnormal{k}}}
\def\ervl{{\textnormal{l}}}
\def\ervm{{\textnormal{m}}}
\def\ervn{{\textnormal{n}}}
\def\ervo{{\textnormal{o}}}
\def\ervp{{\textnormal{p}}}
\def\ervq{{\textnormal{q}}}
\def\ervr{{\textnormal{r}}}
\def\ervs{{\textnormal{s}}}
\def\ervt{{\textnormal{t}}}
\def\ervu{{\textnormal{u}}}
\def\ervv{{\textnormal{v}}}
\def\ervw{{\textnormal{w}}}
\def\ervx{{\textnormal{x}}}
\def\ervy{{\textnormal{y}}}
\def\ervz{{\textnormal{z}}}

% Random matrices
\def\rmA{{\mathbf{A}}}
\def\rmB{{\mathbf{B}}}
\def\rmC{{\mathbf{C}}}
\def\rmD{{\mathbf{D}}}
\def\rmE{{\mathbf{E}}}
\def\rmF{{\mathbf{F}}}
\def\rmG{{\mathbf{G}}}
\def\rmH{{\mathbf{H}}}
\def\rmI{{\mathbf{I}}}
\def\rmJ{{\mathbf{J}}}
\def\rmK{{\mathbf{K}}}
\def\rmL{{\mathbf{L}}}
\def\rmM{{\mathbf{M}}}
\def\rmN{{\mathbf{N}}}
\def\rmO{{\mathbf{O}}}
\def\rmP{{\mathbf{P}}}
\def\rmQ{{\mathbf{Q}}}
\def\rmR{{\mathbf{R}}}
\def\rmS{{\mathbf{S}}}
\def\rmT{{\mathbf{T}}}
\def\rmU{{\mathbf{U}}}
\def\rmV{{\mathbf{V}}}
\def\rmW{{\mathbf{W}}}
\def\rmX{{\mathbf{X}}}
\def\rmY{{\mathbf{Y}}}
\def\rmZ{{\mathbf{Z}}}

% Elements of random matrices
\def\ermA{{\textnormal{A}}}
\def\ermB{{\textnormal{B}}}
\def\ermC{{\textnormal{C}}}
\def\ermD{{\textnormal{D}}}
\def\ermE{{\textnormal{E}}}
\def\ermF{{\textnormal{F}}}
\def\ermG{{\textnormal{G}}}
\def\ermH{{\textnormal{H}}}
\def\ermI{{\textnormal{I}}}
\def\ermJ{{\textnormal{J}}}
\def\ermK{{\textnormal{K}}}
\def\ermL{{\textnormal{L}}}
\def\ermM{{\textnormal{M}}}
\def\ermN{{\textnormal{N}}}
\def\ermO{{\textnormal{O}}}
\def\ermP{{\textnormal{P}}}
\def\ermQ{{\textnormal{Q}}}
\def\ermR{{\textnormal{R}}}
\def\ermS{{\textnormal{S}}}
\def\ermT{{\textnormal{T}}}
\def\ermU{{\textnormal{U}}}
\def\ermV{{\textnormal{V}}}
\def\ermW{{\textnormal{W}}}
\def\ermX{{\textnormal{X}}}
\def\ermY{{\textnormal{Y}}}
\def\ermZ{{\textnormal{Z}}}

% Vectors
\def\vzero{{\bm{0}}}
\def\vone{{\bm{1}}}
\def\vmu{{\bm{\mu}}}
\def\vtheta{{\bm{\theta}}}
\def\va{{\bm{a}}}
\def\vb{{\bm{b}}}
\def\vc{{\bm{c}}}
\def\vd{{\bm{d}}}
\def\ve{{\bm{e}}}
\def\vf{{\bm{f}}}
\def\vg{{\bm{g}}}
\def\vh{{\bm{h}}}
\def\vi{{\bm{i}}}
\def\vj{{\bm{j}}}
\def\vk{{\bm{k}}}
\def\vl{{\bm{l}}}
\def\vm{{\bm{m}}}
\def\vn{{\bm{n}}}
\def\vo{{\bm{o}}}
\def\vp{{\bm{p}}}
\def\vq{{\bm{q}}}
\def\vr{{\bm{r}}}
\def\vs{{\bm{s}}}
\def\vt{{\bm{t}}}
\def\vu{{\bm{u}}}
\def\vv{{\bm{v}}}
\def\vw{{\bm{w}}}
\def\vx{{\bm{x}}}
\def\vy{{\bm{y}}}
\def\vz{{\bm{z}}}
\def\valpha{{\bm{\alpha}}}
\def\vbeta{{\bm{\beta}}}

% Elements of vectors
\def\evalpha{{\alpha}}
\def\evbeta{{\beta}}
\def\evepsilon{{\epsilon}}
\def\evlambda{{\lambda}}
\def\evomega{{\omega}}
\def\evmu{{\mu}}
\def\evpsi{{\psi}}
\def\evsigma{{\sigma}}
\def\evtheta{{\theta}}
\def\eva{{a}}
\def\evb{{b}}
\def\evc{{c}}
\def\evd{{d}}
\def\eve{{e}}
\def\evf{{f}}
\def\evg{{g}}
\def\evh{{h}}
\def\evi{{i}}
\def\evj{{j}}
\def\evk{{k}}
\def\evl{{l}}
\def\evm{{m}}
\def\evn{{n}}
\def\evo{{o}}
\def\evp{{p}}
\def\evq{{q}}
\def\evr{{r}}
\def\evs{{s}}
\def\evt{{t}}
\def\evu{{u}}
\def\evv{{v}}
\def\evw{{w}}
\def\evx{{x}}
\def\evy{{y}}
\def\evz{{z}}

% Matrix
\def\mA{{\bm{A}}}
\def\mB{{\bm{B}}}
\def\mC{{\bm{C}}}
\def\mD{{\bm{D}}}
\def\mE{{\bm{E}}}
\def\mF{{\bm{F}}}
\def\mG{{\bm{G}}}
\def\mH{{\bm{H}}}
\def\mI{{\bm{I}}}
\def\mJ{{\bm{J}}}
\def\mK{{\bm{K}}}
\def\mL{{\bm{L}}}
\def\mM{{\bm{M}}}
\def\mN{{\bm{N}}}
\def\mO{{\bm{O}}}
\def\mP{{\bm{P}}}
\def\mQ{{\bm{Q}}}
\def\mR{{\bm{R}}}
\def\mS{{\bm{S}}}
\def\mT{{\bm{T}}}
\def\mU{{\bm{U}}}
\def\mV{{\bm{V}}}
\def\mW{{\bm{W}}}
\def\mX{{\bm{X}}}
\def\mY{{\bm{Y}}}
\def\mZ{{\bm{Z}}}
\def\mBeta{{\bm{\beta}}}
\def\mPhi{{\bm{\Phi}}}
\def\mLambda{{\bm{\Lambda}}}
\def\mSigma{{\bm{\Sigma}}}

% Tensor
\DeclareMathAlphabet{\mathsfit}{\encodingdefault}{\sfdefault}{m}{sl}
\SetMathAlphabet{\mathsfit}{bold}{\encodingdefault}{\sfdefault}{bx}{n}
\newcommand{\tens}[1]{\bm{\mathsfit{#1}}}
\def\tA{{\tens{A}}}
\def\tB{{\tens{B}}}
\def\tC{{\tens{C}}}
\def\tD{{\tens{D}}}
\def\tE{{\tens{E}}}
\def\tF{{\tens{F}}}
\def\tG{{\tens{G}}}
\def\tH{{\tens{H}}}
\def\tI{{\tens{I}}}
\def\tJ{{\tens{J}}}
\def\tK{{\tens{K}}}
\def\tL{{\tens{L}}}
\def\tM{{\tens{M}}}
\def\tN{{\tens{N}}}
\def\tO{{\tens{O}}}
\def\tP{{\tens{P}}}
\def\tQ{{\tens{Q}}}
\def\tR{{\tens{R}}}
\def\tS{{\tens{S}}}
\def\tT{{\tens{T}}}
\def\tU{{\tens{U}}}
\def\tV{{\tens{V}}}
\def\tW{{\tens{W}}}
\def\tX{{\tens{X}}}
\def\tY{{\tens{Y}}}
\def\tZ{{\tens{Z}}}


% Graph
\def\gA{{\mathcal{A}}}
\def\gB{{\mathcal{B}}}
\def\gC{{\mathcal{C}}}
\def\gD{{\mathcal{D}}}
\def\gE{{\mathcal{E}}}
\def\gF{{\mathcal{F}}}
\def\gG{{\mathcal{G}}}
\def\gH{{\mathcal{H}}}
\def\gI{{\mathcal{I}}}
\def\gJ{{\mathcal{J}}}
\def\gK{{\mathcal{K}}}
\def\gL{{\mathcal{L}}}
\def\gM{{\mathcal{M}}}
\def\gN{{\mathcal{N}}}
\def\gO{{\mathcal{O}}}
\def\gP{{\mathcal{P}}}
\def\gQ{{\mathcal{Q}}}
\def\gR{{\mathcal{R}}}
\def\gS{{\mathcal{S}}}
\def\gT{{\mathcal{T}}}
\def\gU{{\mathcal{U}}}
\def\gV{{\mathcal{V}}}
\def\gW{{\mathcal{W}}}
\def\gX{{\mathcal{X}}}
\def\gY{{\mathcal{Y}}}
\def\gZ{{\mathcal{Z}}}

% Sets
\def\sA{{\mathcal{A}}}
\def\sB{{\mathcal{B}}}
\def\sC{{\mathcal{C}}}
\def\sD{{\mathcal{D}}}
\def\sE{{\mathcal{E}}}
\def\sF{{\mathcal{F}}}
\def\sG{{\mathcal{G}}}
\def\sH{{\mathcal{H}}}
\def\sI{{\mathcal{I}}}
\def\sJ{{\mathcal{J}}}
\def\sK{{\mathcal{K}}}
\def\sL{{\mathcal{L}}}
\def\sM{{\mathcal{M}}}
\def\sN{{\mathcal{N}}}
\def\sO{{\mathcal{O}}}
\def\sP{{\mathcal{P}}}
\def\sQ{{\mathcal{Q}}}
\def\sR{{\mathcal{R}}}
\def\sS{{\mathcal{S}}}
\def\sT{{\mathcal{T}}}
\def\sU{{\mathcal{U}}}
\def\sV{{\mathcal{V}}}
\def\sW{{\mathcal{W}}}
\def\sX{{\mathcal{X}}}
\def\sY{{\mathcal{Y}}}
\def\sZ{{\mathcal{Z}}}

% Entries of a matrix
\def\emLambda{{\Lambda}}
\def\emA{{A}}
\def\emB{{B}}
\def\emC{{C}}
\def\emD{{D}}
\def\emE{{E}}
\def\emF{{F}}
\def\emG{{G}}
\def\emH{{H}}
\def\emI{{I}}
\def\emJ{{J}}
\def\emK{{K}}
\def\emL{{L}}
\def\emM{{M}}
\def\emN{{N}}
\def\emO{{O}}
\def\emP{{P}}
\def\emQ{{Q}}
\def\emR{{R}}
\def\emS{{S}}
\def\emT{{T}}
\def\emU{{U}}
\def\emV{{V}}
\def\emW{{W}}
\def\emX{{X}}
\def\emY{{Y}}
\def\emZ{{Z}}
\def\emSigma{{\Sigma}}

% entries of a tensor
% Same font as tensor, without \bm wrapper
\newcommand{\etens}[1]{\mathsfit{#1}}
\def\etLambda{{\etens{\Lambda}}}
\def\etA{{\etens{A}}}
\def\etB{{\etens{B}}}
\def\etC{{\etens{C}}}
\def\etD{{\etens{D}}}
\def\etE{{\etens{E}}}
\def\etF{{\etens{F}}}
\def\etG{{\etens{G}}}
\def\etH{{\etens{H}}}
\def\etI{{\etens{I}}}
\def\etJ{{\etens{J}}}
\def\etK{{\etens{K}}}
\def\etL{{\etens{L}}}
\def\etM{{\etens{M}}}
\def\etN{{\etens{N}}}
\def\etO{{\etens{O}}}
\def\etP{{\etens{P}}}
\def\etQ{{\etens{Q}}}
\def\etR{{\etens{R}}}
\def\etS{{\etens{S}}}
\def\etT{{\etens{T}}}
\def\etU{{\etens{U}}}
\def\etV{{\etens{V}}}
\def\etW{{\etens{W}}}
\def\etX{{\etens{X}}}
\def\etY{{\etens{Y}}}
\def\etZ{{\etens{Z}}}

% The true underlying data generating distribution
\newcommand{\pdata}{p_{\rm{data}}}
% The empirical distribution defined by the training set
\newcommand{\ptrain}{\hat{p}_{\rm{data}}}
\newcommand{\Ptrain}{\hat{P}_{\rm{data}}}
% The model distribution
\newcommand{\pmodel}{p_{\rm{model}}}
\newcommand{\Pmodel}{P_{\rm{model}}}
\newcommand{\ptildemodel}{\tilde{p}_{\rm{model}}}
% Stochastic autoencoder distributions
\newcommand{\pencode}{p_{\rm{encoder}}}
\newcommand{\pdecode}{p_{\rm{decoder}}}
\newcommand{\precons}{p_{\rm{reconstruct}}}

\newcommand{\laplace}{\mathrm{Laplace}} % Laplace distribution

\newcommand{\Ls}{\mathcal{L}}
\newcommand{\emp}{\tilde{p}}
\newcommand{\lr}{\alpha}
\newcommand{\reg}{\lambda}
\newcommand{\rect}{\mathrm{rectifier}}
\newcommand{\softmax}{\mathrm{softmax}}
\newcommand{\sigmoid}{\sigma}
\newcommand{\softplus}{\zeta}
\newcommand{\standarderror}{\mathrm{SE}}
\newcommand{\cl}{\text{cl}}
% \newcommand{\1}[1]{\mathds{1}\left[#1\right]}


% Wolfram Mathworld says $L^2$ is for function spaces and $\ell^2$ is for vectors
% But then they seem to use $L^2$ for vectors throughout the site, and so does
% wikipedia.
\newcommand{\normlzero}{L^0}
\newcommand{\normlone}{L^1}
\newcommand{\normltwo}{L^2}
\newcommand{\normlp}{L^p}
\newcommand{\normmax}{L^\infty}

\newcommand{\parents}{Pa} % See usage in notation.tex. Chosen to match Daphne's book.
%
%\DeclareMathOperator*{\argmax}{arg\,max}
%\DeclareMathOperator*{\argmin}{arg\,min}
%
%\DeclareMathOperator{\sign}{sign}
% \DeclareMathOperator{\Tr}{Tr}
\let\ab\allowbreak

\newcommand{\seq}{\subseteq}

\renewcommand{\wt}{\mathsf{weight}}

\mathchardef\mhyphen="2D
\newcommand{\val}{\text{val}}
\newcommand{\calP}{\mathcal{P}}
\newcommand{\calH}{\mathcal{H}}
\newcommand{\calM}{\mathcal{M}}
\newcommand{\calB}{\mathcal{B}}
\newcommand{\calE}{\mathcal{E}}
\newcommand{\calQ}{\mathcal{Q}}
\newcommand{\calR}{\mathcal{R}}
\newcommand{\calD}{\mathcal{D}}
\newcommand{\calW}{\mathcal{W}}
\newcommand{\calS}{\mathcal{S}}
\newcommand{\calT}{\mathcal{T}}
\newcommand{\calI}{\mathcal{I}}

% shape terms
\newcommand\numberthis{\addtocounter{equation}{1}\tag{\theequation}}
\newcommand{\midshape}{\mathsf{MidShape}}
\newcommand{\leftshape}{\mathsf{LeftShape}}
\newcommand{\rightshape}{\mathsf{RightShape}}
\newcommand{\dsos}{\mathsf{dsos}}
\newcommand{\mul}{\text{mul}}
\newcommand{\ind}[1]{1_{{#1}}}
\newcommand{\din}{d_{\text{in}}}
\newcommand{\dout}{d_{\text{out}}}
\newcommand{\sbm}{D_{\text{SBM}}}
\newcommand{\cycle}{\ind{\calE}}
\newcommand{\bdd}{\ind{\calF}}
\newcommand{\exc}{\mathsf{excess}}
\newcommand{\cc}{\mathsf{\#cc}}
%\newcommand{\vtxcoef}{\mathsf{vtxcoef}}
\newcommand{\float}{\mathsf{float}}
\newcommand{\fs}{\mathsf{float \mhyphen stick}}
\newcommand{\fc}{\mathsf{float\mhyphen cyc}}
\newcommand{\cyclen}{\mathsf{cyclen}}
\newcommand{\mvs}{\text{MVS}}
\newcommand{\smvs}{\text{SMVS}}
\newcommand{\dang}{\mathsf{dangling}}
\newcommand{\gam}{\gamma}
\newcommand{\Id}{\text{Id}}
\newcommand{\trivial}{\text{trivial}}
\newcommand{\cval}[1]{\mathsf{Colorval}[#1]}	
\newcommand{\crosscol}{\text{cross-color}}
%\renewcommand{\int }{\text{Int}}
\newcommand{\slack}{\text{slack}}
\newcommand{\Glyph}{\mathcal{G}}
\newcommand{\glyph}{\tex{glyph}}
\newcommand{\yes}{\text{Yes}}
\newcommand{\no}{\text{No}}
\newcommand{\mix}{\text{Mix}}
\renewcommand{\cp}{\text{CP}}
\newcommand{\pure}{\text{pure-color}}
\newcommand{\depth}{\text{depth}}
\newcommand{\res}{\text{res}}
\newcommand{\girth}{\text{girth}}
\newcommand{\vap}{\text{vap}}
\newcommand{\boundary}{\text{boundary}}
\newcommand{\chunk}{\text{chunk}}
\newcommand{\cm}{\mathsf{M} }
\renewcommand{\phantom}{\text{phantom}}
\newcommand{\exwedge}{\text{exc-wedge}}
\newcommand{\cnorm}{c_{\mathsf{norm}}}
\newcommand{\ceta}{c_\eta}
\newcommand{\edgeco}{\left(\eta\cdot\sqrt{\frac{p}{1-p}} \right)}
\newcommand{\vtxco}{\left(\frac{1}{k}  \right)}
\newcommand{\sval}{\text{SMVS-val}}
\newcommand{\fantome}{\text{distinct phantom}}
% \newcommand{\iso}{\mathsf{\# isolated}}
\newcommand{\isocc}{\mathsf{\#isocc}}
\newcommand{\cnt}{\mathsf{count}}
\newcommand{\gp}{{\gamma'}}
\renewcommand{\cp}{\mathscr{P}}
\newcommand{\col}{\text{color}}
\newcommand{\aut}{\mathsf{Aut}}
\newcommand{\colaut}{\mathsf{ColAut}}
\newcommand{\conn}{\mathsf{conn}}
\newcommand{\corr}{\mathsf{correct}}
\newcommand{\rbcorr}{\mathsf{Rb}}
\newcommand{\rbp}{\mathsf{Rb}_{\mathsf{pinned}}}
\newcommand{\calG}{\mathcal{G}}
\newcommand{\fp}{\mathsf{P}}
\newcommand{\bti}{\mathsf{BacktrackingInt}}
\newcommand{\nbi}{\mathsf{NonBacktrackInt}}
\newcommand{\cmp}{\widetilde{\cm}}



\newcommand{\Erdos}{Erd\H{o}s\xspace}
\newcommand{\Renyi}{R\'enyi\xspace}
\newcommand{\Lovasz}{Lov\'asz\xspace}
\newcommand{\Juhasz}{Juh\'asz\xspace}
\newcommand{\Bollobas}{Bollob\'as\xspace}
\newcommand{\Komlos}{Koml\'os\xspace}
\newcommand{\Luczak}{\L uczak\xspace}
\newcommand{\Szemeredi}{Szemer\'edi\xspace}
\newcommand{\Hastad}{H{\aa}stad\xspace}
\newcommand{\Holder}{H\"{o}lder\xspace}
\newcommand{\Brandao}{Brand\~ao\xspace}
\newcommand{\Bezout}{B\'ezout\xspace}
\newcommand{\Turan}{Tur\'{a}n\xspace}



% Cicle
\renewcommand*{\circle}[1]{\scalebox{0.85}{\footnotesize
    \tikz[baseline=(char.base)]{
        \node[shape=circle,draw,inner sep=0.5pt, minimum size=14pt](char) {   \ifx&#1&
        \color{white} $i$
        \else
        $#1$
        \fi};
    }}}

% Square
\renewcommand*{\square}[1]{\scalebox{0.85}{\footnotesize
    \tikz[baseline=(char.base), square/.style={regular polygon,regular polygon sides=4}]{
        \node[draw,square, inner sep=0pt, minimum size=18pt](char) {
        \ifx&#1&
        \color{white} $t$
        \else
        $#1$
        \fi};
    }}}
    
% Triangle
% \renewcommand*{\triangle}[1]{\scalebox{0.85}{\footnotesize
%     \tikz[baseline=(char.base), square/.style={regular polygon,regular polygon sides=3}]{
%         \node[draw,square, inner sep=-3pt, minimum size=18pt](char) {
%         \ifx&#1&
%         \color{white} $t$
%         \else
%         $#1$
%         \fi};
%     }}}

% Hexagon
\newcommand*{\hexagon}[1]{\scalebox{0.85}{\footnotesize
    \tikz[baseline=(char.base), square/.style={regular polygon,regular polygon sides=6}]{
        \node[draw,square, inner sep=-3pt, minimum size=18pt](char) {
        \ifx&#1&
        \color{white} $t$
        \else
        $#1$
        \fi};
    }}}

\usepackage[notes=true,later=false,camera=false]{dtrt}
\usepackage{algorithm}
\usepackage{algorithmic}
\usepackage{amsmath}
\DeclareMathOperator{\arccosh}{arccosh}
\DeclareMathOperator{\arcsinh}{arcsinh}
\DeclareMathOperator{\arctanh}{arctanh}
\usepackage{tikz}
\newcommand{\chebyerror}{\mathsf{ChebyshevCorrection}}
 %
\newcommand{\todo}[1]{\textcolor{red}{(#1)}}
\newcommand{\etal}{et.~al.\xspace}
\allowdisplaybreaks
\newcommand{\calL}{\mathcal{L}}
\newcommand{\q}{\mathfrak q}
\newcommand{\p}{\mathfrak P}
\newcommand{\RemoveAuthor}{0}
\usepackage{comment}

\usepackage{mathpazo}
\usepackage{microtype}

%%%%%%%%%%%%%%% Author Notes
\ifnum\showauthornotes=1
\newcommand{\Authornote}[3]{{\sf\footnotesize\color{#3}{[#1: #2]}}}
\else
\newcommand{\Authornote}[3]{}
\fi


\newcommand{\jnote}[1]{\Authornote{Jeff}{#1}{blue}}
\newcommand{\anote}[1]{\Authornote{Aaron}{#1}{cyan}}
\newcommand{\mnote}[1]{\Authornote{Madhur}{#1}{red}}
\newcommand{\snote}[1]{\Authornote{Sofia}{#1}{orange}}





\begin{document}
\title{Sharp Phase Transition for Ellipsoid Fitting}%
\author{
Sofia de la Cerda
\thanks{{University of Chicago}. \textit{sofiasl@uchicago.edu} }
\and 
Aaron Potechin 
\thanks{{University of Chicago}. \textit{potechin@uchicago.edu}}
\and
Madhur Tulsiani
\thanks{{Toyota Technological Institute at Chicago}. \textit{madhurt@ttic.edu}}
\and
Jeff Xu
\thanks{{Toyota Technological Institute at Chicago}.\textit{jeffxusichao@ttic.edu}}
}
\maketitle
% \begin{abstract}
% 	We %resolve 
%     \snote{prove} the ellipsoid fitting conjecture of Saunderson,
% Chandrasekaran, Parrilo, and Willsky. \cite{SCPW12} %up to a $o(d^2)$ factor.
% %  Namely, for \(m\le (1-o_d(1))d^2/4\) independent Gaussian
% %points in \(\mathbb{R}^d\), we show that with high probability there exists an
% %ellipsoid passing through all \(m\) points. At the same time, for $m\geq (1+o_d(1)) d^2/4$, we show that no such ellipsoid exists. 
% Concretely, for $m=\alpha d^2+o(d^2)$ independent Gaussian points in dimension $d$, we show that with high probability as $d\to\infty$,
% \begin{enumerate}
% 	\item for $\alpha < 1/4$, there exists a centered ellipsoid passing through all $m$ points;
% 	\item for $\alpha > 1/4$, no such ellipsoid exists.
% \end{enumerate}
% %Therefore, we pin down the phase transition to be precisely $d^2/4$ in the leading order.
% That is, the ellipsoid fitting problem has a sharp phase transition at $\alpha=\frac 14$.

% We build on and extend the iterative scheme of \cite{PotechinXu2026Theta} for constructing a solution of the SDP for the \Lovasz theta function of a random $G(n,1/2)$ graph with value $(1+o_n(1))\sqrt{n}$ (which is optimal up to the vanishing $o_n(1)$ factor). %, introduced in the context of the \Lovasz theta function of random graphs, to construct basic SDP solutions at sharp thresholds. 
% Our analysis proceeds via graph matrix decomposition, a method for analyzing random matrices with correlated entries that are polynomials of the underlying input. In our analysis, we obtain free independence as well as sharp spectral bounds for a family of graph matrices we call backbone-dangling shapes. %studying correlated random matrices: we obtain free independence as well as sharp spectral bounds for a broad family of random matrices with correlated entries that are polynomials of the underlying Gaussian input.

% \end{abstract}. 
\begin{abstract}
	We resolve the ellipsoid fitting conjecture of Saunderson,
Chandrasekaran, Parrilo, and Willsky~\cite{SCPW12} up to a vanishing
factor.
%  Namely, for \(m\le (1-o_d(1))d^2/4\) independent Gaussian
%points in \(\mathbb{R}^d\), we show that with high probability there exists an
%ellipsoid passing through all \(m\) points. At the same time, for $m\geq (1+o_d(1)) d^2/4$, we show that no such ellipsoid exists. 
Concretely, for $m$ independent Gaussian points in dimension $d$, we show that with high probability,\begin{enumerate}
	\item for $m \leq (1-o_d(1)) \cdot d^2/4$, there exists a centered ellipsoid passing through all $m$ points;
	\item for $m\geq (1+o_d(1) )\cdot d^2/4$, no such ellipsoid exists.
\end{enumerate}
This confirms that the ellipsoid fitting problem has a sharp phase transition at $d^2/4$.
%Consequently, we determine the phase transition at the sharp threshold $d^2/4$ up to lower-order terms.
%
%{\color{orange}Sofia: Concretely, there exists a function $f(d)=o(d^2)$ such that for $m$ independent Gaussian points in dimension $d$, with high probability:
%\begin{enumerate}
%	\item if $m < d^2/4-f(d)$, there exists a centered ellipsoid passing through all $m$ points;
%	\item if $m > d^2/4+f(d)$, no such ellipsoid exists.
%\end{enumerate}
%Consequently, we determine the phase transition at the sharp threshold $d^2/4$ up to lower-order terms.} 
%
%\anote{I concur with dropping the paragraph below as having a short abstract is nice when experts are already familiar with the problem and the importance of the result. I edited the previous sentence.}
%\snote{I also agree to drop the paragraph below.}
%\jnote{one option is to drop the following paragraph..not sure..}
%We build on and extend the iterative scheme of \cite{PotechinXu2026Theta} for constructing a solution of the SDP for the \Lovasz theta function of a random $G(n,1/2)$ graph with optimal value $(1+o_n(1))\sqrt{n}$. %, introduced in the context of the \Lovasz theta function of random graphs, to construct basic SDP solutions at sharp thresholds. 
%Our analysis proceeds via graph matrix decomposition, a method for analyzing random matrices with correlated entries that are polynomials of the underlying input. In our analysis, we obtain free independence as well as sharp spectral bounds for a family of graph matrices we call backbone-dangling shapes. %studying correlated random matrices: we obtain free independence as well as sharp spectral bounds for a broad family of random matrices with correlated entries that are polynomials of the underlying Gaussian input.
%\jnote{Sofia i reverted back to the prior abstract cuz some changes are not in snote...also lets not use $m=\al d^2$ because $\al$ is reserved for graph matrix throughout. Let's use $m = c_{\nu} d^2$ if needed.}




%\jnote{Sofia, please try not to comment out what's already written..feel free to draft an alternative version and leave as it is}

\end{abstract}
%\snote{Note: I changed the formatting of the emails}
\clearpage
\newpage
%\newpage
\thispagestyle{empty}
\setcounter{page}{0}


\thispagestyle{empty}
\setcounter{page}{0}
%\newpage
\thispagestyle{empty}
\setcounter{page}{0}
\tableofcontents
%\clearpage\newpage
\clearpage \newpage

\section{Introduction}
%\todo{1.diagonal error terms 2. verification of the inner product value for refutation}

A basic question at the intersection of convex geometry, random matrix theory,
and semidefinite programming is the following: given a sample of size $m$ from the standard normal distribution \(v_1,\ldots,v_m\in \mathbb R^d\),
%independent random points
% \(v_1,\ldots,v_m\in \mathbb R^d\) sampled from $\mathcal N(0, I_d)$, 
when does there exist an (origin-centered, possibly degenerate) ellipsoid passing
through all of them?  Equivalently, for what values of \(m=m(d)\) is there,
with high probability, a positive semidefinite matrix \(\Lambda\succeq 0\) such that  $v_i^\top \Lambda v_i = 1$ for all $i\in [m]$?
%\begin{definition}[Ellipsoid Matrix]
%Given points $(v_1,\ldots,v_m\in \mathbb R^d)$, an ellipsoid fitting the given points is a symmetric matrix $\Lambda \in \R^{d \times d}  $ s.t.  such that
%\(
%    v_i^\top \Lambda v_i = 1
%\) for all $i\in [m]$ ,
%and $\Lambda \succeq 0 $. 
%\end{definition}
%Following the convention of \cite{PTVW22,KD22, HKPX23}, throughout this paper we work
%with the normalization under which \(\mathbb E\|v_i\|_2^2=1\).  

The ellipsoid fitting problem was introduced by Saunderson \cite{Saunderson11}
and studied by Saunderson, Chandrasekaran,
Parrilo, and Willsky \cite{SCPW12,SPW13}. The latter authors conjectured that the
problem exhibits a sharp phase transition at
\[
    m \approx \frac{d^2}{4}.
\]
%More precisely, for every fixed \(\varepsilon>0\), if
%\(m\le (1-\varepsilon)d^2/4\), then \(v_1,\ldots,v_m\) should have the
%ellipsoid fitting property with probability \(1-o_d(1)\), while if
%\(m\ge (1+\varepsilon)d^2/4\), then the property should fail with probability
%\(1-o_d(1)\). 
% The upper side of this threshold is consistent with the
%dimension of the linear span of the cone of \(d\times d\) positive semidefinite
%matrices, which is \(\binom{d+1}{2}\sim d^2/2\). From this perspective, the conjecture predicts that
%the positivity constraint forces an additional factor-two loss, lowering the
%feasibility threshold from \(d^2/2\) to \(d^2/4\).
%\paragraph{Negative Result for Ellipsoid Fitting} 

%\snote{For $m\leq \binom{d+1}2$ the matrices $(v_iv_i^\top)_{i\in [m]}$ are linearly independent, and we may interpret the existence of a symmetric (not necessarily PSD) matrix $\Lambda$ such that $v_i^\top \Lambda v_i=1$ for all $i$ as a system of $m$ linear equations in $\binom{d+1}2$ unknowns. For $m=\binom{d+1}2$ linear independence implies almost surely there's exactly one solution, so for $m>\binom{d+1}2$ there is no symmetric solution.} \jnote{i think its ok to call this dimension-counting :)}
The natural dimension-counting argument %suggests impossiblity result below the threshold of 
shows that this is impossible when $m >  \binom{d+1}{2} \sim d^2/2$, the dimension of the space of $d\times d$ symmetric matrices. From this perspective, the conjecture predicts that the positive semidefiniteness constraint forces an additional factor-two loss, lowering the feasibility threshold from \(d^2/2\) to \(d^2/4\). To the best of our knowledge, this na\"{i}ve dimension-counting argument remains the only known result—and hence the state of the art—for the impossibility direction of ellipsoid fitting prior to this work.

%\paragraph{Positive Result for Ellipsoid Fitting}

A long line of work has progressively improved the lower bound on the feasible
regime.  Early results established feasibility for
\(m\le O(d^{6/5-\varepsilon})\) \cite{SPW13}. 
%\snote{In the past few years, there has been a revival of interest in this problem, with connections to degree-two Sum-of-Squares and random matrix constructions,
%obtaining bounds of the form \(m\le O(d^{3/2-\varepsilon})\) \cite{GJJPR20} and then
%\(m\le d^2/\operatorname{polylog}(d)\) \cite{PTVW22,KD22}.}
In the past few years, there has been a revival of interest in this problem, with connections to degree-two Sum-of-Squares and random matrix constructions,
obtaining bounds of the form \(m\le O(d^{3/2-\varepsilon})\) \cite{GJJPR20} and then
\(m\le d^2/\operatorname{polylog}(d)\) \cite{PTVW22,KD22}. Most recently, a line of independent works \cite{HKPX23,BMMP24,TW25} has obtained bounds that are tight up to an absolute constant factor, establishing feasibility for \(m\le c d^2\) for some universal constant \(c>0\).\footnote{The coefficient is not explicitly tracked in the existing analysis, though we believe a safe estimate of $c=1/1000$ would have sufficed in \cite{HKPX23}, albeit with a substantial gap from the conjectured threshold of $1/4$.} These works use substantially different techniques, but the underlying construction is essentially the same identity-perturbation ansatz analyzed in prior works. In parallel with this line of progress, \cite{MD24, MB25} provided non-rigorous statistical-physics evidence pointing to \(d^2/4\) as the transition threshold.



%At the same time with this line of progress, \cite{MB25} also provides non-rigorous evidence from statistical physics hinting upon $\frac{d^2}{4}$ as the transition threshold.

%  More
%recently,  \cite{HKPX23} analyzed the
%identity-perturbation construction using graphical matrix decomposition and
%proved feasibility for
%\(
%    m \le c d^2
%\)
%for some small constant \(c>0\), resolving the feasible direction of the SCPW
%conjecture up to a constant factor. We note that related constant-factor results were obtained independently and concurrently by Tulsiani and Wu \cite{TW25} and by Bandeira et al.  via different techniques.

%\snote{Despite this progress, ...}
Despite this active progress, the sharp constant \(1/4\) remains open. We emphasize that the
remaining gap is not merely technical in the analysis, but more importantly, conceptual: empirical evidence,
together with heuristic predictions from statistical physics, suggests that the
previously studied constructions already fail well below the conjectured
threshold, around \(m\approx d^2/10\). Thus, even setting aside the difficulty of
rigorously analyzing a construction near the threshold, a central challenge is to
find, or even identify, a plausible candidate construction capable of reaching
the \(d^2/4\) regime.  

%Among prior works on the ellipsoid fitting conjecture, the closest to ours is \cite{HKPX23}. In particular, we also adopt the language and techniques of graph matrices to construct and analyze the random matrices arising in the problem. A key new ingredient in our work is the application and extension of techniques from the recent result \cite{PotechinXu2026Theta}, which obtains a sharp bound for the \Lovasz-theta function of a $G(n,1/2)$ random graph. We describe the connection between these works and our main technical ideas in detail in the following section.

%\snote{Despite the line of active progress, the sharp constant \(1/4\) remains open. We emphasize that the remaining gap is not merely technical for the analysis, but more importantly, conceptual: empirical evidence, together with heuristic predictions from statistical physics, suggests that the previously studied identity perturbation construction already fails well below the conjectured threshold, around \(m\approx d^2/10\). Thus, even setting aside the difficulty of rigorously analyzing a construction near the threshold, a central challenge is to find, or even identify, a plausible candidate construction capable of reaching the \((1/4-\epsilon)d^2\) regime.}
%%Despite the line of active progress, the sharp constant \(1/4\) remains open. We emphasize that the remaining gap is not merely technical for the analysis, but more importantly, conceptual: empirical evidence, together with heuristic predictions from statistical physics, suggests that the previously studied constructions already fail well below the conjectured threshold, around \(m\approx d^2/10\). Thus, even setting aside the difficulty of rigorously analyzing a construction near the threshold, a central challenge is to find, or even identify, a plausible candidate construction capable of reaching the \(d^2/4\) regime.
%%
%%\paragraph{Out Results}
%%
%%In this work, we resolve the conjecture by showing both sides of the conjecture:
%%\begin{enumerate}
%%	\item \textbf{Positive Result:} it is possible to fit an ellipsoid all the way up to the $m< d^2/4$ threshold in the leading order;
%%	\item \textbf{Negative Result:} it is impossible fit an ellipsoid once we exceed the  $m> d^2/4$ threshold.
%%\end{enumerate}
%%
%% %Of the line of works on the ellipsoid fitting conjecture, the paper most similar to this paper is~\cite{HKPX23}
%%Concretely, we prove the following theorems.
%%
%%
%%\begin{theorem} \label{thm:ellipsoid-thm}
%%For \(m\le (1-o_d(1))\frac{d^2}{4}\), let \(v_1,\dots,v_m\sim \mathcal N(0,I_d/d)\) be independent \(d\)-dimensional Gaussian points. Then, with probability $1-o_d(1)$, there exists an ellipsoid passing through each point.
%%\end{theorem}
%%
%%\begin{theorem} \label{thm:ellipsoid-refutation-thm}
%%For \(m\ge (1+o_d(1))\frac{d^2}{4}\), let \(v_1,\dots,v_m\sim \mathcal N(0,I_d/d)\) be independent \(d\)-dimensional Gaussian points. Then, with probability $1-o_d(1)$, there does not exist an ellipsoid passing through each point.
%%\end{theorem}
%%
%%Of the line of works on the ellipsoid fitting conjecture, the paper most similar to this paper is~\cite{HKPX23}: in particular, we also appeal to the techniques and language of graph matrices to construct and analyze the random matrices that arise in this problem. Most importantly, our work applies and extends the techniques from the recent result \cite{PotechinXu2026Theta} that obtains sharp bound for \Lovasz-Theta bound for $G(n,\frac{1}{2})$. We describe in details of the connection between these works and our main techniques in the following section.
%
%%
%
%
%{\color{orange}Sofia: \paragraph{Our Results}
%
%In this work, we resolve the ellipsoid fitting conjecture by establishingh a sharp phase transition at \(m=d^2/4\), up to a $o(f(d))$ term. With high probability:
%\begin{enumerate}
%    \item \textbf{Feasibility below the threshold:} when $m<\frac 14 d^2-f(d)$ there exists a fitting ellipsoid;
%    
%    \item \textbf{Infeasibility above the threshold:} when $\alpha>\frac 14d^2+f(d)$ no such ellipsoid exists.
%  
%\end{enumerate}
%
%More concretely, we prove the following theorem.
%
%\begin{theorem}[Main theorem]
%\label{thm:ellipsoid-thm}
%Let $d\to\infty$, $m=m(d)$ and $x_1, \dots, x_m\sim \mathcal N(0, I_d)$ be a sample of size $m$ from the standard normal distribution in $\R^d$. The probability that there exists an ellipsoid passing through all the points in the sample satisfies
%    $$\lim_{d\to\infty}\mathbb P\bigg[\exists\: \Lambda\in \textup{Sym}_d: \Lambda \textup{ is PSD}, \: \left<\Lambda x_i, x_i\right>=1\:\forall\: i\in [m]\bigg]=\begin{cases}
%        1 & m<\frac 14 d^2-f(d) \\
%        0 & m>\frac 14 d^2+f(d).
%    \end{cases}$$
%\end{theorem}
%}

\paragraph{Our Results}
%
In this work, we resolve both directions of the ellipsoid fitting conjecture and establish a sharp phase transition at \(m=d^2/4\), up to lower-order terms:
\begin{enumerate}
    \item \textbf{Feasibility below the threshold:} an ellipsoid can be fitted through the points throughout the regime
    \(
        m < \frac{d^2}{4}\,;
    \)
    
    \item \textbf{Infeasibility above the threshold:} no such ellipsoid exists once
   \(
        m > \frac{d^2}{4}\,.
    \)
  
\end{enumerate}

More concretely, we prove the following two theorems.

\begin{theorem}[Main theorem for the positive side]
\label{thm:ellipsoid-thm}
Let
\(
    m \leq \bigl(1-o_d(1)\bigr) \cdot d^2/4\)
and let \(v_1,\dots,v_m\sim \mathcal N(0,I_d/d)\) be independent \(d\)-dimensional Gaussian vectors. Then, with probability \(1-o_d(1)\), there exists an ellipsoid passing through all \(m\) points.
\end{theorem}
%
\begin{theorem}[Main theorem for the negative side]
\label{thm:ellipsoid-refutation-thm}
Let
\(
    m \geq \bigl(1+o_d(1)\bigr) \cdot d^2/4
\)
and let \(v_1,\dots,v_m\sim \mathcal N(0,I_d/d)\) be independent \(d\)-dimensional Gaussian vectors. Then, with probability \(1-o_d(1)\), there does not exist any ellipsoid passing through all \(m\) points.
\end{theorem}

% \begin{remark}
% 	Our results apply with a slack of \(
% 	(1\pm \varepsilon) \cdot \frac{d^2}{4} 
% 	 \)
% 	 with $\varepsilon = 1/\poly\log(d) $ while we do not optimize over this dependence.
% \end{remark}
% \snote{
% \begin{remark}
%     These theorems establish the threshold at $\frac{d^2}{4}$ and hold with $o_d(1)$ replaced by $\varepsilon = 1/\poly\log(d)$. We make no attempt to optimize this polylogarithmic factor.
% \end{remark}
% }
\begin{remark}
    These theorems establish the threshold at $\frac{d^2}{4}$ and hold with $o_d(1)$ instantiated as $\varepsilon = 1/\poly(\log \log d)$. We make no attempt to optimize this factor.
\end{remark}

We prove both theorems by exhibiting and analyzing explicit SDP witnesses. The two solutions are drastically different in their constructions while their analyses share a strikingly similar structure.  Among prior works on the ellipsoid fitting conjecture, the closest to ours is~\cite{HKPX23}: we also make use of the language and techniques of graph matrices to construct and analyze the correlated random matrices arising in the problem. A key new ingredient is the application and extension of techniques from the recent work~\cite{PotechinXu2026Theta}, which establishes a sharp bound for the \Lovasz-theta function of \(G(n,1/2)\). In the following section, we explain the connection between these results and describe the main technical ideas underlying our proof.

\paragraph{Roadmap of Our Work}

%, so 
%while being familiar with \cite{PotechinXu2026Theta} is highly recommended as the \Lovasz Theta function setting is somewhat simpler, it is not required. %no prior familiarity with \cite{PotechinXu2026Theta} is needed. 
%
%Both of results for ellipsoid and refutation are proven via explicit SDP witnesses.
%Throughout technical overview section, we focus primarily on the ellipsoid fitting.
%We give a brief overview of our candidate ellipsoid matrix construction.
%Unlike prior works, which use a one-shot construction for the candidate matrix,
%our construction is iterative, following the approach of
%\cite{PotechinXu2026Theta}.  To fully specify the construction in a form
%suitable for our technical analysis, we first give a self-contained overview of
%graph matrices specialized to our setting in \cref{sec:graph-matrix}.  With
%this graph-matrix framework in place, we then complete the description of our
%iterative construction in \cref{sec:iterative-construction}.  
%
%After describing most ideas for our ellipsoid fitting, we describe how our iterative scheme can be modified to obtain a dual witness for refutation as well.
%
%The remaining sections highlight the main technical challenges specific to our work for the ellipsoid fitting problem. The remainder of our work is organized as follows.

Both our feasibility and refutation results are established via explicit SDP solutions.

In the next section, we provide an overview of our techniques and highlight the main technical challenges that arise in the construction and the analysis. We focus exclusively on the positive side of the conjecture, namely, the existence of a fitting ellipsoid below the threshold. We then explain in \cref{sec:refutation} how these arguments can be adapted to establish the corresponding refutation result above the threshold.  As a preview, all of the technical ingredients developed for the primal program will carry over to the refutation argument with only minor modifications in the setup of the iterative procedure. 

The remainder of the paper is organized as follows.

\begin{enumerate}
	\item In~\cref{sec:refutation}, we show that an analogous procedure produces an SDP dual witness that refutes the existence of a fitting ellipsoid for $m>d^2/4$.
	\item The Gram matrix $M=\calL\calL^\top$ and its inverse play a key role in our construction. In~\cref{sec:MP-sec}, we give an explicit inversion via orthogonal polynomials and graph matrices.
%that if $m = c_\nu \cdot d^2$ for some $c_{\nu} < \frac{1}{4}$ then the limiting distribution of the eigenvalues of $M$ as $d \to \infty$ is the Marchenko-Pastur distribution with parameter $\gamma = \frac{2m}{d^2} = 2c_{\nu} $. We then decompose $M^{-1}$ in terms of the orthogonal polynomials for the Marchenko-Pastur distribution and show how these polynomials can be interpreted in terms of graph matrices. %formally establish the inverse of $M$ via polynomial expansion over the Marchenko-Pastur distribution of the corresponding parameter $\gamma$. 
This justifies the correction mechanism of our iterative scheme and, more importantly, enables a decomposition of $M^{-1}$ in the graph-matrix basis that is additionally amenable to a stability analysis.  \item In~\cref{sec:inner-matrix-Q}, we establish the free independence of backbone-dangling shapes as well as the variance of the inner matrix $Q$ in the limit.
\item Finally, in~\cref{sec:formal-truncation}, we incorporate details for the truncation procedure to instantiate the construction,
    and complete the proof of the main theorem.
\end{enumerate} 
%


\paragraph{Notation.}
Throughout this work, all probabilistic statements hold with probability
\(1-o_d(1)\).  We write \(\|\cdot \|_{sp}\) for the spectral norm of a matrix. 
% We work with the parameter $\gam = \frac{2m}{d^2}$ as a proxy for the threshold.
%Additionally, we encode the affine constraints using
%\[
%\calL(X)\coloneqq (v_t^\top Xv_t)_{t\in[m]},
%\qquad
%\calL^*(w)\coloneqq \sum_{t\in[m]} w[t]v_tv_t^\top,
%\qquad
%M\coloneqq \calL\calL^*.
%\]
We use \(\tau\) as a generic placeholder for an arbitrary shape; accordingly, \(\cm_\tau\) denotes the graph matrix associated with \(\tau\).  The symbols \(\alpha\) and \(\beta\)
are reserved for specific shapes that play distinguished roles in our analysis. We additionally use the following notations:
\begin{enumerate}
	\item $I_d$ denotes the identity matrix of dimension $d\times d$. We ignore the dependence on $d$ when the dimension is clear;
	\item  $\gam = \frac{2m}{d^2}$ is a proxy for the threshold;
	\item $\calL$ evaluates the constraint values, $\calL^*$ lifts coefficient vectors to linear combinations of the rank-one forms $v_tv_t^\top$, and $M =\calL\calL^\top$
is their Gram matrix. \[
\calL(X)\coloneqq (v_t^\top Xv_t)_{t\in[m]},
\qquad
\calL^*(w)\coloneqq \sum_{t\in[m]} w[t]v_tv_t^\top,
\qquad
M\coloneqq \calL\calL^*.
\]
\item  
Let $\calS:=\operatorname{Im}(\mathcal L^*)
=\operatorname{span}\{v_iv_i^\top:i\in[m]\}$ and let  \[ \Pi \coloneqq \Pi_S  = \calL^*M^{-1}\calL \] denote the orthogonal projection to $\calS$, and $\Pi_S^\perp$ the projection to its orthogonal complement.

\item   $\eta_\tau$ denotes the deviation vector $\eta_\tau = \calL(\cm_\tau)$ for a shape $\tau$. When the subscript is left out, we use $\eta = \eta_\emptyset \coloneqq \calL(I_d) - \1_m $ as well as $\eta_0$ interchangeably to denote the raw deviation of the identity perturbation.
\item We use the letter $\calR$ to denote the perturbation matrix from the identity-perturbation construction, \[ \calR \coloneqq  \calL^* M^{-1} \eta =  \calL^*M^{-1} (\calL (I_d) -\1_m)\,.   \] In other words, $I_d -\calR $ satisfies the linear constraints. 
\item We define the variance of a $d\times d$ matrix $X$  as \(
\Var(X) \coloneqq \frac{1}{d }\cdot\E \Tr[ XX^T ]\,. \) With some abuse of notation,  we also write $\Var(\tau) = \Var(\cm_\tau)$ for the graph matrix of shape $\tau$.
\end{enumerate}
%\snote{What does the dimension of a matrix mean here?}
%\jnote{added 'symmetric"}
%
%Thus, $\calL$ evaluates the constraint values, $\calL^*$ lifts coefficient vectors to linear combinations of the rank-one forms $v_tv_t^\top$, and $M =\calL\calL^\top$
%is their Gram matrix. We use $\eta_\tau$ to denote deviation vector $\eta_\tau = \calL(\cm_\tau)$ for a shape $\tau$. When the subscript is left out, we use $\eta = \eta_\emptyset \coloneqq \calL(\1_m) - \1_m $ as well as $\eta_0$ interchangeably to denote the raw deviation of identity perturbation.




\paragraph{Concurrent and Independent Work} While preparing this manuscript, we became aware of two concurrent and independent works~\cite{KS26} and~\cite{MW26} that also obtain comparable results.

% \todo{add discussion}
% \todo{add: our results also apply to hold under universality}
\paragraph{Acknowledgments} We are grateful to Sidhanth Mohanty for bringing the independent, concurrent work to our attention. We also thank Frederic Koehler and Youngtak Sohn for coordinating the preparation of the manuscripts.

  
\paragraph{AI Disclosure}  We used generative AI tools to help polish the writing. All the conceptual contributions are our own, and we take full responsibility for any errors. We also note that the connection to the prior work of \cite{kogan2025extremal} was brought to the attention of one of the authors by generative AI after we completed our proof via graph matrices. %had initially completed the written proof.

%\paragraph{Connection to \Lovasz-Theta Function} In particular, the current state of the ellipsoid fitting problem closely parallels that of the \Lovasz-Theta function on random graphs $G(n,\frac{1}{2})$, with similarity going beyond the fact that they both concern the behavior of basic SDP on random input. Several structural and technical parallels are worth highlighting:\begin{enumerate}
%	\item The input distributions of i.i.d. Gaussian samples and $G(n,\frac{1}{2})$ are highly similar, or even identical, from a universality perspective;
%%	\item The stage of our standing is essentially the same (especially prior to \cite{PotechinXu2026Theta}) where we have an explicit witness construction down to the final constant factor. Moreover, similar techniques have been employed on both fronts. 
%	\item The state of progress is also strikingly similar, especially prior to \cite{PotechinXu2026Theta}: in both settings, explicit witness constructions were known that reached the correct scaling but fell short of the conjectured sharp constant. Moreover, seeking a candidate construction for the SDP dual witness is an interesting challenge on its own;
%	\item Prior best bounds are (can be) essentially obtained by the same set of techniques by studying correlated random matrices with entries that are polynomials of the underlying random sample;%	\item The threshold of $m\approx \frac{d^2}{4}$ also has a natural interpretation in the other context: for Theta function, the constraints require each each entry of the dual witness in pursuit to be $0$ on each non-edge in the graph sample, which in turn corresponds to $\frac{1}{2} \cdot  \binom{d}{2} =\frac{d^2}{4}$ many entrywise constraints (as opposed to affine constraints) of the dual witness.
%\item The threshold \(m\approx d^2/4\) also mirrors the constraint count in the
%\Lovasz-Theta problem. In the Theta function, the desired witness matrix
%is constrained on the non-edges, which account for roughly half of the
%\(\binom n2\) off-diagonal entries, i.e.,
%\(
%    \frac12\binom n2 \sim \frac{n^2}{4}.
%\)
%Thus the \(d^2/4\) threshold in ellipsoid fitting is naturally analogous to the
%number of entrywise non-edge constraints in the random-graph theta problem.
%	 	 \end{enumerate}
%

%\paragraph{Connection to the \Lovasz-Theta Function.}
%The current state of the ellipsoid fitting problem closely parallels that of the
%\Lovasz-Theta function on random graphs \(G(n,1/2)\). The analogy goes
%beyond the fact that both problems concern basic SDPs on natural random inputs.
%Several structural and technical parallels are worth highlighting:
%\begin{enumerate}
%    \item The underlying input distributions---i.i.d. Gaussian samples in the
%    ellipsoid fitting problem and the random graph \(G(n,1/2)\) in the
%    \Lovasz-Theta problem---are highly analogous, and often behave
%    similarly from a universality perspective.
%
%    \item The state of progress is also strikingly similar, especially prior to
%    \cite{PotechinXu2026Theta}: in both settings, explicit witness
%    constructions were known that achieved the correct scaling but fell short
%    of the conjectured sharp constant. Moreover, identifying the right dual SDP
%    witness is itself a central challenge in both problems.
%
%    \item The best prior bounds in the two settings can be viewed as arising
%    from a common set of techniques: both can be done via analyzing correlated random
%    matrices whose entries are low-degree polynomials of the underlying random
%    input.
%
%    \item The threshold \(m\approx d^2/4\) also mirrors the constraint count in
%    the Theta problem. In the theta SDP, the desired witness matrix
%    is constrained on the non-edges; for \(G(n,1/2)\), these account for
%    roughly half of the \(\binom n2\) off-diagonal entries, namely
%    \(
%        \frac12\binom n2 \sim \frac{n^2}{4}.
%    \)
%    Thus, the \(d^2/4\) threshold in ellipsoid fitting is naturally analogous
%    to the number of entrywise non-edge constraints in the random-graph theta
%    problem.
%\end{enumerate}
%
%Having highlighted these parallels, it is natural to revisit ellipsoid fitting in light of the recent progress on the \Lovasz-Theta problem. The work of \cite{PotechinXu2026Theta} suggests an iterative construction scheme that ultimately exploits graph matrices, orthogonal polynomials, and free-independence-type behavior  to produce and analyze sharp SDP witnesses than those obtained by more direct constructions. Motivated by this perspective, we propose to adapt the overall strategy of \cite{PotechinXu2026Theta} as a natural next step toward the ellipsoid fitting problem.
\section{Technical Overview of Ellipsoid Fitting}
%We give an overview of our techniques and highlight the technical challenges along our way. In this section, we focus exclusively on our result for the positive side of the conjecture, i.e., fitting an ellipsoid below the threshold. We will then show how the similar techniques can be applied to obtain our result for refutation as well subsequently in~\cref{sec:refutation}. 
% 
% As a heads-up, essentially all technical ingredients in the analysis can be recycled from that for the primal program.

In this section, we provide an overview of our techniques and highlight the main technical challenges for the positive side of the conjecture.
%\paragraph{Roadmap of Our Work}
%
%We give an overview of our techniques in the subsequent section. At a high level, our notation and setup (especially regarding graph matrices) closely follow \cite{HKPX23}.  
%
%% Part of our motivation for revisiting the ellipsoid fitting problem comes from recent progress on the \Lovasz-Theta function. We discuss the connection between these two problems in detail in \cref{sec:connection-theta}. At the same time, our presentation is self-contained: the construction and analysis are developed directly in the ellipsoid fitting setting, so no prior familiarity with \cite{PotechinXu2026Theta} is needed.
%% 
%% Next, we give a brief overview for our candidate ellipsoid matrix construction. Unlike the prior works where we have an one-shot construction for the candidate matrix, we follow an iterative process to construct our candidate matrix, in the same fashion of \cite{PotechinXu2026Theta}. To fully instantiate our construction to enable our technical analysis, we give a self-contained overview for graph matrices specialized to our setting in \cref{sec:graph-matrix}. Finally, with the basics about graph matrix installed, we finish the full description of our iterative construction in \cref{sec:iterative-construction}.
%% 
%% In the subsequent sections, we highlight the key challenges in our PSD analysis. These are the main challenges unique to our setting of ellipsoid problem.
%
%
%One motivation for revisiting the ellipsoid fitting problem comes from recent progress on the \Lovasz-theta function. We discuss the connection between these two problems in detail in \cref{sec:connection-theta}. At the same time, our presentation is self-contained: the construction and analysis are developed directly in the ellipsoid fitting setting, so no prior familiarity with \cite{PotechinXu2026Theta} is needed.
%
%Next, we give a brief overview of our candidate ellipsoid matrix construction.
%Unlike prior works, which use a one-shot construction for the candidate matrix,
%our construction is iterative, following the approach of
%\cite{PotechinXu2026Theta}.  To fully specify the construction in a form
%suitable for our technical analysis, we first give a self-contained overview of
%graph matrices specialized to our setting in \cref{sec:graph-matrix}.  With
%this graph-matrix framework in place, we then complete the description of our
%iterative construction in \cref{sec:iterative-construction}.
%
%The remaining sections highlight the main technical challenges specific to our work for the ellipsoid fitting problem. The remainder of our work is organized as follows.
%\begin{enumerate}
%	\item 
%In~\cref{sec:MP-sec}, we formally establish the inverse of $M$ via polynomial expansion over Marchenko-Pastur distribution of the corresponding parameter $\gamma$.  This justifies the correction mechanism and, more
%    importantly, expresses the correction terms in the graph-matrix basis.  \item In~\cref{sec:inner-matrix-Q}, we establish the freely independence of the backbone-dangling shapes as well as the variance of the inner matrix $Q$ in the limit.
%\item Finally, in~\cref{sec:formal-truncation}, we incorporate the truncation
%    details needed to instantiate the construction
%    and complete the proof of the main theorem.
%\end{enumerate} 
%%
%
%
%\paragraph{Notation.}
%Throughout this work, all probabilistic statements hold with probability
%\(1-o_d(1)\).  We write \(\|M\|_{sp}\) for the spectral norm of a matrix \(M\),
%and reserve the notation \(\cm\) for graph matrices.  We use \(\tau\) as a
%generic placeholder for an arbitrary shape, so that \(\cm_\tau\) denotes the
%graph matrix associated with that shape.  The symbols \(\alpha\) and \(\beta\)
%are reserved for specific shapes that play distinguished roles in our analysis.
%
%
%\paragraph{AI Disclosure.}  We used generative AI tools to help polish the writing. We take full responsibility for any error. We also note that the connection to the prior work of \cite{kogan2025extremal} was brought to the attention of one of the authors only after we had initially completed the contained proof via trace moment method


%\paragraph{Roadmap for this Section}
%Part of our motivation for revisiting the ellipsoid fitting problem comes from recent progress on the \Lovasz-Theta function. We discuss the connection between these two problems in detail in \cref{sec:connection-theta}. At the same time, our presentation is self-contained: the construction and analysis are developed directly in the ellipsoid fitting setting, so the reader does not need prior familiarity with \cite{PotechinXu2026Theta}.

%\subsection{Recap of Prior Approaches: the identity-perturbation construction.} We now briefly recap the prior approaches based on identity-perturbation.
%A natural starting point is the observation that the identity matrix almost
%satisfies the interpolation constraints:
%\[
%    v_i^\top I_d v_i = \|v_i\|_2^2 \approx 1.
%\]
%The identity-perturbation construction therefore searches for a witness of the
%form
%\[
%    \Lambda
%    =
%    I_d - \sum_{i=1}^m w_i v_i v_i^\top.
%\]
%Let
%\[
%    \eta[i] := \|v_i\|_2^2-1,
%    \qquad
%    M[i,j] := \langle v_i,v_j\rangle^2.
%\]
%The constraints \(v_i^\top \Lambda v_i=1\) are equivalent to the linear system
%\(
%    M w = \eta.
%\)
%Thus, provided \(M\) is invertible, the identity-perturbation candidate is
%\[
%    \Lambda = I_d - R,
%    \qquad
%    R := \sum_{i=1}^m (M^{-1}\eta)_i v_i v_i^\top.
%\]
%The linear constraints are then automatically satisfied, and the entire problem
%reduces to proving
%\(
%    \|R\|_{\mathrm{sp}} < 1.
%\)
%This is the route taken in \cite{KD22,PTVW22,HKPX23}.
%
%The limitation of this one-shot construction is that the affine constraints
%are enforced exactly at the cost of creating a delicate global spectral norm
%problem.  The matrix \(R\) depends on the inverse of the highly correlated
%random matrix \(M\), and the positivity of \(\Lambda=I_d-R\) hinges on
%controlling the full operator norm of this correction.  This motivates a more
%flexible approach that separates the two tasks: enforcing positive
%semidefiniteness and correcting the affine constraints.
%
\begin{comment}
\subsection{Iterative Construction from Non-linear Activation}
In this work, we follow the recipe of designing SDP solution for obtaining sharp-in-the-constant bound for \Lovasz-Theta function of random graph. On a high level, the approach proceeds as follows,
\begin{enumerate}
	\item Fix a non-negative function $F(\cdot)$ to be chosen that will act as a spectral map;
	\item Construct an inner matrix $Q$ s.t. $F(Q)$ satisfies the affine constraints from the specific problem. For technical convenience, the specific spectral map to be applied is chosen as $\frac{1}{C_F} \cdot F(Q_i)$ for any $Q_i$.
\end{enumerate}

To be concrete, we consider the non-negative function $F$ as follows,
\begin{definition}[Nonnegative function \(F\) and its Chebyshev expansion] \label{def:F-function}
Let \(\mu_{\mathrm{sc}}\) be the semicircle distribution on \([-2,2]\), with
density
\[
    f_{\mathrm{sc}}(x)=\frac{\sqrt{4-x^2}}{2\pi}.
\]
We define
\[
    F(x):=2x_+
    =
    \begin{cases}
        2x, & x\in[0,2],\\
        0,  & x\in[-2,0].
    \end{cases}
\]
Let
\[
    P_j(x):=U_j\left(\frac{x}{2}\right),
\]
where \(U_j\) is the \(j\)-th Chebyshev polynomial of the second kind.  Then
\(\{P_j\}_{j\ge0}\) is an orthonormal basis for
\(L^2([-2,2],\mu_{\mathrm{sc}})\), with \(P_0=1\) and \(P_1=x\).

We write the Chebyshev expansion of \(F\) as
\[
    F(x)=C_F+x+\sum_{j\ge2}b_j \cdot P_j(x),
\]
where
\[
    C_F=b_0=\mathbb E_{X\sim\mu_{\mathrm{sc}}}[F(X)]
    =\frac{8}{3\pi},
    \qquad
    b_1=1,
\]
and, for \(j\ge2\),
\[
    b_j\coloneqq \E_{x\sim \mu_{\mathrm{sc}}}[F(x)]
\]
\end{definition}

With the spectral map described, much of the difficulty lies upon finding an inner matrix $Q$ such that $F(Q)$ gives a feasible solution. Towards this end, we construct the inner matrix $Q$ via an iterative process that can be summarized as follows.
\begin{enumerate}
	\item At any iteration $i$, there exists some \emph{correction} term $\Delta_i$ s.t. $\frac{1}{C_F}\cdot  F(Q_i) + \Delta_i $ satisfies the affine constraints;
	\item Then we update $Q_{i+1} \leftarrow Q_i +  \Delta_i $. 
\end{enumerate}
To unpack this process, our approach replaces this one-shot correction in identity-perturbation by an iterative correction
scheme.  The goal is to separate the two tasks which are coupled in the
identity-perturbation argument: positive semidefiniteness and satisfaction of
the affine constraints.  At each level \(i\), we maintain an inner matrix
\(Q_i\).  From \(Q_i\), we first produce a positive semidefinite candidate by
applying a nonnegative spectral function.  Then we measure the constraint
error and find some corresponding correction term to satisfy the linear constraints, albeit at the cost of damaging the fragile positivity that we have started with.

So far, the process as described is verbatim the exact approach employed for theta function. Indeed, the problem-specific component only arises from how to correct for the linear constraints. And this is where the original construction of identity-perturbation kicks in, as it gives us a systematic approach for correcting the deviation of affine constraints, except that we no longer correct with respect to the identity matrix!

 In particular, consider the Gram matrix $M$ defined as earlier, for any iteration $i$, there is some vector $w_i \in \R^m$ we can find s.t. \[
 \frac{1}{C_F} F( Q_i)  + \underbrace{\sum_{t\in [m]} w_i[t] \cdot  v_t  \cdot v_t^T}_{\coloneqq  \Delta_i}
 \]
 satisfies the affine constraints. Concretely, let $\eta_i \in \R^{m}$ denote the deviation of $\frac{1}{C_F} F(C_F \cdot Q_i)$ at each constraint, the coefficient vector $w$ is then taken as \[ 
 w_i = M^{-1} \cdot \eta_i  
 \]
 assuming the matrix $M$ is invertible. 
 

 

\paragraph{First Attempt with the Base Iteration} At this point, we execute the above strategy for the base matrix $Q_0$ to motivate our iterative procedure. Let us start by describing our initialization and the first iteration. For starters, consider the following initialization \[ 
Q_0 \coloneqq {C_F}\cdot \sum_{t\in [m]} w_{0}[t] \cdot v_t\cdot v_t^T\,.
\]
Next, consider $\frac{1}{C_F} F(Q_0)$, and its polynomial expansion via the orthogonal polynomials for $\mu_{SC}$ as defined above, we have \[ 
\frac{1}{C_F }\cdot  F(Q_0) = I + \frac{1}{C_F} \cdot  Q_0 + \underbrace{\sum_{j\ge 2} b_j \cdot P_j(Q_0)}_{\text{deviation for affine constraints} }
\]
where we use $P_1(Q_0) = Q_0$. Plugging in the definition of $Q_0$, in particular, \(
\sum_{t\in[m]} w_0[t] \cdot v_t \cdot v_t^\top 
\)
is the desired perturbation to ensure the identity matrix satisfies the affine constraints, we observe that the deviation term for the affine constraints is only incurred by $ \sum_{j\ge 2} b_j \cdot P_j(Q_0)
$. 

Therefore, to correct for the troublesome terms, we will consider the updated deviation vector \[
\eta_1[t] =  v_t^\top (  \sum_{j\ge 2} b_j \cdot P_j(Q_0) ) v_t
 \] 
 for any $t\in [m]$, and subsequently solve for the corresponding inverse of the linear system. This culminates at the following initialization and first level of iteration.

 \begin{mdframed}[linewidth=0.4pt]
\textbf{Summary of Initialization and the First Iteration:}

\begin{enumerate}
\item \textbf{Iteration level $i=0$.}  
Define
\[
Q_0 \coloneqq {C_F}\cdot \sum_{t\in [m]} w_{0}[t] \cdot v_t\cdot v_t^\top 
\]
Let $\calB(Q_0)$ be the collection of shapes in $Q_0$, in graph matrix representation, we have \[
Q_0 = \sum_{\tau \in \calB(Q_0)} c(\tau) \cdot \cm_\tau\,. 
 \]
 
%The base shape collection for $Q_0$ is the collection of graph matrices obtained by 

\item \textbf{Iteration level $i=1$.}  
When we take $F(Q_0)$, we have the terms $\sum_{j \geq 2}{ b_j {C_F^j} \cdot P_j(Q_0) }$ which need to be corrected for the affine constraints. To correct these terms, we measure the deviation in the affine constraints due to these extra terms by $\eta_1 \in \R^m$ with \[ 
\eta_1[i] = v_i^T ( \sum_{j \geq 2}{ b_j {C_F^j} \cdot P_j(Q_0) } )v_i,  \qquad \forall t\in [m]
\]
and define $w_1 \in \R^m$ s.t. 
\[
M w_1 = -\eta_1 
 \]
 and update our inner matrix as \[
Q_1 \coloneqq Q_0 + \sum_{t\in [m]} w_1[t] \cdot  v_t v_t^T.
\]
\end{enumerate}

\end{mdframed}

 \paragraph{Graph Matrix POV} To further shed light on our iterative process, it is useful to describe a core lemma of our analysis, as that immediately gives a intuitive description to the evolution of the summand in our inner matrix $Q$.
 
 \begin{proposition}[Informal]
 	For an inner matrix $Q$ that is a linear combination with shapes from a ground-set $\calB(Q)$,
 	\[ 
 	P_j(Q_i) = \sum_{\tau: \text{j-fold concatenation of } \calB(Q) } \cm_\tau \,.
 	\]
 \end{proposition}
This gives us a handy description to study $P_j(Q)$ via clean graph matrices. In particular, it inspires for the following identification of deviation terms for the subsequent iterations $i>1$. 
 For each $j \geq 2$, consider the set of $j$-way concatenations of base shapes in $Q_i$:
\[
\calP_j(Q_i) :=
\{\tau_1 \circ \dots \circ \tau_j : \tau_t \in \calB(Q_i) \forall i \in [j]\}.
\]
Partition $\calP_j(Q_i)$ depending on whether the concatenation uses a newly introduced base shape:
\[
\textsf{Old-}\calP_j(Q_i)
=
\{\beta \in \calP_j(Q_i) :
\tau_t \in \calB(Q_{i-1}) \text{ for all } t \in [j]\},
\]
\[
\textsf{New-}\calP_j(Q_i)
=
\{\beta \in \calP_j(Q_i) :
\exists t \text{ such that }
\tau_t \in \calB(Q_i)\setminus \calB(Q_{i-1})\}.
\]
We observe the following invariant,
\begin{proposition}[Iteration Invariant] \label{prop:iterative-invar}  For any $i\geq 1$,we have
 \[ 
\frac{1}{C_F} F(Q_i) =  I + Q_i +  \sum_{j\ge 2} \left(\sum_{\tau \in  \textsf{Old-}\calP_j(Q_i)} c_\tau \cdot \cm_\tau +\sum_{\tau \in  \textsf{New-}\calP_j(Q_i)} c_\tau \cdot \cm_\tau \right)
\]
and that \[ 
I + Q_i +  \sum_{j\ge 2} \sum_{\tau \in  \textsf{Old-}\calP_j(Q_i)} c_\tau \cdot \cm_\tau 
\]
satisfies the affine constraints.	
\end{proposition}
\begin{remark}
	This applies for the base iteration of $i=0$ if we treat ``Old'' as empty-set.
\end{remark}
We defer its proof to~\cref{sec:defer-construction}. This leaves the final ``new'' as the sole source for deviations, and that is the correction target for the subsequent iteration.

%Next, for the recursive step, consider the following.
\paragraph{Summary of the Iterative Construction}

Our general recursive step is as follows:
\begin{mdframed}[linewidth=0.4pt]
\textbf{Recursive Update ($i \to i+1$):}

\begin{enumerate}
\item For each $j \geq 2$, consider the set of $j$-way concatenations of base shapes in $Q_i$:
\[
\calP_j(Q_i) :=
\{\tau_1 \circ \dots \circ \tau_j : \tau_t \in \calB(Q_i) \forall i \in [j]\}.
\]
\item Partition $\calP_j(Q_i)$ depending on whether the product uses a newly introduced base shape:
\[
\textsf{Old-}\calP_j(Q_i)
=
\{\beta \in \calP_j(Q_i) :
\tau_t \in \calB(Q_{i-1}) \text{ for all } t \in [j]\},
\]
\[
\textsf{New-}\calP_j(Q_i)
=
\{\beta \in \calP_j(Q_i) :
\exists t \text{ such that }
\tau_t \in \calB(Q_i)\setminus \calB(Q_{i-1})\}.
\]
\item For each $\beta=\tau_1\circ\dots\circ\tau_j \in \textsf{New-}\calP_j(Q_i)$, note that $\cm_{\beta}$ appears in $F(Q_i)$ with coefficient
\[
c_\beta  \coloneqq b_j \cdot  \prod_{t=1}^{j} c({\tau_t}).
\]
and $Q_i$ does not contain a correction for this term.
%which is well-defined since each $\tau_t$ already carries a coefficient from previous levels.

\item To correct for these terms, %Let $\tau' \coloneqq \corr(\tau)$ denote the backbone-corrected shape (see \cref{def:backbone-correction}), and set
we define the deviation error for the new level $\eta_{i+1} \in \R^m $ as \[ 
\eta_{i+1}[t]  = v_t^T (\sum_{j\geq 2} \sum_{\tau \in \textsf{New-}\calP_j(Q) }  c(\tau) \cdot \cm_\tau   ) v_t \,, \qquad \forall i\in [m]\,.
\]
and solve $w_{i+1}$ such that \[ 
Mw_{i+1} = -\eta _{i+1}\,.
\]
Equivalently, this can be viewed as for any fixed shape $\tau \in \textsf{New-}\calP_j(Q_i)$, we define a vector $w_\tau$ to correct for the error term incurred by a new term $\tau$, and we have \[ 
w_{i+1} = \sum_{j\geq 2} \sum_{\tau\in\textsf{New-}\calP_j(Q_i)} w_{\tau} \,.
\]


Finally, we update the inner matrix as 
%\jnote{double check the $-1$ here}
%\item Define the update
\[
Q_{i+1}-Q_i
\coloneqq
\sum_{t\in [m]} w_{i+1}[t] \cdot  v_t v_t^T = \sum_{t\in [m]} \sum_{j\geq 2} \sum_{\tau \in\textsf{New-}\calP_j(Q_i)} w_\tau[t] \cdot v_tv_t^T \,.
\]
Equivalently, the base shape collection evolves as
\[
\calB(Q_{i+1})
\coloneqq
\calB(Q_i)
\cup
\{ \text{Correction Shape for } \eta_\tau : \tau \in \textsf{New-}\calP_j(Q_i) \forall j\geq 2 \}
\]
%where for each $\beta \in \textsf{New-}\calP_j(Q_i)$, $c({\bar{\beta}}) = b_j\prod_{t=1}^{j} c(\tau_t)$.
%and $Q_{i+1}$ is the linear combination of shapes in $\calB(Q_{i+1})$ with coefficients $c(\cdot)$, i.e.,
%\[ 
%Q_{i+1}  = \sum_{\tau \in \calB(Q_{i+1})} c_\tau
%\cdot \cm_{\tau}\,.
%\]
\end{enumerate}

\end{mdframed}
\begin{remark}
    In our formal analysis, we will consider an inner-truncation such that for any $Q_i$ during the iterative process, we only consider $P_j(Q_i)$ for any $j\leq D$.
\end{remark}

For our final $Q$, we will take $Q = Q_{t}$ for some truncation parameter $t$ and then handle the remaining small correction terms directly. This is addressed in~\cref{sec:real-world-proof}. That said, for the purposes of this overview, it suffices for us to consider $Q$ be the limit of $Q_i$ as $i \to \infty$.
\jnote{to fill in}
\end{comment}

%\subsection{Parallel with the \Lovasz-Theta Function on \texorpdfstring{$G(n,\frac{1}{2})$}{G(n,1/2)}} 
\label{sec:connection-theta} 
% \jnote{drop this? its a bit too long}\anote{I'd lean towards having a paragraph highlighting the connection to [PX26], though I think it could use some editing. I took a stab at a potential replacement paragraph.}
%Potential replacement paragrah: Our starting point is the recent breakthrough of Potechin and Xu \cite{PotechinXu2026Theta} which showed that with high probability, the \Lovasz-Theta function of a random $G(n,1/2)$ graph has value $(1 \pm{o_n(1)})\sqrt{n}$. This paper gave an intricate iterative construction of a solution to the SDP for the \Lovasz-Theta function and analyzed it by exploiting connections between graph matrices, orthogonal polynomials, and free independence. This was promising as it gave a novel solution for this SDP which was explicit and optimal. Moreover, the construction of \cite{PotechinXu2026Theta} matched a distinctive feature of the statistical physics predictions of https://arxiv.org/abs/2310.01169 [MK24] for ellipsoid fitting. In particular, https://arxiv.org/abs/2310.01169 [MK24] predicts that for ellipsoid fitting, nuclear norm minimization gives a solution where half of the eigenvalues are zero. For the construction of \cite{PotechinXu2026Theta}, it also turns out that half of the eigenvalues are zero! 
%jnote{im not sure about leaving this in overview or intro.}
%
%Our starting point is the recent progress for settling the tight threshold for the \Lovasz-Theta function in \cite{PotechinXu2026Theta}. In retrospect, these two problems share a striking similarity in both the input distributions from a universality point of view and in the state of progress. \cite{PotechinXu2026Theta} suggests an iterative construction scheme that exploits graph matrices, orthogonal polynomials, and free-independence-type behavior to produce and analyze SDP witnesses at sharp thresholds, yielding phase-transition bounds beyond those achieved by more direct constructions. 

Our starting point is the recent breakthrough of Potechin and Xu \cite{PotechinXu2026Theta} which showed that with high probability, the \Lovasz-theta function of a random $G(n,1/2)$ graph has value $(1 \pm{o_n(1)})\sqrt{n}$. This paper gives an intricate iterative construction of a solution to the SDP for the \Lovasz-theta function and analyzes it by exploiting connections between graph matrices, orthogonal polynomials, and free independence. This is promising as it gives a novel solution for this SDP which is explicit and optimal. Moreover, the construction of \cite{PotechinXu2026Theta} matched a distinctive feature of the statistical physics predictions of \cite{MD24} for ellipsoid fitting. In particular, \cite{MD24} predicts that for ellipsoid fitting, nuclear norm minimization gives a solution where half of the eigenvalues are zero. For the construction of \cite{PotechinXu2026Theta}, it also turns out that half of the eigenvalues are zero!

%Motivated by this intriguing connection, we employ the overall strategy developed therein as a natural starting point for a new approach to the ellipsoid fitting problem.
% Motivated by this perspective, we take the overall strategy developed there as a natural starting point for a new approach to the ellipsoid fitting problem.
%
%The current state of the ellipsoid fitting problem prior to this work closely parallels that of the
%\Lovasz-Theta function on random graphs \(G(n,1/2)\). The analogy goes
%beyond the fact that both problems concern basic SDPs on natural random inputs.
%Several structural and technical parallels are worth highlighting:
%\begin{enumerate}
%    \item The underlying input distributions---i.i.d. Gaussian samples in the
%    ellipsoid fitting problem and the random graph \(G(n,1/2)\) in the
%    \Lovasz-Theta problem often behave
%    similarly from a universality perspective in random matrix theory.
%
%    \item The state of progress is also strikingly similar: in both settings, explicit witness constructions
%were known that achieved the correct scaling but fell short of the conjectured
%sharp constant. Moreover, identifying the right dual SDP witness is itself a
%central challenge in both problems. In particular, although recent advances
%have led to increasingly sharp analyses of correlated random matrices, there remained a fundamental obstacle: without a plausible candidate
%construction near the conjectured threshold, there was little to analyze in the
%first place. 
%
%The recent work~\cite{PotechinXu2026Theta} overcame this obstacle
%for the theta function by introducing a new iterative construction, explicitly suggesting
%that a similar paradigm may also be fruitful for ellipsoid fitting.

%    \item The best prior bounds in the two settings can be viewed as arising
%    from a common set of techniques: both can be done via analyzing correlated random
%    matrices whose entries are low-degree polynomials of the underlying random
%    input.
%
%    \item The threshold \(m\approx d^2/4\) also mirrors the constraint count in
%    the Theta problem. In the theta SDP, the desired witness matrix
%    is constrained on the non-edges; for \(G(n,1/2)\), these account for
%    roughly half of the \(\binom n2\) off-diagonal entries, namely
%    \(
%        \frac12\binom n2 \sim \frac{n^2}{4}.
%    \)
%    Thus, the \(d^2/4\) threshold in ellipsoid fitting is naturally analogous
%    to the number of entrywise non-edge constraints in the theta function problem.\end{enumerate}

%Having highlighted these parallels, it is natural to revisit ellipsoid fitting in light of the recent progress on the \Lovasz-Theta problem. The work of \cite{PotechinXu2026Theta} suggests an iterative construction scheme that ultimately exploits graph matrices, orthogonal polynomials, and free-independence-type behavior  to produce and analyze sharp SDP witnesses, yielding sharp phase transition bounds than those obtained by direct constructions previously. It is even explicitly mentioned therein that a similar paradigm may also be fruitful for ellipsoid fitting, while admittedly there are obvious technical challenges ahead.  Motivated by this perspective, we revisit the overall strategy developed there as a natural starting point for a new attack on the ellipsoid fitting problem.
  
%  Having highlighted these parallels, it is natural to revisit ellipsoid fitting in light of the recent progress on the \Lovasz-Theta problem.
%   \cite{PotechinXu2026Theta} suggests an iterative construction scheme that exploits graph matrices, orthogonal polynomials, and free-independence-type behavior to produce and analyze SDP witnesses at sharp thresholds, yielding phase-transition bounds beyond those achieved by more direct constructions. Motivated by this perspective, we employ the overall strategy developed therein as a natural starting point for a new approach to the ellipsoid fitting problem.
  
\subsection{Overview of the Iterative Construction }
At a high level, there are two components in the construction:
\begin{enumerate}
    \item choose a nonnegative function \(F\), which will be applied spectrally;
    % \item \snote{choose a basis-invariant function \(F:\textup{Sym}_d\to\textup{Sym}_d\) which is nonnegative (i.e., all elements in its image are PSD);}
    \item construct an inner matrix \(Q\) such that the spectrally transformed
    matrix satisfies the affine constraints of the SDP. Moreover, the inner matrix $Q$ is constructed by an iterative process.
\end{enumerate}

In this section we describe the iterative construction of \(Q\). We start by introducing the non-negative function $F$ and its polynomial expansion. This is exactly the same choice as that for the prior work of \cite{PotechinXu2026Theta} for the \Lovasz-theta function.
\begin{definition}[Non-negative Function \(F\) and its Chebyshev Expansion]
\label{def:F-function}
Let \(\mu_{\mathrm{sc}}\) be the semicircle distribution on \([-2,2]\), with density
\(
    f_{\mathrm{sc}}(x)
    =
    \frac{\sqrt{4-x^2}}{2\pi}.
\)
Define
%\[
%    F(x):=2x_+
%    =
%    \begin{cases}
%        2x, & x\in[0,2],\\
%        0,  & x\in[-2,0].
%    \end{cases}
%\]
\[F(x)=\textup{max}(2x,0)=\begin{cases}
        2x, & x>0,\\
        0,  & x\leq 0.
    \end{cases}\]
Let
\(
    P_j(x):=U_j\left(\frac{x}{2}\right),
\)
where \(U_j\) is the \(j\)-th Chebyshev polynomial of the second kind. Then
\(\{P_j\}_{j\ge0}\) is an orthonormal basis for
\(L^2([-2,2],\mu_{\mathrm{sc}})\), with \(P_0=1\) and \(P_1=x\).

We write
\[
    F(x)
    =
    C_F+x+\sum_{j\ge2} b_j P_j(x),
\]
where
\(
    C_F
    =
    \mathbb E_{x\sim\mu_{\mathrm{sc}}}[F(x)]
    =
    \frac{8}{3\pi},
    \mathbb E_{x\sim\mu_{\mathrm{sc}}}[F(x)P_1(x)]=1,
\)
and, for \(j\ge2\),
\(
    b_j\coloneqq
    \E_{x\sim\mu_{\mathrm{sc}}}[F(x)P_j(x)].
\)
\end{definition}

  We use the following normalization to ensure the identity term has coefficient $1$. Once the $F$ function is fixed, the main difficulty is to find an inner matrix \(Q\) for which the spectral
matrix \(\frac1{C_F}F(  Q)\) satisfies the affine constraints. 
%
% \anote{This is a bit confusing. If I remember correctly from last week's meetings, it should be that if we take $F_{adj}(Q) = \frac{1}{C_F}F({C_F}Q)$ then 
%\[
%    F_{adj}(Q) = \frac{1}{C_F}F({C_F}Q)
%    =
%    I+Q+\sum_{j\ge2} \frac{b_j}{C_F} P_j({C_F}Q)
%\]
%and $F_{adj}(\frac{Q}{C_F}) = \frac{1}{C_F}F(Q)$.} \jnote{indeed there was some inconsistency earlier...fixed now?}\anote{This looks better.}
\[
    \frac{1}{C_F}F( Q)
    =
    I+\frac{1}{C_F}\cdot Q +\sum_{j\ge2}   \frac{b_j}{C_F}\cdot  P_j(  Q).
\]
%This convention keeps the identity term is normalized in the
%form needed for the affine constraints as we consider the identity perturbation construction as a starting point.


 Towards that end, we introduce an iterative process to construct \(Q\), following the same recipe %in 
 as \cite{PotechinXu2026Theta}:  
\begin{enumerate}
	\item At level \(i\), we maintain an inner matrix \(Q_i\);
	\item Apply the non-negative spectral map to obtain a PSD matrix $X_i \leftarrow \frac{1}{C_F} F(Q_i) $;
	\item $X_i$ may not satisfy the affine constraints exactly, and we update the inner matrix to $Q_{i+1}$ to correct for the deviation. \end{enumerate} 
%   Applying the
%nonnegative spectral map then gives a positive semidefinite candidate
%\(
%    \frac1{C_F}F( Q_i),
%\)
% while it may not satisfy the affine
%constraints.  That said, we then identify the constraint error and add a correction to the
%%inner matrix. This gives to the iterative task where we alternately apply activation to an inner matrix to obtain a PSD matrix, and correct for the deviation at the affine constraints by the resulting matrix.
%We now elaborate on the concrete update term. Concretely, we consider the following primitive map that is also heavily used in the prior works for ellipsoid fitting in order,
%\[
%    L(X)_t := v_t^\top Xv_t,
%    \qquad
%    L^*(w):=\sum_{t\in[m]}w[t]\,v_t v_t^\top,
%\]
%and
%\[
%    M:=LL^*,
%    \qquad
%    M[s,t]=\langle v_s,v_t\rangle^2.
%\]
%Given a current candidate \(X\), any affine-constraint error can be corrected by
%a matrix of the form \(L^*(w)\) assuming $M$ is invertible.  Indeed, picking
%\(
%    \eta[t]=v_t^\top Xv_t-1,
%\)
%then choosing \(w\) so that
%\(
%    Mw=-\eta
%\)
%ensures
%\[
%    L(X+L^*(w))-\mathbf 1
%    =
%    \eta+Mw
%    =
%    0.
%\]
%Thus, assuming \(M\) is invertible, the linear correction is
%\(
%    w=-M^{-1}\eta.
%\) This is precisely the main idea behind the identity perturbation construction from the prior works.
%\jnote{i think the current version is ok. There's no need to elaborate upon this.}
%{\color{orange}Sofia:
%\paragraph{The Identity Perturbation Ansatz}
%
%Previous papers used an ansatz of the form 
%\[I_d+\sum_{i=1}^m c_iv_iv_i^\top.\]
%For $m\leq \binom{d+1}{1}$, the rank-1 matrices $(v_iv_i^\top)_{i\in [m]}$ are linearly independent so their Gram matrix $M=(\left<v_i,v_j\right>^2)_{i,j\in [m]}$ is invertible. This implies that there's exactly one choice of constants $(c_i)_{i\in [m]}$ that allow this matrix to satisfy the linear constraints. We can conceptualize this using the orthogonal projection to the affine subspace of $\textup{Sym}_d$ of matrices satisfying the linear constraints of ellipsoid fitting. This projection can be written as
%\[\Lambda\mapsto X_{LS}+\pi(\Lambda)\]
%where $\pi=I-\calL^* M^{-1}\calL$ and $X_{LS}$ is the least-squares solution to the linear constraints. Indeed, it is clear that $\pi^2=\pi$ and it fixes $\textup{im}(\calL^*)$ and kills $\textup{im}(\calL^*)^\perp=\ker(\calL)$.
%
%%For $\alpha<1/2$ there is exactly one choice of constants $(c_i)_{i\in [m]}$ that allow this matrix satisfy the linear constraints, a fact which is a consequence of the invertibility of the Gram matrix $M=(\left<v_i,v_j\right>^2)_{i,j\in [m]}$ and the linearly independence of the rank-1 matrices $(v_iv_i^\top)_{i\in [m]}$. We can conceptualize this as a particular case of the following construction: let $\textup{LC}=\bigcap_{i\in [m]}\{M:\left<Mv_i, v_i\right>=1\}$ be the affine subspace of matrices satisfying the linear constraints. Let $\pi$ be the orthogonal projection to $\textup{LC}$. The identity perturbation ansatz is then $\pi(I_d)$ and it is a solution to the ellipsoid fitting problem whenever it is PSD, a sufficient condition to which is $\lVert \pi(I_d)-I_d\rVert_{sp}<1$. The papers \cite{KD22,PTVW22,HKPX23} prove precisely this bound.
%
%We can write
%\[\pi(\Lambda)=\Lambda-\sum_{i=1}^m (M^{-1}\calL(\Lambda))[i]v_iv_i^\top;\quad X_{LS}+\pi(\Lambda)=\Lambda-\sum_{i=1}^m (M^{-1}\calL(\Lambda)-\1)[i]v_iv_i^\top.\]
%We call the term $-\sum_{i=1}^m (M^{-1}\calL(\Lambda)-\1)[i]v_iv_i^\top$ a correction such that when it is added to $\Lambda$ the sum satisfies the linear constraints. By choosing $\Lambda$ positive it becomes sufficient to prove a spectral bound on this correction, so that the resulting matrix is both positive and satisfies linear contraints. To that end, it becomes necessary to understand the spectrum of $M^{-1}$.
%
%The previous literature \cite{KD22,PTVW22,HKPX23} used precisely this bound to prove that the choice of $\Lambda=I_d$ yields a solution when $m$ is small enough.
%%We thus seek a PSD matrix $\Lambda$ such that $w(\Lambda)$ is small enough so that $\pi(\Lambda)$ remains PSD. To that end, it becomes necessary to understand the spectrum of $M^{-1}$.
%}

\paragraph{Identity Perturbation as Initialization}
We now briefly recap the prior approaches based on identity-perturbation as that in turn gives us an initialization for the inner matrix $Q_0$.
The identity-perturbation construction searches for a witness of the
form
\(
    \Lambda
    =
    I_d - \calR  \) with\[
    \eta \coloneqq \calL(I)-\1_m ,
    \qquad
    \calR \coloneqq \calL M^{-1} \eta\,.
%    M[i,j] := \langle v_i,v_j\rangle^2, \qquad R\coloneqq \sum_{i=1}^m w[i] v_i v_i^\top
    \]
%The constraints \(v_i^\top \Lambda v_i=1\) are equivalent to setting the linear system
%\(
%    M w_0 = \eta_0.
%\)
%Thus, provided \(M\) is invertible, the identity-perturbation candidate is
%\[
%    \Lambda = I_d + R,
%    \qquad
%    R := \sum_{i=1}^m (M^{-1}\eta)_i v_i v_i^\top.
%\]
The linear constraints are then automatically satisfied, and the entire problem
reduces to proving
\(
    \|\calR\|_{\mathrm{sp}} < 1.
\)
This is precisely the route taken in \cite{KD22,PTVW22,HKPX23}. The matrix \(\calR\) depends on the inverse of the highly correlated
random matrix \(M\), and the positivity of \(\Lambda=I_d-\calR\) hinges on
controlling the operator norm of this correction. 

However, it is known that $\calR$ is not of small norm close to the threshold, and this is where the non-negative spectral map $F$ comes to our rescue as it produces a positive albeit large norm matrix as a result. The remaining challenge is then to correct its affine constraint violations %hopefully 
without destroying positivity.

\jnote{im debating on whether to elaborate on the specific correction term here or simply keep it abstract. i'm also ok with the current version in which we fully specify the correction primitive, and mention we are gonna tell u why later.}
\paragraph{Iteration in Action}
Motivated by the prior approaches via identity perturbation, our starting point is %We iniialize starting point as
\(
    Q_0
    \coloneqq - C_F \cdot \calR\,,
\)
%be the perturbation component from the identity-perturbation construction 
which is precisely the perturbation component from the identity-perturbation construction with the scaling chosen so that
\(
    I+   Q_0/C_F \)
satisfies the affine constraints. The particular choice of scaling comes from the Chebyshev expansion of $F$: 
\[
    \frac1{C_F}F(  Q_0)  = I+ \frac{1}{C_F} Q_0 + \sum_{j>1} \frac{b_j}{C_F} \cdot P_j(Q_0) \,.
   \]
Next, observe that the first two terms in the equation already satisfy the affine constraints, and the only affine deviation
comes from the unscaled higher-order terms
\(
 H_{0} \coloneqq \sum_{j\ge2} b_j \cdot   P_j( Q_0)\) crucially without the $1/C_F$ scaing. 
Define the (unscaled) first residual vector by \[ \eta_1= \calL(\sum_{j\ge2}  b_j \cdot   P_j(  Q_0)  ) = \calL(H_{0}  ) \in \R^m\]
%, in other words
%\[
%    \eta_1[t]
%    =
%    v_t^\top
%    \left(
%\sum_{j\ge2}  b_j \cdot   P_j(  Q_0)     \right)  
% v_t ,
%    \qquad t\in[m].
%\]
Therefore, the higher-degree component $ \sum_{j>1} \frac{b_j}{C_F} \cdot P_j(Q) 
$ in the resulting function of $ \frac{1}{C_F} F(Q_0) $ would incur a deviation of $\frac{\eta_1}{C_F}$ when measured with respect to the affine constraints in this notation. 

\paragraph{Picking a Freely Independent Correction Mechanism}
Analogous to the identity perturbation construction, one straightforward choice 
%is to solve for
%\[
%    Mw_1=-\eta_1
%\]
%and update
%\[
%    Q_1
%    =
%    Q_0+
%    \sum_{t\in[m]}w_1[t]\,v_t v_t^\top.
%\]
%Equivalently, 
%for any shape $\tau$ from the higher-order term, let $\eta_\tau = L(-\cm_\tau) $ denote deviation in affine constraints it incurs. 
is to pick the correction term as \(
\mathcal{L}^{*}(M^{-1}(-\eta_1)) 
\)
as this would then allow us to group the correction term with the previously troublesome term as \[ 
H_{0} +  \mathcal{L}^*(M^{-1}(-\eta_1) )
\]
to correct for the prior deviation term incurred by $H_0$ since the affine constraints evaluated with these two terms are simply \[ 
\calL\left(H_{0} + \mathcal{L}^*(M^{-1}(-\eta_1) )\right) = \eta_{1} + \calL^*\calL(M^{-1}(-\eta_1)) = \eta_1-\eta_1 = 0\,
\] 
since $\calL^*\calL = M$. \jnote{add a remark saying how this can be viewed as a projection to the image of rank-$1$ forms as \[ 
\calL^*M^{-1} \calL 
\]
is a projection matrix 
}
 However, as we show in~\cref{sec:cooking-correction}, surprisingly this is not the ``right'' correction mechanism to be employed. For our machinery to apply, it is crucial that we can view each summand in the inner matrix as a ``free'' matrix with semicircular spectrum. The naive choice above does not suffice for this purpose: it is indeed true that it contains various matrices with semicircular spectrum; however, it also contains an extra copy of $\gam \cdot H_{0} $ once we examine its polynomial expansion.
 
 To salvage this correction machinery, we consider an alternative correction mechanism given by 
 \[ 
 \Delta_0  = \corr(H_0) \coloneqq \frac{1}{1-\gamma}\left(\calL^*M^{-1} \calL (-H_0) + \gam \cdot  H_{0}  \right)
 \]
It is straightforward to verify that $\Delta_1$ achieves the same task of removing prior deviations when combined with $ H_{0}$:
\[
\calL(\Delta_0) = \frac{1}{1-\gamma}  \calL\calL^*M^{-1}(-\calL(H_0) )+\frac{\gam }{1-\gamma} \calL( H_{0})\, = -\frac{1}{1-\gamma}\eta_1  + \frac{\gamma}{1-\gamma} \eta_1 =   -\eta_1\,.
 \]

Notice that the scaling is picked so %such 
that when we apply $\frac{1}{C_F} F(Q_1)$ to $Q_1 = Q_0 + H_0$, the degree-$1$ term from the update gives exactly \[
\frac{1}{C_F} (Q_1-Q_0) = \frac{\Delta_0}{C_F} =  \frac{\corr(H_0) }{C_F} \,, \] the desired correction term for $H_0$ in the polynomial expansion for $F(Q_0)/C_F$.
%\anote{It's probably better to write $(Q_1 - Q_0)$ rather than $P_1(Q_1 - Q_0)$.} \jnote{okii}
%\[ 
%\frac{1}{C_F} (Q_1-Q_0) = \frac{1}{C_F} \cdot \sum_{t\in [m]} w_1[t]v_tv_t^\top \,. 
%\]
%chosen to as the desired correction for the previous deviation of $\frac{1}{C_F} \cdot H_{Q_0}$. Let's now unpack this: for any constraint $i\in [m]$, grouping the prior deviation terms with the degree-1 term of the update gives us  \[ 
%v_i^\top \left(\frac{1}{C_F} \cdot \Delta_1  \right)v_i  + v_i^\top \left(\frac{H_{Q_0}}{C_F} \right) v_i = -\frac{\eta_1}{C_F} + \frac{\eta_1}{C_F} = 0\,,
%\]
%hence the update corrects for the affine deviations of the prior iteration, albeit at a cost of potentially creating new deviations since we need to further address \[ 
%H_{1} \coloneqq \sum_{j\ge 2} b_j \cdot \left( P_j(Q_1) - P_j(Q_0) \right)\,,
%\]
%We would like to cautious the reader that this shall be contrasted from the other natural object of \(
% \sum_{j\ge 2} b_j \cdot  P_j(Q_1)\). The particular choice is chosen because $P_j(Q_0)$ from the prior level has been grouped with the newly added correction terms in $P_1(Q_1-Q_0)$. 

We emphasize that in general, especially for $i\ge 1$, \(H_i\) should not be confused with the seemingly more natural quantity %(concretely consider $i=1$)  
%\anote{I think $Q_1$ and $Q_0$ should be replaced by $Q_i$ and $Q_{i-1}$}
\[ 
\sum_{j\ge 2} b_j P_j(Q_i)
\]
The difference is important: the contribution \(P_j(Q_i)\) from the previous level has already been absorbed into the newly added linear correction term \(P_1(Q_i-Q_{i-1})\).  Concretely, in the case of $i=1$, \(H_1\) is therefore defined to include only the genuinely new higher-order deviation created by the update from \(Q_0\) to \(Q_1\).


At this point, one may hope to show that the iterative scheme works by showing
that the newly created deviations become progressively smaller across subsequent
iterations. Before instantiating the precise correction term to be chosen and turning to the analysis, we first give a full summary of the iterative process. We caution, however, that even though the following summary completely specifies the construction, the resulting object is arguably quite opaque. 
To make the analysis tractable and transparent, we will later introduce the graph-matrix viewpoint. 
% Indeed, when phrased in this form, it is far from clear even to the authors how one could analyze the iteratively constructed object directly.
%This viewpoint organizes the terms that appear in the construction and gives a clean combinatorial description of the matrices arising from \(P_j(Q_i)\) in full generality. With this language in place, we will then return to the full iterative procedure and describe the later iterations in detail in~\cref{sec:iterative-construction}.

%\paragraph{Summary of the Initialization and Iterative Procedure in Abstract}
%We give the summary below.
\begin{comment}
\begin{mdframed}[linewidth=0.2pt]
\textbf{Summary of the Initialization and Iterative Procedure in Abstract}
\begin{enumerate}
    \item \textbf{Initialization.}
    Choose \(Q_0\) so that
    \[
        I+Q_0/C_F
    \]
    satisfies the affine constraints.     \item \textbf{First correction.}
    The higher-order part of the spectral expansion is
    \[
        \sum_{j\ge2}  \frac{b_j}{C_F}   \cdot  P_j(Q_0).
    \]
    Define the unscaled deviation vector as
    \[
        \eta_1[t]
        =
        v_t^\top
        \left(
          \sum_{j\geq 2}    b_j \cdot  P_j(Q_0).        \right)
        v_t,
        \qquad t\in[m].
    \]
    The higher-order part of the polynomial expansion from $F(Q_0)/C_F $ incurs an error of $ \eta_1/C_F$ with respect to the affine constraints. Next, we solve for
    \[
        Mw_1=-\eta_1,
    \]
    and set
    \[
        Q_1
        =
        Q_0+
        \sum_{t\in[m]}w_1[t]\,v_t v_t^\top.
    \]
    \item \textbf{Subsequent Iterations $i\rightarrow i+1$. } 
    We choose an update \(
    \Delta_{i+1} \in \R^{d\times d}
    \)
    such that\[ 
    \frac{1}{C_F} \cdot F(Q_i) + \frac{1}{C_F} \cdot  \Delta_i  
    \]
    satisfies the affine constraints, i.e., $v_t^\top  (\frac{1}{C_F} \cdot F(Q_i) + \frac{1}{C_F} \cdot  \Delta_i  ) v_t =1$ for any $t\in [m]$.
    
    Moreover, we pick $\Delta_{i+1}$ specifically of the following form \[ 
    \Delta_{i+1} = \sum_{t\in [m]} w_{i+1}[t] v_t v_t^\top 
    \]
    where $w_{i+1}\in \R^{m}$ is obtained from \[ 
    Mw_{i+1} = -\eta_{i+1}
    \]
    with the updated deviation vector $\eta_{i+1}\in \R^m$ defined as \[
    \eta_{i+1}[t] = v_t^\top (\frac{1}{C_F} \cdot F(Q_i) ) v_t\,.
     \]
     
     Finally, we set $Q_{i+1} \leftarrow Q_i +\Delta_{i+1}$.
    \end{enumerate}
\end{mdframed}
\end{comment}
%\jnote{this is a condensed version. i commented out the original version that is more clear to me but may appear verbose...feel free to pick and choose }
\begin{mdframed}[linewidth=0.2pt]\textbf{Summary of the Initialization and Iterative Procedure.} 
\\\\
\textbf{Initialization:} Pick \(Q_0 = -C_F \cdot \calR\) so that \(I+Q_0/C_F\) satisfies the affine constraints.
\\\\
\textbf{Iterative Update $i\rightarrow i+1$:}
 Given \(Q_i\), 
%define the affine deviation \[ \widetilde{\eta}_{i+1}\coloneqq \calL(\frac{1}{C_F} \cdot F(Q_i))-\1_m, \]
%  %For convenience of our analysis, 
%We introduce the unscaled deviation of \[ 
%  \eta_{i+1}  \coloneqq C_F \cdot \widetilde{\eta}_{i+1}\,, 
%  \]
  define \[ 
  H_{i} \coloneqq \sum_{j\ge 2} b_j \cdot \left(P_j(Q_{i}) - P_j(Q_{i-1})  \right) \numberthis \label{eq:primal-corr}
  \]
  for $i\ge 1$, with \( 
  H_{0} \coloneqq \sum_{j\ge 2} b_j \cdot P_j(Q_0)\,.
  \) 
  
  Set the update as 
  \[ 
  \Delta_{i}  = \corr(H_i) \coloneqq \frac{1}{1-\gamma} \left(\calL^*M^{-1} \calL(-H_{i} ) + \gam \cdot H_{i}     \right)
   \]
%  
%  \[
 
% \Delta_{i+1} \coloneqq L^*(w_{i+1}) = \sum_{t\in [m]} w_{i+1}[t]v_tv_t^\top \,.
% \] 
 Finally, we set $Q_{i+1} \coloneqq Q_i+ \Delta_{i}$.
%   Equivalently, \[ L\!\left(\frac{1}{C_F}F(Q_i)+ \frac{\Delta_{i+1}}{C_F} \right)=\1_m . \] 
\end{mdframed}
%

	Geometrically, our update defined above can be viewed as \[ 
	\corr(H_i) =   \frac{1}{1-\gamma} (-\Pi_S +\gam I) H_i =  \frac{1}{1-\gamma} \Pi_S^\perp H_i - H_i \,,	\]
	by observing $ \Pi_S = \calL^*M^{-1} \calL  $ and  \[ 
	-\Pi_S + \gam I =   (I-\Pi_S)- (1-\gam )I = \Pi_S^\perp - (1-\gamma) \cdot  I\,.
	\]
	where we remind the reader $\Pi_S=\calL^*M^{-1}\calL$ and $\Pi_S^{\perp}=I-\Pi_S$ are orthogonal projections to $\calS = \operatorname{Im}(\mathcal L^*)$ and its orthogonal complement respectively. 

%{\color{orange} Sofia:
\begin{remark}
    The update admits a closed form
    \begin{align*}
        Q_{i+1} &= Q_0+\frac{1}{1-\gamma}\Pi_S^\perp(F(Q_i)-Q_i-C_FI)-(F(Q_i)-Q_i-C_FI)\,.
    \end{align*}
  \end{remark}
\jnote{i think your current version is correct but it might require a few lines of proof...add a proof to the appendix? its likely commented out below somewhere..}
%}
%\jnote{define $|\cdot |$ operation with some justification and add to the main body}
%\jnote{Sofia- feel free to add the closed form here if you think thats helpful, and then put the justification for the closed form in the appendix + a pointer :)}
%\snote{Indeed we have $\pi=I-\calL^* M^{-1}\calL$ and $\sum_{j\geq 2} b_j P_j(Q_i)=F(Q_i)-C_FI-Q_i$ so rearranging finds this equivalent expression.}
% \jnote{so lets see....\[ Q_1 =Q_0 + \corr(H_0) = Q_1 + \frac{1}{1-\gam} \pi(H_0) - H_0  \]
% in your language, it seems we have} \[
% Q_1^{(B)}
% =
% Q_0
% +
% \frac{1}{1-\gamma}
% \pi\!\left(\lvert Q_0\rvert-C_F I\right)
% -
% \left(\lvert Q_0\rvert-C_F I\right).
% \]

% But
% \[
% \begin{aligned}
% \lvert Q_0\rvert-C_F I
% &=
% \frac{1}{2}F(Q_0)-C_F I\\
% &=
% \frac{1}{2}\left(C_F I+Q_0+H_0\right)-C_F I\\
% &=
% \frac{1}{2}\left(Q_0+H_0-C_F I\right).
% \end{aligned}
% \]

% Therefore,
% \[
% \begin{aligned}
% Q_1^{(B)}
% &=
% Q_0
% +
% \frac{1}{2(1-\gamma)}
% \pi\!\left(Q_0+H_0-C_F I\right)\\
% &\qquad
% -
% \frac{1}{2}\left(Q_0+H_0-C_F I\right).
% \end{aligned}
% \]

% Using \(\pi(Q_0)=0\), we obtain
% \[
% \boxed{
% \begin{aligned}
% Q_1^{(B)}
% &=
% \frac{1}{2}Q_0
% +
% \frac{1}{2(1-\gamma)}\pi(H_0)
% -
% \frac{1}{2}H_0\\
% &\qquad
% -
% \frac{C_F}{2(1-\gamma)}\pi(I)
% +
% \frac{C_F}{2}I.
% \end{aligned}
% }
% \] 

%To see why this works, observe that for all $i \geq 0$, 
%\begin{enumerate}
%\item $L\!\left(\frac{1}{C_F}(Q_{i+1} - Q_i)\right) = \frac{1}{C_F}LL^{*}(w_{i+1}) = Mw_{i+1} = -\frac{1}{C_F}\eta_{i+1} = -\widetilde{\eta}_{i+1}$.
%\item $\widetilde{\eta}_{i+2} = L(\frac{1}{C_F} \cdot F(Q_{i+1})) - \1_m = \widetilde{\eta}_{i+1} + L(\frac{1}{C_F} \cdot (F(Q_{i+1}) - F(Q_i)))$ and 
%\[
%L\left(\frac{1}{C_F} \cdot (F(Q_{i+1}) - F(Q_i))\right) = L\!\left(\frac{1}{C_F}(Q_{i+1} - Q_i)\right) + \sum_{j=2}^{\infty}{\frac{b_j}{C_F}(P_j(Q_{i+1}) - P_j(Q_i))}
%\]
%\end{enumerate}
%so \[\widetilde{\eta}_{i+2} =\frac{1}{C_F} \calL\left( \sum_{j=2}^{\infty} b_j (P_j(Q_{i}) - P_j(Q_{i-1})) \right)
%=\frac{1}{C_F} \cdot H_{Q_{i+1}}
%\]
%consists of terms in $P_j(Q_{i+1})$ which don't already appear in $P_j(Q_i)$. As $i$ increases, these terms will have smaller and smaller norms which allows us to truncate the process once the norms of these terms become sufficiently small.
%\paragraph{Graph-matrix viewpoint.} 
%Before describing the subsequent iterations, we first introduce the graph-matrix
%viewpoint that organizes the terms appearing in the construction.  This
%viewpoint gives a clean combinatorial description of the matrices obtained by
%applying \(P_j\) to \(Q_i\). 
%
% Informally, if \(Q_i\) is a linear combination of
%shapes from a ground set \(\calB(Q_i)\), then \(P_j(Q_i)\) is approximated by the
%weighted sum of all \(j\)-fold proper concatenations of shapes from
%\(\calB(Q_i)\), up to lower-order intersection terms.
%\anote{It may be good to describe the general iteration here. One way to do this would be to write the expression for $\frac{1}{C_F}(F(Q_{i+1}) - F(Q_i))$ which leads to the next deviation which needs to be corrected for.}
%\jnote{check below}

At the heart of the iterative process is the following invariant with an upshot that the needed correction term becomes progressively smaller which %that 
ultimately allows us to terminate the iteration.% at some point.

\begin{proposition}[Invariant for Iterative Process] \label{prop:invariant-primal}
	For any $i \ge 1$, we have
	\[
 \calL( \frac{1}{C_F } F(Q_i) )  - \1_m = \frac{1}{C_F} \cdot  \calL(H_{i}).
\]
%\snote{is this hyphen/dash supposed to be here?} %\jnote{yah, but either is fine w. me}
In words, at the end of each iteration $i$, the entire affine deviation is incurred by the new
higher-order term \(H_{i}\).
\end{proposition}
In particular, this invariant is agnostic to the specific correction primitive: it holds for any update rule \(\Delta(H_i)\) satisfying
\(
    \calL\bigl(\Delta(H_i)+H_i\bigr)=0.
\)
The subsequent analysis, however, relies crucially on the particular choice of \(\corr(\cdot )\) defined above.

\begin{proof}[Proof of \cref{prop:invariant-primal}]
%First, using \(M=\calL\calL^*\), the correction satisfies
%\[
%\begin{aligned}
%    \calL(\Delta_i)
%    &=
%    \frac{1}{1-\gamma}
%    \left(
%        \calL\calL^*M^{-1}\calL(-H_i)
%        +\gamma\calL(H_i)
%    \right)\\
%    &=
%    \frac{-\calL(H_i)+\gamma\calL(H_i)}{1-\gamma}
%    =
%    -\calL(H_i).
%\end{aligned}
%\]

For the base case $i=0$, initialization gives
\(
    \calL\!\left(I+\frac{Q_0}{C_F}\right)=\1_m.
\)
%Since
%\[
%    F(Q_0)=C_FI+Q_0+H_0,
%\]
%we obtain
%\[
%    \calL\!\left(\frac{1}{C_F}F(Q_0)\right)-\1_m
%    =
%    \frac{1}{C_F}\calL(H_0).
%\]
%
Now suppose the claim holds at level \(i\). Since
\(
    Q_{i+1}-Q_i=\Delta_i
\)
and
\(
    F(Q_{i+1})-F(Q_i)=\Delta_i+H_{i+1},
\)
we have
\[
\begin{aligned}
    \calL\!\left(\frac{1}{C_F}F(Q_{i+1})\right)-\1_m
    &=
    \frac{1}{C_F}\calL(H_i)
    +\frac{1}{C_F}\calL(\Delta_i)
    +\frac{1}{C_F}\calL(H_{i+1})\\
    &=
    \frac{1}{C_F}\calL(H_{i+1}).
\end{aligned}
\]
This completes the induction.
\end{proof}

One should view the iterative process described here as an ``idealized'' process. In~\cref{sec:concrete-primal-graph-mat}, we give a concrete procedure that simplifies the above process by additionally discarding negligible norm components throughout the iterative process to streamline the analysis.

\begin{comment}
	
=======
\snote{
Recall the closed forms
\begin{align*}
    Q_{i}&=Q_0+\frac{1}{1-\gamma}\pi(|Q_i|-C_FI)-(|Q_i|-C_FI)\\
    H_i&=|Q_i|-C_FI-(|Q_{i-1}|-C_FI)=|Q_i|-|Q_{i-1}|.
\end{align*}
We can write 
\[\frac{1}{C_F}F(Q_i)=I+\frac{Q_i}{C_F}+\frac{|Q_i|-C_FI}{C_F}\]
and rearranging the closed form for $Q_i$,
\[I+\frac{Q_i}{C_F}+\frac{|Q_{i-1}|-C_FI}{C_F}=I+\frac{Q_0}{C_F}+\frac{\pi(|Q_{i-1}-C_FI|)}{C_F(1-\gamma)}.\]
now applying \(\calL\) to both sides and using $\calL\circ\pi=0$,
\begin{align*}
    \calL\left(I+\frac{Q_i}{C_F}+\frac{|Q_{i-1}|-C_FI}{C_F}\right)&=\calL\left(I+\frac{Q_0}{C_F}\right)=\1_m\\
    \1_m-\calL\left(\frac{Q_i+|Q_i|}{C_F}\right)&=\calL\left(\frac{|Q_i|-|Q_{i-1}|}{C_F}\right)\\
    \1_m-\calL\left(\frac{F(Q_i)}{C_F}\right)&=\frac{1}{C_F}\calL(H_i).
\end{align*}
}
\snote{Somehow this is off by a factor of $C_F$?}

>>>>>>> a45f63eda9de99bc8e21fc7bbbeec22d86e166e4
The base case is verified in our prior discussion, and it suffices for us to consider the inductive step. 
Fix \(i\ge 0\), and suppose inductively that
\[
\eta_{i+1}=\calL(H_{Q_i}).
\]
Using \(M=\calL\calL^*\) and the definition of \(\Delta_{i+1}\), we have
\[
\begin{aligned}
\calL\!\left(\frac{1}{C_F}(Q_{i+1}-Q_i)\right)
&=
\frac{1}{C_F}\calL(\Delta_{i+1})
%&=
%\frac{1}{C_F(1-\gamma)}
%\left(
%-MM^{-1}\eta_{i+1}
%+\gamma\calL(H_{Q_i})
%\right)\\
%&=
%\frac{-\eta_{i+1}+\gamma\eta_{i+1}}
%     {C_F(1-\gamma)}\\
=
-\frac{1}{C_F}\eta_{i+1}
=
-\widetilde{\eta}_{i+1}.
\end{aligned}
\]
Hence, the linear part of the update exactly cancels the current affine
deviation. It remains to identify the deviation created at the next iteration. We have
\[
\begin{aligned}
\widetilde{\eta}_{i+2}
&=
\calL\!\left(\frac{1}{C_F}F(Q_{i+1})\right)-\1_m\\
&=
\widetilde{\eta}_{i+1}
+
\calL\!\left(
\frac{1}{C_F}\bigl(F(Q_{i+1})-F(Q_i)\bigr)
\right).
\end{aligned}
\]
By the polynomial expansion of \(F\),
\[
\begin{aligned}
\calL\!\left(
\frac{1}{C_F}\bigl(F(Q_{i+1})-F(Q_i)\bigr)
\right)
&=
\calL\!\left(
\frac{1}{C_F}(Q_{i+1}-Q_i)
\right)\\
&\quad+
\frac{1}{C_F}
\calL\!\left(
\sum_{j\ge 2}b_j
\bigl(P_j(Q_{i+1})-P_j(Q_i)\bigr)
\right).
\end{aligned}
\]
The first term cancels \(\widetilde{\eta}_{i+1}\), and therefore
\[
\begin{aligned}
\widetilde{\eta}_{i+2}
&=
\frac{1}{C_F}
\calL\!\left(
\sum_{j\ge 2}b_j
\bigl(P_j(Q_{i+1})-P_j(Q_i)\bigr)
\right)=
\frac{1}{C_F}\calL(H_{Q_{i+1}}).
\end{aligned}
\]
Multiplying by \(C_F\) completes the induction.

%As \(i\) increases, we will show these
%new contributions become progressively smaller, allowing us to truncate the
%iteration once their norms are sufficiently small.
\end{comment}
\subsection{Graph Matrices:
 Visualizing Orthogonal Polynomials} \label{sec:graph-matrix}

%\snote{ To demystify the above construction and its analysis, we now introduce our main tool: graph matrices.}
 To demystify the above construction and its analysis, we now introduce our main tool---graph matrices.
 
\paragraph{Overview of Graph Matrices.}
Graph matrices provide a versatile framework for analyzing structured random matrices
arising in average-case problems. They have played a central role in proving
Sum-of-Squares lower bounds for a wide range of average-case inference problems
\cite{BHKKMP16, PR20, pang:LIPIcs.CCC.2021.26, JPRTX, JPRX23, KPX24, NGCA24, Xu25, PX25,
PotechinXu2026Theta}. Beyond the SoS literature, graph matrices have also been used
to analyze power-sum decompositions of polynomials \cite{BHKX22}, the ellipsoid
fitting conjecture \cite{PTVW22, HKPX23}, and first-order iterative algorithms,
including belief propagation and approximate message passing \cite{JP24}.


For completeness, we begin with a brief introduction to graph-matrix techniques specialized to our setting, with particular emphasis on their nontrivial connection to the orthogonal polynomials that will appear later. Readers familiar with graph matrices may skip the discussion through the subsection ``Visualizing Orthogonal Polynomials.'' 

Throughout this section, the underlying random input is a collection of
Gaussian vectors
\(G =
  \{v_1,\ldots,v_m\} 
\) in dimension $d$ 
drawn independently from the distribution specified above. Equivalently, we write
\(G\in \mathbb R^{m\times d}\) for the matrix whose \(i\)-th row is \(v_i\).
The graph-matrix framework applies much more generally, for instance to random
graphs and  non-product distributions; we refer the reader to the references above
for a more thorough introduction.



\begin{definition}[Shape]
A shape \(\tau\) is a tuple
\(
    \tau=(V(\tau),U_\tau,V_\tau,E(\tau))
\)
associated with a multigraph \((V(\tau),E(\tau))\). Each vertex in \(V(\tau)\) is
assigned a vertex type, which determines the range from which its labels are drawn.
Each edge \(e\in E(\tau)\) is assigned a Fourier index \(t(e)\in \mathbb N\).
The subsets \(U_\tau,V_\tau\subseteq V(\tau)\) are called the left and right
boundaries of the shape, respectively.
\end{definition}

We emphasize that \(V_\tau\) denotes the right boundary of the shape, whereas
\(V(\tau)\) denotes the full vertex set of the underlying graph.

\begin{definition}[Mapping of a shape]
Given a shape \(\tau\), a mapping of \(\tau\) is a function
\(
    \sigma:V(\tau)\to \mathbb N
\)
such that:
\begin{enumerate}
    \item each vertex is assigned a label from the range specified by its vertex type;
    \item \(\sigma\) is injective on vertices of the same type.
\end{enumerate}
\end{definition}

In \cref{fig:M_alpha_beta}, we illustrate the shapes corresponding to the matrices
\(M_\alpha\) and \(M_\beta\) defined in \cref{eq:M-decomposition}, which will be a
central focus of our analysis. These shapes have two vertex types: square vertices
take labels in \([m]\), while circle vertices take labels in \([d]\). The two ovals
indicate the left and right boundaries \(U_\tau\) and \(V_\tau\).

\begin{figure}[ht!]
    \centering
    \begin{subfigure}[b]{0.31\textwidth}
        \centering
        \includegraphics[width=0.85\textwidth]{diagrams/GOE.pdf}
        \caption{\(\mathsf{GOE}\), zero diagonal.}
        \label{fig:GOE}
    \end{subfigure}
    \hfill
    \begin{subfigure}[b]{0.31\textwidth}
        \centering
        \includegraphics[width=1.2\textwidth]{diagrams/M_alpha.pdf}
        \caption{\(M_{\alpha}\).}
        \label{fig:M_alpha}
    \end{subfigure}
    \hfill
    \begin{subfigure}[b]{0.31\textwidth}
        \centering
        \includegraphics[width=\textwidth]{diagrams/M_beta.pdf}
        \caption{\(M_{\beta}\).}
        \label{fig:M_beta}
    \end{subfigure}

    \captionsetup{width=.9\linewidth}
    \caption{Graph-matrix representations of a \(d\times d\) \(\mathsf{GOE}\)
    matrix with zero diagonal, and of the \(m\times m\) matrices \(M_\alpha\) and
    \(M_\beta\) from \cref{eq:M-decomposition}. Square vertices take labels in
    \([m]\), and circle vertices take labels in \([d]\). The two ovals indicate the
    left and right boundaries \(U_\tau,V_\tau\). If an edge \(e\) is not explicitly
    labeled with an index, then \(t(e)=1\) by default.}
    \label{fig:M_alpha_beta}
\end{figure}

We now describe how a shape gives rise to a matrix once the underlying random input
\(G\) is fixed.

\begin{definition}[Graph matrix of a shape]
\label{def:graph-matrix-for-shape}
Given a shape \(\tau\), its graph matrix \(M_\tau\) is indexed by boundary labelings.
For boundary labelings \(S\) of \(U_\tau\) and \(T\) of \(V_\tau\), define
\[
    \cm_{\tau}[S,T]
    =
    \sum_{\substack{
        \sigma:\,V(\tau)\to \mathbb N \\
        \sigma \text{ is a mapping of } \tau \\
        \sigma(U_\tau)=S,\ \sigma(V_\tau)=T
    }}
    \prod_{e\in E(\tau)}
    \chi_{t(e)}\bigl(G[\sigma(e)]\bigr).
\]
Here, if \(e=\{x,y\}\) is an edge between a square vertex and a circle vertex, then
\(G[\sigma(e)]\) denotes the corresponding entry \(G[\sigma(x),\sigma(y)]\).
\end{definition}

For each entry \(\cm_\tau[S,T]\), the boundary labels are fixed by the constraints
\(\sigma(U_\tau)=S\) and \(\sigma(V_\tau)=T\). Thus the entry is a sum over labelings
of the internal vertices
\(
    V(\tau)\setminus (U_\tau\cup V_\tau).
\)
For example, consider the shapes in \cref{fig:M_alpha_beta}. Suppose
\(G\in \mathbb R^{m\times d}\), with square vertices labeled by elements of \([m]\)
and circle vertices labeled by elements of \([d]\). Then for \(i\neq j\in[m]\),
\begin{align*}
    \cm_{\alpha}[i,j]
    &=
    \sum_{a\neq b\in[d]}
    \chi_1(G[i,a])\chi_1(G[i,b])
    \chi_1(G[j,a])\chi_1(G[j,b]), \\
    \cm_{\beta}[i,j]
    &=
    \sum_{a\in[d]}
    \chi_2(G[i,a])\chi_2(G[j,a]).
\end{align*}
These two matrices will come up again when we analyze the system of affine constraints.
\begin{remark}
The nonzero entries of \(M_\alpha\) may equivalently be written as
\[
    \cm_{\alpha}[i,j]
    =
    2\sum_{a<b\in[d]}
    \chi_1(G[i,a])\chi_1(G[i,b])
    \chi_1(G[j,a])\chi_1(G[j,b]).
\]
The factor of \(2\) will be important in our later analysis, though it can be
safely ignored on a first reading.
\end{remark}

Since the mapping \(\sigma\) is injective on vertices of the same type, and since
\(U_\tau\neq V_\tau\) in both examples in \Cref{fig:M_alpha_beta}, there is no valid
mapping with \(\sigma(U_\tau)=\sigma(V_\tau)\). Consequently, both \(M_\alpha\) and
\(M_\beta\) have zero diagonal.

Following the notation of \cite{HKPX23}, we specialize the graph-matrix framework to
the following setting:
\begin{itemize}
    \item \(G\in\mathbb R^{m\times d}\) is a random Gaussian matrix whose rows are
    \(v_1,\ldots,v_m\sim N(0,\frac1d I_d)\).
    \item The Fourier characters \(\{\chi_t\}_{t\in\mathbb N}\) are the appropriately
    scaled Hermite polynomials.
    \item For all graph matrices appearing in our analysis:
    \begin{itemize}
        \item the left and right boundary labelings satisfy \(|S|=|T|=1\);
        \item there are two vertex types: square vertices take labels in \([m]\), and
        circle vertices take labels in \([d]\).
    \end{itemize}
\end{itemize}


\paragraph{Shape Concatenation and Intersection Terms}

We give a brief overview of the multiplication of graph matrices to shed light on the non-trivial connection with orthogonal polynomials to be introduced. To start with, given shapes $\tau_1, \tau_2$, we call them \emph{composable} if $V_{\tau_1}$ and $U_{\tau_2}$ have matching boundaries, specifically that they have the same number of both square and circle vertices. For any composable pair $\tau_1,\tau_2$, we define their properly concatenated shape as simply the underlying shape obtained by gluing $\tau_1,\tau_2$ together along the boundary (specifically, the right boundary of $\tau_1$ and equivalently the left boundary of $\tau_2$). Formally it is defined as follows, 
%provided that they have matching boundaries, specifically that they have the exact number of both square and circle vertices. We give its formal definition as follows,


\begin{definition}[Proper Shape Concatenation]\label{def:shape-concate}
Given shapes $\tau_1$ and $\tau_2$ that are composable, we take the concatenation $\tau_1 \circ \tau_2$ of $\tau_1$ and $\tau_2$ to be the shape such that:
\begin{enumerate}
		\item $V(\tau_1 \circ \tau_2) = V(\tau_1) \cup V(\tau_2)$ where the vertices in $V_{\tau_1}$ are identified with the vertices in $U_{\tau_2}$.
        \item $U_{\tau_1 \circ \tau_2} = U_{\tau_1}$ and $V_{\tau_1 \circ \tau_2} = V_{\tau_2}$.
		\item $E(\tau_1 \circ \tau_2) = E(\tau_1) \cup E(\tau_2)$. 
\end{enumerate}
%In other words, $\ \circ \beta$ is the shape obtained by gluing $\alpha$ and $\beta$ together along $V_{\alpha}$ and $U_{\beta}$, keeping all other vertices distinct, and taking all of the edges in both $\alpha$ and $\beta$.
	
This definition extends naturally to a sequence of $k$ shapes $\tau_1,\ldots,\tau_k$ such that for each $j \in [k-1]$, $|V_{\tau_j}| = |U_{\tau_{j+1}}|$. We call this a $k$-way (or $k$-fold) concatenation.  We can obtain the concatenation $\tau_1 \circ \cdots \circ \tau_k$ by concatenating these shapes together one by one (it is not hard to check that the order of the concatenations does not matter).
\end{definition}

We illustrate the concatenation of $\al\circ \beta$ in the following diagram, and we use the dotted oval to denote the concatenation boundary of $V_\al = U_\beta$.

\begin{figure}[ht!]
    \centering
    \begin{subfigure}[b]{0.45\textwidth}
        \centering
        \includegraphics[width=\textwidth]{diagrams/concat.pdf}
        \caption{Proper concatenation \(\alpha\circ\beta\).}
        \label{fig:concat-proper}
    \end{subfigure}
    \quad
    \begin{subfigure}[b]{0.45\textwidth}
        \centering
        \includegraphics[width=\textwidth]{diagrams/concat-int.pdf}
        \caption{Concatenation with an intersection.}
        \label{fig:concat-intersection}
    \end{subfigure}
    \centering
	 \begin{subfigure}[b]{0.8\textwidth}
        \centering
        \includegraphics[width=\textwidth]{diagrams/p4.pdf}
        \caption{Proper $4$-fold concatenation \(\alpha\circ\beta\circ \al \circ \beta \).}
%        \label{fig:concat-proper}
    \end{subfigure}
    \caption{Examples of shape concatenations.}
    \label{fig:concat}
\end{figure}

With the proper concatenation introduced, we now return to the motivation for
introducing this operation. For concreteness, consider two shapes \(\alpha\) and
\(\beta\) with compatible boundaries. The reason concatenation is central is that
it captures the leading contribution in the product of the corresponding graph
matrices. Namely, when we multiply \(\cm_\alpha\) and \(\cm_\beta\) using ordinary
matrix multiplication, one might hope in an idealized setting that
\[
    \cm_{\alpha}\cdot \cm_{\beta}
    =
    \cm_{\alpha\circ\beta}.
\]
This identity, however, is not exact. The product on the left generally produces
additional terms, known in the graph-matrix literature as \emph{intersection
terms}. 

The source of the discrepancy is the difference between global and local
injectivity constraints. The graph matrix \(\cm_{\alpha\circ\beta}\) is defined
using labelings that are injective across the entire concatenated shape
\(\alpha\circ\beta\). By contrast, the product \(\cm_\alpha \cdot \cm_\beta\) only
enforces injectivity separately within the copy of \(\alpha\) and within the copy
of \(\beta\). As a result, vertices from \(\alpha\) that are not identified
through the boundary may nevertheless collide with vertices from \(\beta\). 


We now formalize this notion through intersection patterns, which record which vertices from the two shapes collide during the matrix product.
%\anote{This should be modified slightly as we have two types of vertices.}
\begin{definition}[Intersection patterns/shapes]
Given two shapes $\tau_1$ and $\tau_2$ that are composable, we define an intersection pattern for $\tau_1$ and $\tau_2$ to be a non-empty set $I \subseteq \left(V(\tau_1) \setminus V_{\tau_1}\right) \times \left(V(\tau_2) \setminus U_{\tau_2}\right)$ of intersection edges where each vertex is incident to at most one edge. Whenever we have an intersection edge $e = \{a,b\}$ (between vertices of the same type), this means that when we consider maps $\sigma: V(\tau_1) \cup V(\tau_2) \to [n]\cup [d] $, we have the restriction that $\sigma(a) = \sigma(b)$.

We define $\tau_{I} = \tau_{I}(\tau_1\cdot \tau_2)$ to be the shape obtained from $\tau_1$ and $\tau_2$ as follows:
\begin{enumerate}
\item Start with the properly concatenated shape $\tau_1\circ \tau_2$;
\item For each intersection edge $e = \{a,b\}$, merge $a$ and $b$ together into one vertex.
\end{enumerate}
\end{definition}
%For convenience of our analysis
To make our analysis more convenient, instead of merging the collided vertices, we will primarily view intersection patterns as shapes except that they additionally come with intersection edges that specify vertex collisions. See the blue dotted edge in \cref{fig:concat-intersection} for an example of an intersection edge, together with the resulting shape corresponding to this
intersection pattern.
\begin{definition}
We define $Int_{\alpha,\beta}$ to be the set of possible intersection patterns for $\alpha$ and $\beta$.
\end{definition}

\begin{definition}[$\circ$ and $\cdot$ notation]
	Given two shapes $\al_1$ and $\al_2$, we use $\circ$ to denote proper concatenation, and $\cdot $ to denote (possibly) improper concatenations. In other words, $\al_1\circ  \al_2$ denotes a single shape obtained by proper concatenation, while $\al_1\cdot \al_2$ denotes a collection of shapes that include both proper concatenations and intersection shapes.
\end{definition}
%\snote{With the intersection terms introduced, we can deduce the graph-matrix analog of ordinary matrix multiplication.}
With the intersection terms introduced, we can deduce the following analog of ordinary matrix multiplication of graph matrices in the diagram language.
\begin{proposition}
For all shapes $\alpha$ and $\beta$ that can be concatenated, \[
{\cm}_{\alpha}{\cm}_{\beta}  = \cm_{\al\cdot \beta}   = {\cm}_{\alpha \circ \beta} + \sum_{I \in Int_{\alpha,\beta}}{{\cm_{\tau_I}}}\,.\]
\end{proposition}


\paragraph{Orthogonal Polynomials as Horizontal Concatenation}
We now introduce another important ingredient of our work: orthogonal polynomials.
The connection between graph matrices and orthogonal polynomials was first made
explicit in \cite{PotechinXu2026Theta}. In their setting, one considers the
orthogonal polynomials \(\{P_t(x)\}_{t\ge 0}\) with respect to the semicircle
distribution \(\mu_{\mathrm{sc}}\) on \([-2,2]\). They show that if \(Q\) is a
linear combination of graph matrices from a specific family called \emph{backbone-correction shapes}, then these individual matrices are all freely independent of each other, and \(P_t(Q)\) admits a clean decomposition via the \(t\)-fold proper concatenations of the shapes appearing in \(Q\) additionally provided that the matrix $Q$ is appropriately normalized.

Returning to our setting, the relevant shapes are no longer the
backbone-correction shapes studied in \cite{PotechinXu2026Theta}. Nevertheless,
it is natural to expect that the underlying phenomenon is not specific to this
particular family. Rather, one may hope that the same graph-matrix
interpretation of orthogonal polynomials applies more broadly, with the
essential condition being an appropriate form of free independence among the
underlying graph matrices.

In our setting, we call the shapes that arise in the analysis
\emph{backbone-dangling shapes}. At a high level, these shapes may be viewed as
generalizations of the graph matrices studied in \cite{HKPX23}; indeed, our
initialization \(Q_0\) is itself based on the identity-perturbation construction
considered therein. For intuition, we now give a few examples of backbone-dangling shapes that arise
in our analysis, deferring their formal definition and classification to later
sections.

 Analogous to \cite{PotechinXu2026Theta}, a key component of our
analysis is to show that the graph matrices associated with backbone-dangling
shapes remain freely independent. Crucially, the free independence of the underlying graph matrices ensures that the orthogonal-polynomial interpretation
via graph matrices continues to hold: applying an orthogonal
polynomial to a linear combination of these graph matrices is captured, up to negligible error terms, by graph matrices associated with concatenated shapes.

\begin{proposition}[(Informal) Equivalence of Chebyshev Polynomials and Concatenated Shapes]
%For an inner matrix \(Q_i\) with shape ground set \(\calB(Q_i)\) such that \[ 
%Q_i = \sum_{\al\in \calB(Q_i)} c_\al \cdot \cm_\al \,,
%\]
For a linear combination of backbone-dangling shapes, $Q = \sum_{\psi \in \calB(Q) : \text{backbone-dangling} } c_\psi  \cdot \cm_\psi$, provided $\Var(Q) =1$,
 for any $j\geq 1$,  we have
\[
    P_j(Q)
    =
    \sum_{\substack{\tau:\text{ \(j\)-fold concatenation of }\calB(Q) \\ \tau =\psi_1\circ \psi_2\circ \ldots \circ \psi_j} }
    c_\tau \cm_\tau + o_d(1),
\]  
up to lower-order error terms of spectral norm $o_d(1)$ and the coefficients are given by \(
c_\tau = \prod_{t\in [j]} c_{\psi_t} \).

In words,   \(P_j(Q)\) is approximated by the
weighted sum of all \(j\)-fold proper concatenations of shapes from
\(\calB(Q)\), up to lower-order error terms.

\end{proposition}

With the graph-matrix framework in place, together with the correspondence between orthogonal polynomials and graph matrices of concatenated shapes, we now describe our analysis of the iterative construction and the technical challenges specific to the ellipsoid fitting problem. %return to the iterative construction.

%
%\subsection{Iterative Construction Made Concrete via Graph Matrices} \label{sec:iterative-construction}
%%This correspondence lets us identify which terms have already been corrected
%%and which terms are new. 
%In this section, we give a full description of the iterative process, which will in turn identify the collection of graph matrices that arise throughout our analysis. Throughout this section, we restrict our attention to proper shapes.
%
%For an inner matrix $Q_i$, we first consider what arises when we apply $F$, more concretely, in the polynomial expansion of $\frac{1}{C_F} F(Q_i)$, which is \[ 
%\frac{1}{C_F} F(Q_i) = I + Q_i/C_F  + \sum_{j\geq 2} \frac{b_j}{C_F} P_j(Q_i)
%\,.\]
%Let's zoom into the higher-degree term. For each \(j\ge2\), define the collection of  shapes in the $j$-way concatenation of shapes in $Q_i$ as
%\[
%    \calP_j(Q_i)
%    :=
%    \left\{
%        \tau_1\circ\cdots\circ\tau_j:
%        \tau_t\in\calB(Q_i)\ \forall t\in[j]
%    \right\}.
%\]
%Moreover, depending on whether the concatenation uses any shape from the most recent update, specifically in $\calB(Q_i)\setminus \calB(Q_{i-1})$,  we can partition the collection of $j$-way concatenation shapes as \[
%    \textsf{Old}\text{-}\calP_j(Q_i)
%    =
%    \left\{
%        \beta=\tau_1\circ\cdots\circ\tau_j\in\calP_j(Q_i):
%        \tau_t\in\calB(Q_{i-1})\ \forall t\in[j]
%    \right\},
%\]
%and
%\[
%    \textsf{New}\text{-}\calP_j(Q_i)
%    =
%    \left\{
%        \beta=\tau_1\circ\cdots\circ\tau_j\in\calP_j(Q_i):
%        \exists t\in[j]\ \text{such that }\
%        \tau_t\in\calB(Q_i)\setminus\calB(Q_{i-1})
%    \right\}.
%\]
%%Here and below, the notation also denotes the corresponding weightedgraph-matrix sums.
%This partition then allows us to make explicit the following invariant that we would like to maintain across the iterations: in particular, the "new" terms are the only source of deviation with respect to the affine constraints.
%\begin{proposition}[Iteration Invariant]
%\label{prop:iterative-invar}
%For any \(i\ge1\), consider the inner matrix \(
%Q_i = \sum_{\tau} c_\tau \cdot \cm_\tau 
%\),
% we have
%\[
%\frac{1}{C_F}F(  Q_i)
%=
%I+Q_i/C_F + \frac{1}{C_F}\cdot 
%\sum_{j\ge2}
%\left(
%    \sum_{\tau\in\textsf{Old}\text{-}\calP_j(Q_i)}
%    c_\tau\cm_\tau
%    +
%    \sum_{\tau\in\textsf{New}\text{-}\calP_j(Q_i)}
%    c_\tau\cm_\tau
%\right).
%\]
%Moreover,
%\[
%    I+Q_i/C_F + \frac{1}{C_F}\cdot 
%    \sum_{j\ge2}
%    \sum_{\tau\in\textsf{Old}\text{-}\calP_j(Q_i)}
%    c_\tau\cm_\tau
%\]
%satisfies the affine constraints.
%\end{proposition}
%
%\begin{remark}
%The same statement applies to the base level \(i=0\) if the old part is treated
%as empty.
%\end{remark}
%
%We defer the proof of \cref{prop:iterative-invar} to
%\cref{sec:defer-construction}.  The invariant shows that, at iteration \(i\),
%the only remaining source of affine-constraint error is the new part
%\(
%    \sum_{j\ge2}
%    \sum_{\tau\in\textsf{New}\text{-}\calP_j(Q_i)}
%    c_\tau\cm_\tau.
%\)
%This is therefore the correction target for the next iteration. One lingering question
%remains: how do we correct these affine deviations?
%
%To answer this, we again draw inspiration from the identity-perturbation
%construction. We seek a correction matrix of the form
%\(
%    \sum_{t\in[m]} w[t]\,v_tv_t^\top
%\)
%whose constraint values cancel the current deviation. Concretely, if
%\[
%    \eta_{\mathrm{new}}[i]
%    :=
%    v_i^\top
%    \left(
%        \sum_{\tau\in \mathsf{New}\text{-}\calP_j(Q_i)}
%        c_\tau \cm_\tau
%    \right)
%    v_i
%\]
%denotes the affine deviation to be corrected, then we want
%\[
%    v_i^\top
%    \left(
%        \sum_{t\in[m]} w[t]\,v_tv_t^\top
%    \right)
%    v_i
%    =
%    -\eta_{\mathrm{new}}[i],
%    \qquad \forall i\in[m].
%\]
%Equivalently, this reduces to solving the linear system
%\(
%    Mw=-\eta_{\mathrm{new}}.
%\)
%Thus, just as in the identity-perturbation construction, the desired correction
%is obtained by applying \(M^{-1}\) to the current affine deviation, and thereby yielding an update to our inner matrix.
%
%
%Importantly, due to to correlated randomness between the rank-$1$ forms and the deviation vector, we carefully choose our correction update as follows for any $\tau \in \mathsf{New-}P_j(Q_i)$:
% \[ 
%    \corr(\tau) \coloneqq -  \frac{1}{1-\gam } \left( \mathcal{L}^*(M^{-1} \eta_{\tau}) - \gam \cdot \tau  \right)\,.
%    \]
%    \jnote{this needs to be highlighted}
%
%%In other words, each iteration corrects the newly produced affine deviation by
%%solving the same constraint matrix \(M\), and then converts the resulting
%%coefficient vector \(w\) into the rank-one correction
%%\[
%%    \sum_{t\in[m]} w[t]\,v_tv_t^\top .
%%\]
%
%\paragraph{Concrete Iterative construction.}
%We summarize our iterative process as follows.
%\begin{mdframed}[linewidth=0.2pt]
%\textbf{Iterative update \(i\to i+1\).}
%\begin{enumerate}
%    \item For each \(j\ge2\), form the set of \(j\)-fold concatenations
%    \[
%        \calP_j(Q_i)
%        =
%        \left\{
%            \tau_1\circ\cdots\circ\tau_j:
%            \tau_t\in\calB(Q_i)\ \forall t\in[j]
%        \right\}.
%    \]
%
%    \item Split \(\calP_j(Q_i)\) into old and new terms:
%    \[
%        \textsf{Old}\text{-}\calP_j(Q_i)
%        =
%        \left\{
%            \tau_1\circ\cdots\circ\tau_j:
%            \tau_t\in\calB(Q_{i-1})\ \forall t\in[j]
%        \right\},
%    \]
%    and
%    \[
%        \textsf{New}\text{-}\calP_j(Q_i)
%        =
%        \left\{
%            \tau_1\circ\cdots\circ\tau_j:
%            \exists t\in[j]\ \text{such that }\
%            \tau_t\in\calB(Q_i)\setminus\calB(Q_{i-1})
%        \right\}.
%    \]
%
%    \item For each new concatenation
%    \[
%        \beta=\tau_1\circ\cdots\circ\tau_j
%        \in \textsf{New}\text{-}\calP_j(Q_i),
%    \]
%    its coefficient is
%    \[
%        c_\beta
%        =
%        b_j \cdot \prod_{t=1}^j c_{\tau_t}.
%    \]
%    in the polynomial expansion of $ F(Q_i)$.
%    This term appears in the spectral expansion of \(F(Q_i)\), but has not yet
%    been corrected inside \(Q_i\).
%
%    \item Define the new unscaled affine residual by
%    \[
%        \eta_{i+1}[t]
%        =
%        v_t^\top
%        \left(
%            \sum_{j\ge2}
%            \sum_{\tau\in\textsf{New}\text{-}\calP_j(Q_i)}
%            c_\tau\cm_\tau
%        \right)
%        v_t,
%        \qquad t\in[m].
%    \]
%    Solve
%    \[
%        Mw_{i+1}=-\eta_{i+1}.
%    \]
%    Equivalently, one may decompose
%    \[
%        w_{i+1}
%        =
%        \sum_{j\ge2}
%        \sum_{\tau\in\textsf{New}\text{-}\calP_j(Q_i)}
%        w_\tau,
%    \]
%    where \(w_\tau\) is the correction vector associated with the new shape
%    \(\tau\). We will set its corresponding correction term as \[ 
%    \corr(\tau) \coloneqq -  \frac{1}{1-\gam } \left( \mathcal{L}^*(M^{-1} \eta_{\tau}) - \gam \cdot \tau  \right)\,.
%    \]
%
%    \item Update the inner matrix by
%    \[
%        Q_{i+1}-Q_i
%        =
%        \sum_{t\in[m]}w_{i+1}[t]\,v_t v_t^\top.
%    \]
%    Equivalently,
%    \[
%        Q_{i+1}-Q_i
%        =
%        \sum_{t\in[m]}
%        \sum_{j\ge2}
%        \sum_{\tau\in\textsf{New}\text{-}\calP_j(Q_i)}
%        w_\tau[t]\,v_t v_t^\top.
%    \]
%    In graph-matrix language, the ground set evolves by adding the correction
%    shapes associated with the newly appearing concatenations:
%    \[
%        \calB(Q_{i+1})
%        =
%        \calB(Q_i)
%        \cup
%        \left\{
%            \text{correction shapes for }
%            \tau:
%            \tau\in\textsf{New}\text{-}\calP_j(Q_i),\ j\ge2
%        \right\}.
%    \]
%\end{enumerate}
%\end{mdframed}
%
%\begin{remark}
%In the formal analysis, we use an inner truncation: for each \(Q_i\), we only
%retain \(P_j(Q_i)\) for \(j\le D\).
%\end{remark}
%
%For the final construction, we take \(Q=Q_t\) for a truncation level \(t\), and
%then handle the remaining error terms directly.  This is addressed in~\cref{sec:formal-truncation}.  For the purposes of the overview, however, it is
%useful to think of \(Q\) as the limit of the iterates \(Q_i\) as \(i\to\infty\). 
%
%\anote{This feels somewhat redundant with Section 2.2. One possibility would be to focus on how to express an iteration using graph matrices, though this would require discussing $M^{-1}$ first.}
%\jnote{comment here - let's move the summary to the appendix}

\subsection{Correction Primitive: Finding the Inverse \texorpdfstring{\(M^{-1}\)}{of M}}
%With the iterative construction of our candidate matrix in place---particularly the construction of the inner matrix \(Q\)---we now turn to the technical challenges specific to ellipsoid fitting. 
At each step, the iterative procedure corrects the current affine deviation by applying \(M^{-1}\) to the deviation
vector \(\eta\), closely mirroring the identity-perturbation construction. The
difficulty is that, near the conjectured threshold, the matrix \(M\) can no
longer be analyzed as a small perturbation of the identity. Consequently, the
Neumann-series-based approach used in previous analyses is no longer sufficient,
and a more refined understanding of \(M^{-1}\) becomes necessary.

Let us first recall the Neumann-series-based approach, a.k.a. matrix Taylor expansion, used in the prior works. We closely follow the notation of \cite{HKPX23}, and use the decomposition 
\[ \numberthis \label{eq:M-decomposition}
    M = A+B,
\]
where
\[
    A
    :=
    \left(1+\frac1d\right)I_m
    +
    M_\alpha+M_\beta+M_D  
\]
and
\[
    B
    :=
    \frac1d J_m
    +
    \frac1d
    \left(
        \mathbf 1_m\eta^\top+\eta\mathbf 1_m^\top
    \right).
\]
Here \(J_m\) is the all-ones matrix, \(M_\alpha\) and \(M_\beta\) are
off-diagonal centered matrices, and \(M_D\) is diagonal and in fact negligible spectrally.  More explicitly, for
\(i\ne j \in [m] \),
\[
    \cm_\alpha[i,j]
    :=
    \sum_{a\ne b\in[d]}
        v_i[a]v_i[b]v_j[a]v_j[b] = 2\cdot \sum_{a<b\in [d]} v_i[a]v_i[b]v_j[a]v_j[b] \numberthis \label{eq:M_al_def}
\]
\[
    \cm_\beta[i,j]
    :=
    \sum_{a\in[d]}
    \left(v_i[a]^2-\frac1d\right)
    \left(v_j[a]^2-\frac1d\right), \numberthis \label{eq:M_beta_def}
\]
and
\[
    \cm_D[i,i]
    :=
    \|v_i\|_2^4-\frac2d\|v_i\|_2^2-1. \numberthis \label{eq:M_D_def}
\]
See \cref{fig:M_alpha} and \cref{fig:M_beta} for their diagrams. The main insight from the prior work is that $B$ is a rank-$2$ matrix, and one can use the Woodbury matrix identity to reduce the inverse of $M$ to that %for 
of $A$. 

For expository purposes, we advise the reader to ignore the negligible norm components in $A$ and focus on the first three terms \(
 \cm_\al +\cm_\beta + I\,.
\)
The omitted components do not affect the analysis in any essential way.
In particular,
in the regime with a sufficiently large constant factor gap, one may write
\[
     {A}^{-1}
    =
    \frac{1}{I+( {A}-I)}
    =
    \sum_{t\geq 0} (-1)^t ( {A}-I)^t,
\]
and the series converges precisely because \(\|A-I\|_{\mathrm{sp}}<1\).
However, this perturbative argument breaks down close to the threshold, even
when we remain a constant factor away from it.
Our key idea is that, although \( A\) is no longer close to the identity,
it remains close in spirit to a Wishart matrix. The only caveat is that the
underlying vectors are not genuinely i.i.d., because of their tensor-product
structure. Momentarily ignoring this distinction, if \(  A\) were an ideal
Wishart matrix with aspect ratio \(\gamma\), then its spectrum would satisfy
\[
    \mathsf{spec}(A)
    \subseteq
    \bigl[(1-\sqrt{\gamma})^2,\,(1+\sqrt{\gamma})^2\bigr] \pm o_d(1).
\]

We now identify the relevant aspect ratio $\gamma$ at the threshold.
Let us view \(A\) as the Gram matrix of the lifted rank-one vectors \(v_i v_i^\top\),
after removing the low-rank directions captured by \(B\). By symmetry, these
lifted vectors do not live in the full \(d^2\)-dimensional space but instead in a subspace of dimension
\(
   N_0= d(d+1)/2
\).
%\jnote{is there a -1?}
Thus, heuristically, \( A\) should behave like a Wishart matrix formed from
\(m\) samples in dimension \(N_0\approx d^2/2\). The corresponding Marchenko--Pastur aspect
ratio is therefore
\[
    \gamma
    =
    \frac{m}{N_0}
    =
    \frac{2m}{d^2}(1+o_d(1)).
\]
Throughout this work, we write simply \(\gamma=2m/d^2\),
suppressing lower-order terms in \(d\).
% For our analysis, once we view $\cm_D$ and $\frac{1}{d}I_d$ as $o_d(1)$ error terms, we can translate from $A$ to the original matrix $A$ straightforwardly.


%In Section~\jnote{add pointer}, via trace moment method, we rigorously show that the spectrum of $A$ indeed follows the Marchenko--Pastur law, even at the spectral edge. Similar result has been obtained in the literature 


\paragraph{Inversion from Marchenko-Pastur.}

The above heuristic next leads us to the following polynomial basis for the
Marchenko--Pastur distribution with parameter \(\gamma\).
\begin{definition}[MP Polynomial Basis] \label{def:MP-basis}
Fix \(0<\gamma<1\). Let \(\{\mathfrak \q_t^{(\gamma)}\}_{t\ge0}\) be the polynomial sequence defined by
\[
    \q_0^{(\gamma)}(x)=1,
    \qquad
   \q_1^{(\gamma)}(x)=x-1,
\]
and, for \(t\ge1\),
\[
    \q_{t+1}^{(\gamma)}(x)
    =
    \bigl(x-(1+\gamma)\bigr)\q_t^{(\gamma)}(x)
    -
    \gamma \cdot \q_{t-1}^{(\gamma)}(x).
\]
When the dependence on \(\gamma\) is clear, we write \(\q_t=\q_t^{(\gamma)}\). 
%For a matrix \(A\), we define \(P_t(A)\) by functional calculus, i.e., by substituting \(A\) for \(x\) in the polynomial \(P_t\).
\end{definition}
%\begin{remark}
%It should be noted that the above sequence polynomials are not orthogonal under MP distribution. 
%\end{remark}

The following generating-function identity gives the desired expansion of
\(x^{-1}\) in this basis. We defer its proof to the appendix.

\begin{claim}\label{clm:inverse-expansion}
For \(x \in [(1-\sqrt{\gamma})^2, (1+\sqrt{\gamma})^2] \) in the Marchenko--Pastur support and \(\gamma<1\),
\(
    \frac{1}{x}
    =
    \frac{1}{1-\gamma}
    \sum_{t\geq 0} (-1)^t \q_t(x).
\)
\end{claim}
%
%\begin{proof}
%Let
%\[
%    G(z,x)
%    :=
%    \sum_{t\geq 0} \q_t(x) z^t
%\]
%be the generating function of the polynomials. Using the recurrence above,
%one obtains
%\[
%    G(z,x)
%    =
%    \frac{1+\gamma z}
%    {1-\bigl(x-(1+\gamma)\bigr)z+\gamma z^2}.
%\]
%Evaluating at \(z=-1\) gives
%\[
%    \sum_{t\geq 0} (-1)^t \q_t(x)
%    =
%    G(-1,x)
%    =
%    \frac{1-\gamma}{x}.
%\]
%Rearranging proves the claim.
%\end{proof}

Formally, this gives us the inverse expansion
\[
    A^{-1}
    =
    \frac{1}{1-\gamma}
    \sum_{t\geq 0} (-1)^t \cdot  \q_t(A) \,. \numberthis \label{eq:poly-expansion-A}
\]
The remaining issue is that we have yet to rigorously justify that the spectrum of $A$ is indeed a Marchenko-Pastur distribution with parameter $\gamma$, and more importantly, up to the edge of the spectrum as well  
since
 this polynomial expansion is meaningful only on the
spectral interval where the Marchenko--Pastur approximation is valid. 
%Instead 
%We establish the following control on the spectrum of $A$, and we also establish $A$ follows MP$(\gamma)$ spectrum en route to controlling its spectral edge via the trace moment method.  

It should be noted that a similar result can also be deduced from the recent work of \cite{kogan2025extremal}, which establishes a more general theorem for kernel random matrices. Thus, our result may be viewed as a rediscovery of their theorem in the special case considered here. At a technical level, both proofs are based on the trace moment method, %while we also believe our proof is distinct from a usual, 
though our proof is different from the usual direct expansion of the expected value of the trace of high powers of the matrix.%high trace.

Instead, our proof for the spectrum of $A$ closely follows our intuitive connection between orthogonal polynomials and graph matrices of concatenated shapes. En route, we generalize the previous equivalence about Chebyshev polynomials of the second kind to the corresponding orthogonal polynomials for Marchenko--Pastur, and show that an analogous equivalence with graph matrix of concatenated shapes holds.%continues to hold.
\begin{restatable}[MP Polynomials and Shape Concatenation]{lemma}{mpPolynomialShapeConcatenation}
\label{lem:mp-polynomial-shape-concatenation}
Consider
\(
    A=\cm_\alpha+\cm_\beta+\left(1+\frac1d\right)I_m+M_D 
\) from the decomposition of $M$ in~\cref{eq:M-decomposition}.
For any   \(t \), we have
\[
    \q_t(A)
    \coloneqq
    \q_t^{(\gamma)}(A)
    =
    \fp_t(A)+o_d(1),
\]
where
\[
    \fp_t(A)
    \coloneqq
    \sum_{\substack{
        \tau=\tau_1\circ\cdots\circ\tau_t\\
        \tau_i\in\mathcal B(A)\ \forall i\in[t]
    }}
    \cm_\tau .
\]
Here \(\q_t^{(\gamma)}\), formally defined in~\cref{def:MP-orthogonal-poly}, denotes the \(t\)-th orthogonal polynomial associated
with the Marchenko--Pastur distribution of parameter \(\gamma\) , and
\(\fp_t(A)\) is the corresponding linear combination of graph matrices of
\(t\)-fold properly concatenated shapes. 

Crucially, the concatenation sum is
taken over the nontrivial graph-matrix shapes in \(\mathcal B(A) = \{\al,\beta \} \), and does not
include the trivial identity shape.
\end{restatable}

Diagrammatically, \(\q_t(A)\) is given by the sum over all \(t\)-fold proper
concatenations whose constituent shapes are \(\alpha\) or \(\beta\). We recall
these two basic shapes below.
\begin{figure}[ht!]
  \centering
      \begin{subfigure}[b]{0.38\textwidth}
        \includegraphics[width=\textwidth]{diagrams/M_alpha.pdf}
        \caption{\(\cm_{\alpha}\).}
%        \label{fig:M_alpha}
    \end{subfigure}
    \quad
    \begin{subfigure}[b]{0.32\textwidth}
        \includegraphics[width=\textwidth]{diagrams/M_beta.pdf}
        \caption{\(\cm_{\beta}\).}
%        \label{fig:M_beta}
    \end{subfigure}
    \captionsetup{width=.9\linewidth}
%    \caption{Graph-matrix representations of a \(d\times d\) \(\mathsf{GOE}\)
%    matrix with zero diagonal, and of the \(m\times m\) matrices \(M_\alpha\) and
%    \(M_\beta\) from \Cref{prop:M-decomposition}. Square vertices take labels in
%    \([m]\), and circle vertices take labels in \([d]\). The two ovals indicate the
%    left and right boundaries \(U_\tau,V_\tau\). If an edge \(e\) is not explicitly
%    labeled with an index, then \(t(e)=1\) by default.}
%    \label{fig:M_alpha_beta}
\end{figure}


%From the equivalence established above, it is then fairly straightforward to deduce the desired control of the spectrum. 
In particular, the graph-matrix and
orthogonal-polynomial equivalence viewpoint allows us to cleanly decouple the main term from the deviation terms in the trace power calculation, and thereby deduce the spectrum of $A$.
% \anote{While the Marchenko-Pastur distribution is somewhat similar to the semicircle distribution, they are diffierent because of the denominator}

\paragraph{Highlight of the Recurrence.}
The following surprising cancellation via graph matrices lies at the heart of our equivalence lemma between orthogonal polynomials and concatenated shapes. We highlight this equality at the basic level of $t=2$. Throughout this section, we will ignore the $o_d(1)$ term while our subsequent discussion shall also shed light on the source of such error terms. Recall that by the MP-recurrence from~\cref{def:MP-orthogonal-poly}, we have \[ 
\q_{2}(A) = (A - (1+\gamma) \cdot I_m )\cdot  \q_1(A) - \gam \cdot \q_0(A).
\]
%Plugging in our boundary condition of $\q_0(A) = I_d$, we want to show \[ 
%\q_2(A) = (A - (1+\gamma) \cdot I_d ) \cdot  \q_1(A) - \gamma \cdot  I_d  = \fp_2(A) +o_d(1) \,.
%\]
Our goal is to show \[ 
\q_2(A) = \fp_2(A) +o_d(1)
\]
where we emphasize that the RHS is the sum of the graph matrices of all four $2$-fold concatenations of $\{\al,\beta\}$, namely $\{\al\circ \al, \al\circ \beta, \beta\circ \al, \beta\circ\beta \}$, and the $o_d(1)$ subsumes negligible components involving $1/d\cdot  I_m$ and $M_D$. We ignore such negligible components throughout the overview.
  
 Recall that we have $A = \cm_\al +\cm_\beta + I_m +o_d(1) $, $\q_0(A) = I_m $ and $\q_1(A) = \fp_1(A) = \cm_\al+\cm_\beta +o_d(1)$, so the LHS simplifies to \[ 
\q_2(A) =  (\cm_\al +\cm_\beta -\gam \cdot I_m) \cdot (\cm_\al +\cm_\beta) - \gam \cdot I_m  \numberthis  \label{eq:q2(A)} \,. 
\]
We now highlight the following key matrix equalities of graph matrices,\begin{enumerate}
	\item \[  \cm_\al \cdot \cm_\al = \cm _{\al \circ \al} +\gam \cdot I_m  +\gam \cdot \cm_\al +o_d(1)  \label{eq:al-al} \numberthis \]
	\item \[\cm_\beta \cdot \cm_\beta =  \cm _{\beta\circ\beta} + \gam \cdot \cm_\beta +o_d(1) \numberthis  \label{eq:beta-beta} \]
	\item \[ \cm_\al \cdot \cm_\beta  =  \cm_{\al\circ \beta} +o_d(1), \qquad \text{and} \qquad  \cm _{\beta}\cdot \cm_\al = \cm_{\beta\circ\al}+o_d(1)\,.  \]
\end{enumerate}
%Observe that given these three equations, the desired equality  \[
%\q_2(A) =   \cm_{\al\circ \al} + \cm_{\al\circ\beta} +\cm_{\beta\circ\beta} + \cm_{\beta\circ\al}
% \]
% is fairly straightforward
% as any terms generated from $-\gam I$ on the LHS of~\cref{eq:q2(A)} exactly cancels with the corresponding $\gam$ terms from~\cref{eq:al-al} and~\cref{eq:beta-beta}. 
% 
 
 The desired equality \[
\q_2(A) =   \cm_{\al\circ \al} + \cm_{\al\circ\beta} +\cm_{\beta\circ\beta} + \cm_{\beta\circ\al}
 \]
 is then immediate given these three equations.
 Therefore, verifying the desired equivalence lemma boils down to showing \cref{eq:al-al} and~\cref{eq:beta-beta}, especially regarding the $\gam$-factors therein. In words, in the multiplication of $\cm_\al\cdot \cm_\al$, and of $\cm_\beta \cdot \cm_\beta$, in addition to the proper concatenation terms that are part of $\fp_2(A)$, there is an additional factor of \(
 \gam (\cm_\al +\cm_\beta+I)
 \) from the intersection terms!
 
To illustrate this, instead of first expanding the algebra of the matrix product, we begin with diagrams of the intersection patterns that drive the two identities. In the diagrams below, each red dotted edge denotes an intersection edge, meaning that the two incident vertices receive the same label. In all five cases, each green edge appears exactly twice: it is of type \(h_1\) in (a)--(d), and of type \(h_2\) in \cref{fig:half_flat}. In \cref{fig:full_diamond_1} and \cref{fig:full_diamond_2}, the same is also true additionally for the purple edges.
  \begin{figure}[H]
    \centering

    \begin{subfigure}[b]{0.4\textwidth}
        \includegraphics[width=\textwidth]{diagrams/full_diamond_1.pdf}
        \caption{Diamond Backtracking-1}
        \label{fig:full_diamond_1}
    \end{subfigure}
    \quad
    \begin{subfigure}[b]{0.4\textwidth}
        \includegraphics[width=0.94\textwidth]{diagrams/full_diamond_2.pdf}
        \caption{Diamond Backtracking-2}
        \label{fig:full_diamond_2}
    \end{subfigure}

    \vspace{1em}

    \begin{subfigure}[b]{0.4\textwidth}
        \includegraphics[width=0.94\textwidth]{diagrams/half_diamond.pdf}
        \caption{Half-Diamond Backtracking-1}
         \label{fig:half_diamond_1}

    \end{subfigure}
     \begin{subfigure}[b]{0.4\textwidth}
        \includegraphics[width=0.94\textwidth]{diagrams/half_diamond_2.pdf}
        \caption{Half-Diamond Backtracking-2}
         \label{fig:half_diamond_2}

    \end{subfigure}
    
	    \vspace{1em}
%
%	\cenring 
        \begin{subfigure}[b]{0.48\textwidth}
        \includegraphics[width=0.94\textwidth]{diagrams/half_flat.pdf}
        \caption{Half-Flat Backtracking}
         \label{fig:half_flat}

            \end{subfigure}
\label{fig:backtracking-int}
\caption{Well-behaved Backtracking Intersections}
    % \captionsetup{width=.9\linewidth}
    % \caption{Examples of graph matrices in the decomposition of $Q$.
    % Recall that square vertices take labels in $[m]$ and circle vertices take labels in $[d]$.
    % Unlabeled edges have Hermite index $1$.}
    % \label{fig:examples}
\end{figure}
Without giving the formal definition of each intersection as it may be transparent from the diagram illustrations alone already, we note that we have three ``important'' intersection patterns that we call \emph{backtracking intersections}. Specifically, they are classified with the specific $m\times m$ matrix defined in the following.
 \begin{enumerate}
	\item \textbf{Diamond Backtracking}:
	 both of the two inner circle vertices appear twice, as well as the square vertex on the boundary. In this case, we have a diagonal matrix
%	 \begin{align*}
%		 \cm_{\mathsf{DiamondInt}}[i,i] \coloneqq   \left(\sum_{a\neq b} \sum_{k\notin \{i,j\}} v_i[a]\cdot v_j[a]\cdot v_i[b]\cdot v_j[b] \right)^2 &= 2\cdot d\cdot (d-1)\cdot (m-1) \cdot \frac{1}{d^4} \\& =  \frac{2m}{d^2} =\gam \,;
%	\end{align*} 
%	 where we use each gaussian random variable has variance $\frac{1}{d}$ in our set-up and we ignore $o_d(1)$ factors. Crucially, the factor of $2$ comes from the fact that for any pairs of $a \neq b$ labels, we have two choice of intersection patterns as illustrated by \cref{fig:full_diamond_1} and \cref{fig:full_diamond_2)} respectively.
%	 
%	 Therefore,we have \[ \cm_{\mathsf{DiamondInt} }  =\gam \cdot I_d\,. \]	 
obtained by considering
\[
    (\cm_\alpha^2)[i,i]
    =
    \sum_{k\neq i}\cm_\alpha[i,k]\cm_\alpha[k,i],
\]
and restricting to the summands in which the two circle labels match as an
unordered pair.

More explicitly, define
\[
\begin{aligned}
\cm_{\mathsf{DiamondInt}}[i,i]
&:=
\sum_{k\neq i}
\sum_{a\neq b\in[d]}
v_i[a]v_i[b]v_k[a]v_k[b]
\sum_{\substack{a'\neq b'\in[d]\\ \{a',b'\}=\{a,b\}}}
v_k[a']v_k[b']v_i[a']v_i[b']
\\
&=
\sum_{a\neq b\in[d]}
v_i[a]^2v_i[b]^2
\left(
    \sum_{k\neq i}2v_k[a]^2v_k[b]^2
\right).
\end{aligned}
\]
 where the factor of \(2\) comes from the two possible orderings for $a',b'$ to match with given $a,b$ as a set, corresponding to the two full-diamond intersection patterns illustrated in \cref{fig:full_diamond_1} and \cref{fig:full_diamond_2}.
	 \footnote{This can be equivalently viewed as coming from $ \E \left(2\sum_{a<b} \sum_{k\notin \{i,j\}} v_i[a]\cdot v_j[a]\cdot v_i[b]\cdot v_j[b] \right)^2 = 4 \binom{d}{2}m/d^4.$ }
	 
	 Notice this is a scalar with mean \[ 
	 d(d-1) \cdot 2(m-1) \cdot \frac{1}{d^4} = 2m/d^2 = \gam 
	 \]
	  where we use that each Gaussian random variable has variance $\frac{1}{d}$ in our setup. This allows us to show that up to $o_d(1)$ error terms, we have \(
	  \cm_{\mathsf{DiamondInt}} = \gam \cdot  I_m \,.
	  \)
\item \textbf{Half-Diamond Backtracking}: 
%\item both of the two inner circle vertices appear twice, while the leftmost and rightmost square vertex appears once exactly. This is equivalent to looking at the entries of
%\begin{align*}
%	\cm_\al \cdot \cm_\al[i,j]& = \sum_{k\notin\{i,j\}} \cm_{\al}[i,k] \cdot \cm_{\al}[k,j]\\ &= \sum_{k\notin\{i,j\}}\sum_{s\neq t\in [d]}   v_i[s]v_i[t]v_k[s]v_k[t]\sum_{s'\neq t'\in [d]} v_k[s']v_k[t']v_j[s']v_j[t']
%\end{align*}
%and restrict our attention to $\{s,t\} = \{s',t'\}$. For any ordered tuple $(s,t)$, there are two choices to arrange $s',t'$ as the order of $s',t'$ does not affect the set equality above. Therefore, the corresponding graph matrix entries for half-diammond backtracking is given by 
In the half-diamond backtracking pattern, the two inner circle vertices are
identified across the two copies, while the left and right boundary square
vertices remain distinct. Equivalently, we consider the product
\[
    (\cm_\alpha^2)[i,j]
    =
    \sum_{k\notin\{i,j\}} \cm_\alpha[i,k]\cm_\alpha[k,j],
\]
and restrict to entries of $i\neq j\in[m]$ with the summands in which
\(
    \{s,t\}=\{s',t'\}.
\)
Expanding gives
\[
\begin{aligned}
\cm_{\mathsf{HalfDiamondInt}}[i,j]
&\coloneqq
\sum_{k\notin\{i,j\}}
\sum_{s\neq t\in[d]}
v_i[s]v_i[t]v_k[s]v_k[t]
\sum_{\substack{s'\neq t'\in[d]\\ \{s',t'\}=\{s,t\}}}
v_k[s']v_k[t']v_j[s']v_j[t']
\\
&=
		 \underbrace{ \sum_{s\neq t\in[d]} v_i[s]v_i[t] v_j[s]v_j[t] }_{=\cm_\al[i,j]}  \cdot  \underbrace{(\sum_{k\notin \{i,j\}} 2 v_k^2[s]v_k^2[t])}_{=  \gam} \end{aligned}
\]
%

% We have a matrix with non-zero entries defined at $i\neq j\in [m]$ as \begin{align*}
%		\cm_{\mathsf{HalfDiamondInt}}[i,j] &\coloneqq\sum_{k\notin\{i,j\}}\sum_{s\neq t\in [d]}   v_i[s]v_i[t]v_k[s]v_k[t]\sum_{\substack{ s'\neq t'\in [d] \\ \{s',t'\} = \{s,t \}}} v_k[s']v_k[t']v_j[s']v_j[t']\\&=
%		 \underbrace{ \sum_{s\neq t\in[d]} v_i[s]v_i[t] v_j[s]v_j[t] }_{\cm_\al[i,j]}  \cdot  \underbrace{(\sum_{k\notin \{i,j\}} 2 v_k^2[s]v_k^2[t])}_{=  \gam} \,.
%\end{align*} 
One key observation is that the  second term is tightly concentrated around its mean independent of the choice of $s\neq t$,\[
\E \sum_{k\notin \{i,j\}} 2 v_k^2[s]v_k^2[t]  = 2 \cdot (m-2) \cdot  (\frac{1}{d})^2 = \frac{2m}{d^2 } =\gam 
 \]
 when we consider any fixed $s\neq t\in [d]$. Therefore, with some algebra after centering with respect to the mean of the second term, we show that up to some error term with $o_d(1)$ spectral norm, we have\[
 \cm_{\mathsf{HalfDiamondInt}}  = \gam \cdot \cm_\al \,.
 \] 
 
 
\item \textbf{Half-Flat Backtracking}: the inner circle vertex of the two
\(M_\beta\)-gadgets is identified across the two copies, while the left and
right boundary square vertices remain distinct. Equivalently, we consider
\[
    (\cm_\beta^2)[i,j]
    =
    \sum_{k\notin\{i,j\}}\cm_\beta[i,k]\cm_\beta[k,j],
\]
and restrict to the summands in which the two circle labels agree.

Recall that
\[
    \cm_\beta[i,j]
    =
    \sum_{a\in[d]} h_2(v_i[a])h_2(v_j[a]).
\]
Thus the half-flat contribution has entries
\[
\begin{aligned}
\cm_{\mathsf{HalfFlatInt}}[i,j]
&:=
\sum_{k\notin\{i,j\}}
\sum_{a\in[d]}
h_2(v_i[a])h_2(v_k[a])
h_2(v_k[a])h_2(v_j[a])
\\
&=
\underbrace{ \sum_{a\in[d]}
h_2(v_i[a])h_2(v_j[a])}_{=\cm_\beta[i,j]}
\underbrace{ \left(
    \sum_{k\notin\{i,j\}} h_2(v_k[a])^2
\right)}_{= \gam} .
\end{aligned}
 \]
 Analogously to the previous two terms, the second term is tightly concentrated around
 \[ 
 (m-2) \cdot \frac{2}{d^2}
 \]
 where the factor $2$ instead comes from $\E[h_2(v_k[a])^2] = \frac{2}{d^2}$.
  \end{enumerate}

 Subsequently in~\cref{sec:MP-sec}, we give a formal analysis of the error terms that we have ignored so far, and generalize this cancellation pattern to higher degrees beyond $\q_2$ as well to establish the equivalence between orthogonal polynomials and graph matrices of concatenated shapes. 
 
 Furthermore, the equivalence allows us to deduce the spectrum 
 of $A$ up to the edge.
% In the half-diamond backtracking pattern, the two inner circle vertices are
%identified across the two copies, while the left and right boundary square
%vertices remain distinct. Equivalently, we consider the product
%\[
%    (\cm_\alpha^2)[i,j]
%    =
%    \sum_{k\notin\{i,j\}} \cm_\alpha[i,k]\cm_\alpha[k,j],
%\]
%and restrict to the summands in which
%\[
%    \{s,t\}=\{s',t'\}.
%\]
%Expanding gives
%\[
%\begin{aligned}
%\cm_{\mathsf{HalfDiamondInt}}[i,j]
%&:=
%\sum_{k\notin\{i,j\}}
%\sum_{s\neq t\in[d]}
%v_i[s]v_i[t]v_k[s]v_k[t]
%\sum_{\substack{s'\neq t'\in[d]\\ \{s',t'\}=\{s,t\}}}
%v_k[s']v_k[t']v_j[s']v_j[t']
%\\
%&=
%\sum_{s\neq t\in[d]}
%v_i[s]v_i[t]v_j[s]v_j[t]
%\left(
%\sum_{k\notin\{i,j\}}2v_k[s]^2v_k[t]^2
%\right).
%\end{aligned}
%\]
%For each fixed ordered pair \((s,t)\), the factor in parentheses has mean
%\[
%    \frac{2(m-2)}{d^2}
%    =
%    \gamma+O(d^{-2}).
%\]
%Throughout this section, we will ignore the $o_d(1)$ term while our subsequent discussion shall also shed light on the source of such error terms.


\begin{restatable}[Spectral radius of \(A\)]{lemma}{spectralradiusA} \label{lem:spectral-radius-A} 
With high probability, \[ \mathsf{spec}(A) \subseteq \bigl[(1-\sqrt{\gamma})^2,\,(1+\sqrt{\gamma})^2\bigr] \pm o_d(1), \] where \(\gamma = \frac{2m}{d^2}. \) \end{restatable}

We emphasize that our purpose in revisiting the graph-matrix proof is
not merely to reprove a result that can also be obtained by more direct moment
calculations. Rather, the matrix \(A\) provides a simpler setting in which to
illustrate the main ideas that will later be needed for the analysis of the
inner matrix \(Q\) arising from the iterative construction, where no prior result
seems to apply in a black-box manner.
% I% a useful blueprint for organizing the
%terms in \(Q\) and for understanding how polynomial evaluations correspond to
%concatenations of shapes.
%
 
% \subsection{Shifting From $A$ to $M$}

%At this point, we have justified the expansion for $A^{-1}$.
Finally, recall from~\cref{eq:M-decomposition} that $M$ can be viewed as a rank-$2$ update on $A$; existing results from the prior works allow us to switch from the inverse of $A$ to that of $M$ via the Woodbury identity.

%
%\paragraph{Inverting $M$ via Woodbury Identity} Recall that the $M$ can be decomposed into $A$ and a rank-2 matrix $B$ (\cref{eq:M-decomposition}), Woodbury identity then allows us to go from the inverse of $A$ to that of $M$. 
%\begin{lemma}[Inverse of $M$]
%	we have \[ 
%M^{-1} = A^{-1} + \frac{1}{s^2-ru}\cdot A^{-1} \cdot (u\cdot \frac{\1_m\1_m^T}{d} -s \frac{\eta\cdot \1_m^T+ \1_m^T \cdot \eta }{d}
% + r \cdot \frac{\eta \cdot \eta^T }{d}) \cdot A^{-1}\,. \numberthis \label{eq:M-inverse}
% \]
% where Let \(r,s,u\) be the scalars defined below\[
%    r\coloneqq \frac{1_m^T A^{-1} 1_m}{d};\]
%\[
%    s \coloneqq 1+ \frac{\eta^T A^{-1} 1_m}{d} ;
%\]
%and
%\[
%    u \coloneqq -1 + \frac{\eta^T A^{-1} \eta }{d} .\]
%\end{lemma}
\begin{restatable}[Explicit Inverse of \(M\)]{lemma}{Minverse} \label{lem:M-inverse} Let \[ r\coloneqq \frac{\1_m^T A^{-1}\1_m}{d}, \qquad s\coloneqq 1+\frac{\eta^T A^{-1}\1_m}{d}, \qquad u\coloneqq -1+\frac{\eta^T A^{-1}\eta}{d}. \] Then \begin{equation} \label{eq:M-inverse} M^{-1} = A^{-1} + \underbrace{ \frac{1}{s^2-ru}\, A^{-1} \left( u\,\frac{\1_m\1_m^T}{d} - s\,\frac{\eta\1_m^T+\1_m\eta^T}{d} + r\,\frac{ \eta\eta^T}{d} \right) A^{-1}}_{M_W \coloneqq}. \end{equation} \end{restatable}

Combining this with our polynomial expansion for $A$, we have now obtained an explicit formula for $M^{-1}$. More importantly, it admits a clean decomposition in the graph matrices, allowing us to ultimately write our inner matrix $Q$ via graph matrices.
\jnote{i moved basically all the technical discussion of Woodbury to the later section, but leave here the expression for M inverse. This looks fairly clean to me now.}

%\begin{fact}[Matrix Invertibility] \label{fact:invertibility}
%    Suppose $A \in \R^{n_1 \times n_1}$ and $C \in \R^{n_2 \times n_2}$ are both invertible matrices, and $U\in \R^{n_1 \times n_2}$ and $V \in \R^{n_2 \times n_1}$ are arbitrary.
%    Then, $A + U C V$ is invertible if and only if $C^{-1} + V A^{-1} U$ is invertible.
%\end{fact}
%
%\begin{fact}[Woodbury matrix identity~\cite{Woo59}]  \label{fact:woodbury}
%    Suppose $A \in \R^{n_1 \times n_1}$ and $C \in \R^{n_2 \times n_2}$ are both invertible matrices, and $U\in \R^{n_1 \times n_2}$ and $V \in \R^{n_2 \times n_1}$ are arbitrary. Then
%    \[ 
%    (A+ UCV)^{-1} = A^{-1} - A^{-1}U\left(C^{-1}+ VA^{-1}U\right)^{-1}VA^{-1}  \,.
%    \]
%\end{fact}
%
%In light of \Cref{fact:woodbury}, we can write $B$ as $B = UCU^T$ where $U = V^T = \frac{1}{\sqrt{d}} \begin{bmatrix} 1_m & \eta \end{bmatrix} \in \R^{m \times 2}$ and
%$C = \begin{bmatrix} 1 & 1 \\ 1 & 0 \end{bmatrix}$,
%and $M = A + UCU^T$.
%Note that $C^{-1} = \begin{bmatrix} 0 & 1 \\ 1 & -1 \end{bmatrix}$, and we have
%\begin{align*}
%    C^{-1} + U^T A^{-1}U  = \begin{bmatrix}
%	\frac{1_m^T A^{-1} 1_m}{d} & 1 + \frac{\eta^T A^{-1} 1_m}{d}\\\\
%	1+ \frac{\eta^T A^{-1} 1_m}{d} & -1 + \frac{\eta^T A^{-1} \eta }{d} 
%    \end{bmatrix}
%    \eqqcolon
%    \begin{bmatrix}
%    	r & s \\
%    	s & u
%    \end{bmatrix}
%    \,. 
%    \numberthis \label{eq:russ}
%\end{align*}
%
%Assuming $A$ is shown to be invertible, from \Cref{fact:invertibility}, we can prove that $M$ is invertible  by showing that the $2\times 2$ matrix $C^{-1} + U^T A^{-1}U$ is invertible, which is in fact equivalent to $ru - s^2 \neq 0$ from the above definitions.
%
%%\begin{claim}[$M$ is invertible]
%%	\[
%%	ru - s^2\enq 0
%%	 \]
%%\end{claim}
%%\todo{fill in the proof - should follow from this subsequent meta-lemma about the scalars }
%
%
%%\begin{lemma} \label{lem:hyper-parameter-bounds}
%%	We give the following bounds for the scalars defined above.\[ 
%%	 r= (1+o_d(1)) \cdot  \frac{m}{d} \cdot \frac{1}{1-\gamma} \,,
%%	\]
%%	\[
%%	s = 1+ o_d(1)\,,
%%	 \]
%%	 and \[
%%	 u = -1 +\gam +o_d(1) \,.
%%	  \]
%%	  \jnote{the extra $\gamma$ here is crucial}
%%\end{lemma}
%
%\begin{restatable}[Scalar bounds]{lemma}{HyperParameterBounds}
%\label{lem:hyper-parameter-bounds}
%Let \(r,s,u\) be the scalars defined above. Then, with high probability,
%\[
%    r\coloneqq \frac{1_m^T A^{-1} 1_m}{d}
%    =
%    (1+o_d(1))\cdot \frac{m}{d}\cdot \frac{1}{1-\gamma},
%\]
%\[
%    s \coloneqq 1+ \frac{\eta^T A^{-1} 1_m}{d}  = 1+o_d(1),
%\]
%and
%\[
%    u \coloneqq -1 + \frac{\eta^T A^{-1} \eta }{d}  = -1+\gamma+o_d(1).
%\]
%%In particular, the leading correction in \(u\) is the \(+\gamma\) term.
%\end{restatable}
%
%\begin{corollary}[$M$ is invertible] W.h.p. \[ 
%ru-s^2 \neq 0\,.
%\]	
%\end{corollary}
%%Given $M$ is also invertible, we may now proceed to expand $M^{-1}$ in the graph-matrix basis. 
%Finally, expanding out Woodbury, we have \[ 
%M^{-1} = A^{-1} + \frac{1}{s^2-ru}\cdot A^{-1} \cdot (u\cdot \frac{\1_m\1_m^T}{d} -s \frac{\eta\cdot \1_m^T+ \1_m^T \cdot \eta }{d}
% + r \cdot \frac{\eta \cdot \eta^T }{d}) \cdot A^{-1}\,. \numberthis \label{eq:M-inverse}
% \]
%Combining with our polynomial expansion for $A$, we have now obtained an explicit formula for $M^{-1}$. More importantly, it admits a clean decomposition in the graph matrices, allowing us to ultimately write our inner matrix $Q$ via graph matrices.

%\begin{proof}[Proof Sketch to \cref{lem:hyper-parameter-bounds}] We defer the formal proof to the appendix, while giving a sketch of that by highlighting the dominant terms. Notice they are all random variables with means given by the following.
%	\paragraph{Analysis for $r$}
%	Recall that $r = \frac{1}{d}  \cdot \1_m^T \cdot A^{-1} \cdot \1_m  $, and \(
%	A^{-1} = \frac{1}{1-\gamma} \sum_{t\geq 0} P_t(A)\,.
%	 \)	
%	 For $t=0$  we have \[
%	 \frac{1}{d} \cdot \1_m^TP_0(A) \1_m = \frac{m}{d}\,,
%	  \] 
%	  and \[ 
%	  \1_m^T P_t(A) \1_m = o_d(1) 
%	  \]
%%	  \todo{fill in their concentration here. - would be great to quantify the $o_d(1)$ slack for example.}
%	  Therefore, we have \[ 
%	  r = \frac{1}{1-\gamma}  \cdot \frac{\1_m^T A^{-1} \1_m }{d} = \frac{1}{1-\gamma} \cdot \frac{m}{d}+o_d(1)\,.
%	  \]
%	  
%	 \paragraph{Analysis for $s$}
%	 Recall that $s = 1+ \frac{1}{d} \cdot \eta^T A^{-1} \1_m$, it suffices for us to bound the second term as $o_d(1)$. Notice for $t=0$, we have \[ 
%	 \E[\eta^T  \cdot  \1_m]  = 0
%	 \]
%%	 \todo{fill in concentration here, and that for higher $t$.}
%	 
%	 \paragraph{Analysis for $u$} Recall that $u = -1 + \frac{1}{d} \eta^T A^{-1}\eta $. For $t=0$, we have \[
%\frac{1}{d} \cdot \eta^\top P_0(A) \eta  = (1+o_d(1)) \cdot \frac{2md }{d^2} = (1+o_d(1)) \gam \,.
%\]
% by analogous calculation for the higher moments. 
%   	
%  	For $t=1$, we have \[
%  	\q_1(A) = \fp_1(A) = \cm_\al + \cm_\beta \,.
%  	 \]
%  	 The crucial observation is $\eta^T \cm_\al \eta = o_d(1)$  while the second term gives $\gamma^2 $ (and importantly with a  minus sign from the alternating sign in the polynomial expansion for $x^{-1}$). 
%%	 \todo{This is done in claim 4.2- would be great to move things over and reorganize }
%	 \end{proof}
 

%\subsection{Visualizing Orthogonal Polynomials with Graph Matrices}
%\paragraph{Overview of Graph Matrices}
%Graph matrices have emerged as a versatile framework for analyzing structured random matrices arising in average-case problems. They have played a central role in establishing Sum-of-Squares lower bounds for a variety of average-case inference problems \cite{BHKKMP16, PR20, JPRTX, JPRX23, KPX24, NGCA24, Xu25, PX25, PotechinXu2026Theta}. Beyond the SoS literature, they have also been applied to the analysis of power-sum decompositions of polynomials \cite{BHKX22}, the ellipsoid fitting conjecture \cite{PTVW22, HKPX23}, and the study of first-order iterative algorithms, including belief propagation and approximate message passing \cite{JP24}.
%
%We begin with a brief overview of graph matrices specialized to the present setting. Throughout this section, we take the underlying random input to be Gaussian vectors \(A = {v_1,\ldots,v_m}\subseteq\mathbb{R}^d\) drawn independently from the distribution specified above. The framework is  more general and applies to other random inputs, such as random graphs, and we refer the reader to the references above for a comprehensive introduction and broader perspective.
%
%\begin{definition}[Shape] A shape $\tau$ is a tuple $(V(\tau), U_\tau, V_\tau, E(\tau))$ associated with a (multi) graph $(V(\tau), E(\tau))$.
%Each vertex in $V(\tau)$ is associated with a vertex-type that indicates the range of the labels for the particular vertex.
%Each edge $e\in E(\tau)$ is also associated with a Fourier index $t(e) \in \N$.
%Moreover, we have $U_\tau, V_\tau \subseteq V(\tau)$ as the left and right boundary of the shape.
%\end{definition}
%We remind the reader that $V_\tau$ should be distinguished from $V(\tau)$, where $V_\tau$ is the right boundary set, while $V(\tau)$ is the set of all vertices in the graph. 
%In the following example,
%\Cref{fig:M_alpha_beta} show the shapes for matrices $M_\al$ and $M_\beta$ defined in  that will be a main focus of our analysis.
%For these shapes, there are two vertex-types (square and circle).
%The two ovals in each shape indicate the left and right boundaries $U_\tau$ and $V_\tau$.
%
%We next describe how to associate a shape to a matrix (given the underlying matrix $G$).
%
%
%\begin{figure}[ht!]
%    \centering
%    \begin{subfigure}[b]{0.25\textwidth}
%        \includegraphics[width=\textwidth]{diagrams/GOE.pdf}
%        \caption{$\mathsf{GOE}$, zero diagonal.}
%        \label{fig:GOE}
%    \end{subfigure}
%    \quad
%    \begin{subfigure}[b]{0.32\textwidth}
%        \includegraphics[width=\textwidth]{diagrams/M_alpha.pdf}
%        \caption{$M_{\alpha}$.}
%        \label{fig:M_alpha}
%    \end{subfigure}
%    \quad
%    \begin{subfigure}[b]{0.32\textwidth}
%        \includegraphics[width=\textwidth]{diagrams/M_beta.pdf}
%        \caption{$M_{\beta}$.}
%        \label{fig:M_beta}
%    \end{subfigure}
%    \captionsetup{width=.9\linewidth}
%    \caption{Graph matrix representation of a $d\times d$ $\mathsf{GOE}$ matrix with zero diagonal, and the $m \times m$ matrices $M_{\alpha}$ and $M_{\beta}$ as defined in \Cref{prop:M-decomposition}.
%    Square vertices take labels in $[m]$ and circle vertices take labels in $[d]$.
%    The two ovals indicate the left and right boundaries $U_\tau, V_\tau$ of the shapes.
%    If an edge $e$ is not labeled with an index, then $t(e) = 1$ by default.}
%    \label{fig:M_alpha_beta}
%\end{figure}
%
%\begin{definition}[Graph matrix for shape] \label{def:graph-matrix-for-shape}
%    Given a shape $\tau$, we define its graphical matrix $M_{\tau}$ to be the matrix indexed by all possible boundary labelings of $S, T$, and for each of its entry, we define
%    \[ 
%    M_{\tau}[S, T] = \sum_{\substack{\sigma: V(\tau) \to \N \\ \sigma(U_\tau) = S,\ \sigma(V_\tau) = T }} \prod_{e\in E(\tau)}\chi_{t(e)}(G[\sigma(e)]) \,.
%    \]
%\end{definition}%Crucially, we emphasize that vertices outside the specially labeled set are indistinguishable with each other, and there is a bijection between the injective map $\sigma$ and the corresponding edge-set provided it satisfies the shape and boundary constraitns.
%%\end{definition}
%% \jnote{maybe move the following to an individual prelim section.}
%
%%Throughout this work, we will restrict our attention to $|S|=|T|=1$, i.e., matrices indexed with a single label (that may be either in $[d]$ or in $[m]$).
%
%
%Observe that for each entry $M_{\tau}[S,T]$, since $\sigma$ must map $U_\tau$ and $V_\tau$ to $S$ and $T$,
%$M_{\tau}[S,T]$ is simply a sum over labelings of the ``middle'' vertices $V(\tau) \setminus (U_\tau \cup V_\tau)$.
%Take \Cref{fig:M_alpha_beta} for example. Suppose $G\in \R^{m \times d}$ and square and circle vertices take labels in $[m]$ and $[d]$ respectively, then we can write out the entries of the matrix: for $i \neq j\in [m]$,
%\begin{align*}
%    M_{\al}[i,j] &= \sum_{a\neq b\in [d]} \chi_1(G[i,a]) \cdot \chi_1(G[i,b]) \cdot \chi_1(G[j,a]) \cdot \chi_1(G[j,b])  \,, \\
%    M_{\beta}[i,j] &= \sum_{a\in [d]} \chi_2(G[i,a]) \cdot \chi_2(G[j,a]) \,.
%\end{align*}
%
%\begin{remark}
%It is important to point out that the non-zero entries of $M_\al$ can be equivalently rewritten as
% \[  M_{\al}[i,j] = 2\sum_{a< b\in [d]} \chi_1(G[i,a]) \cdot \chi_1(G[i,b]) \cdot \chi_1(G[j,a]) \cdot \chi_1(G[j,b]) \,.\]
% The factor $2$ is crucial is our analysis.	\end{remark}
%
%Note also that since $\sigma$ must be injective for vertices of the same type and $U_\tau \neq V_\tau$ in both examples, there is no mapping such that $\sigma(U_\tau) = \sigma(V_\tau)$.
%Thus, by \Cref{def:graph-matrix-for-shape}, both matrices have zeros on the diagonal.
%
%Following the notation from \cite{HKPX23},  we specialize to the following setting:
%\begin{itemize}
%    \item $G \in \R^{m\times d}$ is a random Gaussian matrix whose rows are $v_1,\ldots, v_m \sim N(0, \frac{1}{d} I_d)$.
%    \item The Fourier characters $\{\chi_t\}_{t\in\N}$ are the (scaled) Hermite polynomials.
%    \item For all graph matrices that arise in our analysis,
%    \begin{itemize}
%        \item $|S| = |T| = 1$,
%        \item There are two vertex-types: square vertices take labels in $[m]$ and circle vertices take labels in $[d]$.
%    \end{itemize}
%\end{itemize}
%
%\subsection{Polish attempt:}
%\subsection{Graph Matrices:
% a Visualization of Orthogonal Polynomials} 
%\paragraph{Overview of Graph Matrices.}
%Graph matrices provide a versatile framework for analyzing structured random matrices
%arising in average-case problems. They have played a central role in proving
%Sum-of-Squares lower bounds for a wide range of average-case inference problems
%\cite{BHKKMP16, PR20, JPRTX, JPRX23, KPX24, NGCA24, Xu25, PX25,
%PotechinXu2026Theta}. Beyond the SoS literature, graph matrices have also been used
%to analyze power-sum decompositions of polynomials \cite{BHKX22}, the ellipsoid
%fitting conjecture \cite{PTVW22, HKPX23}, and first-order iterative algorithms,
%including belief propagation and approximate message passing \cite{JP24}.
%
%We begin with a brief overview of graph matrices, specialized to the setting of this
%paper. Throughout this section, the underlying random input is a collection of
%Gaussian vectors
%\[
%    G = \{v_1,\ldots,v_m\}\subseteq \mathbb R^d,
%\]
%drawn independently from the distribution specified above. Equivalently, we write
%\(G\in \mathbb R^{m\times d}\) for the matrix whose \(i\)-th row is \(v_i\).
%The graph-matrix framework applies much more generally, for instance to random
%graphs and even non-product distributions; we refer the reader to the references above
%for a broader introduction.
%
%\begin{definition}[Shape]
%A shape \(\tau\) is a tuple
%\(
%    \tau=(V(\tau),U_\tau,V_\tau,E(\tau))
%\)
%associated with a multigraph \((V(\tau),E(\tau))\). Each vertex in \(V(\tau)\) is
%assigned a vertex type, which determines the range from which its labels are drawn.
%Each edge \(e\in E(\tau)\) is assigned a Fourier index \(t(e)\in \mathbb N\).
%The subsets \(U_\tau,V_\tau\subseteq V(\tau)\) are called the left and right
%boundaries of the shape, respectively.
%\end{definition}
%
%We emphasize that \(V_\tau\) denotes the right boundary of the shape, whereas
%\(V(\tau)\) denotes the full vertex set of the underlying graph.
%
%\begin{definition}[Mapping of a shape]
%Given a shape \(\tau\), a mapping of \(\tau\) is a function
%\[
%    \sigma:V(\tau)\to \mathbb N
%\]
%such that:
%\begin{enumerate}
%    \item each vertex is assigned a label from the range specified by its vertex type;
%    \item \(\sigma\) is injective on vertices of the same type.
%\end{enumerate}
%\end{definition}
%
%In \Cref{fig:M_alpha_beta}, we illustrate the shapes corresponding to the matrices
%\(M_\alpha\) and \(M_\beta\) defined in \Cref{prop:M-decomposition}, which will be a
%central focus of our analysis. These shapes have two vertex types: square vertices
%take labels in \([m]\), while circle vertices take labels in \([d]\). The two ovals
%indicate the left and right boundaries \(U_\tau\) and \(V_\tau\).
%
%\begin{figure}[ht!]
%    \centering
%    \begin{subfigure}[b]{0.25\textwidth}
%        \includegraphics[width=\textwidth]{diagrams/GOE.pdf}
%        \caption{\(\mathsf{GOE}\), zero diagonal.}
%        \label{fig:GOE}
%    \end{subfigure}
%    \quad
%    \begin{subfigure}[b]{0.32\textwidth}
%        \includegraphics[width=\textwidth]{diagrams/M_alpha.pdf}
%        \caption{\(M_{\alpha}\).}
%        \label{fig:M_alpha}
%    \end{subfigure}
%    \quad
%    \begin{subfigure}[b]{0.32\textwidth}
%        \includegraphics[width=\textwidth]{diagrams/M_beta.pdf}
%        \caption{\(M_{\beta}\).}
%        \label{fig:M_beta}
%    \end{subfigure}
%    \captionsetup{width=.9\linewidth}
%    \caption{Graph-matrix representations of a \(d\times d\) \(\mathsf{GOE}\)
%    matrix with zero diagonal, and of the \(m\times m\) matrices \(M_\alpha\) and
%    \(M_\beta\) from \Cref{prop:M-decomposition}. Square vertices take labels in
%    \([m]\), and circle vertices take labels in \([d]\). The two ovals indicate the
%    left and right boundaries \(U_\tau,V_\tau\). If an edge \(e\) is not explicitly
%    labeled with an index, then \(t(e)=1\) by default.}
%    \label{fig:M_alpha_beta}
%\end{figure}
%
%We now describe how a shape gives rise to a matrix once the underlying random input
%\(G\) is fixed.
%
%\begin{definition}[Graph matrix of a shape]
%\label{def:graph-matrix-for-shape}
%Given a shape \(\tau\), its graph matrix \(M_\tau\) is indexed by boundary labelings.
%For boundary labelings \(S\) of \(U_\tau\) and \(T\) of \(V_\tau\), define
%\[
%    M_{\tau}[S,T]
%    =
%    \sum_{\substack{
%        \sigma:\,V(\tau)\to \mathbb N \\
%        \sigma \text{ is a mapping of } \tau \\
%        \sigma(U_\tau)=S,\ \sigma(V_\tau)=T
%    }}
%    \prod_{e\in E(\tau)}
%    \chi_{t(e)}\bigl(G[\sigma(e)]\bigr).
%\]
%Here, if \(e=\{x,y\}\) is an edge between a square vertex and a circle vertex, then
%\(G[\sigma(e)]\) denotes the corresponding entry \(G[\sigma(x),\sigma(y)]\).
%\end{definition}
%
%For each entry \(M_\tau[S,T]\), the boundary labels are fixed by the constraints
%\(\sigma(U_\tau)=S\) and \(\sigma(V_\tau)=T\). Thus the entry is a sum over labelings
%of the internal vertices
%\(
%    V(\tau)\setminus (U_\tau\cup V_\tau).
%\)
%For example, consider the shapes in \Cref{fig:M_alpha_beta}. Suppose
%\(G\in \mathbb R^{m\times d}\), with square vertices labeled by elements of \([m]\)
%and circle vertices labeled by elements of \([d]\). Then for \(i\neq j\in[m]\),
%\begin{align*}
%    M_{\alpha}[i,j]
%    &=
%    \sum_{a\neq b\in[d]}
%    \chi_1(G[i,a])\chi_1(G[i,b])
%    \chi_1(G[j,a])\chi_1(G[j,b]), \\
%    M_{\beta}[i,j]
%    &=
%    \sum_{a\in[d]}
%    \chi_2(G[i,a])\chi_2(G[j,a]).
%\end{align*}
%
%\begin{remark}
%The nonzero entries of \(M_\alpha\) may equivalently be written as
%\[
%    M_{\alpha}[i,j]
%    =
%    2\sum_{a<b\in[d]}
%    \chi_1(G[i,a])\chi_1(G[i,b])
%    \chi_1(G[j,a])\chi_1(G[j,b]).
%\]
%The factor of \(2\) is important in our analysis.
%\end{remark}
%
%Since the mapping \(\sigma\) is injective on vertices of the same type, and since
%\(U_\tau\neq V_\tau\) in both examples in \Cref{fig:M_alpha_beta}, there is no valid
%mapping with \(\sigma(U_\tau)=\sigma(V_\tau)\). Consequently, both \(M_\alpha\) and
%\(M_\beta\) have zero diagonal.
%
%Following the notation of \cite{HKPX23}, we specialize the graph-matrix framework to
%the following setting:
%\begin{itemize}
%    \item \(G\in\mathbb R^{m\times d}\) is a random Gaussian matrix whose rows are
%    \(v_1,\ldots,v_m\sim N(0,\frac1d I_d)\).
%    \item The Fourier characters \(\{\chi_t\}_{t\in\mathbb N}\) are the appropriately
%    scaled Hermite polynomials.
%    \item For all graph matrices appearing in our analysis:
%    \begin{itemize}
%        \item the left and right boundary labelings satisfy \(|S|=|T|=1\);
%        \item there are two vertex types: square vertices take labels in \([m]\), and
%        circle vertices take labels in \([d]\).
%    \end{itemize}
%\end{itemize}
%
%
%\subsection{Main Technical Lemmas}
%
%\begin{restatable}[Variance of inner matrix \(Q\)]{lemma}{varianceInnerQ} \label{lem:variance-inner-Q} We have \[ \lim_i \Var(Q_i) = 1, \] where \[ \Var(Q_i) \coloneqq \frac{1}{d}\,\E\!\left[\Tr(Q_i Q_i^\top)\right]. \] \end{restatable} 
%
%\begin{restatable}[Spectral radius of \(Q\)]{theorem}{spectralRadiusQ} \label{thm:spectral-radius-Q} With high probability, \[ \mathsf{spec}(Q) \subseteq [-2,2] \pm o_d(1). \] \end{restatable}
%
%\begin{restatable}[
%    (Semicircle) Orthogonal Polynomials and Shape Concatenation
%]{lemma}{semicircleShapeConcatenation}
%\label{lem:semicircle-shape-concatenation}
%Consider a linear combination of backbone-dangling shapes from a ground set $\calB(\cdot )$,
%\[
%    M=\sum_{\tau\in\mathcal B(Q)} c_\tau \cdot \cm_\tau
%\]
%%Let
%%\(\{P^{\mathsf{sc}}_t\}_{t\ge 0}\) denote the orthogonal polynomials for the
%%semicircle law on \([-2,2]\), normalized so that
%%\[
%%    P^{\mathsf{sc}}_0(x)=1,\qquad
%%    P^{\mathsf{sc}}_1(x)=x,\qquad
%%    P^{\mathsf{sc}}_{t+1}(x)
%%    =
%%    xP^{\mathsf{sc}}_t(x)-P^{\mathsf{sc}}_{t-1}(x).
%%\]
%%Equivalently, \(P^{\mathsf{sc}}_t(x)=U_t(x/2)\), where \(U_t\) is the
%%Chebyshev polynomial of the second kind.
%Then, for any $t>0 $,
%\[
%    P^{\mathsf{sc}}_t(M)
%    =
%    \sum_{\substack{
%        \tau=\tau_1\circ\cdots\circ\tau_t\\
%        \tau_i\in\mathcal B(M)\ \forall i\in[t]
%    }}
%    \left(\prod_{i=1}^{t} c_{\tau_i}\right)\cm_\tau
%    + o_d(1),
%\]
%where the sum is over proper \(t\)-fold concatenations of shapes from
%\(\mathcal B(M)\), and the error is in spectral norm, provided \[ 
%\Var(M) = 1 \,.
%\]
%
%% For \(t=0\), the empty
%%concatenation is interpreted as the identity matrix.
%%When the dependence on the semicircle law is clear, we abbreviate
%%\(P_t=P^{\mathsf{sc}}_t\,.\)
%\end{restatable}
%
%
%An analogous statement holds for the matrix \(A\) arising from the correction
%for linear constraints. In this case, the relevant polynomial basis is the
%family of orthogonal polynomials associated with the Marchenko--Pastur law on
%\[
%    \bigl[(1-\sqrt{\gamma})^2,\,(1+\sqrt{\gamma})^2\bigr].
%\]
% To distinguish from the orthogonal polynomials for $\mu_{sc}$, we use \(\q_t\) to denote the orthogonal
%polynomials for Marchenko-Pastur from \cref{def:MP-basis}. Analogous to our result for Chebyshev polynomials, the orthogonal polynomials for Marchenko-Pastur also admits a clean decomposition in graph matrix via concatenated shapes.
%
%\begin{lemma}[MP Polynomials and Shape Concatenation]
%Consider
%\[
%    A=\cm_\alpha+\cm_\beta+\left(1+\frac1d\right)I_d + M_D\,,
%\]  for any $t>0$,
%\[
%    \q_t(A) \coloneqq \q^{(\gam)}_t(A) 
%    =  \underbrace{    \sum_{\substack{
%        \tau=\tau_1\circ\cdots\circ\tau_t\\
%        \tau_i\in\mathcal B(A)\ \forall i\in[t]
%    }} 
%    \cm_\tau}_{\fp_t(A) \coloneqq }
%    +o_n(1).
%\]
%where the LHS is the orthogonal polynomial evaluated on $A$ while the RHS is a linear combination of graph matrix of concatenated shapes that crucially do not include the trivial shape which corresponds to the identity matrix.
%%Here \(\{q_t\}_{t\ge0} = \{q^{\mathsf{MP}}_t \}\) denotes the orthogonal polynomials associated with the
%%Marchenko--Pastur distribution of parameter \(\gamma\) from~\cref{def: MP-basis}.
%\end{lemma}
%Thus, unlike the one-shot identity-perturbation
%construction, where one only needs to understand \(M^{-1}\eta\), here one must
%understand the action of \(M^{-1}\) on a family of residual vectors generated
%by the spectral map \(Q_t\mapsto F(C_FQ_t)\).  The central technical object is
%therefore the correction operator
%\[
%    \mathcal C := L^*M^{-1}.
%\]
%A useful measure of the size of the \(t\)-th correction is
%\[
%    \|w^{(t)}\|_2^2
%    =
%    \|M^{-1}r_t\|_2^2
%    =
%    r_t^\top M^{-2}r_t.
%\]
%Controlling the iteration therefore requires a precise expansion of \(M^{-1}\)
%and norm bounds for the associated matrix-valued correction
%\[
%    L^*M^{-1}r_t.
%\]
\subsection{The Rise of Backbone-Dangling Shapes: Iterations Made Concrete}
We now zoom in on the combinatorial structure of the shapes that arise in the iteratively constructed matrix \(Q\). We begin by recalling the graph-matrix decomposition used in prior work, which naturally leads to a distinguished family of graph matrices, called \emph{backbone-dangling shapes}. We follow the naming from prior work, while substantially generalizing the definition to accommodate the shapes generated by our iterative construction.
%For starters, we focus on the initialization $Q_0$ and revisit the construction in \cite{PTVW22,HKPX23}.

%\subsection{Graph Matrix Decomposition}
\paragraph{Recap of Identity Perturbation: Decomposition at the Base Level}
We start by recapping the decomposition into graph matrices as presented in the previous work that exploits graph-matrix decomposition \cite{PTVW22, HKPX23}. By and large, we follow the notation of \cite{HKPX23} most closely, while also emphasizing the distinctions from the prior work that appear in our decomposition.

%\todo{i copy pasted the following from HKPX23' - plz help rephrase}
%By writing $A^{-1}$ as a truncated series (\Cref{def:truncated-A-inverse}) and the decomposition of $\calR$ (\Cref{prop:R-decomposition}), we may view $\calR_1$ and $\calR_2$ as linear combinations of $d \times d$ matrices of the forms
%\begin{align*}
%    \sum_{i\in[m]}  \left(Q_1 Q_2 \cdots Q_k \eta\right)[i] \cdot v_i v_i^T
%    \quad \text{and} \quad 
%    \sum_{i\in[m]}  \left(Q_1 Q_2 \cdots Q_k 1_m\right)[i] \cdot v_i v_i^T
%    \numberthis \label{eq:R-terms}
%\end{align*}
%respectively, where $Q_1,\dots,Q_k \in \{M_\al, M_\beta, M_D, \frac{1}{d} I_m\}$ are $m \times m$ matrices.
Let's recall that \[ 
Q_0 = -C_F \cdot  \mathcal{L}^*(M^{-1}\eta)= - C_F \cdot \sum_{t\in [m]} w_0[t] \cdot v_i \cdot v_i^\top
\]
with \( w_0 = M^{-1} \cdot \eta  \) and \(\eta = \eta_0 =   \calL(I_d) -\1_m\) denoting the affine deviation of the identity matrix.



The focus of \cite{PTVW22, HKPX23} boils down to studying the following collection of shapes called \emph{dangling shapes}. These are the matrices in the decomposition of the base matrix $Q_0$ that appear in the off-diagonal. The expansion also contains nonzero diagonal terms; however, their contribution is negligible, so we suppress them throughout the discussion. These matrices can be viewed as arising from a $3$-way concatenation of a backbone-path, a dangling $A^{-1}$-path, and a final $h_2$ deviation attachment:
    \begin{enumerate}
        \item \textbf{Backbone Path:} the off-diagonal of component $\sum_{t\in [m]} v_tv_t^\top $ is represented by the shape shown in \cref{fig:R-base} which can be viewed as a backbone path; this shape serves as
the base gadget for the other matrices appearing in the inner matrix $Q$.
%        \item For diagonal terms, we have a double edge (that shall be distinguished from an $h_2$ edge as the double edge here stands for $v_i[a]^2$) connected to a middle square vertex that receives a label in $[m]$.
	\item \textbf{Dangling $A^{-1}$ Path} 
    For each term, we specify a length-$k\geq 0$ dangling path that starts from the middle square vertex such that
    \begin{itemize}
        \item for any $k>0$, each step comes from one of the following gadgets in $\{ M_\al, M_\beta, M_D, \frac{1}{d}I_m\}$;
        \item The dangling path is a (possibly improper) concatenation of shapes along the path.    \end{itemize}
    \item \textbf{Final $h_2$ Deviation Attachment:} Since $w_0 = M^{-1} \eta$, each $A^{-1}$ term is hit with an additional $\eta$ factor; we therefore attach an $h_2$ gadget (\cref{fig:gadgets}) to the end of the dangling path.
    We call this the ``final $h_2$ attachment gadget''.

    \end{enumerate}
See \cref{fig:R-A} for an example. It is important to point out that we make two simplifications here: 
\begin{enumerate}
	\item These are the restrictions of $Q_0$ to proper concatenations ignoring possible intersections from the $3$-way vertical concatenation. Controlling all such intersection patterns is a nontrivial issue and necessitates the lengthy charging arguments developed in prior work. For the moment, we set this complication aside. %Momentarily, we set this complication aside, and 
    We will return to it when choosing the correction terms.
	\item Moreover, we use the proxy of $M^{-1}\approx A^{-1}$ by ignoring the extra rank-$2$ component. This is less of a concern, as it can be straightforwardly shown to be lower-order terms, or to give a constant multiple of another $A^{-1}$ term. \anote{This sentence is a bit confusing} \jnote{ok to drop - i wanted to say the rank-2 component for $Q_0$ is not by itself negligible but it gives an $A^{-1}$ term as well.}
\end{enumerate}

  \jnote{add pointer}.

There is, however, one aspect of this analysis that simplifies substantially in our setting. In the prior analyses of \cite{PTVW22,HKPX23}, the expansion of \(A^{-1}\) is expressed through ordinary powers of \(A-I\). Consequently, the dangling path representing \(A^{-1}\) may contain numerous collisions among the gadgets placed along the path. By contrast, our expansion in orthogonal polynomials naturally produces proper, and hence injective, concatenations of shapes. This eliminates most of the collision patterns that arise from the inverse expansion itself.


%\paragraph{Step 1: Decomposing $A^{-1}$ in Graph Matrices}
%To expand \(Q_0\) in the graph-matrix basis, we make the following observations. The expansion also contains nonzero diagonal terms; however, their contribution is negligible, so we suppress them throughout the discussion.
%\begin{enumerate}
%\item \textbf{Backbone Path in the Off-Diagonal:} the off-diagonal of component $\sum_{t\in [m]} v_tv_t^\top $ is represented by the shape shown in \cref{fig:R-base} which can be viewed as a backbone path; this shape serves as
%the base gadget for the other matrices appearing in $Q$
%	\item \textbf{Expressing $w_0$}: The equivalence between the Chebyshev expansion of \(A^{-1}\) and
%the graph-matrix expansion in terms of concatenated shapes in
%\(\q_t(A)\) already gives a convenient diagrammatic representation of
%\(A^{-1}\) via proper concatenations of $\al$ and $\beta$ gadgets;
%\item \textbf{Deviation Attachment at the End:} The $A^{-1}$-decomposition can be viewed as attaching a dangling path from the middle square vertex in the backbone path. In the case of $Q_0$, $\eta$ is attached to the end of $A^{-1}$ path as well.
%\end{enumerate}
%Each matrix in $Q_0$ can be represented by attaching dangling-path (\Cref{fig:gadgets}) to the square vertex in the middle.
%%
%%
%%First, the equivalence between the Chebyshev expansion of \(A^{-1}\) and
%%the graph-matrix expansion in terms of concatenated shapes in
%%\(\q_t(A)\) already gives a convenient diagrammatic representation of
%%\(A^{-1}\) via proper concatenations of $\al$ and $\beta$ gadgets. 
%%
%%Next, we combine this representation with the base component of
%%\(Q\). In particular, the off-diagonal part of the matrix
%%\(
%%    \sum_{i=1}^m v_i v_i^\top
%%\)
%%is represented by the shape shown in \cref{fig:R-base}; this shape serves as
%%the base gadget for the other matrices appearing in $Q$
%%
%Moreover, since $\eta_0[j] = \|v_j\|_2^2 - 1 = \sum_{a=1}^d h_2(v_j[a])$ for $j\in [m]$ , each shape has an extra $h_2$ attached at the end.
%\Cref{fig:R-A} shows an example of such a shape.


\begin{figure}[ht]
    \centering
    \begin{subfigure}[b]{0.28\textwidth}
        \includegraphics[width=\textwidth]{diagrams/path.pdf}
        \caption{Off-diagonal part of $\sum_{i=1}^m v_i v_i^T$.}
        \label{fig:R-base}
    \end{subfigure}
    \quad
    \begin{subfigure}[b]{0.32\textwidth}
        \includegraphics[width=0.94\textwidth]{diagrams/gadgets.pdf}
        \caption{Left: $M_{\alpha}$ gadget. Middle: $M_{\beta}$ gadget. Right: final $h_2$ gadget.}
        \label{fig:gadgets}
    \end{subfigure}
    \quad
    \begin{subfigure}[b]{0.28\textwidth}
        \includegraphics[width=0.94\textwidth]{diagrams/dangling.pdf}
        \caption{Off-diagonal part of $\sum_{i=1}^m (M_\al M_\beta \eta)[i] \cdot v_i v_i^T$.}
        \label{fig:R-A}
    \end{subfigure}
    % \quad
    % \begin{subfigure}[b]{0.35\textwidth}
    %     \includegraphics[width=\textwidth]{Figures/dangling_r.pdf}
    %     \caption{Example of a dangling shape in $R_s$. Dotted line: a jump across $\frac{1}{d}\eta\eta^T$.}
    %     \label{fig:R-s}
    % \end{subfigure}
    % \quad
    \captionsetup{width=.9\linewidth}
    \caption{Examples of graph matrices in the decomposition of $Q$ adapted from \cite{HKPX23}.
    Recall that square vertices take labels in $[m]$ and circle vertices take labels in $[d]$.
    Unlabeled edges have Hermite index $1$.}
    \label{fig:examples}
\end{figure}

%In the following, we recall the definition of these shapes, which were coined
%\emph{dangling shapes} in prior work, and adapt them to our setting.
%\begin{definition}[Gadgets]
%We call each shape in $\{M_\al, M_\beta, M_D, \frac{1}{d}I_m\}$ a \emph{gadget} in the dangling shape.
%\end{definition}


%
%The focus of \cite{PTVW22, HKPX23} boils down to studying the following collection of shapes coined \emph{backbone-dangling shapes}. These are the matrices in the decomposition of the base matrix $Q_0$ split into diagonal and off-diagonal terms. 
%    \begin{enumerate}
%        \item For the off-diagonal terms, we start with a path from $U$ to $V$ (each containing a circle vertex receiving labels in $[d]$) passing through a middle square vertex that receives a label in $[m]$.
%        \item For diagonal terms, we have a double edge (that shall be distinguished from an $h_2$ edge as the double edge here stands for $v_i[a]^2$) connected to a middle square vertex that receives a label in $[m]$.
%    \end{enumerate}
%
%
%    For each term, we specify a length-$k\geq 0$ dangling path that starts from the middle square vertex such that
%    \begin{itemize}
%        \item for any $k>0$, each step comes from one of the following gadgets in $\{ M_\al, M_\beta, M_D, \frac{1}{d}I_m\}$;
%        \item The dangling path is a (possibly improper) concatenation of shapes along the path.    \end{itemize}
%    For matrices with an additional $\eta$ factor, we attach an $h_2$ gadget (\cref{fig:gadgets}) to the end of the dangling path .
%    We call this the ``final $h_2$ attachment gadget''.
    % At the end of the path, a coordinate $j\in [m]$ has already been specified,
    % \begin{itemize}
    %     \item (Final $H_2$ attachment) we additionally pick up $\eta_j$: we have one vertex choice $a\in [d]$ and an $H_2$ edge $v_j[a]^2-\frac{1}{d}$ connected to it. This appears with a scaling constant of $\frac{r+s}{s^2-ru}$.
    %     \item ($1_m$ term) we pick up a factor $1$ with the scaling constant $-\frac{u+s}{s^2-ru}$.
    % \end{itemize}
%That said, for correction terms iteratively constructed, shapes may have multiple $A^{-1}$ paths, and collision between vertices among different paths, as well as those on the main backbone path, may still happen. This is captured by the following definition.
%
%\begin{definition}[$A^{-1}$-Path Intersection]
%	For a shape $\tau$ with multiple $A^{-1}$ paths, each viewed as a proper concatenation, we call vertex intersection between gadgets on different paths an $A^{=1}$-path intersection.
%\end{definition}

\paragraph{Dangling Path Once Again: Vertical Concatenation} 

 We now turn to the matrices that arise from the correction terms in the
iterative construction. For concreteness, we describe the first iteration; the
same decomposition applies to later iterations essentially verbatim. Again we focus on terms that arise from proper concatenation while deferring the discussion on intersection terms.

%Instead of reasoning with the entire $H_{0}$ term, %altogether, 
%let's consider $ H_{0}$ in its 
Consider the graph matrix decomposition of $H_{0}$ restricted to the properly concatenated terms that are on the off-diagonal, 
\[ 
\mathsf{Proper} H_{0} = \sum_{\substack{\tau \in \calB(H_{0})\\ \text{proper off-diagonal}}} c(\tau) \cdot \cm_\tau \,,
\]
and consider a fixed term $\tau \in \calB(H_{0})$. We ignore the diagonal terms as they are negligible. Modulo its coefficient $c(\tau)$, its assigned correction term in the iterative process is given by \[ 
\corr(\tau) \coloneqq \frac{1}{1-\gamma} \left(- \calL^*M^{-1}\calL(\cm_\tau)  + \gam \cdot \cm_\tau  \right)\,.
\]

Throughout this section, we focus solely on the first term of $  \calL^*M^{-1}\calL(\cm_\tau)$, and most importantly, the proper shapes therein: the remaining term of $\gam \cdot H_{0}$ is picked precisely to cancel with certain intersection terms that also arise from the sum of rank-$1$ terms. We defer the discussion to~\cref{sec:cooking-correction}.

\paragraph{Upgrade at the Terminal Attachment}
 We call such an operation \emph{vertical concatenation} from a graph matrix point of view: it is again a $3$-way concatenation of the \emph{backbone-path}, \emph{dangling $A^{-1}$-path} and \emph{deviation attachment} from the prior works except the final deviation attachment is no longer an $h_2$-gadget. Instead, the
terminal square vertex may be attached to a newly generated shape
\(\tau \in H_{0} \) and in general from the most recently generated higher-degree term $H_i$ in the Chebyshev expansion, with the two endpoints of \(\tau\)
connected to the same square vertex via two $h_1$ edges. We call the final attachment a \emph{deviation attachment}, and call the two edges connecting boundary vertices of $\tau$ to the square vertex \emph{attachment edges}. Correspondingly, we also call the terminal square vertex a \emph{deviation attachment vertex}.
 
 We showcase the process of vertical concatenation from a given shape $\tau \in H_{0}$ from $(a)$ to $(c)$ in~\cref{fig:vertical-concate}.


%\begin{enumerate}
%	\item This can be viewed as starting with the familiar backbone path given by $\sum_{t} v_iv_i^\top$;
%	\item The particular coefficient $w_1$ is obtained from concatenating $M^{-1}$ (and thereby $A^{-1}$) as a dangling path from the backbone-path given by $\q_t(A)$ for any $t\ge 0$;
%	\item The end of the dangling path is further concatenated with $\eta_1$. However, to contrast with the base case in \cref{fig:R-A}, the attachment is no longer an $h_2$-gadget.
%\end{enumerate}
\begin{figure}[h]
    \centering
    \begin{subfigure}[b]{0.4\textwidth}
        \includegraphics[width=\textwidth]{diagrams/p2_Q.pdf}
        \caption{Shape $\tau \in P_2(Q_0)$}
        \label{fig:p2(Q)}
    \end{subfigure}
    \quad
    \begin{subfigure}[b]{0.4\textwidth}
        \includegraphics[width=0.94\textwidth]{diagrams/eta_P2_Q.pdf}
        \caption{Deviation  w.r.t. $ \cm_\tau $}
        \label{fig:eta(p2(Q))}
    \end{subfigure}

    \vspace{1em}

    \begin{subfigure}[b]{0.48\textwidth}
        \includegraphics[width=0.94\textwidth]{diagrams/inner_1.pdf}
        \caption{Backbone-dangling shape: $\corr(\tau) \in Q_1\setminus Q_0$}
        \label{fig:corr(tau)}
    \end{subfigure}
    \hfill
    \begin{subfigure}[b]{0.48\textwidth}
        \includegraphics[width=0.94\textwidth]{diagrams/inner_2.pdf}
        \caption{Backbone-dangling shape in $ Q_1\setminus Q_0$}
        \label{fig:corr(tau)-2}
      
%        \label{fig:corr(tau)-2} 
    \end{subfigure}
%	    \vspace{1em}
	\caption{Vertical Concatenation and Backbone-Dangling Shapes}
	\label{fig:vertical-concate}
%	\cenring 
%        \begin{subfigure}[b]{0.48\textwidth}
%        \includegraphics[width=0.94\textwidth]{diagrams/old-3.pdf}
%        \caption{Backbone-dangling shape in $ Q_1\setminus Q_0$}
%            \end{subfigure}
%
    % \captionsetup{width=.9\linewidth}
    % \caption{Examples of graph matrices in the decomposition of $Q$.
    % Recall that square vertices take labels in $[m]$ and circle vertices take labels in $[d]$.
    % Unlabeled edges have Hermite index $1$.}
    % \label{fig:examples}
\end{figure}

%	For intuition, it suffices for us to ignore the low-rank component in $M^{-1}$ and view the dangling path as an $A^{-1}$ path, thus we may focus on $M^{-1}\approx A^{-1}$ for most of the analysis. The only place where such distinction warrants extra care  is the analysis for the variance of the base case $Q_0$. 

% In particular, we focus on the proper shape arising from these terms, as the remaining term of $\gam \cdot H_{Q_0}$ is picked to cancel with certain intersection terms that also arise from the sum of rank-$1$ terms. We will defer the discussion, and it suffices for us to focus on the proper shapes in the concatenation through this section.

%Thus the correction matrix is obtained by applying \(M^{-1}\) to the residual
%vector \(\eta\), and then lifting the resulting coefficients back to matrix
%space via the map \(w\mapsto \sum_t w_t v_t v_t^\top\). 
%
%On the one hand,  the factor \(v_t v_t^\top\) gives the familiar base component, while the coefficient
%\((M^{-1}\eta)_t\) contributes once again a dangling path attached at the square vertex
%\(t\). Expanding \(M^{-1}\) through the MP-polynomial expansion of \(A^{-1}\)
%and the rank-two Woodbury correction therefore yields a decomposition into
%base gadgets with dangling concatenated paths, together with the error terms
%from the orthogonal polynomials and graph matrix equivalence.

%Let's now zoom into the distinction from the base case- the final attachment at the end of $A^{-1}$ path: the terminal attachment is no longer necessarily the plain
%\(h_2\)-attachment that appears in the base case of \(Q_0\). Instead, the
%terminal square vertex may be attached to a newly generated shape
%\(\tau \in H_{0} \) and in general from the most recently generated higher-degree term $H_i$ in Chebyshev expansion, with the two endpoints of \(\tau\)
%connected to the same square vertex via two $h_1$ edges. We call the final attachment a \emph{deviation attachment}, and call the two edges connecting boundary vertices of $\tau$ to the square vertex \emph{attachment edges}. Correspondingly, we also call the terminal square vertex a \emph{deviation attachment vertex}.

Finally, as we apply such construction iteratively across the iteration levels, this gives rise to a recursive terminal attachment process defined as follows.
\begin{definition}[Recursive Terminal Attachment]
Let \(s\) be a square vertex. A recursive terminal attachment at \(s\) is
generated by the following rules.

\begin{enumerate}
    \item \textbf{Base case: \(h_2\)-attachment.}
    The attachment may be a single \(h_2\)-edge from \(s\) to a circle vertex
    \(a\):
    \[
        s \xleftrightarrow{h_2} a .
    \]

    \item \textbf{Concatenated-path attachment.}
    For any \(t > 1\), let \(\tau_1,\ldots,\tau_t\) be backbone-dangling
    shapes with compatible circle boundaries, and let    \[
        \tau =\tau_t \circ \tau_{t-1}\circ \cdots \circ \tau_1
    \]
    denote their proper concatenation. $\tau$ is then attached by adding two $h_1$ edges connecting $s$ to $U_\tau$ and $V_\tau$ (assuming $U_\tau\neq V_\tau$).
%    
%    If \(U_\tau \) and \(V_\calR\) are the left
%    and right circle boundary vertices of \(\calR\), then \(\calR\) may be attached
%    to \(s\) by adding two \(h_1\)-edges connecting \(s\) to \(U_\calR\) and
%    \(V_\calR\). That is, we add
%    \[
%        U_\calR \xleftrightarrow{h_1} s \xleftrightarrow{h_1} V_\calR .
%    \]
%    When \(U_\calR=V_\calR\), we use the diagonal convention and attach \(s\) to
%    this single circle vertex by two parallel \(h_1\)-edges. \anote{I don't think we want to allow this as if $U_\calR=V_\calR$ then we either don't have a proper concatenation or $t = 0$.}
\end{enumerate}
\end{definition}

Notice we also single out the off-diagonal shapes appearing in the subsequent levels in \(Q\), since these are
the main objects of our analysis---the diagonal terms will be treated separately
as small-norm error terms and bounded at the end.

 Finally, we wrap this section up with the formal definition of \emph{backbone-dangling} shapes as follows: in summary, %crux,
 these are the shapes that appear in the recursive application of the terminal-attachment process if the entire resulting shape is vertex injective.
\begin{definition}[Backbone-Dangling Shape]\label{def:backbone-dangling-shape}
A shape is called a backbone-dangling shape if it satisfies the following
properties.

\begin{enumerate}
    \item It contains an off-diagonal backbone path from a circle vertex
    \(U_\tau\) to another circle vertex \(V_\tau\), passing through exactly \emph{one} square
    vertex via two \(h_1\)-edges:
    \[
        U_\tau \xleftrightarrow{h_1} s \xleftrightarrow{h_1} V_\tau .
    \]

    \item Attached to this square vertex \(s\) is a dangling path arising from
    the \(A^{-1}\)-expansion.

    \item At the end of the dangling path, the terminal square vertex is
    equipped with a recursive terminal attachment in the sense of the preceding
    definition;
    \item It is a \emph{proper} shape, i.e., there is no vertex intersection.
\end{enumerate}
\end{definition}

Intuitively, the shapes we study can be viewed as generalizing the backbone shapes in~\cite{HKPX23} both vertically and horizontally. 
%
%Before we give examples, it should be noted that there is indeed a discrepancy when we restrict our attention to the proper shapes at the end of terminal attachment. That said, we show that unless we are in the base case of $\eta_0$, each non-trivial intersection is negligible, and the only dominant term is indeed those from proper concatenation throughout (both horizontally when we extend to $P_j(\cdot)$ as well as vertically when we apply terminal attachment).
%
%There are two main distinctions in the shapes we study:
%\begin{enumerate}
%	\item Our backbone-dangling shape is a generalization from those in \cite{HKPX23}, in that we may have multiple vertical levels (corresponding to levels in the iterative process). In contrast,  \cite{HKPX23} only considers the depth-$1$ structure.
%	\item For each vertical level, we additionally have horizontal extension corresponding to concatenation from Chebyshev polynomials. In contrast, \cite{HKPX23} only studies the width-$1$ variant.
%\end{enumerate}



%\paragraph{Example of Backbone-Dangling Shapes} 
%We now give examples of backbone-dangling shapes to give an intuitive illustration. For convenience, we color the edges in our illustration to identify where each edge comes from. For starters, we showcase a shape from $P_2(Q)$ in \cref{fig:p2(Q)}: in particular, we use orange edges to highlight the edges part of the backbone path, and blue edges to highlight those from $M^{-1}$ (and more precisely, $A^{-1}$). The final red edges correspond the final $h_2$-attatchment for each dangling path. \Cref{fig:eta(p2(Q))} then illustrates the corresponding deviation vector, and finally, \cref{fig:corr(tau)} showcases the corresponding backbone-dangling shape to be added into the inner matrix for the subsequent iteration.
%Intuitively, the shapes we study can be viewed as generalizing the backbone shapes in~\cite{HKPX23} both vertically and horizontally. We give several examples of backbone-dangling shapes to build intuition. For convenience, the edges in the illustrations are colored according to their origin. We begin in \cref{fig:p2(Q)} with a shape arising from \(P_2(Q)\): the orange edges indicate the backbone path, while the blue edges indicate the contribution from \(M^{-1}\), more precisely from \(A^{-1}\). The final red edges correspond to the terminal \(h_2\)-attachment on each dangling path. Next, \cref{fig:eta(p2(Q))} illustrates the associated deviation vector. Evaluating $M^{-1}$ at the corresponding deviation is equivalent to adding an $M^{-1}$ path on top. Finally, \cref{fig:corr(tau)} shows the corresponding backbone-dangling shape that is added to the inner matrix in the next iteration.

%\begin{figure}[ht]
%    \centering
%    \begin{subfigure}[b]{0.4\textwidth}
%        \includegraphics[width=\textwidth]{diagrams/p2(Q).pdf}
%        \caption{Shape $\tau \in  P_2(Q_0)$}
%        \label{fig:p2(Q)}
%    \end{subfigure}
%    \quad
%    \begin{subfigure}[b]{0.4\textwidth}
%        \includegraphics[width=0.94\textwidth]{diagrams/eta(P2(Q)).pdf}
%        \caption{Deviation w.r.t. $\tau$: $\eta^\top \cm_\tau \eta  $ }
%        \label{fig:eta(p2(Q))}
%    \end{subfigure}
%    \quad
%    \begin{subfigure}[b]{0.5\textwidth}
%        \includegraphics[width=0.94\textwidth]{diagrams/inner_1.pdf}
%        \caption{Backbone-Dangling Shape: $\corr(\tau) \in Q_1\setminus Q_0$}
%        \label{fig:corr(tau)}
%    \end{subfigure}
%    \quad
%    \begin{subfigure}[b]{0.5\textwidth}
%        \includegraphics[width=0.94\textwidth]{diagrams/inner_2.pdf}
%        \caption{Backbone-Dangling Shape: $\corr(\tau) \in Q_1\setminus Q_0$}
%        \label{fig:corr(tau)}
%    \end{subfigure}
%
%    % \quad
%    % \begin{subfigure}[b]{0.35\textwidth}
%    %     \includegraphics[width=\textwidth]{Figures/dangling_r.pdf}
%    %     \caption{Example of a dangling shape in $R_s$. Dotted line: a jump across $\frac{1}{d}\eta\eta^T$.}
%    %     \label{fig:R-s}
%    % \end{subfigure}
%    % \quad
%%    \captionsetup{width=.9\linewidth}
%%   \caption{Examples of graph matrices in the decomposition of $Q$.
%%    Recall that square vertices take labels in $[m]$ and circle vertices take labels in $[d]$.
%%    Unlabeled edges have Hermite index $1$.}
%%    \label{fig:examples}
%\end{figure}


%\subsubsection{Why are Backbone-Dangling Shapes Freely Independent?}
%Next, with the combinatorial description of backbone dangling shapes in place, we can proceed to elaborate on why we expect their graph matrices to be freely independent.
%First, we recall that each backbone-dangling shape can be viewed as consisting of two components: the backbone-path corresponding to the $\sum_{t} v_t v_t^\top$ component, and the dangling path attached to the square vertex  corresponding to $M^{-1}\eta $.  
%
%There are two main factors governing the free independence. Consider the trace calculations of $Q$ and in particular the dominants terms in $\E[\Tr(QQ^\top)^q] $, we can again view the walks as having two components as well: backbone and dangling components. On a high level, it suffices for us to discuss each component individually:
%\begin{enumerate}
%	\item Restricted to the backbone path (from the component $\sum_{t\in [m]} v_tv_t^\top ) $, we observe that the dominant term requires each labeled square vertex (from $[m]$) to appear twice exactly. This shall be contrasted from circle vertex that do not have such constraints. This discrepancy from the stark contrast between the combinatorial factor of the underlying square and circle vertices: each square vertex's label contributes a cost of $m\approx d^2 \gg d$ where $d$ is the contribution factor of a circle vertex. Therefore, in the dominant walks, we would like to avoid extra repetition of the square vertex where the twice-appearance is the minimum appearance requirement due to mean-$0$ random variables involved.
%	\item 
%	For readers familiar with the separator language for graph matrix norm bounds, this is equivalent to   
%\end{enumerate}

\subsection{Semicircular Spectrum of the Inner Matrix \texorpdfstring{$Q$}{Q}}
With a combinatorial characterization of the summands in $Q$, we have two final questions to settle in order to establish that the spectrum of $Q$ follows a semicircle distribution in $[-2,2]$. In turn, this will ultimately ensure that our choice of Chebyshev polynomials of the second kind for the polynomial expansion is valid. Concretely, we show the following:
  \begin{enumerate}
	\item \textbf{Semicircular Spectrum}: each constituent matrix in $Q$ is freely independent of each other;
	\item \textbf{Variance Normalization}: the variance of the inner matrix is $1$ (in the limit).
\end{enumerate}

Towards identifying the spectral property of the inner matrix, we set out to understand the combinatorial property of each shape arising in the expansion of $Q$. 

%\subsubsection{The Rise of Backbone-Dangling Shapes}
%We now zoom into the combinatorial characterizations of shapes in $Q$ constructed iteratively from the above process. We will start by recapping the construction and analysis from prior works via graph matrices that point to a specific collection of graph matrices - coined \emph{backbone-dangling shapes}. We follow the previous notation while extensively generalize its definition.
\begin{comment}
\subsubsection{The Rise of Backbone-Dangling Shapes}

We now zoom in on the combinatorial structure of the shapes that arise in the iteratively constructed matrix \(Q\). We begin by recalling the graph-matrix decomposition used in prior work, which naturally leads to a distinguished family of graph matrices, called \emph{backbone-dangling shapes}. We follow the naming from prior work, while substantially generalizing the definition to accommodate the shapes generated by our iterative construction.
%For starters, we focus on the initialization $Q_0$ and revisit the construction in \cite{PTVW22,HKPX23}.

%\subsection{Graph Matrix Decomposition}
\paragraph{Recap of Identity Perturbation: Decomposition at the Base Level}
We start by recapping the decomposition into graph matrices as presented in the previous work that exploits graph-matrix decomposition \cite{PTVW22, HKPX23}. By and large, we follow the notation of \cite{HKPX23} most closely, while also emphasizing the distinctions from the prior work that appear in our decomposition.

%\todo{i copy pasted the following from HKPX23' - plz help rephrase}
%By writing $A^{-1}$ as a truncated series (\Cref{def:truncated-A-inverse}) and the decomposition of $\calR$ (\Cref{prop:R-decomposition}), we may view $\calR_1$ and $\calR_2$ as linear combinations of $d \times d$ matrices of the forms
%\begin{align*}
%    \sum_{i\in[m]}  \left(Q_1 Q_2 \cdots Q_k \eta\right)[i] \cdot v_i v_i^T
%    \quad \text{and} \quad 
%    \sum_{i\in[m]}  \left(Q_1 Q_2 \cdots Q_k 1_m\right)[i] \cdot v_i v_i^T
%    \numberthis \label{eq:R-terms}
%\end{align*}
%respectively, where $Q_1,\dots,Q_k \in \{M_\al, M_\beta, M_D, \frac{1}{d} I_m\}$ are $m \times m$ matrices.
Let's recall that \[ 
Q_0 = C_F \mathcal{L}^*(M^{-1}\eta) =  C_F \cdot \sum_{t\in [m]} w_0[t] \cdot v_i \cdot v_i^\top
\]
with \( w_0 = M^{-1} \cdot \eta  \). We remind the reader that $\eta= \eta_0$ denotes the original affine deviation of the identity matrix.
%\paragraph{Step 1: Decomposing $A^{-1}$ in Graph Matrices}

First, the equivalence between the Chebyshev expansion of \(A^{-1}\) and
the graph-matrix expansion in terms of concatenated shapes in
\(\q_t(A)\) already gives a convenient diagrammatic representation of
\(A^{-1}\) via proper concatenations of $\al$ and $\beta$ gadgets. 

Next, we combine this representation with the base component of
\(Q\). In particular, the off-diagonal part of the matrix
\(
    \sum_{i=1}^m v_i v_i^\top
\)
is represented by the shape shown in \cref{fig:R-base}; this shape serves as
the base gadget for the other matrices appearing in $Q$
Each matrix in $Q_0$ can be represented by attaching ``gadgets'' (\Cref{fig:gadgets}) to the square vertex in the middle.
For the case of $Q_0$, since $\eta_0[j] = \|v_j\|_2^2-1 = \sum_{a=1}^d h_2(v_j[a])$ for $j\in [m]$ , each shape has an extra $h_2$ attached at the end.
\Cref{fig:R-A} shows an example of such a shape.


\begin{figure}[ht]
    \centering
    \begin{subfigure}[b]{0.28\textwidth}
        \includegraphics[width=\textwidth]{diagrams/path.pdf}
        \caption{Off-diagonal part of $\sum_{i=1}^m v_i v_i^T$.}
        \label{fig:R-base}
    \end{subfigure}
    \quad
    \begin{subfigure}[b]{0.32\textwidth}
        \includegraphics[width=0.94\textwidth]{diagrams/gadgets.pdf}
        \caption{Left: $M_{\alpha}$ gadget. Middle: $M_{\beta}$ gadget. Right: final $h_2$ gadget.}
        \label{fig:gadgets}
    \end{subfigure}
    \quad
    \begin{subfigure}[b]{0.28\textwidth}
        \includegraphics[width=0.94\textwidth]{diagrams/dangling.pdf}
        \caption{Off-diagonal part of $\sum_{i=1}^m (M_\al M_\beta \eta)[i] \cdot v_i v_i^T$.}
        \label{fig:R-A}
    \end{subfigure}
    % \quad
    % \begin{subfigure}[b]{0.35\textwidth}
    %     \includegraphics[width=\textwidth]{Figures/dangling_r.pdf}
    %     \caption{Example of a dangling shape in $R_s$. Dotted line: a jump across $\frac{1}{d}\eta\eta^T$.}
    %     \label{fig:R-s}
    % \end{subfigure}
    % \quad
    \captionsetup{width=.9\linewidth}
    \caption{Examples of graph matrices in the decomposition of $Q$ adapted from \cite{HKPX23}.
    Recall that square vertices take labels in $[m]$ and circle vertices take labels in $[d]$.
    Unlabeled edges have Hermite index $1$.}
    \label{fig:examples}
\end{figure}

In the following, we recall the definition of these shapes, which were coined
\emph{dangling shapes} in prior work, and adapt them to our setting.
\begin{definition}[Gadgets]
We call each shape in $\{M_\al, M_\beta, M_D, \frac{1}{d}I_m\}$ a \emph{gadget} in the dangling shape.
\end{definition}



The focus of \cite{PTVW22, HKPX23} boils down to studying the following collection of shapes - called dangling shapes. In our terms,  these are the matrices in the decomposition of the base matrix $Q_0$  into diagonal and off-diagonal terms: 
    \begin{enumerate}
        \item For the off-diagonal terms, we start with a path from $U$ to $V$ (each containing a circle vertex receiving labels in $[d]$) passing through a middle square vertex that receives a label in $[m]$.
        \item For diagonal terms, we have a double edge (that shall be distinguished from an $h_2$ edge as the double edge here stands for $v_i[a]^2$) connected to a middle square vertex that receives a label in $[m]$.
    \end{enumerate}


    For each term, we specify a length-$k\geq 0$ dangling path that starts from the middle square vertex such that
    \begin{itemize}
        \item for any $k>0$, each step comes from one of the following gadgets in $\{ M_\al, M_\beta, M_D, \frac{1}{d}I_m\}$;
        \item The dangling path is a (possibly improper) concatenation of shapes along the path.    \end{itemize}
    For matrices with an additional $\eta$ factor, we attach an $h_2$ gadget (\cref{fig:gadgets}) to the end of the dangling path .
    We call this the ``final $h_2$ attachment gadget''.
    % At the end of the path, a coordinate $j\in [m]$ has already been specified,
    % \begin{itemize}
    %     \item (Final $H_2$ attachment) we additionally pick up $\eta_j$: we have one vertex choice $a\in [d]$ and an $H_2$ edge $v_j[a]^2-\frac{1}{d}$ connected to it. This appears with a scaling constant of $\frac{r+s}{s^2-ru}$.
    %     \item ($1_m$ term) we pick up a factor $1$ with the scaling constant $-\frac{u+s}{s^2-ru}$.
    % \end{itemize}
In the prior analyses of \cite{PTVW22, HKPX23}, the expansion of \(A^{-1}\)
proceeds through powers of \(A-I\). As a result, the dangling path
representing \(A^{-1}\) may contain many potential collisions among the
gadgets appearing along the path. In our setting, this difficulty is largely
avoided: after expanding \(A^{-1}\) in the MP polynomial basis, the dangling
path is represented by a proper concatenation of shapes. Crucially, the backtracking intersection terms are canceled out and the remaining intersection terms are separated out and absorbed into the error terms
controlled by the equivalence lemma. 

%That said, for correction terms iteratively constructed, shapes may have multiple $A^{-1}$ paths, and collision between vertices among different paths, as well as those on the main backbone path, may still happen. This is captured by the following definition.
%
%\begin{definition}[$A^{-1}$-Path Intersection]
%	For a shape $\tau$ with multiple $A^{-1}$ paths, each viewed as a proper concatenation, we call vertex intersection between gadgets on different paths an $A^{=1}$-path intersection.
%\end{definition}

\paragraph{Dangling Path Once Again: Decomposition for Iterative Constructions}
 We now turn to the matrices that arise from the correction terms in the
iterative construction. For concreteness, we describe the first iteration; the
same decomposition applies to later iterations essentially verbatim.

The correction matrix has the form
\(
    \sum_{t\in[m]} w_t\, v_t v_t^\top,
\)
where
\(
    w=M^{-1}\eta.
\)
Here \(\eta\in\mathbb R^m\) is the vector of constraint violations generated
by the newly introduced terms in the spectral expansion of \(Q\). More
precisely, for each \(i\in[m]\),
\[
    \eta[i]
    =
    v_i^\top
    \left(
        \sum_{j\ge2}
        \sum_{\tau\in \textsf{New-}\calP_j(Q)}
            c(\tau)\,\cm_\tau
    \right)
    v_i .
\]
Thus the correction matrix is obtained by applying \(M^{-1}\) to the residual
vector \(\eta\), and then lifting the resulting coefficients back to matrix
space via the map \(w\mapsto \sum_t w_t v_t v_t^\top\).

On the one hand,  the factor \(v_t v_t^\top\) gives the familiar base component, while the coefficient
\((M^{-1}\eta)_t\) contributes once again a dangling path attached at the square vertex
\(t\). Expanding \(M^{-1}\) through the MP-polynomial expansion of \(A^{-1}\)
and the rank-two Woodbury correction therefore yields a decomposition into
base gadgets with dangling concatenated paths, together with the error terms
from the orthogonal polynomials and graph matrix equivalence.

On the other hand, the terminal attachment is no longer necessarily the plain
\(h_2\)-attachment that appears in the base case of \(Q_0\). Instead, the
terminal square vertex may be attached to a newly generated shape
\(\tau \in \textsf{New-}\mathcal P_j(Q)\), with the two endpoints of \(\tau\)
connected to the same square vertex.

\begin{definition}[Recursive Terminal Attachment]
Let \(s\) be a square vertex. A recursive terminal attachment at \(s\) is
generated by the following rules.

\begin{enumerate}
    \item \textbf{Base case: \(h_2\)-attachment.}
    The attachment may be a single \(h_2\)-edge from \(s\) to a circle vertex
    \(a\):
    \[
        s \xrightarrow{h_2} a .
    \]

    \item \textbf{Concatenated-path attachment.}
    For any \(t \ge 1\), let \(\tau_1,\ldots,\tau_t\) be backbone-dangling
    shapes with compatible circle boundaries, and let
    \[
        \calR=\tau_t \circ \tau_{t-1}\circ \cdots \circ \tau_1
    \]
    denote their proper concatenation. If \(U_\calR\) and \(V_calR\) are the left
    and right circle boundary vertices of \(\calR\), then \(\calR\) may be attached
    to \(s\) by adding two \(h_1\)-edges connecting \(s\) to \(U_\calR\) and
    \(V_\calR\). That is, we add
    \[
        U_\calR \xrightarrow{h_1} s \xrightarrow{h_1} V_\calR .
    \]
    When \(U_\calR=V_\calR\), we use the diagonal convention and attach \(s\) to
    this single circle vertex by two parallel \(h_1\)-edges. \anote{I don't think we want to allow this as if $U_\calR=V_\calR$ then we either don't have a proper concatenation or $t = 0$.}
\end{enumerate}
\end{definition}

We now single out the off-diagonal shapes appearing in \(Q\), since these are
the main objects of our analysis. The diagonal terms will be treated separately
as small-norm error terms and bounded at the end. Finally, we wrap up this section with the formal definition of \emph{backbone-dangling} shapes as follows: in crux, these are the shapes that appear in the recursive application of terminal-attachment process if the entire resulting shape is vertex injective.
\begin{definition}[Backbone-Dangling Shape]\label{def:backbone-dangling-shape}
A shape is called a backbone-dangling shape if it satisfies the following
properties.

\begin{enumerate}
    \item It contains an off-diagonal backbone path from a circle vertex
    \(U_\tau\) to another circle vertex \(V_\tau\), passing through exactly \emph{one} square
    vertex via two \(h_1\)-edges:
    \[
        U_\tau \xrightarrow{h_1} s \xrightarrow{h_1} V_\tau .
    \]

    \item Attached to this square vertex \(s\) is a dangling path arising from
    the \(A^{-1}\)-expansion.

    \item At the end of the dangling path, the terminal square vertex is
    equipped with a recursive terminal attachment in the sense of the preceding
    definition;
    \item It is a \emph{proper} shape.
\end{enumerate}
\end{definition}

Before we give examples, it should be noted that there is indeed a discrepancy when we restrict our attention to the proper shapes at the end of terminal attachment. That said, we show that unless we are in the base case of $\eta_0$, each non-trivial intersection is negligible, and the only dominant term is indeed those from proper concatenation throughout (both horizontally when we extend to $P_j(\cdot)$ as well as vertically when we apply terminal attachment).
\jnote{in construction}

There are two main distinctions in the shapes we study:
\begin{enumerate}
	\item Our backbone-dangling shape is a generalization from those in \cite{HKPX23}, in that we may have multiple vertical levels (corresponding to levels in the iterative process). In contrast,  \cite{HKPX23} only considers the depth-$1$ structure.
	\item For each vertical level, we additionally have horizontal extension corresponding to concatenation from Chebyshev polynomials. In contrast, \cite{HKPX23} only studies the width-$1$ variant.
\end{enumerate}



\paragraph{Example of Backbone-Dangling Shapes} 
%We now give examples of backbone-dangling shapes to give an intuitive illustration. For convenience, we color the edges in our illustration to identify where each edge comes from. For starters, we showcase a shape from $P_2(Q)$ in \cref{fig:p2(Q)}: in particular, we use orange edges to highlight the edges part of the backbone path, and blue edges to highlight those from $M^{-1}$ (and more precisely, $A^{-1}$). The final red edges correspond the final $h_2$-attatchment for each dangling path. \Cref{fig:eta(p2(Q))} then illustrates the corresponding deviation vector, and finally, \cref{fig:corr(tau)} showcases the corresponding backbone-dangling shape to be added into the inner matrix for the subsequent iteration.

We now give several examples of backbone-dangling shapes to build intuition. For convenience, the edges in the illustrations are colored according to their origin. We begin in \cref{fig:p2(Q)} with a shape arising from \(P_2(Q)\): the orange edges indicate the backbone path, while the blue edges indicate the contribution from \(M^{-1}\), more precisely from \(A^{-1}\). The final red edges correspond to the terminal \(h_2\)-attachment on each dangling path. Next, \cref{fig:eta(p2(Q))} illustrates the associated deviation vector. Evaluating $M^{-1}$ at the corresponding deviation is equivalent to adding an $M^{-1}$ path on top. Finally, \cref{fig:corr(tau)} shows the corresponding backbone-dangling shape that is added to the inner matrix in the next iteration.

%\begin{figure}[ht]
%    \centering
%    \begin{subfigure}[b]{0.4\textwidth}
%        \includegraphics[width=\textwidth]{diagrams/p2(Q).pdf}
%        \caption{Shape $\tau \in  P_2(Q_0)$}
%        \label{fig:p2(Q)}
%    \end{subfigure}
%    \quad
%    \begin{subfigure}[b]{0.4\textwidth}
%        \includegraphics[width=0.94\textwidth]{diagrams/eta(P2(Q)).pdf}
%        \caption{Deviation w.r.t. $\tau$: $\eta^\top \cm_\tau \eta  $ }
%        \label{fig:eta(p2(Q))}
%    \end{subfigure}
%    \quad
%    \begin{subfigure}[b]{0.5\textwidth}
%        \includegraphics[width=0.94\textwidth]{diagrams/inner_1.pdf}
%        \caption{Backbone-Dangling Shape: $\corr(\tau) \in Q_1\setminus Q_0$}
%        \label{fig:corr(tau)}
%    \end{subfigure}
%    \quad
%    \begin{subfigure}[b]{0.5\textwidth}
%        \includegraphics[width=0.94\textwidth]{diagrams/inner_2.pdf}
%        \caption{Backbone-Dangling Shape: $\corr(\tau) \in Q_1\setminus Q_0$}
%        \label{fig:corr(tau)}
%    \end{subfigure}
%
%    % \quad
%    % \begin{subfigure}[b]{0.35\textwidth}
%    %     \includegraphics[width=\textwidth]{Figures/dangling_r.pdf}
%    %     \caption{Example of a dangling shape in $R_s$. Dotted line: a jump across $\frac{1}{d}\eta\eta^T$.}
%    %     \label{fig:R-s}
%    % \end{subfigure}
%    % \quad
%%    \captionsetup{width=.9\linewidth}
%%   \caption{Examples of graph matrices in the decomposition of $Q$.
%%    Recall that square vertices take labels in $[m]$ and circle vertices take labels in $[d]$.
%%    Unlabeled edges have Hermite index $1$.}
%%    \label{fig:examples}
%\end{figure}

\begin{figure}[ht]
    \centering

    \begin{subfigure}[b]{0.4\textwidth}
        \includegraphics[width=\textwidth]{diagrams/p2(Q).pdf}
        \caption{Shape $\tau \in P_2(Q_0)$}
        \label{fig:p2(Q)}
    \end{subfigure}
    \quad
    \begin{subfigure}[b]{0.4\textwidth}
        \includegraphics[width=0.94\textwidth]{diagrams/eta(P2(Q)).pdf}
        \caption{Deviation  w.r.t. $ \cm_\tau $}
        \label{fig:eta(p2(Q))}
    \end{subfigure}

    \vspace{1em}

    \begin{subfigure}[b]{0.48\textwidth}
        \includegraphics[width=0.94\textwidth]{diagrams/inner_1.pdf}
        \caption{Backbone-dangling shape: $\corr(\tau) \in Q_1\setminus Q_0$}
        \label{fig:corr(tau)}
    \end{subfigure}
    \hfill
    \begin{subfigure}[b]{0.48\textwidth}
        \includegraphics[width=0.94\textwidth]{diagrams/inner_2.pdf}
        \caption{Backbone-dangling shape in $ Q_1\setminus Q_0$}
        \label{fig:corr(tau)-2}
      
%        \label{fig:corr(tau)-2} 
    \end{subfigure}
%	    \vspace{1em}
%
%	\cenring 
%        \begin{subfigure}[b]{0.48\textwidth}
%        \includegraphics[width=0.94\textwidth]{diagrams/old-3.pdf}
%        \caption{Backbone-dangling shape in $ Q_1\setminus Q_0$}
%            \end{subfigure}
%
    % \captionsetup{width=.9\linewidth}
    % \caption{Examples of graph matrices in the decomposition of $Q$.
    % Recall that square vertices take labels in $[m]$ and circle vertices take labels in $[d]$.
    % Unlabeled edges have Hermite index $1$.}
    % \label{fig:examples}
\end{figure}
	For intuition, we advise to reader to ignore the rank-$2$ component for subsequent iterations beyond $Q_0$ as they turn out to be negligible, and thus focus on $M^{-1}\approx A^{-1}$ for most of the analysis. The only place where such distinction becomes crucial is the analysis for the variance of the base case $Q_0$. 

%\subsubsection{Why are Backbone-Dangling Shapes Freely Independent?}
%Next, with the combinatorial description of backbone dangling shapes in place, we can proceed to elaborate on why we expect their graph matrices to be freely independent.
%First, we recall that each backbone-dangling shape can be viewed as consisting of two components: the backbone-path corresponding to the $\sum_{t} v_t v_t^\top$ component, and the dangling path attached to the square vertex  corresponding to $M^{-1}\eta $.  
%
%There are two main factors governing the free independence. Consider the trace calculations of $Q$ and in particular the dominants terms in $\E[\Tr(QQ^\top)^q] $, we can again view the walks as having two components as well: backbone and dangling components. On a high level, it suffices for us to discuss each component individually:
%\begin{enumerate}
%	\item Restricted to the backbone path (from the component $\sum_{t\in [m]} v_tv_t^\top ) $, we observe that the dominant term requires each labeled square vertex (from $[m]$) to appear twice exactly. This shall be contrasted from circle vertex that do not have such constraints. This discrepancy from the stark contrast between the combinatorial factor of the underlying square and circle vertices: each square vertex's label contributes a cost of $m\approx d^2 \gg d$ where $d$ is the contribution factor of a circle vertex. Therefore, in the dominant walks, we would like to avoid extra repetition of the square vertex where the twice-appearance is the minimum appearance requirement due to mean-$0$ random variables involved.
%	\item 
%	For readers familiar with the separator language for graph matrix norm bounds, this is equivalent to   
%\end{enumerate}
\end{comment}

\paragraph{Why are Backbone-Dangling Shapes Freely Independent?}

With the combinatorial description of backbone-dangling shapes in place, we now
explain why one should expect their associated graph matrices to exhibit
free-independence. Recall that each backbone-dangling shape consists of
two components: a backbone path, coming from the matrix
\(\sum_{t\in[m]} v_tv_t^\top\), and a dangling path attached to a square vertex,
coming from the correction term \(M^{-1}\eta\).

The free-independence heuristic is governed by two largely separate
considerations, corresponding to these two components. Consider the trace moment
calculation for \(Q\), and in particular the dominant contributions to
\(
    \mathbb E\operatorname{Tr}\bigl((QQ^\top)^q\bigr).
\)
A walk contributing to this trace can itself be viewed as having a backbone
component and a dangling component. At a high level, we would like to view these two
components separately.

\paragraph{The Backbone Component}
 Restricting attention to the backbone path, which comes from the matrix
    \(
        \sum_{t\in[m]} v_tv_t^\top,
    \)
    the dominant trace contributions require each labeled square vertex, whose
    label lies in \([m]\), to appear exactly twice, provided $m\gg d$. This should be contrasted
    with circle vertices, whose labels lie in \([d]\), and which do not obey the
    same combinatorial constraint.

    The reason is the sharp difference between the label costs of square and
    circle vertices. A square vertex contributes a factor of \(m\approx d^2\),
    whereas a circle vertex contributes only a factor of \(d\). Thus,
    repetitions of square vertices are much more wasteful in the combinatorial counting of the trace walks.
    In the dominant walks, one therefore wants to avoid any unnecessary
    repetition of square vertices. The twice-appearance condition is the minimal
    requirement forced by the mean-zero random variables appearing along the
    backbone path.
    
    
    For readers familiar with the graph matrix norm bound language, this can be equivalently viewed by contrasting the norm bound value with the square vertex being the separator as opposed to the circle vertex, which accounts for a gap of $\tilde{O}(\sqrt{d})$ factor in the norm bound.
    
    \paragraph{The Dangling Component} Once we observe that each square vertex makes exactly two appearances, we can further deduce in the dominant term that each square vertex must be paired with the same dangling component (in both summands from $Q$, and its dangling edge-set). Therefore, this allows us to further bundle steps of the trace walk that correspond to the same underlying shape (and labeled edge-set),  \[ 
    \Tr[(Q\cdot Q^\top )^{q}] =  \Tr \sum_{\tau_1,\tau_2,\dots,\tau_{2q}} \cm_{\tau_1} \cdot \cm_{\tau_2} \cdot \dots \cdot \cm_{\tau_{2q}}\,.
    \]
    At this point, analogous to usual trace moment calculations for Wigner matrices,  it is well-anticipated that the partition for the dominant term should proceed in a non-crossing manner, giving rise to the semicircle shape of the resulting spectrum of $Q$. 
    
    We give a formal verification of this strategy via the equivalence with orthogonal polynomials in the later section. It culminates in establishing the following lemma.
    \begin{restatable}[Orthogonal Polynomials and Shape Concatenation for Semicircular Distributions]{lemma}{semicircleShapeConcatenation}
\label{lem:semicircle-shape-concatenation}
Consider a linear combination of backbone-dangling shapes 
\[
    K=\sum_{\tau\in\mathcal B(K)} c_\tau \cdot \cm_\tau
\]
from a ground set $\calB(K)$.
Then, for any $t>0 $,
\[
    P^{\mathsf{sc}}_t(K)
    =
    \sum_{\substack{
        \tau=\tau_1\circ\cdots\circ\tau_t\\
        \tau_i\in\mathcal B(K)\ \forall i\in[t]
    }}
    \left(\prod_{i=1}^{t} c_{\tau_i}\right)\cm_\tau
    + o_d(1),
\]
where the sum is over proper \(t\)-wise concatenations of shapes from
\(\mathcal B(K)\), and the error is in spectral norm, provided \[ 
\Var(K) = 1 \,.
\]

\end{restatable}

    
 \paragraph{Controlling the Variance of \texorpdfstring{$Q$}{Q}}   
 With the shape of the spectrum settled, it still remains for us to identify the radius of the semicircle distribution, which boils down to the following matrix variance quantity that we consider, \[ 
 \Var(Q) = \frac{1}{d}\cdot  \E[\Tr(QQ^\top)] = \sum_{\tau \in \calB(Q)} \Var(c_\tau \cdot \cm_\tau) +o_d(1)\,.  \] 
 where the second equality follows from free independence of backbone-dangling shapes. For our ansatz construction to apply, it is crucial that $Q$ has variance $1$. We now illustrate how we keep track of the variance across the iterative process.
 
 
%  Furthermore, instead of opening up the trace directly by the definition, it turns out crucial for us to exploit the alternating sign in the polynomial expansion for $A^{-1}$. 
 
By design of our iterative construction, we observe that once a shape $\tau$ is added to $Q$ at some level of $Q_i$, its coefficient stays invariant for the subsequent iterations by~\cref{prop:no-reappearance-old-shapes}. Hence, it suffices for us to partition $\calB(Q)$ by the level in which each shape gets added to the inner matrix.

\paragraph{Variance for Base Level $Q_0$}


 Let's start from our base case.  % Recall that we define the variance of the matrix $X$ of dimension  as \(
%\Var(X) \coloneqq \frac{1}{\dim(X) }\cdot\E \Tr[ XX^T ]\,, \) and $\eta = \calL(\1_d)$ denotes the raw deviation of identity. 
 \begin{restatable}[Decomposition of $Q_0$ via Graph Matrix and its Variance]{lemma}{qzerographvariance}
\label{lem:q0-graph-variance} Recall that $Q_0 \coloneqq -C_F \cdot \calR$,
  	we have 
  	\[ Q_0 = \sum_{\tau \in\calB(Q_0)} c(\tau) \cdot \cm_\tau + o_d(1) \]
  	where $\calB(Q_0)$ is a collection of (proper) \emph{backbone-dangling} shapes with variance
  	\[ 
  \Var(Q_0) = \sum_{\tau\in\calB(Q_0)} \Var(c_{\tau} \cdot \cm_{\tau}) +o_d(1)= (C_F)^2 \cdot \frac{\gam }{1-\gam } +o_d(1)\,.
  \]
\end{restatable} 
%
%\snote{Our formal proof in the subsequent section proceeds via a graph matrix analysis and it involves certain intersections that warrant extra care.}
Our formal proof in the subsequent section proceeds via a graph matrix analysis, and it involves certain intersections that warrant extra care. The direct proof via graph matrix language additionally verifies the structural property that each non-negligible term is a backbone-dangling shape. That said, to illustrate the variance bound, we give a separate argument that is considerably more intuitive, without having to appeal to graph matrix language in depth.
%
%Instead of explicitly computing the coefficient of $c_\tau$ for each shape $\tau$ in the decomposition of $Q_0$, we give an indirect argument through considering the race of $Q_0$ directly to give an intuitive explanation for variance bound. For notational convenience, we abbreviate $Q'\coloneqq \frac{1}{C_F} Q_0$. 
%

Ignoring the extra scaling, it suffices for us to show $\Var(\calR)=\frac{\gamma}{1-\gamma} $. Observe that the trace quantity can be rearranged as \begin{align*}
 	\Tr(\calR \cdot \calR^T) = \sum_{i,j \in [m]} w[i] \cdot w[j] \cdot  \langle v_i, v_j\rangle^2 = w^T M w
 	&= (M^{-1} \eta)^T \cdot M \cdot  (M^{-1}) \eta   \\
 	&= \eta^T \cdot  M^{-1} \cdot  \eta 
 \end{align*}
Finally, it can then be verified by the following scalar calculation whose proof is deferred to~\cref{sec:def-scalar}. \begin{claim}[Base Variance] \label{clm:base-scalar-variance} 
	\[ 
	 \frac{1}{d} \E[\eta^\top M^{-1} \cdot \eta] = \frac{\gamma}{1-\gamma} \,.
	\]
\end{claim}
%
%Recall that $M^{-1}$, as well as $M$, has two components. We bound the quadratic form of $\eta$ with respect to each component separtely as follows. We record the following scalar concentration estimates, deferring their proofs to the appendix. For the variance calculation, we state only their first-moment versions, although the same arguments also yield \(1+o_d(1)\) concentration around the corresponding means with high probability.
%
% \begin{restatable}{claim}{EtaAInvEtaConcentration} 
%\label{clm:concentration-eta-A-eta}
%We have
%\[
%    \frac{1}{d}\cdot \E\!\left[\eta^\top A^{-1}\eta\right]
%    =
%    \gamma + o_d(1).
%\]
%\end{restatable}
%% 	 and analogously for the second term from Woodbury expansion, \[ 
%%\frac{1}{d} \cdot  \E\bigg[\eta^T \left( \frac{1}{s^2-ru}\cdot A^{-1} \cdot (u\cdot \frac{J_m}{d} -s \frac{\eta\cdot \1_m^T+ 1_m^T \cdot \eta }{d}
%% + r \cdot \frac{\eta \cdot \eta^T }{d}) \cdot A^{-1}\right) \eta  \bigg ] = \frac{\gam^2}{1-\gamma}\,. 
%%\]
% 
%\begin{restatable}{claim}{WoodburyCorrectionConcentration} From the rank-$2$ update in Woodbury, 
%\label{clm:concentration-woodbury}
%we have
%\[
%\frac{1}{d} \cdot
%\E\!\left[
%\eta^\top
%\left(
%\frac{1}{s^2-ru}\,
%A^{-1}
%\left(
%u\cdot \frac{J_m}{d}
%-
%s\cdot \frac{\eta \mathbf 1_m^\top+\mathbf 1_m \eta^\top}{d}
%+
%r\cdot \frac{\eta\eta^\top}{d}
%\right)
%A^{-1}
%\right)
%\eta
%\right]
%=
%\frac{\gamma^2}{1-\gamma}.
%\]
%\end{restatable}

%\jnote{below is old}
%By our choice of the coefficient vector $w_0$ obtained from $Mw_0 = \eta_0$, we observe  that after normalizing by $C_F$, the variance term can be rearranged as \[ 
% \Var(Q_0) = C_F^2 \cdot \frac{1}{d}\cdot   \E[\eta_0^\top M^{-1} \eta_0] \,,
% \]
%In particular, we show that \[ 
% \frac{1}{d}\cdot   \E[\eta_0^\top M^{-1} \eta_0]  = \frac{\gamma}{1-\gamma} +o_d(1)\,,
% \]
% giving us $\Var(Q_0) = \frac{\gam}{1-\gam} \cdot (C_F)^2 $ where we ignore $o_d(1)$ terms throughout.
 
 \paragraph{Variance Evolution} 
% \jnote{double check if $New-P_j$ has been introduced by this point.}
Next, we derive a scalar recursion for the variance across iterations. The key point is that the graph-matrix and orthogonal-polynomial structure allows us to track variance shape-by-shape. In particular, it boils down to the following key identities. For any shape $\tau$ that appears in one of the ``new'' terms from the prior level, in~\cref{lem:vertical-concatenation}, we show that its corresponding correction term $\corr(\tau)$ to be added in the next level has variance \[ 
\Var(\corr(\tau)) =  \frac{\gamma}{1-\gamma} \Var(\tau) +o_d(1)\,. \]

%From here, we observe that the similar variance evolution recurrence from \cite{PotechinXu2026Theta} applies identically
We are now ready to deduce the scalar recurrence. Let $S_i \coloneqq \Var(Q_i)$ denote the variance at the $i$-th level, and let \(\textsf{New-}\mathcal P_j(Q_i)\) denote the collection of newly added terms at level-$i$ and degree-$j$, i.e., degree-$j$ terms in $H_i$ that are additionally multi-way proper concatenations from Chebyshev expansions.% the variance is precisely governed by the following recurrenc:\[ 
%S_{i+1} - S_i = \frac{\gam }{1-\gam } \cdot  \sum_{j = 2}^{\infty}{b_j^2(S_i^j - S_{i-1}^j)} \,.
%\]
%
%Letting $S_{\infty} = \lim_{k \to \infty}{S_k}$ assuming this limit exists, and plugging in the $\gam= \frac{1}{2}$ at the interested threshold,  we have that 
%\[S_{\infty} =\frac{\gam}{1-\gam} \left(  C_F^2 + \sum_{j=2}^{\infty}{{b_j^2}S_{\infty}^j}\right)\,.\] 
%Note that for $m = \frac{d^2}{4}$, i.e., $\gamma=1/2$, this is exactly the same recurrence from its counterpart for the Theta function. Finally, since $\sum_{j=2}^{\infty}{b_j^2} = 1 - C_F^2$, $S_{\infty} = 1$ is a solution of this equationThe discussion above culminates at the following lemma,  
%\begin{restatable}[Variance of the Inner Matrix \(Q\)]{lemma}{innermatrixvariance}
%\label{lem:inner-matrix-variance}
%Let \(Q \coloneqq \lim_{i\rightarrow \infty} Q_i\). We have
%\[
%    \Var(Q)=1\,.
%\]
%\end{restatable}
%
%

\begin{align*}
S_{i+1} - S_i &= \sum_{\tau\in \calQ_{i+1}\setminus \calQ_{i}} c_\tau^2  \cdot \Var(\tau) \tag{definition of \(S_i\)} \\&= 
 \sum_{j = 2}^{\infty}{b_j^2\sum_{\beta = \tau_1 \circ \cdots \circ \tau_j \in \textsf{New-}P_j(Q_i)}}  c_\beta^2  \cdot \Var(\corr(\beta)) \tag{each term is a correction} \\ 
 &=  \sum_{j = 2}^{\infty}{b_j^2\sum_{\beta = \tau_1 \circ \cdots \circ \tau_j \in \textsf{New-}P_j(Q_i)}}  c_\beta^2   \cdot \frac{\gamma}{1-\gam } \cdot \Var(\beta)  \tag{\cref{lem:vertical-concatenation}}\\ 
&=\frac{\gam}{1-\gam } \cdot  \sum_{j = 2}^{\infty}{b_j^2\sum_{\beta = \tau_1 \circ \cdots \circ \tau_j \in \textsf{New-}P_j(Q_i)} \prod_{t=1}^j c_{\tau_t}^2} \cdot \Var(\tau_t) \tag{via free independence }\\
&= \frac{\gam }{1-\gam } \cdot  \sum_{j = 2}^{\infty}{b_j^2(S_i^j - S_{i-1}^j)} 
\end{align*}
where we take $S_0 =  C_F^2$  and $S_{-1} = 0$. In the last line, we used that all \(j\)-fold concatenations from \(\mathcal B(Q_i)\) contribute \((S_i)^j\), while those using only shapes from \(\mathcal B(Q_{i-1})\) contribute \((S_{i-1})^j\). Their difference is therefore exactly the contribution from \(\textsf{New-}\mathcal P_j(Q_i)\).


 Summing this expression from $i = 0$ to $k$ gives that 
\[
S_{k+1} = \frac{\gam}{1-\gamma}( S_0 +  \sum_{j=2}^{\infty}{{b_j^2} \cdot S_k^j}) \,.
\]
Letting $S_{\infty} = \lim_{k \to \infty}{S_k}$ (assuming this limit exists), and plugging in $\gam= \frac{1}{2}$ at the ideal threshold,  we have that 
\[S_{\infty} = C_F^2 + \sum_{j=2}^{\infty}{{b_j^2}S_{\infty}^j}\,.\] Since $\sum_{j=2}^{\infty}{b_j^2} = 1 - C_F^2$, $S_{\infty} = 1$ is a solution of this equation, yielding the following lemma.
\begin{lemma}[Ideal Variance of the Inner Matrix $Q$ at the Threshold]
	Let $Q \coloneqq \lim_{i\rightarrow \infty} Q_i$. We have 
	\[ 
	\Var(Q) = 1\,.
	\]
\end{lemma} 
This ultimately gives us the desired semicircle of radius $1$ for the spectrum of $Q$ in the limit ignoring various truncation procedures. 

%\snote{"cooking" invokes the idiom "cooking the books" with pejorative connotations, I'd rather skip it}
\subsection{Finding the Correction Term} \label{sec:cooking-correction}
%\jnote{not sure if its worth highlighting- there is a way to incorporate this into the description of correction term without describing the buggy heursitic.}
We now elaborate on how we arrive at the counter-intuitive correction term per~\cref{eq:primal-corr}. For simplicity, let's consider a fixed term $\tau$ in the graph expansion of $H_i$ for the most recent level-$i$, and its correction term is given by \[ 
\corr(\tau) \coloneqq \frac{1}{1-\gamma} (\calL^*M^{-1}\calL(-\cm_\tau) -\gam \cdot \cm_\tau) 
\]
for any high-degree term $\tau$ that arises in the Chebyshev expansion $H(Q_i)$. We assume throughout that $\tau$ is a properly concatenated shape as otherwise it would be a negligible term.  In particular, we justify why we do not opt for the more natural correction term via \[
\corr_{natural}(\tau) \coloneqq \calL^*M^{-1} \calL (-\cm_\tau)\,.
\]
%To get started, it should be noted that our starting point is indeed the more natural option of $\corr_{alt}(\tau)$.
It should be noted that we arrive at the particular choice of correction by starting with the more natural option of $\corr_{alt}(\tau)$. In fact, there is sufficient heuristic justification for such a term being the ``right'' correction primitive beyond being the most natural one. Let us elaborate upon this.

\paragraph{A ``Buggy'' Heuristic for the Natural Choice of $\corr_{natural}(\tau)$:}
% The biggest support for this correction term is that as evidenced in the above variance calculation, for our machinery to apply, it is crucial that we may view the inner matrix $Q$ as ultimately as a sum of free matrices - matrices with semicircular spectrum. More concretely, 
 From the previous calculation for the variance of the inner matrix $Q$, we would like to show the added correction term (regardless of its concrete options) $\bar{\tau}$ has variance \[ 
\Var(\bar{\tau} ) = \frac{\gamma}{1-\gamma} \Var(\tau)
\]
\snote{use of the word ludicrous} \jnote{i just removed it}
with the key parameter of $\frac{\gamma}{1-\gamma}$ popping up that allows us to deduce a semicircular matrix of variance-$1$ in the limit. From this perspective, there is indeed a reason one may expect $\corr_{alt}$ to be the correct object to consider.
%
Let $\eta_\tau \coloneqq \calL(\eta_\tau)$ be the deviation incurred by $\tau$. By straightforward moment calculation for the vector, one can verify that \[ 
  \E\|\eta_\tau\|_2^2 = \gam \cdot \Var(\tau)\,.
\]
%The idea is that compared to $\tau$, the moment calculation of $\eta_\tau$ picks up two extra $h_1$ edges (which additionally has $2$ pairs of matchings in the second moment) and an additional square-vertex, that combines to a total factor of \[ 
%\frac{1}{d^2} \cdot 2 \cdot m =\gam \,.
%\]
Next, recall from the calculation for $\Var(\calR)$ that we can conveniently rewrite the variance of the correction matrix as \[
\Var(\calL^*M^{-1}\eta_\tau) = \frac{1}{d} \cdot \E[\eta_\tau^\top \cdot M^{(-1)} \cdot \eta_\tau ] \,.
 \]
At this point, suppose we ignore the correlation of underlying randomness between $\eta_\tau$ and $M^{-1}$; we may then heuristically view $\eta_\tau$ as a genuinely random vector with respect to $M^{-1}$, and we would expect \[ 
\E[\eta_\tau^\top \cdot M^{-1} \cdot \eta_\tau ] \approx_{\text{heuristic}} \E[\|\eta_\tau\|_2^2] \cdot \frac{1}{m} \E[\Tr M^{-1}]
\]
assuming $\eta_\tau$ is in isotropic position relative to the eigenbasis of $M^{-1}$. 
 \jnote{i guess its the same but which one is more natural to put here? $M$ or $M^{-1}$?...}  Since we have shown $M$ follows the spectrum of Marchenko-Pastur with parameter $\gamma$ (even up to the edge), we have \( 
\frac{1}{m} \E \Tr M^{-1}  = \frac{1}{1-\gamma}\,.
\)
Combining the above, we would indeed arrive at \[ 
\Var(\corr_{natural}(\tau) ) = \frac{\gam}{1-\gamma} \Var(\tau)\,
\]
as desired! 

That said, as suggested by the paragraph title, this heuristic is flawed because of the correlated randomness.%not directly valid and the correlated randomness poses a substantial risk. 
\paragraph{Graph Matrices Strike Back: Removal of Non-Free Terms}
To address this, we appeal to graph matrices, which also in turn give us an intuitive illustration of how the correlation of underlying randomness kicks in. One intuitive way to see the issue is that there are non-trivial intersection terms that arise in the expansion of $\corr_{alt}(\tau)$ in the graph matrix basis - it is no longer merely proper concatenations that matter.

Informally, in~\cref{lem:proper-decomp-deviation-attach}, we show that \[ 
\corr_{natural}(\tau) = \calL^*M^{-1}\calL(\cm_\tau) = \gam\cdot \cm_\tau + \text{desired properly concatenated shapes} 
\]
up to negligible errors. The presence of $\cm_\tau$ in turn flags another concern for us: we would like the inner matrix $Q$ to consist of ``free'' matrices with semicircular spectrum. And it is clear that $\cm_\tau$ cannot be such a matrix - it arises from $P_j(G)$ for some $j>1$ and some putatively free matrix $G$. One intuitive example would be to consider proper terms in $\fp_2(\calR)$ as showcased in~\cref{fig:p2(Q)}.

At this point, we arrive at $\corr(\tau)$ as it is obtained by precisely removing the troublesome non-free term while being rescaled to have the desired variance at the same time! %And m
Most importantly, $\corr(\tau)$ achieves the same job in terms of satisfying the affine constraints.
%
%
%That said, suppose we may ignore the correlation of randomness between $M^{-1}$ and $\eta_\tau$,  
%
%
%
%
% there are two underlying reasons:\begin{enumerate}
%	\item For our matrix analysis, in particular the spectrum estimate for our final inner matrix $Q$, we would like each summand in $Q$ to be a "free" matrix with semicircular spectrum;
%	\item 
%\end{enumerate}



    %
%\begin{enumerate}
%    \item \textbf{The backbone component.}
%    Restricting attention to the backbone path, which comes from the matrix
%    \[
%        \sum_{t\in[m]} v_tv_t^\top,
%    \]
%    the dominant trace contributions require each labeled square vertex, whose
%    label lies in \([m]\), to appear exactly twice. This should be contrasted
%    with circle vertices, whose labels lie in \([d]\), and which do not obey the
%    same combinatorial constraint.
%
%    The reason is the sharp difference between the label costs of square and
%    circle vertices. A square vertex contributes a factor of \(m\approx d^2\),
%    whereas a circle vertex contributes only a factor of \(d\). Thus,
%    repetitions of square vertices are much more wasteful in the combinatorial counting of the trace walks.
%    In the dominant walks, one therefore wants to avoid any unnecessary
%    repetition of square vertices. The twice-appearance condition is the minimal
%    requirement forced by the mean-zero random variables appearing along the
%    backbone path.
%    
%    
%    For reader familiar with the graph matrix norm bound language, this can be equivalently viewed by contrasting the norm bound value with the square vertex being the separator as opposed to the circle vertex, which accounts for a gap of $\tilde{O}(\sqrt{d})$ factor in the norm bound.
%    
%  \item 
%\end{enumerate}


%\subsection{Main Technical Lemmas}
%
%\begin{restatable}[Variance of inner matrix \(Q\)]{lemma}{varianceInnerQ} \label{lem:variance-inner-Q} We have \[ \lim_i \Var(Q_i) = 1, \] where \[ \Var(Q_i) \coloneqq \frac{1}{d}\,\E\!\left[\Tr(Q_i Q_i^\top)\right]. \] \end{restatable} 
%
%\begin{restatable}[Spectral radius of \(Q\)]{theorem}{spectralRadiusQ} \label{thm:spectral-radius-Q} With high probability, \[ \mathsf{spec}(Q) \subseteq [-2,2] \pm o_d(1). \] \end{restatable}
%
%\begin{restatable}[
%    (Semicircle) Orthogonal Polynomials and Shape Concatenation
%]{lemma}{semicircleShapeConcatenation}
%\label{lem:semicircle-shape-concatenation}
%Consider a linear combination of backbone-dangling shapes from a ground set $\calB(\cdot )$,
%\[
%    M=\sum_{\tau\in\mathcal B(Q)} c_\tau \cdot \cm_\tau
%\]
%
%Then, for any $t>0 $,
%\[
%    P^{\mathsf{sc}}_t(M)
%    =
%    \sum_{\substack{
%        \tau=\tau_1\circ\cdots\circ\tau_t\\
%        \tau_i\in\mathcal B(M)\ \forall i\in[t]
%    }}
%    \left(\prod_{i=1}^{t} c_{\tau_i}\right)\cm_\tau
%    + o_d(1),
%\]
%where the sum is over proper \(t\)-fold concatenations of shapes from
%\(\mathcal B(M)\), and the error is in spectral norm, provided \[ 
%\Var(M) = 1 \,.
%\]
%
%\end{restatable}
%
%
%An analogous statement holds for the matrix \(A\) arising from the correction
%for linear constraints. In this case, the relevant polynomial basis is the
%family of orthogonal polynomials associated with the Marchenko--Pastur law on
%\[
%    \bigl[(1-\sqrt{\gamma})^2,\,(1+\sqrt{\gamma})^2\bigr].
%\]
% To distinguish from the orthogonal polynomials for $\mu_{sc}$, we use \(\q_t\) to denote the orthogonal
%polynomials for Marchenko-Pastur from \cref{def:MP-basis}. Analogous to our result for Chebyshev polynomials, the orthogonal polynomials for Marchenko-Pastur also admits a clean decomposition in graph matrix via concatenated shapes.
%
%\begin{lemma}[MP Polynomials and Shape Concatenation]
%Consider
%\[
%    A=\cm_\alpha+\cm_\beta+\left(1+\frac1d\right)I_d + M_D\,,
%\]  for any $t>0$,
%\[
%    \q_t(A) \coloneqq \q^{(\gam)}_t(A) 
%    =  \underbrace{    \sum_{\substack{
%        \tau=\tau_1\circ\cdots\circ\tau_t\\
%        \tau_i\in\mathcal B(A)\ \forall i\in[t]
%    }} 
%    \cm_\tau}_{\fp_t(A) \coloneqq }
%    +o_n(1).
%\]
%where the LHS is the orthogonal polynomial evaluated on $A$ while the RHS is a linear combination of graph matrix of concatenated shapes that crucially do not include the trivial shape which corresponds to the identity matrix.
%%Here \(\{q_t\}_{t\ge0} = \{q^{\mathsf{MP}}_t \}\) denotes the orthogonal polynomials associated with the
%%Marchenko--Pastur distribution of parameter \(\gamma\) from~\cref{def: MP-basis}.
%\end{lemma}
%


%\input{refutation_overview}


\section{Refutation for Ellipsoid Fitting}
\label{sec:refutation}
In this section, we show how we can extend the iterative construction for ellipsoid fitting discussed above to give a refutation, i.e., the impossibility of ellipsoid fitting, when the number of constraints exceeds the threshold $m \ge d^2/4$.

\subsection{Overview of Refutation}
Recall that we let
\[
\mathcal L(X)_i= v_i^\top Xv_i ,
\qquad
\mathcal S:=\operatorname{Im}(\mathcal L^*)
=\operatorname{span}\{v_iv_i^\top:i\in[m]\},
\]
and let \(\calR\in\mathcal S\) be the identity perturbation satisfying
\(
\mathcal L(I-\calR)=\mathbf 1.
\)
Let \[ \Pi \coloneqq \Pi_S  = \calL^*M^{-1}\calL \] denote the orthogonal projection to $\calS$, and $\Pi^\perp$ the projection to its orthogonal complement.

We establish our refutation result by giving a dual witness for the ellipsoid fitting SDP. To get started, we record the dual SDP as follows,
\begin{proposition}\label{def:refutation-sdp}
	The dual of ellipsoid fitting is given by \(
	\Lambda \succeq 0
	\)
	such that \[ 
	\Lambda \in \calS \coloneqq \text{span}\{v_iv_i^\top: i\in [m] \}, \qquad \langle \Lambda, I-R\rangle < 0
	\]
where we remind the reader that we define the matrix inner-product as $\langle A,B\rangle = \Tr(A\cdot B)$.
\end{proposition}
As a sanity check, notice that such a matrix serves as a refutation for the primal program, as follows.
\begin{proposition}
Such a solution $\Lambda$ refutes the existence of a fitting ellipsoid.	
\end{proposition}
\begin{proof}
	Suppose there exists a solution $X$ for the primal, i.e., $\calL(X)=\1_m$ and $X\succeq 0$ and write $\Lambda = \calL^*(y)$. Then we have \[
  \langle y, \1 \rangle = \langle y, \calL(X)\rangle = \langle \calL^*(y), X \rangle =   \langle \Lambda, X\rangle =\Tr \left( X^{1/2} \Lambda X^{1/2}\right) \geq 0
\]
 by the PSDness of $\Lambda$ and $X$. On the other hand, since $\calL(I-\calR)=\1_m$,  we have \[
\langle y, \1_m \rangle = \langle y, \calL(I-\calR)\rangle = \langle \calL^*(y), I-\calR\rangle <0 
 \]
 yielding a contradiction.
\end{proof}

Next, we show the SDP dual can be reduced to seeking the following target matrix.
\begin{definition}[Target Dual Witness for Refutation] \label{def: target-refutation-witness}
	We call the following matrix our desired witness for refutation.
	 \[ 
 \Lambda = I_d+ c_R \cdot \calR - \Pi^\perp \cdot I_d  + Z = (I_d-   \Pi^\perp I)+ c_R \cdot \calR + Z  \succeq 0
 \] 
for some scalar $c_R > \frac{1-\gamma}{\gamma}$ and some matrix $Z \in \calS$ such that \( |\langle \calR, Z \rangle |= o(d)\) and $\Tr(Z) = o(d) $.
\end{definition}


To see why this is sufficient, and particularly how we arrive at the choice of $c_R >\frac{1-\gamma}{\gamma}$, we make the following observations:\begin{enumerate} 
	\item \jnote{check what condition on $\Pi^\perp I_d$ is most convenient} $\|\Pi^{\perp} I_d \|_{Frob} = \tilde{O}(1) = O(\frac{1}{\sqrt{d}}) $  since $I$ is almost entirely in the range of the rank-1 forms by considering the all-$1$ combination once appropriately scaled (concretely, $ d/m \cdot \1_m$). For intuition, this is a negligible term.
	\item $\Tr(\calR) = o(d)$;
% 	\item $\langle \calR, \calR\rangle = d \cdot \Var(R) = d \cdot \frac{\gam}{1-\gam }  $;
 	\item  $\langle \calR, \Pi^\perp I \rangle = 0   $ since $\calR\in \cal S$.\end{enumerate}
% 
% Let's now unpack the SDP-dual by rewriting as follows, \[ 
% \Lambda = I_d+ c_R \cdot R - \Pi^\perp \cdot I_d  +W = (I_d-   \Pi^\perp I)+ c_R \cdot R +W  \succeq 0
% \] 
% for some $W \in \calS $. Next, we make the following observations and assumptions about $W$,\begin{enumerate}
% 	\item $\|\Pi^{\perp} I\| =o(d) $  since $I$ is almost entirely in the range of the rank-1 forms by considering the all-$1$ combination once appropriately scaled;
% 	\item $\Tr(R) = o(d)$;
% 	\item $\langle R, R\rangle = d \cdot \Var(R) = d \cdot \frac{\gam}{1-\gam }  $;
% 	\item  $\langle R, \Pi^\perp I \rangle = 0   $ since $R\in \cal S$;
% 	\item Additionally, we assume \[
% 	|\langle W, \calR\rangle| = o(d)\,. 
% 	 \] 
% 	 \end{enumerate}
%These observations allow us to 
% further simplify the objective value as \[
% \langle \Lambda , I - \calR\rangle = \Tr(\Lambda)  - c_R \cdot \langle \calR, \calR\rangle + \Tr(\calR) -\langle \calR, W\rangle =  \Tr(\Lambda)  - c_R \cdot  \langle \calR, \calR\rangle+o(d)\,.
%  \]
%  Plugging in $\Tr(W)=o(d)$, we have \begin{align*}
% 	\langle \Lambda , I-R\rangle &=  \Tr(\Lambda) - \langle \Lambda,R\rangle \\
% 	&= d \left(  1 -     c_R\cdot \Var(\calR)  + o(1) \right)
% \end{align*} 
% since $\Var(\calR) = \frac{1}{d} \langle \calR,\calR \rangle$ by definition. Finally, recall from the calculation in the primal in~\cref{lem:q0-graph-variance}, we have \( 
% \Var(\calR) = \frac{\gam}{1-\gam }
% \)
% so we need to pick $c_R$ such that \[ 
% c_R \cdot \frac{\gam }{1-\gam} >1
% \]
% up to $o_d(1)$ terms.
This culminates in the following proposition.
\begin{proposition}
	The target dual matrix from \cref{def: target-refutation-witness} implies an SDP dual witness from~\cref{def:refutation-sdp}.
\end{proposition}

\begin{proof}
It remains only to verify that
\(\langle\Lambda,I_d-R\rangle<0\).
%Because \(R\in\calS\) and \(\Pi^\perp I_d\in\calS^\perp\),
%\[
%    \langle \Pi^\perp I_d,R\rangle=0.
%\]
%Furthermore, orthogonality of \(\Pi I_d\) and \(\Pi^\perp I_d\)
%gives
%\(
%    \langle \Pi^\perp I_d,I_d\rangle
%    =
%    \|\Pi^\perp I_d\|_{\mathrm F}^2.
%\)
%Therefore,
We have
\begin{align*}
    \langle\Lambda,I_d-\calR\rangle
    &=
    \left\langle
        I_d-\Pi^\perp I_d+c_R\cdot  \calR+W',\,
        I_d-R
    \right\rangle                                                    \\
    &=
    d-\|\Pi^\perp I_d\|_{\mathrm F}^2
      +(c_R-1)\Tr(\calR)
      +\Tr(W')
      -c_R \cdot \langle \calR,\calR\rangle
      -\langle W',\calR\rangle                                             \\
    &\le
    d-c_R \cdot \langle \calR,\calR\rangle+o(d)                                      \\
    &=
    d\left(
        1-c_R \cdot \frac{\gamma}{1-\gamma}+o_d(1)
    \right).
\end{align*}
Since
\[
    c_R>\frac{1-\gamma}{\gamma},
\]
the final expression is negative for all sufficiently large \(d\).
Thus
\(\langle\Lambda,I_d-\calR\rangle<0\), and \(\Lambda\) is a valid dual
witness.
\end{proof}

\jnote{below is old stuff that is no longer needed i believe...commenting out for now..}
%\begin{comment}
%This calculation motivates the following ansatz:\[
%\Lambda = \frac{1}{C_F} F(Q) 
% \] for some inner matrix $Q$. Concretely, we expect it to satisfy the following assumptions,
% \begin{definition}[Our Target Inner Matrix for Refutation]
% 	We seek a matrix $Q$ that satisfies the following assumptions:
% 	\begin{enumerate}
% 		\item Image constraint: $\frac{1}{C_F} \cdot F(Q) \in \calS     $;
% 		\item Semicircular constraint: $Q  $ is a matrix with semicircular spectrum of variance-1;
%% 		\[ 
%% 		Q_0 = C_F \cdot c_R \cdot R 
%% 		\]
% 		\item \[\langle  F(Q)/C_F , \calR\rangle = \langle I + Q_0/C_F  , \calR \rangle +o_d(1) = \langle I, \calR\rangle + \langle c_R \cdot \calR ,\calR\rangle +o_d(1)\,.  \]
% 	\end{enumerate}
% \end{definition}
% Now we show that given such a target dual witness matrix, we can take $ \Lambda = \frac{1}{C_F} F(Q)$ as our final solution. It suffices for us to verify the negative correlation constraint:
% \begin{align*}
% 	\langle \Lambda , I-R\rangle &= \langle \Lambda, I\rangle - \langle \Lambda,R\rangle \\
% 	&= d \cdot 1 - d\cdot (o_d(1) + c_R\cdot \Var(R) )
% \end{align*} 
% where we use $\langle I, R\rangle =o(d)$ and $\Var(X)\coloneqq \frac{1}{d} \Tr(X\cdot X^\top)$. Thus it suffices for us to pick $c_R$ such that \(
% c_R \cdot \Var(R) >1.
% \)
% Finally, recall from the calculation in the primal in~\cref{lem:q0-graph-variance}, we have \[ 
% \Var(\calR) = \frac{\gam}{1-\gam }
% \]
% so we need to pick $c_R$ such that \[ 
% c_R \cdot \frac{\gam }{1-\gam} >1\,.
% \]
% Next, we consider the choice of $c_R$. For convenience, we denote $\rho \coloneqq \frac{1-\gam}{\gam}$, the above condition is equivalent to verifying \[ 
% c_R > \rho\,.
% \]
% \end{comment}

\subsection{Dual Iterative Construction}
\snote{We will refer to our inner matrix for the dual refutation as $W$ to distinguish it from the matrix $Q$ in the primal program.}
To distinguish from our iterative process for the primal program, we will refer to our inner matrix for the dual refutation as $W$ throughout. Prompted by our desired target matrix, we start with \[ 
W_0 = C_F \cdot c_R \cdot \calR 
\]
for some constant $c_R$ to be chosen. We now apply the same spectral map $F$ to our inner matrix. 

\paragraph{Correction for Image Constraint} Throughout the iterative construction for the dual program, in contrast to satisfying the affine constraints as prescribed in the primal program, we would like to ensure instead $\frac{1}{C_F} F(W) \in \calS $---the image of the rank-$1$ forms. 

 Considering the first iteration $\frac{1}{C_F} F(W_0)$ and again its polynomial expansion, we have \[
\frac{1}{C_F} F(W_0) = I_d + c_R \cdot \calR + \frac{1}{C_F} \underbrace{ \sum_{j\ge 2} b_j\cdot  P_j(C_F \cdot c_R \cdot \cal R ) }_{H_0}\,.
 \]
 Since $I_d$ is in the desired image up to $o(d)$ error, we will view this as a negligible error term throughout, and correct such error alongside truncation error. Therefore, to ensure the resulting matrix is in the desired image, it suffices for us to correct for $H_0$. Again, it is by no means automatic that we should expect $H_0\in \calS$.
 
 Let's now consider how such correction might be achieved. By and large, we would like to find $\corr(H_0)$ such that \[
 H_0 + \corr(H_0) \in \calS\,.
 \] 
 Intuitively, one natural choice is to consider the orthogonal projection of $H_0$, and remove its orthogonal complement $(I_d - \Pi_\calS) H_0$. However, as we will discuss in ***, analogously to our correction for the primal program, the correlated randomness between $\Pi_S$ and $H_0$ produces a non-trivial fraction of $H_0$ in the above choice, posing a barrier to our ansatz solution as we would like to view the inner matrix as a sum of \textit{free} matrices. To address this issue, we alternatively consider \[ 
 \corr(H_0) \coloneqq \frac{1}{\gam }     \Pi_{\calS}   H_0- H_0 \,.
 \] 
It is straightforward to verify that this less natural choice also satisfies the image constraint as \[ 
H_0 + \corr(H_0) = \frac{1}{\gamma} \Pi_{\calS} H_0 \in \calS \,.
\]
Replacing the correction primitive in the primal iterative process gives us the following analogous process for the dual.

\begin{mdframed}[linewidth=0.2pt]
\textbf{Summary of the Initialization and Iterative Procedure for Refutation.}
%\\
%Let
%\[
%\calS\coloneqq \operatorname{Im}(\calL^*)
%\]
%and denote by
%\[
%\Pi
%\coloneqq
%\calL^*M^{-1}\calL
%\qquad\text{and}\qquad
%\Pi^\perp\coloneqq I-\Pi
%\]
%the orthogonal projections onto \(\calS\) and \(\calS^\perp\), respectively.
%\
\\\\
\textbf{Initialization:}
Set
\[
W_0\coloneqq C_Fc_R\calR,
\]
where \(c_R>0\) is later chosen so that the limiting inner matrix has
variance \(1\) and yields the desired negative correlation.
%
%Define the initial image-constraint deviation by
%\[
%W_0
%\coloneqq
%C_F I+\sum_{j\ge 2}b_jP_j(Q_0).
%\]
%\\
\\\\
\textbf{Iterative Update \(i\rightarrow i+1\):}
For \(i\ge 1\), define the newly created image-constraint deviation
\[
H_{i}
\coloneqq
\sum_{j\ge 2}b_j
\left(P_j(W_i)-P_j(W_{i-1})\right).
\]
Set
\[
\Delta_{i+1} = \corr(H_{i})
\coloneqq
\left(\frac{1}{\gamma} \Pi_{\calS}H_{i}  -   H_{i} \right) 
=
\frac{1}{\gamma} \left(\calL^*M^{-1}\calL H_{i}- \gamma  H_{i} \right),
\]
and update
\[
W_{i+1}\coloneqq W_i+\Delta_{i+1}.
\]
\\
\textbf{Output:}
For the limiting inner matrix \(Q\), take the dual witness
\[
\Lambda_R\coloneqq \frac{1}{C_F}F(W) - \Pi_{\calS}^\perp I_d.
\]
Note that $Z = \frac{1}{C_F}F(W) - I_d - c_R\calR$ for \cref{def: target-refutation-witness}.
\end{mdframed}

\begin{proposition}[Invariant for the Dual Iterative Process]
\label{prop:invariant-dual}
For every \(i\geq 0\),
\[
    \Pi_{\calS}^{\perp}
    \left(
        \frac{1}{C_F}F(W_i)-\Pi_{\calS}^{\perp}I_d
    \right)
    =
    \frac{1}{C_F}\Pi_{\calS}^{\perp}H_i.
\]
Equivalently,
\[
    \frac{1}{C_F}F(W_i)-I_d-\frac{1}{C_F}H_i
    \in \calS.
\]
Thus, at the end of iteration \(i\), the entire remaining violation of the
image constraint is incurred by the newly created higher-order term \(H_i\).
\end{proposition}

\begin{proof}
By definition of the correction,
\[
    H_i+\Delta_{i+1}
    =
    \frac{1}{\gamma}\Pi_{\calS}H_i
    \in\calS,
\]
and hence
\[
    \Pi_{\calS}^{\perp}\Delta_{i+1}
    =
    -\Pi_{\calS}^{\perp}H_i.
\]

For the base case, \(W_0=C_Fc_R\calR\in\calS\), and
\[
    F(W_0)=C_FI_d+W_0+H_0.
\]
Therefore,
\[
    \Pi_{\calS}^{\perp}
    \left(
        \frac{1}{C_F}F(W_0)-\Pi_{\calS}^{\perp}I_d
    \right)
    =
    \frac{1}{C_F}\Pi_{\calS}^{\perp}H_0.
\]

Now suppose the claim holds at level \(i\). Since
\[
    F(W_{i+1})-F(W_i)
    =
    \Delta_{i+1}+H_{i+1},
\]
we obtain
\begin{align*}
    \Pi_{\calS}^{\perp}
    \left(
        F(W_{i+1})-C_FI_d
    \right)
    &=
    \Pi_{\calS}^{\perp}H_i
    +\Pi_{\calS}^{\perp}\Delta_{i+1}
    +\Pi_{\calS}^{\perp}H_{i+1}\\
    &=
    \Pi_{\calS}^{\perp}H_{i+1}.
\end{align*}
Dividing by \(C_F\) completes the induction.
\end{proof}


\paragraph{Variance Calculation: Reverse Engineering for $c_R$} With the iterative process described above, we illustrate how to pick $c_R >0 $ as promised so that $c_R > \frac{1-\gamma}{\gamma}$ as in \cref{def: target-refutation-witness}. Recall that this is necessary to obtain a negative correlation with $\calR$.

To demonstrate how we pick $c_R$, let us first point to the similarity between our primal and dual corrections. In particular, for any $H_i$ from the $i$-th level of iteration, the primal correction is given by \[ 
\corr_{\mathsf{primal}}(H_i) =   \frac{1}{1-\gamma} \left(\calL^*M^{-1} \calL(-H_{i} ) + \gam \cdot H_{i}     \right) 
\]
recalling from~\cref{eq:primal-corr}. In the meanwhile, the dual update is given by \[ 
\corr_{\mathsf{dual}}(H_i) = \frac{1}{\gamma} \left(\calL^*M^{-1}\calL H_{i}- \gamma  H_{i} \right)  = -\frac{1-\gamma}{\gamma}  \corr_{\mathsf{primal}}(H_i)\,.  \numberthis \label{eq:primal-dual-correction} \]
Moreover, the identical calculation applies term-wise to each shape in the decomposition as well.

%Our final choice of $c_R$ involves incorporating truncation terms. Hence, instead of achieving genuinely negative correlation of $\langle \Lambda_R, \calR \rangle <0$, we start by showing 
Therefore, following the variance calculation in the primal iteration, we obtain the following analogous variance recurrence for the dual program. To distinguish it from the primal variance, define $SW_i \coloneqq \Var(W_i)$ to denote the variance of the dual inner matrix $W_i$.

The same calculation from the primal applies to give us the update equation
\begin{align*}
SW_{i+1} - SW_i &= \sum_{\tau\in \calQ_{i+1}\setminus \calQ_{i}} c_\tau^2  \cdot \Var(\tau) \tag{definition of \(SW_i\)} \\&= 
 \sum_{j = 2}^{\infty}{b_j^2\sum_{\beta = \tau_1 \circ \cdots \circ \tau_j \in \textsf{New-}P_j(Q_i)}}  c_\beta^2  \cdot \Var(\corr_{\mathsf{dual}}(\beta)) \tag{each term is a correction} \\ 
 &= (\frac{1-\gamma}{\gamma})^2 \cdot  \sum_{j = 2}^{\infty}{b_j^2\sum_{\beta = \tau_1 \circ \cdots \circ \tau_j \in \textsf{New-}P_j(Q_i)}}  c_\beta^2  \cdot \Var(\corr_{\mathsf{primal}}(\beta)) \tag{\cref{eq:primal-dual-correction}}
 \\&=   (\frac{1-\gamma}{\gamma})^2 \cdot \sum_{j = 2}^{\infty}{b_j^2\sum_{\beta = \tau_1 \circ \cdots \circ \tau_j \in \textsf{New-}P_j(Q_i)}}  c_\beta^2   \cdot \frac{\gamma}{1-\gam } \cdot \Var(\beta)  \tag{\cref{lem:vertical-concatenation}}
 \\&= (\frac{1-\gamma}{\gamma})^2 \cdot \frac{\gamma}{1-\gam } \cdot \sum_{j = 2}^{\infty}b_j^2 \left( (SW_i)^j - (SW_{i-1})^j\right) 
 % &=  \sum_{j = 2}^{\infty}{b_j^2\sum_{\beta = \tau_1 \circ \cdots \circ \tau_j \in \textsf{New-}P_j(Q_i)}}  c_\beta^2   \cdot \frac{\gamma}{1-\gam } \cdot \Var(\beta)  \tag{\cref{lem:vertical-concatenation}}\\ 
%&=\frac{\gam}{1-\gam } \cdot  \sum_{j = 2}^{\infty}{b_j^2\sum_{\beta = \tau_1 \circ \cdots \circ \tau_j \in \textsf{New-}P_j(Q_i)} \prod_{t=1}^j c_{\tau_t}^2} \cdot \Var(\tau_t) \tag{via free independence }\\
%&= \frac{\gam }{1-\gam } \cdot  \sum_{j = 2}^{\infty}{b_j^2(S_i^j - S_{i-1}^j)} 
\end{align*}
where the last line follows identically to that for the primal recurrence. Applying the telescoping sum again gives us \[ 
SW_{k+1} = \frac{\gam}{1-\gamma} SW_0 + \frac{1-\gamma}{\gamma}  \sum_{j=2}^{\infty}{{b_j^2} \cdot (SW_k)^j} \,.
\] 

For the variance to converge to $1$ in the limit, it is convenient to consider \(
B(s) \coloneqq \sum_{j\ge 2} b_j^2 s^j\,,
\) 
and the variance recurrence then can be rewritten in the form \(
SW_{k+1} =SW_0 + \frac{1-\gamma}{\gamma} B(SW_{k})\,.
\)
Since we additionally have $B(1) = 1-C_F^2$ by Parseval, and \(
\Var(SW_{0}) = C_F^2 \cdot c_R^2 \cdot \frac{\gamma}{1-\gam} \,,
\) requiring $1$ to be a fixed point gives \[ 
1= C_F^2 \cdot c_R^2 \cdot \frac{\gamma}{1-\gamma} + (1-C_F^2)\cdot \frac{1-\gam}{\gam}\,.
\] 
Therefore, we have \[ 
c_R^2 = \frac{1-\gamma}{\gamma} \cdot \frac{1}{C_F^2} \cdot \left(1- \frac{1-\gam}{\gamma} (1-C_F^2) \right)\,.
\]
Finally, we take $c_R$ to be the positive square root of the above. It remains for us to verify $c_R > \frac{1-\gamma}{\gamma}$.
\begin{claim} 
	For $c_R = \sqrt{\frac{1-\gamma}{\gamma} \cdot \frac{1}{C_F^2} \cdot \left(1- \frac{1-\gam}{\gamma} (1-C_F^2) \right)}$, we have $c_R >\frac{1-\gamma}{\gamma}$ provided $\gam >\frac{1}{2}$ which is exactly our desired threshold for refutation.
\end{claim}

\begin{proof} For convenience, write $\rho = \frac{1-\gam}{\gam}$.
	A direct calculation gives
\[
\begin{aligned}
c_R^2-\rho^2
&=
\frac{\rho\bigl(1-\rho(1-C_F^2)\bigr)}{C_F^2}
-\rho^2
=
\frac{\rho(1-\rho)}{C_F^2} >0,
\end{aligned}
\]
since $0<\rho<1$ for $\gam>\frac{1}{2}$.
\end{proof}

%
%
% Summing this expression from $i = 0$ to $k$ gives that 
%\[
%S_{k+1} = \frac{\gam}{1-\gamma}( S_0 +  \sum_{j=2}^{\infty}{{b_j^2} \cdot S_k^j}) \,.
%\]
%Letting $S_{\infty} = \lim_{k \to \infty}{S_k}$ (assuming this limit exists), and plugging in the $\gam= \frac{1}{2}$ at the interested threshold,  we have that 
%\[S_{\infty} = C_F^2 + \sum_{j=2}^{\infty}{{b_j^2}S_{\infty}^j}\,.\] Since $\sum_{j=2}^{\infty}{b_j^2} = 1 - C_F^2$, $S_{\infty} = 1$ is a solution of this equation, yielding the following lemma.
\begin{comment}
	
\paragraph{Choice and variance of the base term.} We now give intuition about how we pick the coefficient $c_R$. Concretely, we will pick $c_R$ such that the inner matrix has variance-$1$, corresponding to setting up \(1\) as a fixed point for the corresponding  recurrence.
%Writing
%\[
%F(x)=C_F+x+\sum_{j\geq2}b_jP_j(x),
%\qquad 
Consider \[
B(s):=\sum_{j\geq2}b_j^2s^j,
\]
we have
\[
B(1)=1-C_F^2.
\]
The variance recurrence takes the form
\[
S_{t+1}=S_0+\rho B(S_t)+o_d(1),
\]
with
\[
S_0=\Var(Q_0)
=
\frac{C_F^2c_R^2}{\rho}+o_d(1).
\]
Requiring \(S=1\) to be a fixed point gives
\[
1
=
\frac{C_F^2c_R^2}{\rho}
+
\rho(1-C_F^2),
\]
and hence
\[
c_R^2
=
\frac{\rho\bigl(1-\rho(1-C_F^2)\bigr)}
{C_F^2}.
\]
%Equivalently,
%\[
%\boxed{
%Q_0
%=
%\sqrt{\rho\bigl(1-\rho(1-C_F^2)\bigr)}\,R.
%}
%\]
%
%We want to pick the base term $Q_0$ to have variance
%\[
%\boxed{
%\Var(Q_0)
%=
%1-\rho(1-C_F^2)+o_d(1).
%}
%\]
%The correction from iterative process contributes the remaining variance
%\[
%\rho B(1)=\rho(1-C_F^2),
%\]
%so that the completed matrix \(Q=Q_0-W\) has variance \(1\). Therefore, we want to set \[ 
%c^2_R \cdot c_F^2 \cdot \Var(R) = 1-\rho(1-C_F^2)
%\]
%
It remains to verify that this rescaling preserves negative correlation.
A direct calculation gives
\[
\begin{aligned}
c_R^2-\rho^2
&=
\frac{\rho\bigl(1-\rho(1-C_F^2)\bigr)}{C_F^2}
-\rho^2\\
&=
\frac{\rho(1-\rho)}{C_F^2}.
\end{aligned}
\]
Recall that $\rho<1$ for $\gam>\frac{1}{2}$, we also have
\[
RHS  >0\,.
\]
Therefore, picking $c_R$ to the positive square-root, we have \[c_R >\rho \,. \]
 Under our assumption,
\[
\langle\Lambda,I-R\rangle
=
d \left( 1-\frac{c_R}{\rho}+o_d(1)\right)<0,
\]
as desired
\end{comment}

\begin{remark}
At the threshold \(\gamma=\frac12\), we have \(\rho=1\) and \(c_R=1\),
so the base term reduces to
\(
W_0=C_F \calR. 
\)	
\jnote{check the sign?}
\end{remark}


\section{Correction Primitive for Affine Constraint Deviations}
\label{sec:MP-sec}

In this section,  we prove that $A$ can indeed be inverted, and moreover, its inverse can be nicely decomposed in the graph matrix basis, via concatenated shapes of $\{\al, \beta \}$. Crucially, even though $A$ contains an identity term (without $1/d$ pre-factor), such a trivial shape does not appear as a ground-element for the concatenation.
\mpPolynomialShapeConcatenation*
In this section, we suggest that the reader momentarily ignore the negligible components in $A$, namely $M_D$ and $\frac{1}{d} I$, with our focus restricted to $\cm_\al + \cm_\beta + I_d $.
Once we establish tight control on the spectral radius of $A$, we then apply polynomial expansion via the orthogonal polynomials for $MP(\gamma)$ to understand $A^{-1}$. In particular, we show that $A^{-1}$ admits a ``nice'' decomposition in the graph matrix basis via concatenated shapes of $\al$ and $\beta$. Finally, as in existing analyses that are tight up to constant factors, we use the Woodbury identity to obtain the inverse for $M$ from $A^{-1}$.  This ultimately allows us to obtain again a nice decomposition of the inner matrix $Q$ in the graph matrix basis.

With the equivalence between concatenated shapes and the orthogonal polynomials, we show that the spectrum of $A$ continues to be MP-like, even at the spectral edge.
\spectralradiusA*
%
%\jnote{i think its ok to completely sack the current 4.1 because the overview already covers this, and we may cut straight to the discussion on norm bound from orthogonal polynomials}
%
%\snote{why "orthogonal polynomial of A"? The polynomials have nothing a priori to do with A. I think we should just say "orthogonal polynomials and concatenated shapes" or "orthogonal polynomials of Marchenko-Pasture and concatenated shapes" but I think the former is better.}
\subsection{MP Orthogonal Polynomial and Concatenated Shapes} For the reader's convenience, we recall the orthogonal polynomial basis
for the Marchenko--Pastur distribution, together with its recurrence.

\begin{definition}[Orthogonal Polynomial for Marchenko-Pastur] \label{def:MP-orthogonal-poly}
Fix \(0<\gamma<1\). Let
\(\{\mathfrak q_t^{(\gamma)}\}_{t\ge 0}\) be the polynomial sequence
defined by
\[
    \mathfrak q_0^{(\gamma)}(x)=1,
    \qquad
    \mathfrak q_1^{(\gamma)}(x)=x-1,
\]
and, for every \(t\ge 1\),
\[
    \mathfrak q_{t+1}^{(\gamma)}(x)
    =
    \bigl(x-(1+\gamma)\bigr)\mathfrak q_t^{(\gamma)}(x)
    -
    \gamma\,\mathfrak q_{t-1}^{(\gamma)}(x).
\]
\end{definition}


\begin{definition}[Concatenated Shapes for $A$]
	Define  $\calB(A) = \{ \al, \beta\}$; we let $\calP_t(A)$ denote the collection of shapes obtained by $t$-fold proper concatenations of shapes in $\calB(A)$, in other words, proper concatenations of $\al$ or $\beta$. For $t=0$, we define $\fp_0(A)=I_d$.
\end{definition}

With these definitions, we introduce our formal lemma for the equivalence between orthogonal polynomials and graph matrices for $A$.
\begin{lemma}[Quantitative version of~\Cref{lem:mp-polynomial-shape-concatenation}]
	\label{lem:mp-error-quant}
	\[ 
	\q_t(A) = \fp_t(A)  + \mathsf{Error}_t(A)\,.
	\]
	with \[ 
	\|\mathsf{Error}_t(A)\|_{sp} = \frac{\poly(t)}{\sqrt{d}} =  o_d(1)\,.
	\]
	provided $t \leq  d^\delta$ for some $\delta>0$ where we recall that $\fp_t(A)$ is the sum of the graph matrices in $\calP_t$, i.e., any $t$-fold concatenations of $\calB(A)$.
	
%	where we emphasize that the left-hand-side is the matrix-valued polynomial with $P_t(X)$ being the orthogonal basis for MP($\gamma$) while the right-hand-side is a linear combination of graph matrices. 
\end{lemma}

We begin by identifying the error incurred when passing to the graph-matrix perspective. 
Recall that $\calP_j(M)$ denotes the collection of proper $j$-way concatenations using base shapes from $\calB(A) = \{\cm_\al ,\cm_\beta \}$. Our goal now is to identify what shapes arise in the error term:  we will formally define their properties, and show they admit a small norm correction $o_d(1)$  and hence can be safely discarded.

Next, we identify the combinatorial properties of the error terms by induction. For the base case $t=1$ (since $t=0$ is trivial as we have identity on both sides), we have \begin{align*}
	LHS = A - I = \cm_\al + \cm_\beta + (1+\frac{1}{d}) I +M_D  - I = \cm_\al + \cm_\beta =  \sum_{\tau\in \calP_1(A)} \cm_\tau +o_d(1)\,.
\end{align*}
In other words, $\mathsf{Error}_1(A)= \frac{1}{d}I_d + M_D = o_d(1) $.
%\paragraph{Preliminaries for Shape Intersections}

\paragraph{Highlight of the Cancellation, and Motivating the Error Term.}
% We now isolate the cancellation mechanism that underlies the correspondence
% between concatenated shapes and the orthogonal polynomials for the
% Marchenko--Pastur distribution. 
%\snote{the next sentence is very long and hard to read}
As highlighted in the technical overview, the key point is that the backtracking
intersections produced by composing two \(\alpha\)-shapes (as well as two $\beta$-shapes) collapse after
summing over labels to exactly match the terms being subtracted in the three-term
recurrence for the MP polynomial basis.

Towards that end, we start by singling out the following three intersection shapes obtained from the multiplications involving $\cm_\al$ and $\cm_\beta$.
\begin{definition}[Half-Diamond Backtracking Intersection]
Let \(A\) and \(B\) be two \(M_\alpha\)-shapes. We say that the composition
\(A\cdot  B\) forms a \emph{half-diamond backtracking intersection} if
\begin{enumerate}
    \item \(U_A \not\equiv V_B\), where \(U_A\) is the left boundary square
    vertex of \(A\) and \(V_B\) is the right boundary square vertex of \(B\);

    \item the two circle vertices of \(A\) receive the same set of labels in $[d]$ as those in $B$.
\end{enumerate}
\end{definition}

Analogously, we also define a full-diamond backtracking intersection except now we have an extra vertex intersection between $U_A$ and $V_B$.
\begin{definition}[Full-Diamond Backtracking Intersection]
Let \(A\) and \(B\) be two \(M_\alpha\)-shapes. We say that the composition
\(A\cdot  B\) forms a \emph{full-diamond backtracking intersection} if
\begin{enumerate}
    \item \(U_A  \equiv V_B\), where \(U_A\) is the left boundary square
    vertex of \(A\) and \(V_B\) is the right boundary square vertex of \(B\);

    \item the two circle vertices of \(A\) receive the same set of labels in $[d]$ as those in $B$.
\end{enumerate}
\end{definition}

\begin{definition}[Half-Flat Backtracking Intersection]
	Let \(A\) and \(B\) be two \(M_\beta \)-shapes. We say that the composition
\(A\cdot B\) forms a \emph{half-flat backtracking intersection} if
\begin{enumerate}
    \item \(U_A \not\equiv V_B\), where \(U_A\) is the left boundary square
    vertex of \(A\) and \(V_B\) is the right boundary square vertex of \(B\);
    \item the circle vertex of \(A\) receives the same label in $[d]$ as that in $B$.
\end{enumerate}
\end{definition}


With these intersection shapes singled out, we first define the error term arising from the intended cancellation from the well-behaved intersections highlighted above. 

\begin{definition}[Backtracking residuals]
\label{def:backtracking-residuals}
Define
\begin{align*}
    \Delta_{\mathsf{HalfDiamondInt}}
    &\coloneqq
    \cm_{\mathsf{HalfDiamondInt}}
    -
    \gamma\cm_\alpha, \\
    \Delta_{\mathsf{HalfFlatInt}}
    &\coloneqq
    \cm_{\mathsf{HalfFlatInt}}
    -
    \gamma\cm_\beta, \\
    \Delta_{\mathsf{DiamondInt}}
    &\coloneqq
    \cm_{\mathsf{DiamondInt}}
    -
    \gamma I_m.
\end{align*}
%We call these matrices the \emph{backtracking residuals}. A
%\emph{backtracking-residual term} is any graph-matrix summand appearing
%when one of these residual matrices is expanded in the graph-matrix basis.
%
\end{definition}
Next, we show that these deviations are indeed negligible. This is the core of the cancellation that enables the equivalence with orthogonal polynomials for the MP distribution.


\begin{proposition}[Half-Diamond Cancellation] \label{prop:half-diamond-cancellation}
We have
\[
    \cm_{\mathsf{HalfDiamondInt}}
        =
    \gamma \cdot \cm_\al  + o_d(1).
\]
and analogously,
\[
    \cm_{\mathsf{HalfFlatInt}}
        =
    \gamma \cdot \cm_\beta   + o_d(1).
\]
where $\gam = \frac{2m}{d^2}$.
\end{proposition}

%\begin{proof}
%Without loss of generality, consider the $(i,j)$-entry of $\cm_{\mathsf{HalfDiamondInt}}$ for $i\neq j\in [m]$, we have
%
%\begin{align*}
%	\cm_{\mathsf{HalfDiamondInt}}[i,j] &= \sum_{k\notin  \{i,j\}}  \sum_{a< b \in [d]} (2 v_i[a]v_i[b]v_k[a]v_k[b]) \cdot  (2 v_k[a]v_k[b]v_j[a]v_j[b]) \\
%	&= 2 \sum_{a<b\in[d]} v_i[a]v_i[b] v_j[a]v_j[b]   \cdot  (\sum_{k\notin \{i,j\}} 2 v_k^2[a]v_k^2[b])\,.
%\end{align*}
%At this point, we may observe that the second term has mean \[ 
%\E \left[ \sum_{k\notin \{i,j\}} 2\cdot v_k^2[a]\cdot  v_k^2[b] \right] = 2(m-2) \cdot \frac{1}{d^2} = \frac{2m}{d^2} - \frac{4}{d^2} =\gamma+o_d(\frac{1}{d} )\,.
%\]
%Thus, we write
%\begin{align*}
%	\cm_{\mathsf{HalfDiamondInt}}[i,j] &= 2 \sum_{a<b\in[d]} v_i[a]v_i[b] v_j[a]v_j[b]   \cdot  (\sum_{k\notin \{i,j\}} 2 v_k^2[a]v_k^2[b])\\
%	&= 2 \gam  \sum_{a<b\in[d]} v_i[a]v_i[b] v_j[a]v_j[b] -  \\&2\sum_{a<b\in[d]} v_i[a]v_i[b] v_j[a]v_j[b] (  (\sum_{k\notin \{i,j\}} 2 v_k^2[a]v_k^2[b]) - \frac{2m}{d^2} + \frac{2}{d^2} )
%\end{align*}
%Zooming into the deviation term, we can rewrite it as \[
%    v_k[a]^2v_k[b]^2-\frac1{d^2}
%    =
%    \frac1d \cdot  h_2(v_k[a])
%    +
%    \frac1d \cdot  h_2(v_k[b])
%    +
%    h_2(v_k[a]) \cdot  h_2(v_k[b]).
%\]
%Since this holds for any non-zero entry of the matrix, we have 
%\[\cm_{\mathsf{HalfDiamondInt}} =\gam \cdot M_\al + \Delta_{\mathsf{HalfDiamondInt}} \]
%with $\Delta_{\mathsf{HalfDiamondInt}} = o_d(1)$.  
%
%% \todo{to fill in - write the deviation term in graph matrix and show that it is small norm.}
%
%Similarly, for $ \cm_{\mathsf{HalfFlatInt}}$, we have \begin{align*}
% \cm_{\mathsf{HalfFlatInt}}[i,j] &= \sum_{k\notin \{i,j\}}  \sum_{a\in[d]} h_2(v_i[a]) \cdot h_2(v_i[k]) \cdot h_2(v_k[a]) \cdot h_2(v_j[a])\\& = \sum_{a\in[d]} h_2(v_i[a]) \cdot h_2(v_j[a]) \cdot 	 \sum_{k\notin \{i,j\}}(h_2(v_k[a]))^2\,.
%\end{align*}
%Again, the second term has mean $m \cdot \frac{2}{d^2} = \gamma $
% where we use $\E[h_2(v_k[a])^2] = \frac{2}{d^2}$. Therefore, up to deviation term with $o_d(1)$ norm, we have \[ 
% \cm_{\mathsf{HalfFlatInt}} = \gam \cdot \cm_\beta \,.
% \]
% 
%% \todo{to fill in - write the deviation term in graph matrix and show that it is small norm.}
%
%\end{proof}

\begin{proposition}[Diamond Cancellation] \label{prop:diamond-cancellation}
We have
\[
    \cm_{\mathsf{DiamondInt}}= \gam \cdot I_m+o_d(1).
\]
\end{proposition}

%\begin{proof}
%The diamond intersection is the fully backtracking configuration obtained when
%the two circle vertices are identified and the two outer square boundary
%vertices are also identified. Thus the composition no longer produces a
%nontrivial backbone shape. 
% \begin{align*}
% \cm_{\mathsf{DiamondInt}}[i,i] = \left(2 \sum_{a<b} \sum_{k\notin \{i,j\}} v_i[a]\cdot v_j[a]\cdot v_i[b]\cdot v_j[b] \right)^2
% \end{align*}
% Next, we observe that this term has mean \[ 
% 4\cdot (m-2) \cdot \binom{d}{2} \cdot \frac{1}{d^2} = \frac{2m}{d^2}=\gamma +o_d(1)
% \]
% where we note that the factor $4$ comes from the scalar-$2$, and $m$ from the summation over the labels of $d$, $\binom{d}{2}$ from the summation of pairs $a<b$ and $\frac{1}{d^2}$ from the variance of the random variables.
% 
%% \todo{fill in the decomposition with deviation term + organization into graph matrix basis}
% \end{proof}

 We verify in~\cref{sec:MP-def-proof} that the deviation terms are of negligible norm. Together, these two identities show that the backtracking pieces generated by
shape concatenation reproduce exactly the correction terms in the
Marchenko--Pastur three-term recurrence. This is the cancellation that allows
the \(t\)-fold concatenation expansion to match \(P_t(A)\), rather than the
ordinary power \(A^t\). We defer the complete proof of \cref{lem:mp-error-quant} regarding the spectral norm bounds to the end of this section.

 Before that, we show that almost as a corollary to the equivalence lemma, we can also obtain control of the spectral edge of $A$ with ease. This is the focus of the subsequent section.

%\paragraph{Formal Identification of Error Terms}
%
%Our key observation is that the error terms can be decomposed into concatenation of local-collision-pieces that are by themselves almost locally vertex-injective up to the final concatenation defined as follows.

%
%\begin{definition}[Local-Collision-Pieces (LCP)]  \label{def:local-collision-pieces}
%We call an (improper) shape $\tau$ obtained from the $t$-wise concatenation of $\tau_1, \tau_2, \ldots, \tau_t$ a \emph{local collision piece} (equivalently, a concatenations of $t$ gadgets) if it satisfies the following properties,\begin{enumerate}
%	\item Each constituent piece is either of shape $\cm_\al$ or $\cm_\beta$;
%	\item There is no vertex intersection in the concatenation from $\tau_1$ to $\tau_{t-1}$. In other words, $\tau = \tau_L \circ  \tau_t$ can be decomposed such that $\tau_L = \tau_1 \circ \ldots \circ \tau_{t-1}$ and it is a proper shape. Additionally, there is a non-trivial vertex intersection between $\tau_t$ and $\tau_L$, and moreover, $\tau_{t}$ does not form a backtracking intersection with $\tau_{t-1}$;
%	\item or there is a 3-way backtracking intersection formed by $\tau_{t-2}, \tau_{t-1}$ and $\tau_t$.
%\end{enumerate}
%Additionally, we define the length of a local-collision-piece as the number of its underlying gadgets. And we call $\tau_1$ and $\tau_t$ the start and end of the local-collision-piece.
%\end{definition}
%
%
%
%
%With this definition in hand, we are now ready to characterize the error term. The up-shot of this characterization is that we do not have immediately backtracking shapes in the concatenation for the error terms.
%\begin{proposition}[Structural Property for Error Terms]   \label{prop:error-term-A-equivalence} 
%	Any term in $P_t(A)$ can be decomposed as \[ 
%\tau_t \cdot \tau_{t-1} \cdot \dots \cdot \tau_{2}\cdot \tau_1
%\]
%with $\tau_i \in \{\al,\beta \}$ for any $i\in[t]$, and moreover, each error term can be viewed as a concatenation of local-collision-pieces with possibly a properly concatenated-shape.
%
%In other words, there is no backtracking-intersection between any adjacent piece in the error terms.
%\end{proposition}
%
%\begin{proof}
%For starters, it should be clarified, $\tau = \tau_t\cdot \tau_{t-1}\cdot \dots\cdot \tau_2\cdot\tau_1$ is not necessarily a proper shape and we allow vertex intersections beyond the boundary across different pieces. However, if restricted to proper shapes such that there is no vertex-intersection beyond the boundary across any pieces, they are exactly the desired terms $\fp_t(A)$. It suffices for us to focus on the error terms.
%
%	We prove this by induction. For $i=2$, each error term is consist of a single local-collision piece of length-$2$. For the inductive step, consider the recurrence of \begin{align*}
%	P_{t+1}(A)  &= (A- (1+\gamma) \cdot I ) \cdot  P_t(A) - \gamma\cdot P_{t-1}(A) \\
%	&= (\cm_\al + \cm_\beta -\gam \cdot I)  \cdot  P_t(A) - \gam   \cdot P_{t-1}(A)  \end{align*}
%%	 and consider the effect on the error terms. There're two families of error terms,
%%	 \begin{enumerate}
%%	 	\item \textbf{New-Error-Term}: this arises from the multiplication of $\calP_t(A)$ with $\cm_\al +\cm_\beta $, and specifically, the non-trivial intersection terms therein. 
%%	 	\item \textbf{Prior-Error-Term}: this arrises from \[ 
%%	 	 (\cm_\al +\cm_\beta -\gam \cdot I ) \cdot \mathsf{Error}_t(A), \qquad -\gam\cdot  \mathsf{Error}_{t-1}(A)
%%	 	\]
%%	 \end{enumerate}
%%	 By induction hypothesis, we can write any term in $P_t(A)$ as a $t$-wise concatenations (albeit with potential non-trivial intersections in the error term), consider any term as $\tau =  \tau_t \cdot  \tau_{t-1}\cdot  \dots \cdot \tau_{2}\cdot \tau_1 $ with $\tau_i\in \{\cm_\al,\cm_\beta\}$ for any $i\in [t]$. 
%Henceforth, any term in 
%\[ 
%	  (\cm_\al + \cm_\beta -\gam \cdot I)  \cdot  P_t(A) 	 \]
%	 can be written as \[ 
%	 \sigma   \cdot \tau = \sigma \cdot \tau_t \cdot  \tau_{t-1}\cdot  \dots \cdot \tau_{2}\cdot \tau_1 \,. 
%	 \] 	 while noticing $\sigma =\{\al, \beta, \emptyset\}$ with $\sigma=\emptyset$ corresponding to the identity matrix and $\tau$ satisfying the prescribed properties by the hypothesis.
%
%	 Next, we consider the following categorizations of such shapes for $\sigma \neq \emptyset $.   \begin{enumerate}
%	 	\item If there is no vertex intersection between any pairs of $\tau_i$, and there is no intersection between $\sigma$ and $\tau$, any such term is in $\fp_t(A)$, i.e., the well-behaved terms as desired;
%%	 	\item For concreteness, consider the case that $\tau_{t-1}$ is not the end of a local-collision-piece. If $\sigma$ and $\tau_t$ form a backtracking intersection, there're two cases depending on whether it is a half-diamond or a diamond intersection. Either case, grouping over such intersections, up to errors in \cref{prop:diamond-cancellation}, we have the following shapes \[ 
%%	 	\al  \cdot  \left(\tau_{t-1}\cdot \dots \cdot \tau_1 \right)\,,
%%	 	\]
%%	 	and \[ \beta \cdot  \left(\tau_{t-1}\cdot \dots \cdot \tau_1 \right)\,,
%%	 	\]
%%	 	from half-diamond intersections, and \[ 
%%	 	\tau_{t-1}\cdot \dots \cdot \tau_1 
%%	 	\]
%%	 	each with an additional factor of $\gamma$ from \cref{prop:diamond-cancellation}.
%	\item If 
%	 \end{enumerate}
%	 
%\end{proof}
%
%\subsection{new attempt}
%
%\paragraph{Formal Identification of Error Terms}
%
%Our key observation is that every error term has a canonical first
%``bad'' collision. Before this collision occurs, the shape is locally proper;
%at the collision step, the new gadget intersects the previously built proper
%prefix in a non-backtracking way. This motivates the following definition.



%\subsection{new identification}

\subsection{Inverting the System of Affine Constraints}

In this section, we bound the spectral radius of $A$, which allows us to justify the polynomial expansion on the corresponding interval up to $o_d(1)$ fluctuation. Once the inverse of $A$ is justified, we appeal to existing results to obtain the inverse of $M$, following essentially the same route as in prior works, except that several scalar concentration bounds used therein require significant sharpening.
%\begin{lemma}[Spectral Radius via Norm Bounds for Centered Matrix] 
%Define $Y \coloneqq \cm_\al+\cm_\beta-\gam I$.
%With high probability,
%	\[
%	\|A - (1+ \gamma ) \cdot I \|_{sp}  = \| Y\|_{sp} \leq (1-o_n(1)) \cdot 2\sqrt{\gamma}\,.
%	\]
%\end{lemma}
%
%As usual, we prove this via trace moment method by bounding \[ 
%\E\left[\Tr(\cm_\al + \cm_\beta -\gamma \cdot I)^{2q} \right] \,,
%\]
%for $q = \poly\log d$,
%while a unique challenge arises here as we crucially rely on the "$-$" sign for cancellation, as well as not all matrix entries have mean-$0$. Our idea is to express $Y^{2q} $ in the orthogonal basis of MP($\gamma$) intended for $A$, i.e., write \[ 
% Y^{2q} = \sum_{t = 0}^{2q} c_t \cdot P_t(A) \,. 
%\]
%where $\{P_t\}_t$ is the orthogonal basis defined in~\cref{fact:MP-basis}.
%
%Next, we can then take the expectation, by the equivalence established above, $P_t(A)$ is given by the concatenated shapes of length-$t$ up to deviation term. Therefore, once we take the trace, up to the deviation term, the only contributing term is simply $P_0(A) = I_d$. It then suffices for us to bound the coefficient $c_0$ as well as the deviation term.
%
%\paragraph{Bounding the Coefficients}
%
%To bound the coefficient, we will inductively multiply with $Y$ and keep track of the coefficients. Let $c_{t}(k)$ be the coefficient of $P_t(A)$ in the expansion of $Y^k$. We observe the following recurrence relation, \[ 
%c_{t}(k) = c_{t+1}(k+1) + \gamma\cdot c_t(k) 
%\]

\paragraph{Inverse of $A$ and its Spectrum}
We now bound the spectrum of $A$, thereby validating our inversion via polynomial expansion in the corresponding orthogonal polynomials. Since the spectrum of $A$ is not symmetric around $0$, we pass through the centered matrix $Y$ as follows.
\begin{lemma}[Spectral Radius via Norm Bounds for the Centered Matrix]
Let
\[
    Y
    \coloneqq
    A-(1+\gamma)I_m
    =
    \cm_\alpha+\cm_\beta+M_D
    +
    \left(\frac1d-\gamma\right)I_m.
\]
Then, with high probability,
\[
    \|A-(1+\gamma)I\|_{\mathrm{sp}}
    =
    \|Y\|_{\mathrm{sp}} 
    \le
    \bigl(1+o_d(1)\bigr)2\sqrt{\gamma}.
\]
\end{lemma}


\begin{proof}
	
%\jnote{somewhere we should put in the expectation, but im not sure where atm...}
We prove the claim by the trace moment method. Fix
\(q=\polylog(d)\), to be chosen sufficiently large. It suffices to show that
\[
    \E\Tr(Y^{2q})
    \le
    \bigl(1+o_d(1)\bigr) \cdot  m \cdot (2\sqrt{\gamma})^{2q}.
\]


The main ingredient is that we should not expand \(Y^{2q}\) in the 
monomial basis. The cancellation coming from the negative diagonal term
\(-\gamma I\) is crucial. Instead, we expand \(Y^{2q}\) in the Marchenko–Pastur orthogonal-polynomial basis \(\{\q_t\}_{t\ge0}\) from~\cref{def:MP-basis}. We have \[
    \q_0(A)=I,
    \qquad
    \q_1(A)=A- I=\cm_\al + \cm_\beta + o_d(1),
\]
and, for \(t\ge1\),
\[
    (A - (1+\gamma) I)  \cdot  \q_t(A)
    =
    \q_{t+1}(A)+\gamma  \cdot \q_{t-1}(A)\,.\]

Write
\[
    Y^k
    =
    \sum_{t=0}^{k} c_t(k) \cdot \q_t(A)  = \sum_{t=0}^k c_t(k) \cdot \left(\fp_t(A)+ \mathsf{Error}_t\right).
\]
Ignoring for the moment the error from approximating $\q_t(A)$ by the $t$-wise proper concatenations in $\fp_t(A)$, the coefficients
satisfy the recurrence
\[
    c_t(k+1)
    =
    c_{t-1}(k)+\gamma c_{t+1}(k),
    \qquad t\ge1,
\]
with
initial condition
\[
    c_0(0)=1,
    c_t(0) =0 \,.
\]
for any $t>0$ as we have \[ 
Y^0 = I = \q_0(A) \,.
\]

%\jnote{these are w.h.p. statements right?...} \jnote{incorporated through the w.h.p. norm bound for the error piece- double check}

We now derive the coefficient recurrence. Since \[ \q_0(A)=I, \qquad \q_1(A)=A-I, \] we have \[ Y\q_0(A) = \bigl(A-(1+\gamma)I\bigr)\q_0(A) = P_1(A)-\gamma \q_0(A). \]
 On the other hand, for \(t\ge1\), the MP recurrence gives us \[ Y\cdot \q_t(A)=\q_{t+1}(A)+\gamma \q_{t-1}(A).  \] 
 Therefore the coefficients satisfy \[ c_0(k+1)=\gamma c_1(k)-\gamma c_0(k), \] \[ c_1(k+1)=c_0(k)+\gamma c_2(k), \] and, for \(t\ge2\), \[ c_t(k+1)=c_{t-1}(k)+\gamma c_{t+1}(k). \]
  We will use the following coefficient bound and spectral norm for the error terms to wrap up the proof. 
 \begin{claim}[Coefficient bound in the MP basis]
\label{clm:MP-coef-bnd} 
  Assume \(0<\gamma\le1\). For every \(k\ge0\) and every \(t\ge0\), \[ |c_t(k)| \le \gamma^{-t/2}(2\sqrt{\gamma})^k. \] 
 In particular, for \(k=2q\), \[ |c_t(2q)| \le \gamma^{-t/2}(2\sqrt{\gamma})^{2q}. \] \end{claim} 
We prove this claim in~\cref{sec:MP-def-proof}. 
On the one hand, for the dominant term, we have  \[ 
\E\Tr c_0(2q) \cdot \q_0 = m \cdot c_0 \leq  m \cdot 2\sqrt{\gamma}^{2q}\,.
\]
On the other hand, for the deviation terms, we have 
\begin{align*}
	\E[ \Tr \sum_{t>1}^{2q} c_t(2q) \cdot (\fp_t(A) + \mathsf{Error}_t)] &\leq \sum_{t=1}^{2q} c_t(2q) \cdot m \cdot \|  \mathsf{Error}_t\|_{sp}\\
	&\leq  m\cdot  2q \cdot \max(c_t(t)\cdot  \frac{\poly(q)}{\sqrt{d}} \cdot \sqrt{\gamma}^{t-1} )\\
	&=m \cdot (2\sqrt{\gamma})^{2q} \cdot \frac{\poly(q,t)}{\sqrt{d}\cdot \sqrt{\gamma}}
\end{align*} 
where we plug in~\cref{lem:LCP-bound} and~\cref{clm:MP-coef-bnd} for the final bound. Since $q = \poly\log (d)$, we have \[ 
\E[\Tr Y^{2q}] \leq  O(1) \cdot   m \cdot 2\sqrt{\gamma}^{2q}\,.
\]
Taking the $1/2q$-th root then gives us the desired norm bound immediately. 
\end{proof}
%
%It follows that
%\[
%    \operatorname{spec}(A)
%    \subseteq
%    \left[
%        (1-\sqrt\gamma)^2-o_d(1),
%        (1+\sqrt\gamma)^2+o_d(1)
%    \right].
%\]
%In particular, when \(\gamma\) is bounded away from \(1\) as in our regime of interest, \(A\) is
%invertible with high probability. By the generating-function identity for
%the MP polynomials,
%\[
%    A^{-1}
%    =
%    \frac{1}{1-\gamma}
%    \sum_{t\ge0}
%        (-1)^t\q_t(A),
%\]
%where the series converges in operator norm. Combining this identity with
%\Cref{lem:mp-error-quant} gives the desired graph-matrix expansion of
%\(A^{-1}\).

\paragraph{Woodbury Identity: Moving Inverse from $A$ to $M$} We next fill in the details concerning the inverse of $M$. This is largely adapted from the corresponding section in \cite{HKPX23} except that we strengthen the concentration bounds for the scalars to remove extra $O(1)$ slack. Concretely, we establish the following lemma.
\Minverse*
\begin{proof}
  Recall that $M$ can be decomposed into $A$ and a rank-2 matrix $B$ (\cref{eq:M-decomposition}); the Woodbury identity then allows us to go from the inverse of $A$ to that of $M$. For completeness, we recall the following facts adapted from \cite{HKPX23}.
  
  \begin{fact}[Matrix Invertibility] \label{fact:invertibility}
    Suppose $A \in \R^{n_1 \times n_1}$ and $C \in \R^{n_2 \times n_2}$ are both invertible matrices, and $U\in \R^{n_1 \times n_2}$ and $V \in \R^{n_2 \times n_1}$ are arbitrary.
    Then, $A + U C V$ is invertible if and only if $C^{-1} + V A^{-1} U$ is invertible.
\end{fact}

\begin{fact}[Woodbury matrix identity~\cite{Woodbury1950}]  \label{fact:woodbury}
    Suppose $A \in \R^{n_1 \times n_1}$ and $C \in \R^{n_2 \times n_2}$ are both invertible matrices, and $U\in \R^{n_1 \times n_2}$ and $V \in \R^{n_2 \times n_1}$ are arbitrary. Then
    \[ 
    (A+ UCV)^{-1} = A^{-1} - A^{-1}U\left(C^{-1}+ VA^{-1}U\right)^{-1}VA^{-1}  \,.
    \]
\end{fact}

In light of \cref{fact:woodbury}, we can write $B$ as $B = UCU^T$ where $U = V^T = \frac{1}{\sqrt{d}} \begin{bmatrix} 1_m & \eta \end{bmatrix} \in \R^{m \times 2}$, % and
$C = \begin{bmatrix} 1 & 1 \\ 1 & 0 \end{bmatrix}$,
and $M = A + UCU^T$.
Note that $C^{-1} = \begin{bmatrix} 0 & 1 \\ 1 & -1 \end{bmatrix}$, and we have
\begin{align*}
    C^{-1} + U^T A^{-1}U  = \begin{bmatrix}
	\frac{1_m^T A^{-1} 1_m}{d} & 1 + \frac{\eta^T A^{-1} 1_m}{d}\\\\
	1+ \frac{\eta^T A^{-1} 1_m}{d} & -1 + \frac{\eta^T A^{-1} \eta }{d} 
    \end{bmatrix}
    \eqqcolon
    \begin{bmatrix}
    	r & s \\
    	s & u
    \end{bmatrix}
    \,. 
    \numberthis \label{eq:russ}
\end{align*}

Assuming $A$ is invertible, from \cref{fact:invertibility}, we can prove that $M$ is invertible  by showing that the $2\times 2$ matrix $C^{-1} + U^T A^{-1}U$ is invertible, which is in fact equivalent to $ru - s^2 \neq 0$ from the above definitions.

%\begin{claim}[$M$ is invertible]
%	\[
%	ru - s^2\enq 0
%	 \]
%\end{claim}
%\todo{fill in the proof - should follow from this subsequent meta-lemma about the scalars }


%\begin{lemma} \label{lem:hyper-parameter-bounds}
%	We give the following bounds for the scalars defined above.\[ 
%	 r= (1+o_d(1)) \cdot  \frac{m}{d} \cdot \frac{1}{1-\gamma} \,,
%	\]
%	\[
%	s = 1+ o_d(1)\,,
%	 \]
%	 and \[
%	 u = -1 +\gam +o_d(1) \,.
%	  \]
%	  \jnote{the extra $\gamma$ here is crucial}
%\end{lemma}

\begin{restatable}[Scalar bounds]{lemma}{HyperParameterBounds}
\label{lem:hyper-parameter-bounds}
Let \(r,s,u\) be the scalars defined above. Then, with high probability,
\[
    r\coloneqq \frac{1_m^T A^{-1} 1_m}{d}
    =
    (1+o_d(1))\cdot \frac{m}{d}\cdot \frac{1}{1-\gamma},
\]
\[
    s \coloneqq 1+ \frac{\eta^T A^{-1} 1_m}{d}  = 1+o_d(1),
\]
and
\[
    u \coloneqq -1 + \frac{\eta^T A^{-1} \eta }{d}  = -1+\gamma+o_d(1).
\]
%In particular, the leading correction in \(u\) is the \(+\gamma\) term.
\end{restatable}

\begin{corollary}[$M$ is invertible] W.h.p. \[ 
ru-s^2 \neq 0\,.
\]	
\end{corollary}
%Given $M$ is also invertible, we may now proceed to expand $M^{-1}$ in the graph-matrix basis. 
Finally, expanding out Woodbury, we have \[ 
M^{-1} = A^{-1} + \frac{1}{s^2-ru}\cdot A^{-1} \cdot (u\cdot \frac{\1_m\1_m^T}{d} -s \frac{\eta\cdot \1_m^T+ \1_m^T \cdot \eta }{d}
 + r \cdot \frac{\eta \cdot \eta^T }{d}) \cdot A^{-1}\,.   \]
and this completes the proof of~\cref{lem:M-inverse}.
\end{proof}

%\todo{the final bound should be able to be wrapped up once we have a spectral norm bound for the deviation term at each polynomial approximation}
%
%\jnote{double check - this suggests that a bound of the form $\sqrt{\gamma}^i$ for the spectral norm of the deviation term would have sufficed for our purpose as long as \[ 
%\gamma^{- i/2} \cdot \mathsf{Bound for Error} \leq \poly(i)
%\]}
%
%
%\jnote{the updated deviation bound does behave as this as we like...check out block-value bound for LCP }
%
%\todo{add links + text for wrapping up the proof}
% \begin{proof}[Proof to~\cref{clm:MP-coef-bnd}]
%  Define \( B_k:=\max_{t\ge0}\gamma^{t/2}|c_t(k)|. \) We prove that \[ B_{k+1}\le 2\sqrt{\gamma}\,B_k. \] First consider \(t\ge2\). Using the recurrence, \[ \begin{aligned} \gamma^{t/2}|c_t(k+1)| &\le \gamma^{t/2}|c_{t-1}(k)| + \gamma^{t/2}\gamma |c_{t+1}(k)| \\ &= \sqrt{\gamma}\,\gamma^{(t-1)/2}|c_{t-1}(k)| + \sqrt{\gamma}\,\gamma^{(t+1)/2}|c_{t+1}(k)| \\ &\le 2\sqrt{\gamma}\,B_k. \end{aligned} \] For \(t=1\), we have \[ \begin{aligned} \gamma^{1/2}|c_1(k+1)| &\le \sqrt{\gamma}|c_0(k)| + \sqrt{\gamma}\gamma |c_2(k)| \\ &\le 2\sqrt{\gamma}\,B_k. \end{aligned} \] Finally, for \(t=0\), we use the boundary recurrence: \[ \begin{aligned} |c_0(k+1)| &\le \gamma |c_1(k)|+\gamma |c_0(k)| \\ &= \sqrt{\gamma}\,\gamma^{1/2}|c_1(k)| + \gamma |c_0(k)|. \end{aligned} \] Since \(0<\gamma\le1\), we have \(\gamma\le \sqrt{\gamma}\), and therefore \[ |c_0(k+1)| \le 2\sqrt{\gamma}\,B_k. \] 
%Combining the three cases gives \[ B_{k+1}\le 2\sqrt{\gamma}\,B_k. \] Since \(B_0=1\), induction yields \[ B_k\le (2\sqrt{\gamma})^k. \] Equivalently, \[ |c_t(k)| \le \gamma^{-t/2}(2\sqrt{\gamma})^k. \] This proves the claim. \end{proof}

%\subsection{Shifting From $A$ to $M$}
%
%\paragraph{Obtaining Inverse of $M$ via Woodbury Identity}
%The decomposition of $M$ into $A$ and a rank-2 matrix $B$ (\cref{eq:M-decomposition}) allows us to apply the Woodbury matrix identity about the inverse of low-rank corrections of invertible matrices.
%
%\begin{fact}[Matrix Invertibility] \label{fact:invertibility}
%    Suppose $A \in \R^{n_1 \times n_1}$ and $C \in \R^{n_2 \times n_2}$ are both invertible matrices, and $U\in \R^{n_1 \times n_2}$ and $V \in \R^{n_2 \times n_1}$ are arbitrary.
%    Then, $A + U C V$ is invertible if and only if $C^{-1} + V A^{-1} U$ is invertible.
%\end{fact}
%
%\begin{fact}[Woodbury matrix identity~\cite{Woo59}]  \label{fact:woodbury}
%    Suppose $A \in \R^{n_1 \times n_1}$ and $C \in \R^{n_2 \times n_2}$ are both invertible matrices, and $U\in \R^{n_1 \times n_2}$ and $V \in \R^{n_2 \times n_1}$ are arbitrary. Then
%    \[ 
%    (A+ UCV)^{-1} = A^{-1} - A^{-1}U\left(C^{-1}+ VA^{-1}U\right)^{-1}VA^{-1}  \,.
%    \]
%\end{fact}
%
%In light of \Cref{fact:woodbury}, we can write $B$ as $B = UCU^T$ where $U = V^T = \frac{1}{\sqrt{d}} \begin{bmatrix} 1_m & \eta \end{bmatrix} \in \R^{m \times 2}$ and
%$C = \begin{bmatrix} 1 & 1 \\ 1 & 0 \end{bmatrix}$,
%and $M = A + UCU^T$.
%Note that $C^{-1} = \begin{bmatrix} 0 & 1 \\ 1 & -1 \end{bmatrix}$, and we have
%\begin{align*}
%    C^{-1} + U^T A^{-1}U  = \begin{bmatrix}
%	\frac{1_m^T A^{-1} 1_m}{d} & 1 + \frac{\eta^T A^{-1} 1_m}{d}\\\\
%	1+ \frac{\eta^T A^{-1} 1_m}{d} & -1 + \frac{\eta^T A^{-1} \eta }{d} 
%    \end{bmatrix}
%    \eqqcolon
%    \begin{bmatrix}
%    	r & s \\
%    	s & u
%    \end{bmatrix}
%    \,. 
%    \numberthis \label{eq:russ}
%\end{align*}
%
%Assuming $A$ is shown to be invertible, next, from \Cref{fact:invertibility}, we can prove that $M$ is invertible  by showing that the $2\times 2$ matrix $C^{-1} + U^T A^{-1}U$ is invertible, which is in fact equivalent to $ru - s^2 \neq 0$ from the above definitions.
%
%%\begin{claim}[$M$ is invertible]
%%	\[
%%	ru - s^2\enq 0
%%	 \]
%%\end{claim}
%%\todo{fill in the proof - should follow from this subsequent meta-lemma about the scalars }
%
%
%%\begin{lemma} \label{lem:hyper-parameter-bounds}
%%	We give the following bounds for the scalars defined above.\[ 
%%	 r= (1+o_d(1)) \cdot  \frac{m}{d} \cdot \frac{1}{1-\gamma} \,,
%%	\]
%%	\[
%%	s = 1+ o_d(1)\,,
%%	 \]
%%	 and \[
%%	 u = -1 +\gam +o_d(1) \,.
%%	  \]
%%	  \jnote{the extra $\gamma$ here is crucial}
%%\end{lemma}
%
%\begin{restatable}[Scalar bounds]{lemma}{HyperParameterBounds}
%\label{lem:hyper-parameter-bounds}
%Let \(r,s,u\) be the scalars defined above. Then, with high probability,
%\[
%    r\coloneqq \frac{1_m^T A^{-1} 1_m}{d}
%    =
%    (1+o_d(1))\cdot \frac{m}{d}\cdot \frac{1}{1-\gamma},
%\]
%\[
%    s \coloneqq 1+ \frac{\eta^T A^{-1} 1_m}{d}  = 1+o_d(1),
%\]
%and
%\[
%    u \coloneqq -1 + \frac{\eta^T A^{-1} \eta }{d}  = -1+\gamma+o_d(1).
%\]
%%In particular, the leading correction in \(u\) is the \(+\gamma\) term.
%\end{restatable}
%
%\begin{corollary}[$M$ is invertible] W.h.p. \[ 
%ru-s^2 \neq 0\,.
%\]	
%\end{corollary}
%
%%Given $M$ is also invertible, we may now proceed to expand $M^{-1}$ in the graph-matrix basis. 
%Expanding out Woodbury, we have \[ 
%M^{-1} = A^{-1} + \frac{1}{s^2-ru}\cdot A^{-1} \cdot (u\cdot \frac{\1_m\1_m^T}{d} -s \frac{\eta\cdot \1_m^T+ \1_m^T \cdot \eta }{d}
% + r \cdot \frac{\eta \cdot \eta^T }{d}) \cdot A^{-1}\,. \numberthis \label{eq:M-inverse}
% \]
%\jnote{chunk of text hidden for sketch for proof of hyperparameter.}
%\subsection{Scalar Concentration for Hyper-Parameters}
%
%\begin{proof}[Proof Sketch to \cref{lem:hyper-parameter-bounds}] We defer the formal proof to the appendix, while giving a sketch of that by highlighting the dominant terms. Notice they are all random variables with means given by the following.
%	\paragraph{Analysis for $r$}
%	Recall that $r = \frac{1}{d}  \cdot \1_m^T \cdot A^{-1} \cdot \1_m  $, and \(
%	A^{-1} = \frac{1}{1-\gamma} \sum_{t\geq 0} P_t(A)\,.
%	 \)	
%	 For $t=0$  we have \[
%	 \frac{1}{d} \cdot \1_m^TP_0(A) \1_m = \frac{m}{d}\,,
%	  \] 
%	  and \[ 
%	  \1_m^T P_t(A) \1_m = o_d(1) 
%	  \]
%%	  \todo{fill in their concentration here. - would be great to quantify the $o_d(1)$ slack for example.}
%	  Therefore, we have \[ 
%	  r = \frac{1}{1-\gamma}  \cdot \frac{\1_m^T A^{-1} \1_m }{d} = \frac{1}{1-\gamma} \cdot \frac{m}{d}+o_d(1)\,.
%	  \]
%	  
%	 \paragraph{Analysis for $s$}
%	 Recall that $s = 1+ \frac{1}{d} \cdot \eta^T A^{-1} \1_m$, it suffices for us to bound the second term as $o_d(1)$. Notice for $t=0$, we have \[ 
%	 \E[\eta^T  \cdot  \1_m]  = 0
%	 \]
%%	 \todo{fill in concentration here, and that for higher $t$.}
%	 
%	 \paragraph{Analysis for $u$} Recall that $u = -1 + \frac{1}{d} \eta^T A^{-1}\eta $. For $t=0$, we have \[
%\frac{1}{d} \cdot \eta^\top P_0(A) \eta  = (1+o_d(1)) \cdot \frac{2md }{d^2} = (1+o_d(1)) \gam \,.
%\]
% by analogous calculation for the higher moments. 
%   	
%  	For $t=1$, we have \[
%  	\q_1(A) = \fp_1(A) = \cm_\al + \cm_\beta \,.
%  	 \]
%  	 The crucial observation is $\eta^T \cm_\al \eta = o_d(1)$  while the second term gives $\gamma^2 $ (and importantly with a  minus sign from the alternating sign in the polynomial expansion for $x^{-1}$). 
%%	 \todo{This is done in claim 4.2- would be great to move things over and reorganize }
%	 \end{proof}
	 
%	 
%
%\subsection{Formal Decomposition of MP Error Terms}
%
% As highlighted in the technical overview, the key point is that the backtracking
%intersections produced by composing two \(\alpha\)-shapes, as well as two $\beta$-shapes collapse, after
%summing over labels, match the terms being subtracted in the three-term
%recurrence for the MP polynomial basis.
%
%\paragraph{Preliminaries and Base Case} Towards that end, we start by isolating out the following three intersection shapes obtained from the multiplications involving $\cm_\al$ and $\cm_\beta$.
%\begin{definition}[Half-Diamond Backtracking Intersection]
%Let \(A\) and \(B\) be two \(M_\alpha\)-shapes. We say that the composition
%\(A\circ B\) forms a \emph{half-diamond backtracking intersection} if
%\begin{enumerate}
%    \item \(U_A \not\equiv V_B\), where \(U_A\) is the left boundary square
%    vertex of \(A\) and \(V_B\) is the right boundary square vertex of \(B\);
%
%    \item the two circle vertices of \(A\) receive the same set of labels in $[d]$ as those in $B$.
%\end{enumerate}
%\end{definition}
%
%Analogously, we also define a full-diamond backtracking-intersection except now we have extra vertex intersection between $U_A$ and $V_B$.
%\begin{definition}[Full-Diamond Backtracking Intersection]
%Let \(A\) and \(B\) be two \(M_\alpha\)-shapes. We say that the composition
%\(A\circ B\) forms a \emph{full-diamond backtracking intersection} if
%\begin{enumerate}
%    \item \(U_A  \equiv V_B\), where \(U_A\) is the left boundary square
%    vertex of \(A\) and \(V_B\) is the right boundary square vertex of \(B\);
%
%    \item the two circle vertices of \(A\) receive the same set of labels in $[d]$ as those in $B$.
%\end{enumerate}
%\end{definition}
%
%\begin{definition}[Half-Flat Backtracking Intersection]
%	Let \(A\) and \(B\) be two \(M_\beta \)-shapes. We say that the composition
%\(A\circ B\) forms a \emph{half-flat backtracking intersection} if
%\begin{enumerate}
%    \item \(U_A \not\equiv V_B\), where \(U_A\) is the left boundary square
%    vertex of \(A\) and \(V_B\) is the right boundary square vertex of \(B\);
%    \item the circle vertex of \(A\) receives the same label in $[d]$ as that in $B$.
%\end{enumerate}
%\end{definition}
%
%
%With these intersection shapes singled out, we are ready to introduce the core for the cancellation that enables the equivalence with orthogonal polynomials for MP-distribution.
%
%\begin{proposition}[Half-Diamond Cancellation] \label{prop:half-diamond-cancellation}
%We have
%\[
%    \cm_{\mathsf{HalfDiamondInt}}
%        =
%    \gamma \cdot \cm_\al  + o_d(1).
%\]
%and analogously,
%\[
%    \cm_{\mathsf{HalfFlatInt}}
%        =
%    \gamma \cdot \cm_\beta   + o_d(1).
%\]
%where $\gam = \frac{2m}{d^2}$.
%\end{proposition}
%\begin{proposition}[Diamond Cancellation] \label{prop:diamond-cancellation}
%We have
%\[
%    \cm_{\mathsf{DiamondInt}}= \gam \cdot I_d+o_d(1).
%\]
%\end{proposition}
%
%
%
%The proof of both propositions follow by bounding the deviation term on top of the discussion in technical overview. We defer their proofs to \cref{sec:MP-def-proof}.
%
%
%
%
%\paragraph{Extension to Higher-Degree}
%To generalize to higher degree, our key observation is that every error term is created by a terminal gadget that incurs
%(non-backtracking) vertex collision with prior gadgets. More precisely, as we build a record of gadgets from right to
%left, the error terms are generated at the first point where the newly
%attached gadget incurs vertex intersection (that is additionally not a backtracking intersection). We enforce
%local injectivity up to this terminal gadget, but we allow the terminal gadget
%to intersect either the current local prefix or a piece that appeared earlier
%in the global record. This flexibility is important for handling three-way
%backtracking/diagonal errors, where a backtracking pair also interacts with
%an earlier part of the concatenation.
%
%%\paragraph{Graph Matrix Basics}
%%
%%In light of this discussion, we first isolate the contribution corresponding to the idealized setting where no unintended vertex collisions occur. This leads to the notion of \emph{proper concatenation}, which captures the clean decomposition of shapes along their shared boundary without introducing any additional intersections. We now formalize this operation.
%%\begin{definition}[(Proper) Shape Concatenation]\label{def:shape-concate}
%%Given shapes $\alpha$ and $\beta$ such that $|V_{\alpha}| = |U_{\beta}|$, we take the concatenation $\alpha \circ \beta$ of $\alpha$ and $\beta$ to be the shape such that:
%%\begin{enumerate}
%%		\item $V(\alpha \circ \beta) = V(\alpha) \cup V(\beta)$ where the vertices in $V_{\alpha}$ are identified with the vertices in $U_{\beta}$.
%%        \item $U_{\alpha \circ \beta} = U_{\alpha}$ and $V_{\alpha \circ \beta} = V_{\beta}$.
%%		\item $E(\alpha \circ \beta) = E(\alpha) \cup E(\beta)$. 
%%\end{enumerate}
%%In other words, $\alpha \circ \beta$ is the shape obtained by gluing $\alpha$ and $\beta$ together along $V_{\alpha}$ and $U_{\beta}$, keeping all other vertices distinct, and taking all of the edges in both $\alpha$ and $\beta$.
%%	
%%This definition extends naturally to a sequence of $k$ shapes $\tau_1,\ldots,\tau_k$ such that for each $j \in [k-1]$, $|V_{\tau_j}| = |U_{\tau_{j+1}}|$. We all this a $k$-way concatenation.  We can obtain the concatenation $\tau_1 \circ \cdots \circ \tau_k$ by concatenating these shapes together one by one (it is not hard to check that the order of the concatenations does not matter).
%%\end{definition}
%%
%%\begin{definition}[Intersection patterns]
%%Given two shapes $\alpha$ and $\beta$ such that $|V_{\alpha}| = |U_{\beta}|$, we define an intersection pattern for $\alpha$ and $\beta$ to be a non-empty set $I \subseteq \left(V(\alpha) \setminus V_{\alpha}\right) \times \left(V(\beta) \setminus U_{\beta}\right)$ of intersection edges where each vertex is incident to at most one edge. Whenever we have an intersection edge $e = \{a,b\}$, this means that when we consider maps $\sigma: V(\alpha) \cup V(\beta) \to [n]$, we have the restriction that $\sigma(a) = \sigma(b)$.
%%
%%We define $\tau_{I}$ to be the shape obtained from $\alpha$ and $\beta$ as follows:
%%\begin{enumerate}
%%\item We start with $\alpha \circ \beta$.
%%\item For each intersection edge $e = \{a,b\}$, we merge $a$ and $b$ together into one vertex.
%%\end{enumerate}
%%\end{definition}
%%For convenience of our analysis, we will view intersection patterns as shapes as well, except it additionally comes with intersection edge that specifies vertex collisions.
%%
%%\begin{definition}
%%We define $Int_{\alpha,\beta}$ to be the set of possible intersection patterns for $\alpha$ and $\beta$.
%%\end{definition}
%%
%%\begin{definition}[$\circ$ and $\cdot$ notation]
%%	Given two shapes $\al_1$ and $\al_2$, we use $\circ$ to denote proper concatenation, and $\cdot $ to denote (possibly) improper concatenations. In other words, $\al_1\circ  \al_2$ denotes a single shape obtained by proper concatenation, while $\al_1\cdot \al_2$ denotes a collections of shapes that include both proper concatenation and intersection shapes.
%%\end{definition}
%%
%%\begin{proposition}
%%For all shapes $\alpha$ and $\beta$ such that $|V_{\alpha}| = |U_{\beta}|$, $
%%{\cm}_{\alpha}{\cm}_{\beta}    = {\cm}_{\alpha \circ \beta} + \sum_{I \in Int_{\alpha,\beta}}{{\cm_{\tau_I}}}$.
%%\end{proposition}
%%
%
%
%We first isolate the backtracking intersections that are canceled by the MP
%recurrence.
%
%\begin{definition}[Well-behaved Backtracking Intersectons]
%Let
%\[
%    \tau_t\cdot \tau_{t-1}\cdot \dots \cdot \tau_1,
%    \qquad \tau_i\in\{\alpha,\beta\},
%\]
%be a gadget record. We say that an adjacent pair
%\(
%    \tau_s\cdot \tau_{s-1}
%\)
%forms an \emph{well-behaved backtracking intersection} if the following hold:
%\begin{enumerate}
%    \item $\tau_s$ and $\tau_{s-1}$ form one of the designated backtracking
%    intersections: half-diamond, half-flat, or full-diamond;
%
%    \item the vertices participating in this backtracking block have no
%    additional vertex intersections with the earlier prefix
%    \[
%        \tau_{s-2}\cdot \dots \cdot \tau_1.
%    \]
%\end{enumerate}
%If the first condition holds but the second condition fails, we call
%$\tau_s\cdot\tau_{s-1}$ a \emph{non-isolated backtracking pair}. These are
%the three-way backtracking/diagonal intersections.
%\end{definition}
%Moreover, we introduce the following "residual" quantity to denote the slack in the desired cancellation showcased above.
%
%\begin{definition}[Backtracking residuals]
%\label{def:backtracking-residuals}
%Define the residual matrices associated with the three designated
%backtracking intersections by
%\begin{align*}
%    \Delta_{\mathsf{HalfDiamondInt}}
%    &\coloneqq
%    \cm_{\mathsf{HalfDiamondInt}}
%    -
%    \gamma\cm_\alpha, \\
%    \Delta_{\mathsf{HalfFlatInt}}
%    &\coloneqq
%    \cm_{\mathsf{HalfFlatInt}}
%    -
%    \gamma\cm_\beta, \\
%    \Delta_{\mathsf{DiamondInt}}
%    &\coloneqq
%    \cm_{\mathsf{DiamondInt}}
%    -
%    \gamma I_m.
%\end{align*}
%We call these matrices the \emph{backtracking residuals}.
%\end{definition}
%
%
%\begin{definition}[Local-Collision Pieces for MP Polynomials (MP-LCP)]
%\label{def:local-collision-pieces}
%Fix a gadget record
%\[
%    \tau_t\cdot \tau_{t-1}\cdot \dots \cdot \tau_1,
%    \qquad \tau_i\in\{\alpha,\beta\}.
%\]
%A consecutive sub-record
%\[
%    R
%    =
%    \tau_b\cdot \tau_{b-1}\cdot \dots \cdot \tau_a
%\]
%is called a \emph{local-collision piece}, abbreviated \emph{LCP}, if it
%satisfies the following conditions:
%\begin{enumerate}
%    \item Each constituent gadget is either of shape $\cm_\alpha$ or
%    $\cm_\beta$.
%
%    \item The concatenation before the terminal gadget is locally proper. That is,
%    \[
%        \tau_{b-1}\circ\tau_{b-2}\circ\cdots\circ\tau_a
%    \]
%    has no vertex intersections beyond the prescribed concatenation
%    boundaries.
%
%    \item The terminal gadget $\tau_b$ has a nontrivial vertex intersection
%    with some earlier gadget $\tau_j$, where $j<b$. The gadget $\tau_j$ may
%    lie either inside the same block, i.e., $a\leq j\leq b-1$, or inside an
%    earlier block of the global decomposition.
%
%    \item If the only intersection involving $\tau_b$ is a well-behaved
%    backtracking intersection with $\tau_{b-1}$, then the leading term is not
%    regarded as an error term; it is canceled by the MP recurrence. We record the residual from \cref{def:backtracking-residuals}.
%    \end{enumerate}
%The length of an LCP is the number of underlying gadgets in its sub-record.
%We call $\tau_a$ the start and $\tau_b$ the end of the LCP.
%\end{definition}
%
%\begin{definition}[LCP-Decomposition]
%Let $\tau$ be an improper shape with underlying concatenation record
%\[
%    \tau_t\cdot \tau_{t-1}\cdot \dots \cdot \tau_1,
%%    \qquad \tau_i\in\{\alpha,\beta\}.
%\]
%We say that $\tau$ admits an \emph{LCP-decomposition} if this record can be
%partitioned into consecutive blocks
%\[
%    \tau_t\cdot \tau_{t-1}\cdot \dots \cdot \tau_1
%    =
%   R_r\cdot R_{r-1}\cdot \dots \cdot R_1
%\]
%such that each block $R_j$ is an
%LCP.
%\end{definition}
%As a remark, we allow vertex intersections between different blocks. The only forbidden
%configuration is an unresolved isolated backtracking pair across the boundary
%between two adjacent blocks. Every isolated backtracking pair is either
%canceled by the MP recurrence or absorbed into the centered residual LCP
%coming from the corresponding backtracking cancellation.
%
%\begin{observation}[Terminal interaction of an LCP]
%\label{obs:lcp-terminal-not-isolated-backtracking}
%Let
%\(
%   R=\tau_b\cdot\tau_{b-1}\cdot\dots\cdot\tau_a
%\)
%be an LCP. Then the terminal gadget \(\tau_b\) is not a purely isolated
%well-behaved backtracking intersection with \(\tau_{b-1}\). More precisely, one of the
%following holds:
%\begin{enumerate}
%    \item \textbf{Non-backtracking Intersection:} \(\tau_b\) has a nontrivial non-backtracking intersection with some
%    earlier gadget in the record; or
%
%    \item \textbf{Three-way Diagonal Intersection:} \(\tau_b\cdot\tau_{b-1}\) forms a non-well-behaved backtracking intersection,
%    i.e., a three-way backtracking/diagonal interaction; or
%
%    \item \textbf{Backtracking Residual:} $\tau_b\cdot \tau_{b-1} $ forms a well-behaved intersection, and \emph{instead of the dominant term} which is canceled by the orthogonal polynomials, we pick up the residual after an isolated
%    backtracking cancellation from~\cref{def:backtracking-residuals}.  \end{enumerate}
%In particular, the leading contribution of a purely isolated backtracking pair
%is never treated as an error term; it is canceled by the MP recurrence.
%\end{observation}
%
%With this definition, we can state the structural property of the error
%terms.
%
%\begin{proposition}[Structural Property for Error Terms]
%\label{prop:error-term-A-equivalence}
%After applying the backtracking cancellations coming from the MP recurrence,
%every summand appearing in $\q_t(A)$ has an underlying concatenation record of the
%form
%\[
%    \tau_t\cdot \tau_{t-1}\cdot \dots \cdot \tau_2\cdot \tau_1,
%    \qquad \tau_i\in\{\alpha,\beta\}.
%\]
%The summands whose records are proper concatenations are exactly the desired
%terms in $\fp_t(A)$. Every remaining summand, i.e., every error term,
%admits an LCP-decomposition.
%%
%%Equivalently, after the MP recurrence cancels all isolated backtracking
%%pairs, every surviving error term is generated by terminal collisions. In
%%particular, there is no unresolved isolated backtracking intersection between
%%adjacent blocks in the LCP-decomposition.
%\end{proposition}
%\begin{comment}
%\begin{proof}
%%First we clarify the bookkeeping. A term may visibly look like a lower-order
%%shape because an adjacent backtracking pair has collapsed through a
%%half-diamond, half-flat, or full-diamond cancellation. Nevertheless, we track
%%the original gadget record. Thus a residual term produced by a two-gadget
%%backtracking block is still counted as having two underlying gadgets. More
%%generally, every summand appearing at level $t$ is assigned a record of
%%length $t$.
%We prove the statement by induction on $t$. The claim is trivial for $t=0$ and $t=1$. 
%\[
%    \q_0(A)=I,
%\]
%and
%\[
%    \q_1(A)=A-I=\cm_\alpha+\cm_\beta.
%\]
%Thus there are no error terms at levels $0$ and $1$.
%
%For $t=2$, every term comes from multiplying two gadgets from
%$\{\alpha,\beta\}$. If the two gadgets concatenate properly, the resulting
%shape belongs to $\mathcal P_2(A)$. If the terminal gadget has a nontrivial
%non-backtracking intersection with the previous gadget, the resulting shape
%is an LCP of length $2$. If the two gadgets form an isolated backtracking
%pair, then its leading contribution is canceled by the lower-order terms in
%the MP recurrence. The only possible survivors are the centered residuals
%from the corresponding backtracking cancellation, and by Definition
%\ref{def:local-collision-pieces} these residuals are recorded as LCPs
%supported on the same two-gadget record. Hence the proposition holds for
%$t=2$.
%
%Now assume the proposition holds for all levels up to $t$. We prove it for
%$t+1$.
%
%Using the MP recurrence and the identity
%\[
%    A-I=\cm_\alpha+\cm_\beta,
%\]
%we have
%\begin{align*}
%    \q_{t+1}(A)
%    &=
%    (A-(1+\gamma)I)\q_t(A)-\gamma \q_{t-1}(A)  \\
%    &=
%    (\cm_\alpha+\cm_\beta-\gamma I)\q_t(A)
%    -\gamma \q_{t-1}(A).
%\end{align*}
%Therefore every new non-scalar multiplication term is obtained by attaching
%a new gadget
%\(
%    \sigma\in\{\alpha,\beta\}
%\)
%to the left of a term from $P_t(A)$. By the induction hypothesis, each term
%of $P_t(A)$ has a gadget record
%\[
%    \tau_t\cdot \tau_{t-1}\cdot \dots \cdot \tau_1,
%    \qquad \tau_i\in\{\alpha,\beta\},
%\]
%and each error term among them admits an LCP-decomposition. After multiplying
%by $\sigma$, the new record is
%\[
%    \sigma\cdot \tau_t\cdot \tau_{t-1}\cdot \dots \cdot \tau_1.
%\]
%
%We classify the possible interactions of $\sigma$ with the previous record.
%
%\medskip
%\noindent
%\textbf{Case 1: $\sigma$ concatenates properly.}
%
%Suppose $\sigma$ has no vertex intersection with the previous shape except
%for the boundary vertex required by concatenation. If the old term was a
%proper $t$-way concatenation, then the new term is a proper $(t+1)$-way
%concatenation, and hence belongs to $\mathcal P_{t+1}(A)$.
%
%If the old term was already an error term, then by induction it admits an
%LCP-decomposition. Since $\sigma$ attaches properly, it introduces no new
%terminal collision. We therefore either attach $\sigma$ to the leftmost
%proper block, or create a new proper block at the left end. The resulting
%term still admits an LCP-decomposition.
%
%\medskip
%\noindent
%\textbf{Case 2: $\sigma$ creates a nontrivial non-backtracking collision.}
%
%Suppose $\sigma$ has a nontrivial vertex intersection with some earlier
%gadget in the record, and this interaction is not an isolated backtracking
%intersection with the adjacent gadget $\tau_t$.
%
%Take the maximal locally proper run immediately preceding $\sigma$. That is,
%choose the smallest index $a$ such that
%\[
%    \tau_t\cdot \tau_{t-1}\cdot \dots \cdot \tau_a
%\]
%is locally proper before adding $\sigma$, while extending further to the
%right would cross an existing LCP boundary. Then
%\[
%    \tau_t\circ \tau_{t-1}\circ \cdots \circ \tau_a
%\]
%is proper by construction.
%
%The terminal gadget $\sigma$ has a nontrivial intersection with some earlier
%gadget in the full record. This earlier gadget may lie inside the locally
%proper run
%\[
%    \tau_t,\tau_{t-1},\ldots,\tau_a,
%\]
%or it may lie inside an earlier LCP block. In either case, the consecutive
%sub-record
%\[
%    \sigma\cdot \tau_t\cdot \tau_{t-1}\cdot \dots \cdot \tau_a
%\]
%satisfies the definition of an LCP: it is locally injective up to the
%terminal gadget $\sigma$, and the terminal gadget intersects something that
%appeared earlier in the global record.
%
%The remaining part of the old record already has an LCP-decomposition by the
%induction hypothesis. Thus adding this new terminal LCP gives an
%LCP-decomposition of the new error term.
%
%\medskip
%\noindent
%\textbf{Case 3: $\sigma$ forms a non-isolated backtracking pair.}
%
%Now suppose $\sigma$ forms a backtracking intersection with the adjacent
%gadget $\tau_t$, but the backtracking block is not isolated from the earlier
%prefix
%\[
%    \tau_{t-1}\cdot \dots \cdot \tau_1.
%\]
%Equivalently, some vertex participating in the backtracking block also
%intersects an earlier gadget. This is the three-way
%backtracking/diagonal case.
%
%This case is an LCP under Definition \ref{def:local-collision-pieces}. Indeed,
%the block
%\[
%    \sigma\cdot \tau_t
%\]
%is locally proper before adding the terminal gadget $\sigma$, since its
%prefix consists only of the single gadget $\tau_t$. The terminal gadget
%$\sigma$ intersects the adjacent gadget $\tau_t$ through the backtracking
%configuration, and because the backtracking pair is non-isolated, the
%corresponding diagonal or identity replacement is not exact relative to the
%earlier prefix. Hence this surviving three-way term is recorded as an LCP
%supported on the non-isolated backtracking block, with the additional
%cross-block vertex intersection allowed by the definition.
%
%The rest of the record is already decomposed by induction, so the whole term
%admits an LCP-decomposition.
%
%\medskip
%\noindent
%\textbf{Case 4: $\sigma$ forms an isolated backtracking pair.}
%
%Finally suppose $\sigma$ and $\tau_t$ form an isolated backtracking pair.
%Then the backtracking cancellation identities apply. The half-diamond and half-flat intersections have leading contributions of
%the form
%\(
%    \gamma \cm_\alpha
%   \) or \(    \gamma \cm_\beta,
%\)
%and these leading terms are canceled by the contribution of
%\(
%    -\gamma P_t(A)
%\)
%in the recurrence. The full-diamond intersection has leading contribution
%\(
%    \gamma I,
%\)
%and this leading term is canceled by the contribution of
%\(
%    -\gamma P_{t-1}(A).
%\)
%
%Thus an isolated backtracking pair never remains as an unresolved
%backtracking interaction in the error expansion. Its leading contribution is
%paired with, and canceled by, the corresponding lower-order term in the MP
%recurrence. The only part that can survive is the centered residual from the
%relevant cancellation identity. By Definition
%\ref{def:local-collision-pieces}, this residual is recorded as an LCP
%supported on the same underlying backtracking pair. Hence the resulting
%summand again admits an LCP-decomposition.
%
%\medskip
%
%The four cases above exhaust all possible interactions between the newly
%attached gadget $\sigma$ and the existing record.
%
%It remains only to account for the scalar recurrence terms
%\[
%    -\gamma \q_t(A)
%    \qquad\text{and}\qquad
%    -\gamma \q_{t-1}(A).
%\]
%These terms introduce no new vertex intersections. They serve only to cancel
%the leading contributions of the isolated backtracking pairs considered in
%Case 4. After grouping each scalar recurrence term with the corresponding
%isolated backtracking contribution, the leading backtracking piece is removed.
%Any leftover contribution is precisely the centered residual already recorded
%as an LCP. If the term being canceled is attached to a summand that already
%contains error pieces, then the existing LCP-decomposition is preserved by the
%induction hypothesis.
%
%Consequently, every surviving non-proper summand in \(\q_{t+1}(A)\) admits an
%LCP-decomposition. Moreover, the construction keeps track of exactly one new
%gadget added at this step, so every surviving summand in \(\q_{t+1}(A)\) has an
%underlying record of exactly \(t+1\) gadgets from \(\{\alpha,\beta\}\).
%
%This completes the induction.
%\end{comment}
%%\end{proof}
%\begin{proof}
%We prove the statement by induction on \(t\). As throughout this
%subsection, we consider the expansion in \(\alpha\)- and \(\beta\)-gadgets;
%all terms containing \(M_D\) or \(\frac{1}{d}I\) are absorbed into
%\(\mathsf{Error}_t(A)\). We also retain the original gadget record when a
%backtracking contribution is replaced by its centered residual. Thus, a
%residual arising from a two-gadget backtracking block continues to carry
%the corresponding two-gadget record.
%
%We first record a simple consequence of
%\Cref{def:local-collision-pieces}. Suppose a record admits an
%LCP-decomposition. Removing its leftmost gadget leaves a record that still
%admits an LCP-decomposition. Indeed, if the leftmost block is proper, then
%the remaining part of that block is still proper. If the leftmost block is
%an LCP, then its leftmost gadget is its terminal gadget, and removing it
%leaves the locally proper prefix guaranteed by the definition of an LCP.
%
%The claim is immediate for \(t=0,1\). Indeed,
%\[
%    \q_0(A)=I,
%\]
%while
%\[
%    \q_1(A)
%    =
%    A-I
%    =
%    \cm_\alpha+\cm_\beta
%    +
%    \left(\frac{1}{d}I+M_D\right),
%\]
%and the parenthesized term is already included in
%\(\mathsf{Error}_1(A)\).
%
%Now fix \(t\geq 1\), and assume the proposition holds at levels \(t\) and
%\(t-1\). The MP recurrence gives
%\begin{align*}
%    \q_{t+1}(A)
%    &=
%    \bigl(A-(1+\gamma)I\bigr)\q_t(A)
%    -
%    \gamma\q_{t-1}(A) \\
%    &=
%    \bigl(\cm_\alpha+\cm_\beta-\gamma I\bigr)\q_t(A)
%    -
%    \gamma\q_{t-1}(A),
%\end{align*}
%up to the already-suppressed terms involving \(M_D\) or
%\(\frac{1}{d}I\).
%
%Before classifying the resulting summands, we group the scalar recurrence
%terms with their corresponding backtracking contributions. Namely, the
%half-diamond and half-flat contributions are grouped with
%\(-\gamma\q_t(A)\), while the full-diamond contributions are grouped with
%\(-\gamma\q_{t-1}(A)\). By
%\Cref{prop:half-diamond-cancellation,prop:diamond-cancellation}, this
%grouping cancels the leading contribution of every isolated
%well-behaved backtracking pair up to diagonal-residue.
%
%Consider now a term obtained by attaching a new gadget
%\(\sigma\in\{\alpha,\beta\}\) to a record
%\[
%    \tau_t\cdot\tau_{t-1}\cdot\dots\cdot\tau_1
%\]
%appearing at level \(t\). By the induction hypothesis, the old record is
%either proper or admits an LCP-decomposition. Moreover, by the observation
%above, the tail
%\[
%    \tau_{t-1}\cdot\dots\cdot\tau_1
%\]
%admits an LCP-decomposition.
%
%There are three possibilities. First, suppose that \(\sigma\) concatenates properly with the entire old
%record. If the old record is proper, then
%\[
%    \sigma\circ\tau_t\circ\dots\circ\tau_1
%\]
%belongs to \(\calP_{t+1}(A)\). Otherwise, we append \(\sigma\) as a proper
%block to the existing LCP-decomposition. In either case, the desired
%structure is preserved.
%
%Second, suppose that \(\sigma\) creates a nontrivial intersection that is
%not a designated backtracking intersection with \(\tau_t\). Then
%\[
%    \sigma\cdot\tau_t
%\]
%is an LCP: its prefix consists of the single, hence proper, gadget
%\(\tau_t\), while its terminal gadget \(\sigma\) intersects a gadget that
%appeared earlier in the global record. That earlier gadget may be
%\(\tau_t\), a gadget in the tail, or a gadget in an earlier block, as
%allowed by \Cref{def:local-collision-pieces}. Appending the
%LCP-decomposition of the tail gives an LCP-decomposition of the full
%record.
%
%Finally, suppose that \(\sigma\cdot\tau_t\) forms a well-behaved backtracking pair. If the pair is isolated, its leading contribution is
%canceled by the corresponding scalar recurrence term, and only the
%residual remains from \cref{def:backtracking-residuals}.
%  By definition, this residual is an LCP
%supported on the same two-gadget record. If the pair is non-isolated, then
%some vertex participating in the backtracking block also intersects the
%earlier tail. The cancellation is therefore no longer exact, and the
%surviving grouped term is precisely a non-isolated
%backtracking/diagonal error; again,
%\(\sigma\cdot\tau_t\) is an LCP. In both cases, combining this terminal
%LCP with the LCP-decomposition of the tail gives the desired decomposition
%of the full record.
%
%These cases exhaust all possible interactions of the newly attached
%gadget with the previous record. Hence every surviving proper record in
%\(\q_{t+1}(A)\) belongs to \(\calP_{t+1}(A)\), while every surviving
%non-proper record admits an LCP-decomposition. By our bookkeeping
%convention, each such record contains exactly \(t+1\) underlying gadgets
%from \(\{\alpha,\beta\}\). This completes the induction.
%\end{proof}
%


\subsection{Formal Decomposition of MP Error Terms}

As highlighted in the technical overview, the key point is that certain
backtracking intersections produced by composing two \(\alpha\)-shapes, or
two \(\beta\)-shapes, collapse after summing over their internal labels.
Their leading contributions agree exactly with the lower-order terms
subtracted in the three-term recurrence for the Marchenko--Pastur
polynomials. The remaining terms will be organized into local-collision
pieces.

%From this point, we will focus ignore the $1/d\cdot I_m$ and $M_D$ component in $M$ as they are negligible in norm.  
%
%\paragraph{Backtracking intersections and their residuals.}
%We begin by recalling the three backtracking intersection patterns that
%arise from products involving \(\cm_\alpha\) and \(\cm_\beta\).
%
%\begin{definition}[Half-Diamond Backtracking Intersection]
%Let \(\tau\) and \(\tau'\) be two \(\alpha\)-shapes. We say that
%\(\tau\circ\tau'\) forms a \emph{half-diamond backtracking intersection}
%if
%\begin{enumerate}
%    \item \(U_\tau\not\equiv V_{\tau'}\);
%    \item the two circle vertices of \(\tau\) receive the same unordered
%    pair of labels in \([d]\) as the two circle vertices of \(\tau'\).
%\end{enumerate}
%\end{definition}
%
%\begin{definition}[Full-Diamond Backtracking Intersection]
%Let \(\tau\) and \(\tau'\) be two \(\alpha\)-shapes. We say that
%\(\tau\circ\tau'\) forms a \emph{full-diamond backtracking intersection}
%if
%\begin{enumerate}
%    \item \(U_\tau\equiv V_{\tau'}\);
%    \item the two circle vertices of \(\tau\) receive the same unordered
%    pair of labels in \([d]\) as the two circle vertices of \(\tau'\).
%\end{enumerate}
%\end{definition}
%
%\begin{definition}[Half-Flat Backtracking Intersection]
%Let \(\tau\) and \(\tau'\) be two \(\beta\)-shapes. We say that
%\(\tau\circ\tau'\) forms a \emph{half-flat backtracking intersection} if
%\begin{enumerate}
%    \item \(U_\tau\not\equiv V_{\tau'}\);
%    \item the circle vertex of \(\tau\) receives the same label in
%    \([d]\) as the circle vertex of \(\tau'\).
%\end{enumerate}
%\end{definition}

%We first define the error term arising from the intended cancellation from the well-behaved intersections highlighted above. 
%
%\begin{definition}[Backtracking residuals]
%\label{def:backtracking-residuals}
%Define
%\begin{align*}
%    \Delta_{\mathsf{HalfDiamondInt}}
%    &\coloneqq
%    \cm_{\mathsf{HalfDiamondInt}}
%    -
%    \gamma\cm_\alpha, \\
%    \Delta_{\mathsf{HalfFlatInt}}
%    &\coloneqq
%    \cm_{\mathsf{HalfFlatInt}}
%    -
%    \gamma\cm_\beta, \\
%    \Delta_{\mathsf{DiamondInt}}
%    &\coloneqq
%    \cm_{\mathsf{DiamondInt}}
%    -
%    \gamma I_m.
%\end{align*}
%We call these matrices the \emph{backtracking residuals}. A
%\emph{backtracking-residual term} is any graph-matrix summand appearing
%when one of these residual matrices is expanded in the graph-matrix basis.
%
%\end{definition}
%
%For accounting purpose, every backtracking-residual term retains the original
%two-gadget record of the backtracking pair from which it arises. We show in \cref{sec:MP-def-proof} that these are negligible components.
% even if its
%visible graph becomes shorter after the leading backtracking contribution
%is removed. In particular, a full-diamond pair reduces at leading order to
%the empty record, represented by \(I_m\), while
%\(\Delta_{\mathsf{DiamondInt}}\) is the diagonal fluctuation that remains
%after the leading term \(\gamma I_m\) is canceled.
%
%\begin{proposition}[Half-Diamond and Half-Flat Cancellation]
%\label{prop:half-diamond-cancellation}
%We have
%\[
%    \cm_{\mathsf{HalfDiamondInt}}
%    =
%    \gamma\cm_\alpha
%    +
%    \Delta_{\mathsf{HalfDiamondInt}},
%    \qquad
%    \left\|
%        \Delta_{\mathsf{HalfDiamondInt}}
%    \right\|_{\mathrm{sp}}
%    =
%    o_d(1),
%\]
%and, analogously,
%\[
%    \cm_{\mathsf{HalfFlatInt}}
%    =
%    \gamma\cm_\beta
%    +
%    \Delta_{\mathsf{HalfFlatInt}},
%    \qquad
%    \left\|
%        \Delta_{\mathsf{HalfFlatInt}}
%    \right\|_{\mathrm{sp}}
%    =
%    o_d(1),
%\]
%where
%\[
%    \gamma=\frac{2m}{d^2}.
%\]
%\end{proposition}
%
%\begin{proposition}[Diamond Cancellation]
%\label{prop:diamond-cancellation}
%We have
%\[
%    \cm_{\mathsf{DiamondInt}}
%    =
%    \gamma I_m
%    +
%    \Delta_{\mathsf{DiamondInt}},
%    \qquad
%    \left\|
%        \Delta_{\mathsf{DiamondInt}}
%    \right\|_{\mathrm{sp}}
%    =
%    o_d(1).
%\]
%\end{proposition}
%
%The proofs of these propositions follow by centering the internal label
%sums around their expectations and bounding the resulting graph matrices.
%We defer the details to \Cref{sec:MP-def-proof}.
%
%The important distinction is that the MP recurrence cancels only the
%leading reduced contribution. For a half-diamond or half-flat pair, the
%two-gadget block reduces to \(\gamma\cm_\alpha\) or
%\(\gamma\cm_\beta\), respectively. For a full-diamond pair, the block
%reduces to \(\gamma I_m\). The corresponding backtracking residual remains
%and is treated as an error term.

\paragraph{Extension to higher degree.}
To pass to higher degree, we build each gadget record from right to left.
A new error term is created at the first point where the newly attached
terminal gadget has an unintended vertex intersection with the previously
constructed record. We enforce local injectivity before this terminal
gadget, while allowing the terminal gadget to intersect either the current
locally proper prefix or an earlier block in the global record. The latter
possibility is needed for non-isolated backtracking configurations, in
which a backtracking pair also interacts with an earlier gadget.

We first distinguish the isolated backtracking pairs canceled by the MP
recurrence from the non-isolated intersections that remain as error terms.

\begin{definition}[Well-Behaved Backtracking Intersections]
\label{def:well-behaved-backtracking}
Let
\[
    \tau_t\cdot\tau_{t-1}\cdot\dots\cdot\tau_1,
    \qquad
    \tau_i\in\{\alpha,\beta, \frac{1}{d} I_m, M_D\},
\]
be a gadget record. We say that an adjacent pair
\(
    \tau_s\cdot\tau_{s-1}
\)
forms a \emph{well-behaved backtracking intersection} if
\begin{enumerate}
    \item \(\tau_s\) and \(\tau_{s-1}\) form one of the designated
    backtracking intersections: half-diamond, half-flat, or full-diamond;
    \item the vertices participating in this backtracking block have no
    additional vertex intersections with the earlier prefix
    \(
        \tau_{s-2}\cdot\dots\cdot\tau_1.
    \)
\end{enumerate}
If the first condition holds but the second fails, we call
\(\tau_s\cdot\tau_{s-1}\) a \emph{non-isolated backtracking pair}.
These are the three-way backtracking/diagonal intersections.
\end{definition}

\begin{definition}[Local-Collision Pieces for MP Polynomials]
\label{def:local-collision-pieces}
Fix a gadget record
\[
    \tau_t\cdot\tau_{t-1}\cdot\dots\cdot\tau_1,
    \qquad
    \tau_i\in\{\alpha,\beta, \frac{1}{d} I_m, M_D\}.
\]
A consecutive sub-record
\(
    R
    =
    \tau_b\cdot\tau_{b-1}\cdot\dots\cdot\tau_a
\)
is called a \emph{local-collision piece}, abbreviated \emph{LCP}, if it
satisfies the following conditions:
\begin{enumerate}
    \item Each constituent gadget is a gadget in $A$, i.e., from $\{\alpha,\beta, \frac{1}{d} I_m, M_D\}$.

    \item The concatenation preceding the terminal gadget is locally
    proper; namely,
    \(
        \tau_{b-1}\circ\tau_{b-2}\circ\cdots\circ\tau_a
    \)
    has no vertex intersections beyond the prescribed concatenation
    boundaries.

    \item The terminal gadget is either $\frac{1}{d}I_m $ or $M_D$, or a gadget \(\tau_b\) that has a nontrivial vertex
    intersection with a gadget that appeared earlier in the global record.
    This earlier gadget may lie inside the same block or inside a preceding
    block of the global decomposition.

    \item If the terminal interaction is an isolated well-behaved
    backtracking pair
    \(\tau_b\cdot\tau_{b-1}\), then its leading reduced contribution is
    canceled by the MP recurrence and is not regarded as an error term.
    Only a backtracking-residual term in the sense of
    \Cref{def:backtracking-residuals} is retained, supported on the same
    two-gadget record
    \(\tau_b\cdot\tau_{b-1}\).
\end{enumerate}
The length of an LCP is the number of underlying gadgets in its record.
We call \(\tau_a\) its start and \(\tau_b\) its terminal gadget. 
\end{definition}

\begin{definition}[LCP-Decomposition]
\label{def:lcp-decomposition}
Let \(\tau\) be an error term with underlying gadget record
\[
    \tau_t\cdot\tau_{t-1}\cdot\dots\cdot\tau_1.
\]
We say that \(\tau\) admits an \emph{LCP-decomposition} if its record can
be written as
\[
    \tau_t\cdot\tau_{t-1}\cdot\dots\cdot\tau_1
    =
    P\cdot R_r\cdot R_{r-1}\cdot\dots\cdot R_1,
\]
where \(P\) is a possibly empty proper concatenation and each \(R_j\) is
an LCP.

Equivalently, after removing a possibly empty leading proper run, the
remaining gadget record is a concatenation of LCPs.
\end{definition}

It should be noted that we allow vertex intersections between different blocks. The only forbidden
configuration is an isolated backtracking pair across the
boundary of two consecutive blocks---such a pair is canceled at leading
order by the MP recurrence, and any surviving backtracking-residual term is
included in an LCP supported on the original two-gadget record.

\begin{observation}[Terminal Interaction of an LCP]
\label{obs:lcp-terminal-not-isolated-backtracking}
Let
\(
    R
    =
    \tau_b\cdot\tau_{b-1}\cdot\dots\cdot\tau_a
\)
be an LCP. Then one of the following holds:
\begin{enumerate}
    \item \textbf{Non-backtracking intersection:}
    \(\tau_b\) has a nontrivial non-backtracking intersection with an
    earlier gadget in the record;

    \item \textbf{Non-isolated backtracking intersection:}
    \(\tau_b\cdot\tau_{b-1}\) is a non-isolated backtracking pair, i.e.,
    a three-way backtracking/diagonal interaction;

    \item \textbf{Backtracking residual:}
    the terminal two-gadget block
    \(\tau_b\cdot\tau_{b-1}\) supports a backtracking-residual term in the
    sense of \Cref{def:backtracking-residuals}.
\end{enumerate}
In particular, the leading reduced contribution of an isolated
well-behaved backtracking pair is never treated as an error term: it is
canceled by the MP recurrence.
\end{observation}

With these definitions, we can state the structural property of the MP
error terms.

\begin{proposition}[Structural Property for Error Terms]
\label{prop:error-term-A-equivalence}
%After placing all terms containing \(M_D\) or \(\frac1d I_m\) into
%\(\mathsf{Error}_t(A)\), and after applying the backtracking cancellations
%from the MP recurrence, every remaining summand in \(\q_t(A)\) has an
%underlying record For 
%<<<<<<< HEAD
For any $t>0$, any term in \(\mathsf{Error}_t(A)\) from the equivalence of MP polynomials and graph matrices of concatenated shapes admits an LCP-decomposition,
%=======
%For any $t>0$, any term in \(\mathsf{Error}_t(A)\) from the equivalence of Chebyshev polynomials and graph matrices of concatenated shapes admits an LCP-decomposition,
%>>>>>>> 258cc9d9278f3d3112b6597fe09243ecb5da8fa1
\[
    \tau_t\cdot\tau_{t-1}\cdot\dots\cdot\tau_2\cdot\tau_1,
    \qquad
    \tau_i\in \{\alpha,\beta, \frac{1}{d} I_m, M_D\}.
\]
%The summands whose records are proper concatenations are exactly the
%desired terms in \(\fp_t(A)\). Every remaining non-proper summand, i.e.,
%every MP-concatenation error term, admits an LCP-decomposition.
\end{proposition}

\begin{proof}
We prove the statement by induction on \(t\). Throughout the proof, terms
containing \(M_D\) or \(\frac1d I_m\) are placed directly into
\(\mathsf{Error}_t(A)\). Moreover, whenever an isolated backtracking pair
is replaced by a backtracking-residual term, that residual retains the
original two-gadget record.

We first note that deleting the leftmost gadget from a record admitting an
LCP-decomposition preserves an LCP-decomposition. Indeed, if the leftmost
block is proper, its remaining sub-record is still proper. If the leftmost
block is an LCP, then its leftmost gadget is its terminal gadget, and
removing it leaves the locally proper prefix from
\Cref{def:local-collision-pieces}.

The claim is immediate for \(t=0,1\). We have
\[
    \q_0(A)=I_m,
\]
and
\[
    \q_1(A)
    =
    A-I_m
    =
    \cm_\alpha+\cm_\beta
    +
    \left(\frac1d I_m+M_D\right),
\]
where the parenthesized term is already included in
\(\mathsf{Error}_1(A)\); both are LCPs of length $1$.

Now fix \(t\geq1\), and assume the proposition holds at levels \(t\) and
\(t-1\). Modulo the already-suppressed terms involving \(M_D\) or
\(\frac1d I_m\), the MP recurrence gives
\begin{align*}
    \q_{t+1}(A)
    &=
    \bigl(A-(1+\gamma)I_m\bigr)\q_t(A)
    -
    \gamma\q_{t-1}(A) \\
    &=
    \bigl(\cm_\alpha+\cm_\beta-\gamma I_m\bigr)\q_t(A)
    -
    \gamma\q_{t-1}(A).
\end{align*}

Thus every new non-scalar multiplication term is obtained by attaching a
new gadget
\[
    \sigma\in \{\alpha,\beta, \frac{1}{d} I_m, M_D\} 
\]
to the left of a record
\(
    \tau_t\cdot\tau_{t-1}\cdot\dots\cdot\tau_1
\)
appearing at level \(t\). By induction, the old record admits an LCP-decomposition, and its tail
\(
    \tau_{t-1}\cdot\dots\cdot\tau_1
\)
also admits an LCP-decomposition. 

We now consider the cases.

\medskip
\noindent
\textbf{Case 0: \(\sigma \notin \{\al,\beta\} \).} 

In this case, we view this as a gadget with negligible contribution from either $M_D$ or $1/d I_m$. This is a terminal gadget by itself, or it forms an LCP as a terminal gadget with a sequence of properly concatenated shapes.

\medskip
\noindent
\textbf{Case 1: \(\sigma\) concatenates properly.}

Suppose that \(\sigma \in \{\al,\beta\} \) has no unintended intersection with the old
record. If the old record is proper, then
\[
    \sigma\circ\tau_t\circ\dots\circ\tau_1
\]
belongs to \(\calP_{t+1}(A)\) (and does not contribute to the error term). Otherwise, we append \(\sigma\) as a proper
block to the existing LCP-decomposition.  

\medskip
\noindent
\textbf{Case 2: \(\sigma\) creates a non-backtracking collision.}

Suppose that \(\sigma\) has a nontrivial intersection with the old record
that is not a designated backtracking intersection with \(\tau_t\). Then
\[
    \sigma\cdot\tau_t
\]
is an LCP: its prefix consists of the single, hence proper, gadget
\(\tau_t\), while its terminal gadget \(\sigma\) intersects a gadget that
appeared earlier in the global record. This earlier gadget may be
\(\tau_t\), a gadget in the tail, or a gadget in an earlier block, as
allowed by \Cref{def:local-collision-pieces}. Combining this terminal LCP
with the LCP-decomposition of the tail gives an LCP-decomposition of the
full record.

\medskip
\noindent
\textbf{Case 3: \(\sigma\cdot\tau_t\) forms a designated backtracking pair.}

First suppose that the pair is well-behaved, and hence isolated from the
earlier tail. If it is a half-diamond or half-flat pair, its leading
two-gadget contribution reduces to
\[
    \gamma\cm_\alpha
    \qquad\text{or}\qquad
    \gamma\cm_\beta,
\]
respectively. After concatenating with the tail, this is exactly the
matching contribution canceled by \(-\gamma\q_t(A)\).

If it is a full-diamond pair, its leading contribution is
\(\gamma I_m\). The two leading \(\alpha\)-gadgets therefore reduce to the
empty record, leaving
\[
    \gamma\,
    \cm_{\tau_{t-1}\cdot\dots\cdot\tau_1},
\]
which is precisely the
matching contribution canceled by \(-\gamma\q_{t-1}(A)\).

Thus, in every isolated backtracking case, the leading reduced
contribution cancels exactly. The only surviving summands are the
backtracking-residual terms from
\Cref{def:backtracking-residuals}. In the full-diamond case, these are the
diagonal backtracking-residual terms. By definition, each residual retains
the original record \(\sigma\cdot\tau_t\), and hence forms a terminal LCP.

Finally, suppose that \(\sigma\cdot\tau_t\) is non-isolated. Then some
vertex participating in the backtracking block also intersects the earlier
tail. This is not a residual from an isolated cancellation; it is a genuine
three-way backtracking/diagonal collision. Consequently,
\(\sigma\cdot\tau_t\) is again a terminal LCP. Combining it with the
LCP-decomposition of the tail gives the desired decomposition of the full
record.

\medskip

The scalar terms in the MP recurrence introduce no new vertex
intersections. Their role in the grouping above is to cancel the leading
contributions of the isolated backtracking pairs. Any pre-existing error
summands appearing in those scalar terms already admit
LCP-decompositions by the induction hypothesis.

The three cases exhaust all possible interactions of the newly attached
gadget with the previous record. Hence every surviving proper record in
\(\q_{t+1}(A)\) belongs to \(\calP_{t+1}(A)\), while every surviving
non-proper \(\alpha/\beta\)-record admits an LCP-decomposition. Under our
bookkeeping convention, each such record contains exactly \(t+1\)
underlying gadgets from \(\{\alpha,\beta\}\). This completes the induction.
\end{proof}



\subsection{Block-Value Bound for Local-Collision-Pieces}

\paragraph{Factor Assignment Scheme}
The crux of our argument is the following factor-assignment scheme. We introduce the following categorization for each step of the walk: we call it an $F/L$ step if it is along an edge appearing for the first/last time, and an $H$-step if neither (i.e., middle appearance).

\begin{mdframed}[frametitle = Combinatorial Factor Assignment for Vertex Factors] \label{prop:vtx-assignment}
 Each vertex $i$ requires a factor $\wt{(i)}$ when it first appears in the walk (namely a factor of $d$ for a circle vertex, and a factor of $m$ for a square vertex) and a subsequent factor of $O(q)$ when it appears as an incoming active-vertex of some $F$ edge in the walk. We assign its corresponding factor  as follows,
\begin{enumerate}
	\item Assign a factor of $\sqrt{\wt(i) } $ for the step if the vertex is appearing for the first time;
	\item Assign a factor of $\sqrt{\wt(i)}$ for the step if the vertex is appearing for the last (final) time;
	\item Assign a factor of $2q \cdot |V(\tau)|$ if the vertex is not incident to the vertex-component of final-appearance vertices connected to the current walk boundary $U_\tau$. (For example, the vertex makes a middle appearance as the destination of an $F$- or $H$-edge.) 
%	 vertex is making a middle appearance, and it is not incident to any \emph{vertex-component} making final appearances connected to the walk boundary $U_\tau$ .	
\end{enumerate}
%Moreover, for completeness, for vertices incident to some $H$-edge, i.e., edge appearing for the middle time, or an $L$-edge but not in ideal-step condition, assign a total cost of $O_r(q)$ to all the (incidental) incoming active vertices.
\end{mdframed}


 
Besides the vertex factor, we also consider the following scheme for edge-factor assignment.
\begin{mdframed}[frametitle = Analytical Factor Assignment for Edges] \label{prop: edge-value-assignment}
Each edge (random variable $g_e$ that appears for an even number of times) contributes an analytical value of $1$, and we assign this value to each individual step that traverses along edge $e$ via factor $B_{ev}(i,e)$---the factor assigned to step-$i$ from edge-$e$ as follows.
\begin{enumerate}
	\item If the appearance of random variable $e$ contains $\geq 2$ copies of an $h_1$ edge, assign the first and final copy a factor of $\frac{1}{\sqrt{d}}$, and assign a factor of $\frac{2q}{\sqrt{d}}$ for any middle appearance;
	\item  If the appearance of random variable $e$ contains $\geq 2$ copies of an $h_2$ edge, assign the first and final copy a factor of $\frac{\sqrt{2}}{d} $;
\end{enumerate}	
\end{mdframed}

\begin{observation}
Let \(x\sim N(0,1/d)\), and let
\[
    h_1(x)=x,\qquad h_2(x)=x^2-\frac1d .
\]
For any \(t_1,t_2\ge 0\) with \(t_1+t_2>0\),
\[
    \E\!\left[h_1(x)^{t_1} h_2(x)^{t_2}\right]\neq 0
\]
only if \(t_1\ge 2\) or \(t_2\ge 2\).
Equivalently, if \(t_1,t_2\in\{0,1\}\) and \(t_1+t_2>0\), then
\[
    \E\!\left[h_1(x)^{t_1} h_2(x)^{t_2}\right]=0.
\]
\end{observation}
 
\begin{lemma}[Block-Value Bound for LCP]  \label{lem:LCP-bound}
	For any $t>0$, and $\tau$ an LCP of length-$t$, we have \[ 
	B(\tau) \leq \frac{\poly(t, q) }{\sqrt{d}} \cdot \sqrt{\gamma}^{t-1} \,.
	\]
\end{lemma}

\begin{proof}[Proof of \cref{lem:LCP-bound}]
	For any LCP $\tau$, we decompose it as $\tau = \tau_I \cdot \tau_R$ where $\tau_I$ is the terminal gadget, and $\tau_R$ is the properly concatenated part. The block-value bound then follows from the next two claims, as we take \[
	B_q(\tau)\leq B_q(\tau_R) \cdot B_q(\tau_I) \leq \frac{\poly(t,q)}{\sqrt{d}} \cdot \sqrt{\gamma}^{t-1}
	\]
\end{proof}


%\begin{claim} For a given LCP $\tau = \tau_{t} \cdot \dots \tau_1$,
%	let $\tau_R = \tau_{t-1} \circ \tau_{t-2} \circ \dots \circ \tau_1$ be the locally proper piece, we have \[ 
%	B_q(\tau_R) \leq  4 \cdot (t-1) \cdot \sqrt{\gamma }^{t-1} 	\]
%\end{claim}
%
%\begin{proof}
%	We apply block-value bound, and for any segment of $t-1$ length, call a vertex a separator vertex if a vertex makes a prior and subsequent appearance beyond the current interval (i.e., a vertex making a middle appearance). Call such vertex a separator vertex. 
%	
%	Observe that there is at least one gadget among the $t-1$ blocks such that any path from the left square vertex to the right square vertex is blocked by the separator vertices. Call this gadget an anchor gadget.
%	
%	We next observe that for any gadget to the left of the "anchor" gadget, no vertex can be making its first appearance, and similarly, any vertex to the right of the anchor gadget cannot be making its final appearance. Therefore, any such $\al$-gadget contributes a block-value of at most \[ 
%	2\cdot \sqrt{\frac{m \binom{d}{2}}{d^2}} = \sqrt{\frac{2m}{d^2}} = \sqrt{\gam }\,.
%	\] 
%	On the other hand, for any $\beta$-gadget, we have \[ 
%	\left(\sqrt{\frac{2}{d^2}}\right)^2 \cdot \sqrt{md} = o_d(1) \,.
%	\]
%	\begin{remark}
%		It is crucial that the $O(1)$ value labeling from $\beta$-gadget requires the local gadget to be an anchor gadget (since the separator vertex has to be in the middle).
%	\end{remark}
%	Since there are at most $t-2$ such blocks other than the gadget block, we have a bound of \[ 
%	B_q( \textsf{Non-Anchor} ) \leq \sqrt{\gamma}^{t-2} 
%	\]
%	from the non-anchor gadgets.
%	
%	For the anchor gadget, we have a block-value of at most \begin{align*}
%	B_q(\textsf{Anchor}(\al))& \leq  (t-1) \cdot (1+o_d(1))\cdot  \left( 2 \cdot 2\cdot \sqrt{\frac{m \binom{d}{2}}{d^2}}  + \frac{m \binom{d}{2}}{d^2} \right) 	  \\
%	&\leq (1+o_d(1)) \cdot (t-1) (2\sqrt{\gamma}+\gamma)	
%	\end{align*}
%	
%	since we have $t-1$ choices for such gadget, and for each $\al$ gadget, we have $2$ sets of vertex separators via a single square vertex, and the choice of taking the $2$ middle circle vertices as the separator vertices (up to choices with $O(\frac{1}{\sqrt{d}})$ gap.)
%	
%	Similarly, for an anchor-gadget being a $\beta$-gadget, we have block-value of \[ 
%	B_q(\textsf{Anchor}(\beta)) \leq (1+o_d(1)) \cdot  (t-1) \cdot \frac{2 \cdot m}{d^2}  = (1+o_d(1)) \cdot  (t-1)\cdot \gamma  
%	\]
%	
%	Combining the above gives us a total block-value bound of \begin{align*} 
%	B_q(\tau_R) &\leq B_q(\textsf{Non-Anchor}) \cdot B_q( \textsf{Anchor} ) \\&\leq  (1+o_d(1))^{t-1}  \sqrt{\gamma}^{t-2} \cdot ((t-1) \cdot (2\sqrt{\gamma } +2\gamma)\\  &\leq 4(t-1)\sqrt{\gamma}^{t-1} 	
%	\end{align*} 
%	assuming $\gamma<1$.
%	\end{proof}

%\subsection{reorg}
\begin{claim}[Block value of a locally proper concatenation]
\label{clm:block-value-proper-prefix}
Let
\[
    \tau=\tau_t\cdot \tau_{t-1}\cdot \dots \cdot \tau_1
\]
be an LCP of length \(t\), and let
\[
    \tau_R
    :=
    \tau_{t-1}\circ \tau_{t-2}\circ \dots \circ \tau_1
\]
be its locally proper concatenation. Assume \(\gamma<1\). Then
\[
    B_q(\tau_R)
    \le
    (1+o_d(1))^t\cdot 4(t-1)\sqrt{\gamma}^{\,t-1}.
\]
for any $q\ll d$.
\end{claim}

\begin{claim}[Block value bound for collision gadget] \label{clm:collision-gadget-block-bound}
	Let $\tau_I$ be a collision gadget (i.e., a terminal gadget in an LCP); we have \[
	B_q(\tau_I) \leq  \frac{\poly(q)}{\sqrt{d}} \cdot \sqrt{\gamma}
	\]
	for   $q\ll d^{\delta}$ for some $\delta>0$
	.
\end{claim}

\begin{proof}[Proof of \cref{clm:block-value-proper-prefix}]
We apply the block-value factor assignment to the proper concatenation
\[
    \tau_R=\tau_{t-1}\circ \tau_{t-2}\circ \dots \circ \tau_1.
\]
For each interval, we call a vertex a \emph{separator vertex} for the gadget if it makes both
a prior and a subsequent appearance outside this interval; in
particular, it is a middle-appearance vertex for the block-value applied to the whole interval of $\tau_R$.
%
We first observe that since each edge appears at least twice in the walk, there exists at least one gadget among the \(t-1\)
gadgets of \(\tau_R\) that has its left boundary disconnected from its right boundary by the separator of $\tau_R$. Call such a gadget an \emph{anchor} gadget, and the rest \emph{non-anchor}.
 



%Indeed, consider the component of final-appearance vertices connected to the
%left square boundary. Since the right square boundary cannot be a
%final-appearance vertex in the local block-step accounting, this component
%cannot contain the right boundary. Therefore its boundary inside the proper
%chain gives a separator. In our shapes, such a separator is realized either by
%a square vertex or by the two circle vertices of an \(\alpha\)-gadget.
%
%Next, consider any non-anchor gadget. A gadget strictly to the left of the
%anchor cannot contain a first-appearance active vertex; otherwise that vertex
%would be a middle-appearance vertex relative to the separator orientation, and
%we would gain an additional gap. Similarly, a gadget strictly to the right of
%the anchor cannot contain a final-appearance active vertex, unless we already
%obtain such a gap. Since these gap terms are smaller, it suffices to bound the
%dominant configurations in which every non-anchor gadget is fully oriented as
%final on the left side or fully oriented as fresh on the right side.

For a non-anchor \(\alpha\)-gadget, the block value is at most
\[
B_q(\al \text{-non-anchor}) =     2\cdot \frac{\sqrt{m\binom d2}}{d^2}
    =
    (1+o_d(1))\sqrt{\frac{2m}{d^2}}
    =
    (1+o_d(1))\sqrt{\gamma}\,.
\]
On the other hand, a non-anchor \(\beta\)-gadget is negligible: its two
\(h_2\)-edges contribute \(2/d^2\), while the dominant vertex contribution is
at most \(\sqrt{md}\). Hence
\[
    B_q(\beta\text{-Non-Anchor})
    \le
    \frac{2\sqrt{md}}{d^2}
    =
    o_d(1),
\]
in the regime \(m=O(d^2)\). In particular, it is bounded by
\((1+o_d(1))\sqrt{\gamma}\). Therefore, since there are at most \(t-2\)
non-anchor gadgets, we have
\[
    B_q(\text{non\mbox{-}anchor})
    \le
    (1+o_d(1))^{t-2}\sqrt{\gamma}^{\,t-2}.
\]

It remains to bound the anchor contribution. There are at most \(t-1\) choices
for the anchor gadget.

First suppose the anchor is an \(\alpha\)-gadget. For each such gadget, there
are two possible square-separator choices, and one possible two-circle
separator choice. The square-separator choices contribute at most
\[
    2\cdot
    \left(
        2\cdot \frac{\sqrt{m\binom d2}}{d^2}
    \right)
    =
    (1+o_d(1))\,2\sqrt{\gamma}.
\]
The two-circle separator contributes at most \(\gamma\), as we have
\[
    (1+o_d(1)) \cdot 2 \cdot \frac{m}{d^2} = (1+o_d(1))  \cdot \gamma\,.
\]
 Hence,
\[
    B_q(\alpha\text{-anchor})
    \le
    (1+o_d(1))(t-1)(2\sqrt{\gamma}+\gamma).
\]

Similarly, if the anchor is a \(\beta\)-gadget, then its middle circle vertex
serves as the separator. The two square-side contributions give
\[
    B_q(\beta\text{-anchor})
    \le
    (1+o_d(1))(t-1)\frac{2m}{d^2}
    =
    (1+o_d(1))(t-1)\gamma.
\]
Combining the two anchor possibilities, we obtain 
\[
    B_q(\text{anchor})
    \le
    (1+o_d(1))(t-1)(2\sqrt{\gamma}+2\gamma).
\]

Putting together the anchor and non-anchor block-value bounds,
\[
\begin{aligned}
    B_q(\tau_R)
    &\le
    B_q(\text{non\mbox{-}anchor})\cdot B_q(\text{anchor}) \\
    &\le
    (1+o_d(1))^{t-1}
    \sqrt{\gamma}^{\,t-2}
    \cdot
    (t-1)(2\sqrt{\gamma}+2\gamma).
\end{aligned}
\]
Since \(\gamma<1\), we have \(\gamma\le \sqrt{\gamma}\), and hence
\[
    2\sqrt{\gamma}+2\gamma
    \le
    4\sqrt{\gamma}.
\]
Therefore
\[
    B_q(\tau_R)
    \le
    (1+o_d(1))^{t-1}\cdot 4(t-1)\sqrt{\gamma}^{\,t-1}.
\]
Weakening \((1+o_d(1))^{t-1}\le (1+o_d(1))^t\), we obtain
\[
    B_q(\tau_R)
    \le
    (1+o_d(1))^t\cdot 4(t-1)\sqrt{\gamma}^{\,t-1}.
\]
This proves the claim.
\end{proof}

\begin{proof}[Proof of~\cref{clm:collision-gadget-block-bound}]
	Recall from the structural property of LCP decomposition and \cref{obs:lcp-terminal-not-isolated-backtracking}, the final collision gadget satisfies one of the following properties,
	\begin{enumerate}
    \item it has a nontrivial non-backtracking intersection with some
    earlier gadget in the record; or

    \item \(\tau_b\cdot\tau_{b-1}\) forms a non-isolated backtracking pair,
    i.e., a three-way backtracking/diagonal interaction; or

    \item  $\tau_b$ is the residual left after an isolated
    backtracking cancellation.
\end{enumerate}

\paragraph{Block value bound for non-backtracking intersection}
In this case, the edge-value is $O(\frac{q^4}{d^2})$ while we have a vertex factor of either \( 
\sqrt{md^2} 
\) or \( \sqrt{m^2} \) provided there is no additional vertex making a middle appearance. However, by the property of non-backtracking intersection, we have at least one extra vertex making a middle appearance, and any such appearance gives a gap of $O(\frac{q}{\sqrt{d}} )$, hence we have \[ 
B_q( \text{Non-backtracking Intersection Gadget} ) \leq O(\frac{t^5}{\sqrt{d}} \cdot \sqrt{\gamma} )\,.
\]

\paragraph{Block value bound for three-way diagonal intersection}
The same bound holds as we observe that we would have a bound of $O(1)$ if every vertex is making a first and final appearance within the backtracking intersection, while any vertex that has additionally made a prior appearance gives a slack of \[ 
\frac{1}{\sqrt{d}} \cdot O(q)\,,
\]
again giving a loose bound of
 \[ 
B_q( \text{Three-way Diagonal Gadget} ) \leq O(\frac{t^5}{\sqrt{d}} \cdot \sqrt{\gamma} )\,.
\]
\paragraph{Block value bound for backtracking residue}
For the half-diamond and half-flat residuals, the expansions in
~\cref{sec:MP-def-proof} show that every  summand either
contains an additional centered \(h_2\)-edge, together with an explicit
\(1/d\)-prefactor, or is an \(O(d^{-2})\)-multiple of
\(\cm_\alpha\) or \(\cm_\beta\). Thus every residual summand gains at
least \(d^{-1/2}\) relative to its uncentered backtracking block. The
full-diamond residual is diagonal and contains a centered \(h_2\)
fluctuation, giving the same gain. Therefore
\[
    B_q( \text{Backtracking Residue})
    \le
    \frac{\poly(t,q)}{\sqrt d}
    (\sqrt\gamma)^2.
\]

\paragraph{Negligible terminal gadgets.}
Finally,
\[
    B_q\!\left(\frac1d I_m\right)=\frac1d,
    \qquad
    B_q(M_D)=\widetilde O(d^{-1/2}),
\]
so these terminal cases satisfy the same claimed bound.

%\todo{fill in block value bounds for the other two cases}
\end{proof}


Finally, we wrap up this section by completing the proof of the spectral norm bound of the error terms when we switch from orthogonal polynomials to graph matrices of concatenated shapes.
\begin{proof}[Proof of \cref{lem:mp-error-quant}]
	We use the factor assignment scheme developed in the series of works for obtaining tight norm bounds for graph matrices \cite{JPRTX, HKPX23, KPX24, Xu25, KX26, PotechinXu2026Theta}, which shows that, for the aforementioned factor assignment scheme, we have \[
	\E\Tr[ \left(\text{Error}_t \cdot \text{Error}_t^\top\right )^q  ] \leq \text{matrix-dimension} \cdot B_q(\text{Error}_t)^{2q}\,.
	 \] 
	 Plugging in $B_q(\text{Error}_t) = O(\frac{\poly(t)}{\sqrt{d}}) $ from the above bound  gives us \[ 
	 \E\Tr[ \left(\text{Error}_t \cdot \text{Error}_t^\top\right )^q  ] \leq  m  \cdot  \left(O(\frac{\poly(t)}{\sqrt{d}}) \right)^{2q}\,.
	  \] 
	  Finally, setting $q= \Theta(\log^4 d)$ gives us the desired bound by~\cref{claim:trace-to-norm-rough}.
	  \end{proof}

%\begin{proof}
%	Recall that a locally-collision-piece of length-$t$ contains a concatenation of $t$ shapes from $\{\al,\beta\}$ as defined in \cref{def:local-collision-pieces}. For starters, consider the block-value restricted to the properly concatenated chunk of $\tau$, and observe that the block-value restricted to this chunk is at most \[ 
%	(2(t-1))^2 \cdot (\sqrt{\gamma})^{t-1}
%	\]
%	To see this, first note that in the dominant case (up to constant dependence), we either have separator vertex being a single square vertex, or two circle vertices from one of the gadgets. Thus we have at most $(2(t-1))^2 $ for determining the locations of the first two separator vertices. Once these two vertices are decided, consider a traversal from these two vertices to orient the edges.
%
%
%	\end{proof}
%
%\begin{claim}
%	
%\end{claim}

%\subsection{Shifting From $A$ to $M$}
%
%\paragraph{Obtaining Inverse of $M$ via Woodbury Identity}
%The decomposition of $M$ into $A$ and a rank-2 matrix $B$ (\cref{eq:M-decomposition}) allows us to apply the Woodbury matrix identity about the inverse of low-rank corrections of invertible matrices.
%
%\begin{fact}[Matrix Invertibility] \label{fact:invertibility}
%    Suppose $A \in \R^{n_1 \times n_1}$ and $C \in \R^{n_2 \times n_2}$ are both invertible matrices, and $U\in \R^{n_1 \times n_2}$ and $V \in \R^{n_2 \times n_1}$ are arbitrary.
%    Then, $A + U C V$ is invertible if and only if $C^{-1} + V A^{-1} U$ is invertible.
%\end{fact}
%
%\begin{fact}[Woodbury matrix identity~\cite{Woo59}]  \label{fact:woodbury}
%    Suppose $A \in \R^{n_1 \times n_1}$ and $C \in \R^{n_2 \times n_2}$ are both invertible matrices, and $U\in \R^{n_1 \times n_2}$ and $V \in \R^{n_2 \times n_1}$ are arbitrary. Then
%    \[ 
%    (A+ UCV)^{-1} = A^{-1} - A^{-1}U\left(C^{-1}+ VA^{-1}U\right)^{-1}VA^{-1}  \,.
%    \]
%\end{fact}
%
%In light of \Cref{fact:woodbury}, we can write $B$ as $B = UCU^T$ where $U = V^T = \frac{1}{\sqrt{d}} \begin{bmatrix} 1_m & \eta \end{bmatrix} \in \R^{m \times 2}$ and
%$C = \begin{bmatrix} 1 & 1 \\ 1 & 0 \end{bmatrix}$,
%and $M = A + UCU^T$.
%Note that $C^{-1} = \begin{bmatrix} 0 & 1 \\ 1 & -1 \end{bmatrix}$, and we have
%\begin{align*}
%    C^{-1} + U^T A^{-1}U  = \begin{bmatrix}
%	\frac{1_m^T A^{-1} 1_m}{d} & 1 + \frac{\eta^T A^{-1} 1_m}{d}\\\\
%	1+ \frac{\eta^T A^{-1} 1_m}{d} & -1 + \frac{\eta^T A^{-1} \eta }{d} 
%    \end{bmatrix}
%    \eqqcolon
%    \begin{bmatrix}
%    	r & s \\
%    	s & u
%    \end{bmatrix}
%    \,. 
%    \numberthis \label{eq:russ}
%\end{align*}
%
%Assuming $A$ is shown to be invertible, next, from \Cref{fact:invertibility}, we can prove that $M$ is invertible  by showing that the $2\times 2$ matrix $C^{-1} + U^T A^{-1}U$ is invertible, which is in fact equivalent to $ru - s^2 \neq 0$ from the above definitions.
%
%%\begin{claim}[$M$ is invertible]
%%	\[
%%	ru - s^2\enq 0
%%	 \]
%%\end{claim}
%%\todo{fill in the proof - should follow from this subsequent meta-lemma about the scalars }
%
%
%%\begin{lemma} \label{lem:hyper-parameter-bounds}
%%	We give the following bounds for the scalars defined above.\[ 
%%	 r= (1+o_d(1)) \cdot  \frac{m}{d} \cdot \frac{1}{1-\gamma} \,,
%%	\]
%%	\[
%%	s = 1+ o_d(1)\,,
%%	 \]
%%	 and \[
%%	 u = -1 +\gam +o_d(1) \,.
%%	  \]
%%	  \jnote{the extra $\gamma$ here is crucial}
%%\end{lemma}
%
%\begin{restatable}[Scalar bounds]{lemma}{HyperParameterBounds}
%\label{lem:hyper-parameter-bounds}
%Let \(r,s,u\) be the scalars defined above. Then, with high probability,
%\[
%    r\coloneqq \frac{1_m^T A^{-1} 1_m}{d}
%    =
%    (1+o_d(1))\cdot \frac{m}{d}\cdot \frac{1}{1-\gamma},
%\]
%\[
%    s \coloneqq 1+ \frac{\eta^T A^{-1} 1_m}{d}  = 1+o_d(1),
%\]
%and
%\[
%    u \coloneqq -1 + \frac{\eta^T A^{-1} \eta }{d}  = -1+\gamma+o_d(1).
%\]
%%In particular, the leading correction in \(u\) is the \(+\gamma\) term.
%\end{restatable}
%
%\begin{corollary}[$M$ is invertible] W.h.p. \[ 
%ru-s^2 \neq 0\,.
%\]	
%\end{corollary}
%
%%Given $M$ is also invertible, we may now proceed to expand $M^{-1}$ in the graph-matrix basis. 
%Expanding out Woodbury, we have \[ 
%M^{-1} = A^{-1} + \frac{1}{s^2-ru}\cdot A^{-1} \cdot (u\cdot \frac{\1_m\1_m^T}{d} -s \frac{\eta\cdot \1_m^T+ \1_m^T \cdot \eta }{d}
% + r \cdot \frac{\eta \cdot \eta^T }{d}) \cdot A^{-1}\,. \numberthis \label{eq:M-inverse}
% \]
%
%
%\begin{proof}[Proof Sketch to \cref{lem:hyper-parameter-bounds}] We defer the formal proof to the appendix, while giving a sketch of that by highlighting the dominant terms. Notice they are all random variables with means given by the following.
%	\paragraph{Analysis for $r$}
%	Recall that $r = \frac{1}{d}  \cdot \1_m^T \cdot A^{-1} \cdot \1_m  $, and \(
%	A^{-1} = \frac{1}{1-\gamma} \sum_{t\geq 0} P_t(A)\,.
%	 \)	
%	 For $t=0$  we have \[
%	 \frac{1}{d} \cdot \1_m^TP_0(A) \1_m = \frac{m}{d}\,,
%	  \] 
%	  and \[ 
%	  \1_m^T P_t(A) \1_m = o_d(1) 
%	  \]
%%	  \todo{fill in their concentration here. - would be great to quantify the $o_d(1)$ slack for example.}
%	  Therefore, we have \[ 
%	  r = \frac{1}{1-\gamma}  \cdot \frac{\1_m^T A^{-1} \1_m }{d} = \frac{1}{1-\gamma} \cdot \frac{m}{d}+o_d(1)\,.
%	  \]
%	  
%	 \paragraph{Analysis for $s$}
%	 Recall that $s = 1+ \frac{1}{d} \cdot \eta^T A^{-1} \1_m$, it suffices for us to bound the second term as $o_d(1)$. Notice for $t=0$, we have \[ 
%	 \E[\eta^T  \cdot  \1_m]  = 0
%	 \]
%%	 \todo{fill in concentration here, and that for higher $t$.}
%	 
%	 \paragraph{Analysis for $u$} Recall that $u = -1 + \frac{1}{d} \eta^T A^{-1}\eta $. For $t=0$, we have \[
%\frac{1}{d} \cdot \eta^\top P_0(A) \eta  = (1+o_d(1)) \cdot \frac{2md }{d^2} = (1+o_d(1)) \gam \,.
%\]
% by analogous calculation for the higher moments. 
%   	
%  	For $t=1$, we have \[
%  	\q_1(A) = \fp_1(A) = \cm_\al + \cm_\beta \,.
%  	 \]
%  	 The crucial observation is $\eta^T \cm_\al \eta = o_d(1)$  while the second term gives $\gamma^2 $ (and importantly with a  minus sign from the alternating sign in the polynomial expansion for $x^{-1}$). 
%%	 \todo{This is done in claim 4.2- would be great to move things over and reorganize }
%	 \end{proof}
%	 
%	 





%\input{variance.tex}
\section{Analysis of the Inner Matrix \texorpdfstring{$Q$}{Q}} \label{sec:inner-matrix-Q}
We analyze the inner matrix in this section.  We start by showing that each term in the iterative construction admits a ``nice'' decomposition into proper shapes.

 There are two more components that we obtain from the decomposition. In the first section, we verify that the variance of the inner matrix is $1$ in the limit at the sharp threshold of $d^2/4$. This allows us to ultimately obtain a positive definite %P.D. 
 matrix when we truncate at finite levels in the next section when we additionally allow some small, vanishing slack to the threshold.

The second part of this section establishes the free independence of backbone-dangling shapes. In particular, it gives a tight norm bound for any linear combination of backbone-dangling shapes.
\begin{comment}
\todo{Major Update: the correction operation should be the following. For any $\tau$ that arises in multi-way concatenation at the most recent level, take its correction as \[ 
\corr(\tau) \coloneqq - \frac{1}{1-\gamma}\left( \mathcal{L}^*(M^{-1} \mathcal{L}(\tau))-   \gam \cdot \tau  \right)  
\] 
In other words, we scale the prior update by $\frac{1}{1-\gamma}$ after subtracting  a factor of $\gam \cdot \tau$. 
This again ensures \[
L(\tau+ \corr(\tau)) =0
 \]
 as we have \[ 
 L(\corr(\tau)) = -\frac{1}{1-\gam} L(\tau) +\frac{\gamma}{1-\gamma} L(\tau) =-L(\tau)
 \]
 as achieved by the prior choice of $L(\tau)$.
}

\snote{
We can define 
$$\corr(\tau)=-\frac{\pi_0(\tau)-\gamma\tau}{1-\gamma}$$
using Madhur's projection $\pi_0$. I still think it's clearer this way, rather than using the $\mathcal{L}^*M^{-1} \mathcal{L}$ composition.
}
\jnote{sure, we can add a remark saying that this is simply the projection to the rank-$1$ forms (with the subscrip-$0$ dropped) }
\end{comment}
%\input{variance.tex}



\subsection{Deviation Attachment and Vertical Concatenation}
In this section, we show that the vertical concatenation - arising from evaluating $M^{-1}$ on the deviation term $\eta_i$ at level-$i$- can indeed be viewed as a proper vertical concatenation where we enforce vertex injectivity across the whole piece. In other words, intersection terms in the vertical attachment are negligible, and the variance of the dominant terms behaves as we anticipate in the technical overview.

 More importantly, we show that it suffices for us to focus on the main term of $A^{-1}$ while the extra low-rank update from Woodbury \cref{fact:woodbury} may be ignored in this process as it corresponds to $o_d(1)$ norm components for $i\geq 1$. 
 
% In addition to verifying the structural property that each term is a backbone-dangling shape, the direct expansion in graph matrix also enables to verify the evolution of the variance across the iterations.
%
% It is crucial to note that such approximation does not hold for the base iteration of $Q_0$, which in part explains the reason we do not compute directly the variance of $Q_0$ via graph matrix decompositions. Fortunately and intriguingly, such complication does not propagate further for the subsequent levels. 
% 
 
 \begin{lemma}[Linearizing Deviation Attachment into Proper Shapes]  \label{lem:vertical-concatenation}
 Given a backbone-dangling shape $\tau$, its deviation vector $\eta_\tau = L(\tau) \in \R^m $, and $w_\tau = M^{-1} \eta_\tau$, write \[ 
\corr(\tau) \coloneqq \frac{1}{1-\gamma} \left( \mathcal{L}^*(w_\tau ) - \gam \cdot \cm_\tau  \right)=   \mathsf{Proper}(\eta_\tau) + \mathsf{Error}(\eta_\tau)\,,
 \]
 where $\mathsf{Proper}(\eta_\tau)$ is a weighted linear combination of proper backbone-dangling shapes,	
 and $\mathsf{Error}$ are the remaining terms. We have
 \[ 
 \Var( \mathsf{Proper}(\eta_\tau))  = \frac{1}{1-\gam}\cdot \Var(\eta_\eta)+o_d(1)=  \frac{\gamma}{1-\gam}\cdot \Var(\tau) + o_d(1) 
 \]
 and
  \[
 \|\mathsf{ErrorShapes}\|_{sp} = o_d(1)\,.
  \]
We remind the reader that for a matrix $X$, we define its variance as $\Var(X)\coloneqq \frac{1}{\dim(X)} \cdot \E[\Tr(XX^\top)] $ and analogously for a vector.
  \end{lemma}
  
  Before we proceed with the analysis, we note that there are three facets in the above lemma:\begin{enumerate}
  	\item The target correction term admits a clean decomposition of graph matrices of proper shapes up to negligible error terms;
  	\item Each proper shape in the decomposition is a backbone-dangling shape except the $\gam \cdot \tau$ component;
  	\item  Altogether the proper shapes have variance precisely scaled by an additional factor of $\frac{1}{1-\gamma}$.
  \end{enumerate}  
 We give a sketch of the analysis here while we note that this is an informal version because it has not incorporated polynomial truncation of $M^{-1}$ (specifically that of $A^{-1}$). We additionally address the truncation issue with quantitative slack subsequently in~\cref{sec:formal-truncation}.

%Recall that there are essentially two terms in $M^{-1}$, one from $A^{-1}$ and the additional term from the low-rank updates. For $\eta_\tau $ arising from iterative process (i.e., $\eta_\tau \neq \eta_0$ from the identity-perturbation initialization), the entire dominant contribution comes from $A^{-1}$, and the subsequent contribution from the low-rank update can be ignored. However, for the initialization term $\eta_0$, extra care is warranted from the low-rank component as well, because the low-rank component precisely aligns with $\eta_0$.

%<<<<<<< HEAD
%\paragraph{Analysis of $A^{-1}$ and Well-behaved Intersections} For the analysis of $A^{-1}$ term, there is one component we would like to highlight for this analysis, and precisely why this lemma warrants a formal treatment: not all the intersection shapes in the vertical concatenation of deviation-attachment is negligible! We isolate out the following collection of intersection terms via the following definitions. 
%=======
\paragraph{Analysis of $A^{-1}$ and Well-behaved Intersections} For the analysis of the $A^{-1}$ term, there is one component we would like to highlight for this analysis, and precisely why this lemma warrants a formal treatment: not all the intersection shapes in the deviation-attachment are negligible! We isolate out the following collection of intersection terms via the following definitions. 
%>>>>>>> 258cc9d9278f3d3112b6597fe09243ecb5da8fa1

%\begin{definition}[Deviation Attachment Edges in $\eta_\tau$ ]
%	Given $\eta_\tau$ a deviation-vector, recall that it has a single square vertex on the boundary attached to two distinct circle vertices. Call these edges \emph{attachment-edges}. 
%\end{definition}

Regardless of potential intersection within the shape, terms in $\mathcal{L}^*(w_\tau)$ are obtained from the following diagram operations:
\begin{enumerate}
	\item Starting with a backbone-path from $U$ to $V$ through a single square vertex in the middle;
	\item Attach the $A^{-1}$-path to the square-vertex on the backbone-path, and consider it as dangling down from the path.
	\item Attach $\eta_\tau$ to the other end (``bottom'') of the dangling $A^{-1}$-path.
	\item Each shape receives an additional coefficient from the coefficient in the polynomial expansion for $A^{-1}$, i.e., $(-1)^t \cdot \frac{1}{1-\gamma}$ for the $\q_t(A)$ term for any $t\geq 0$.  
\end{enumerate}

In particular, it should be emphasized here that as we have two concatenation operations in a $3$-way concatenation, and there are $2$ types of intersections that may appear:
\begin{enumerate}
	\item Intersections from the attachment of $A^{-1}$-Path;
	\item Intersections from the attachment of $\eta_\tau$ to the end of the dangling $A^{-1}$ path;
\end{enumerate}

This prompts us to define the following well-behaved intersection. For convenience, we call $s_{bp}$ the square vertex on the backbone-path where the $A^{-1}$-path is attached, and $s_{\eta}$ the square vertex at the end of the $A^{-1}$-path where $\tau$ is attached via two edges connecting the two circle vertices of the endpoints of $\tau$.

\begin{definition}[Well-behaved Top, Bottom, Top+Bottom Intersection $\mathcal{T}, \mathcal{B}, \mathcal{TB}$] \label{def:well-behaved-intersection-vertical}
We define the following three categories of well-behaved intersection patterns in the three-way concatenation of backbone-path, dangling $A^{-1}$-path, and deviation attachment $\eta_\tau$:

\begin{enumerate}
	\item \textbf{Top-Intersection $\mathcal{T}$}: an intersection such that $s_{bp}$ is incident to $2$ distinct vertices. In other words, the pair of edges that $s_{bp}$ is incident to in the backbone path matches the pair of edges that $s_{bp}$ is incident to in the dangling $A^{-1}$ path. Moreover, there is no other vertex intersection in the three-way concatenation.
	\item \textbf{Bottom-Intersection $\mathcal{B}$}: an intersection such that $s_{\eta}$ is incident to $2$ distinct vertices. In other words, the pair of edges that $s_{\eta}$ is incident to in the $A^{-1}$ path matches the pair of edges that $s_{bp}$ is incident to in the deviation attachment of $\eta_\tau$. Moreover, there is no other vertex intersection in the three-way concatenation.
	\item \textbf{Top+Bottom-Intersection $\mathcal{TB}$ }: it involves the well-behaved intersection around $s_{bp}$ and $s_{\eta}$, and there is no other vertex intersection.
\end{enumerate}
\end{definition}

We give diagram illustrations as follows.
\begin{figure}[ht!]
    \centering
    \begin{subfigure}[b]{0.45\textwidth}
        \centering
        \includegraphics[width=\textwidth]{diagrams/int_original.pdf}
        \caption{original shape $\tau$ to be corrected.}
%        \label{fig:concat-proper}
    \end{subfigure}
    \quad
    \begin{subfigure}[h]{0.45\textwidth}
        \centering
        \includegraphics[width=\textwidth]{diagrams/tint.pdf}
        \caption{$\mathcal{T}$-intersection}
%        \label{fig:concat-intersection}
    \end{subfigure}
    \centering
	 \begin{subfigure}[h]{0.45\textwidth}
        \centering
        \includegraphics[width=\textwidth]{diagrams/tbint.pdf}
        \caption{$\mathcal{TB}$-intersection }
%        \label{fig:concat-proper}
    \end{subfigure}
    \begin{subfigure}[h]{0.45\textwidth}
        \centering
        \includegraphics[width=\textwidth]{diagrams/tint_linearized.pdf}
        \caption{Linearization of $\mathcal{T}$-intersection}
%        \label{fig:concat-proper}
    \end{subfigure}

    \caption{Examples of intersections.}
    \label{fig:concat}
\end{figure}




%\begin{definition}[Well-behaved Intersection Terms for $\mathcal{L}^*(w_\tau)$.]
%We call a term from the following shape concatenation from the above process a well-behaved intersection if \begin{enumerate}
%	\item the the square vertex on the backbone path is incident to $2$ pairs of the same edges: the two circle vertices on the backbone-path collide in pairs with the two circle vertices in the first injective gadget from $A$ (that is additionally an $\al$-gadget)
%	\item There is no additional vertex intersection in the concatenation of the backbone-path, dangling $A^{-1}$-path, and $\eta_\tau$.
%\end{enumerate}
%		
%	
%\end{definition}
\begin{remark}
	Implicitly, the first and final gadget of the dangling $A^{-1}$-path must be an $\al$-gadget to form a well-behaved intersection.
\end{remark}
% \snote{is "incidental" not supposed to be "incident"?}
With these definitions, we observe that the well-behaved intersection term admits an immediate linearization by removing the square-vertex on the backbone-path as well as incident edges. Grouping the intersection terms that reduce to the same underlying shape gives us the following.

\snote{definition 5.4 and proposition 5.5 are somewhat redundant with each other}
\jnote{dropped previous 5.4}
%\begin{definition}[Linearization of Well-behaved Intersection]\label{def:linearization-well-behaved}
%		Let $\tau_I$ be a well-behaved intersection in the sense of \cref{def:well-behaved-intersection-vertical}; we define the resulting linearized shape $R(\tau_I)$ as the same shape with each $s\in \{s_{bp},s_{\eta}\}$ and the incidental edges removed if $s$ has a well-behaved intersection matched up in its incidental edges.
%\end{definition}
%A well-behaved intersection can then be reduced to a proper shape as follows,
\begin{proposition}[Linearization of Well-behaved Intersection] \label{prop:linearization-well-behaved-vert}
	Let $\tau_I$ be a well-behaved intersection in the sense of \cref{def:well-behaved-intersection-vertical}; we define its linearized shapes as follows:
	\begin{enumerate}
		\item $R(\tau_I)$ has the same vertex boundary as $\tau_I$;
		\item $R(\tau_I)$ is obtained from $\tau_I$ via the following operation: for any square vertex $s \in \{s_{bp}, s_{\eta}\}$ incident to well-behaved intersection (i.e., its incident edges match up), remove the square vertex and its incident edges.
	\end{enumerate}
	Then we have \[
	\cm_{\tau_I} = \frac{m}{d^2}\cdot  \cm_{R(\tau_I)} + \tilde{O}(1/d)\,.
	 \]
\end{proposition}
%
%\begin{proposition}\label{prop:linearization-well-behaved}
%Let $\tau_R$ be any proper shape from the deviation attachment (i.e., a shape arising from the above attachment vertex without any vertex exception), and let $\tau_{I(R)}$ be collection of well-behaved intersection shape that reduce to $\tau_R$ under the removal of the initial square-vertex and its incidental edges,  we have \[ 
%	\cm_{\tau_{I(R)}} = (1+o_d(1)) \cdot  \gam \cdot  \cm_{\tau_R} \,.
%	\]
%	Moreover, $\tau_R$ is also a backbone-dangling shape.
%\end{proposition}
%\begin{proof}
%The crucial observation is that for any $\tau_R$, there are two well-behaved intersection shapes that reduce to a fixed $\tau_R$. We showcase in the following diagrams. The two shapes on the second line both reduce to the top shape under the removal process.  Since each linearization corresponds to a factor of $(1+o_d(1))\cdot m/d^2$, we have $2m/d^2 = \gam$ in total.
%\end{proof}
%
With the linearization operation defined, we first identify the properties of the shapes arising from linearization of well-behaved intersections.
\begin{claim} For terms in $\calL^*(w_\tau)$, the linearization of any well-behaved intersection is either $\tau$ (the original shape from horizontal concatenation in the prior level), $\tau^\top$, or a backbone-dangling shape. \label{clm:linearized-property-vertical}
\end{claim}
\begin{proof}
	We prove this by cases depending on the length $t$ of the $A^{-1}$-path. The crucial property to check is that we have only one single square vertex on the backbone path in the resulting shape.
\begin{enumerate}
	\item For the case of $t=0$, there is no $\mathcal{T}$ intersection (and hence neither $\mathcal{TB}$) while there is a $\mathcal{B}$ intersection. Such intersection recovers $\tau$ and $\tau^\top$ (up to flipping of the boundary) as the backbone edges linearize with the deviation-attachment edges. We remind the reader that $\tau$ is not a \emph{backbone-dangling} shape as it has $>1$ square vertex on the backbone path for $\tau\neq \emptyset$ beyond the base iteration.
	\item For the case of $t=1$, linearization of $\mathcal{TB}$ gives $\tau$ and $\tau^\top$. 
	\item For the other cases, consider the backbone-path as depth-$0$, and the horizontal-$\tau$-level as depth $t+1$. Any other linearization has a single vertex on the backbone path as this property is only violated when the horizontal $\tau$-level becomes depth-$0$ after linearization. It is straightforward to verify that each $\calT$ and $\calB$ linearization decreases the depth by $1$ (hence $\mathcal{TB}$ by $2$).
\end{enumerate}	
This completes the proof of our claim.
	\end{proof}
 From the linearization, we obtain an explicit decomposition into proper shapes by combining the proper shapes in the decomposition with the linearization of the well-behaved intersection shapes. Let $\mathsf{Proper}(\eta_\tau)$ be defined to be the weighted combination of properly concatenated shapes in this process as well as the well-behaved intersection terms formally as follows,

\begin{definition}	For any deviation vector $\eta_\tau$ such that $\tau\neq \emptyset$, 	we define $
	\mathsf{Proper} \mathcal{L}^*(w_\tau)$ to be the weighted linear combination of the following terms in the $3$-way concatenation in $\frac{1}{1-\gamma}\cdot  \calL^*(w_\tau)$:
	\begin{enumerate}
		\item Proper concatenations such that each non-trivial gadget in $A^{-1}$-path is an $\al$-gadget;
		\item Linearization of well-behaved intersections such that each non-trivial gadget in $A^{-1}$-path is an $\al$-gadget.
	\end{enumerate}
	It is important to distinguish between $\mathsf{Proper}(\eta_\tau)$ and $\mathsf{Proper}\mathcal{L}^*(w_\tau)$.
\end{definition}

\begin{figure}[ht!]
  \centering
      \begin{subfigure}[b]{0.60\textwidth}
        \includegraphics[width=\textwidth]{diagrams/reduce.png}
%        \label{fig:M_alpha}
    \end{subfigure}
            \caption{From $2$ Intersection Terms to $1$ Proper Shape}

%    \quad
%    \begin{subfigure}[b]{0.32\textwidth}
%        \includegraphics[width=\textwidth]{diagrams/M_beta.pdf}
%        \caption{\(\cm_{\beta}\).}
%%        \label{fig:M_beta}
%    \end{subfigure}
%    \captionsetup{width=.9\linewidth}
%    \caption{Graph-matrix representations of a \(d\times d\) \(\mathsf{GOE}\)
%    matrix with zero diagonal, and of the \(m\times m\) matrices \(M_\alpha\) and
%    \(M_\beta\) from \Cref{prop:M-decomposition}. Square vertices take labels in
%    \([m]\), and circle vertices take labels in \([d]\). The two ovals indicate the
%    left and right boundaries \(U_\tau,V_\tau\). If an edge \(e\) is not explicitly
%    labeled with an index, then \(t(e)=1\) by default.}
%    \label{fig:M_alpha_beta}
\end{figure}


\begin{lemma}[Proper Decomposition for the Scaled Deviation Attachment]\label{lem:proper-decomp-deviation-attach}
	For any horizontally concatenated term $\tau$ to be corrected in the iterative process, 	 we have \[ 
	\mathsf{Proper}\mathcal{L}^*(w_\tau) = \frac{\gam}{1-\gam} \cdot \tau+ \sum_{\substack{\psi \in \mathsf{Proper Shapes}(\eta_\tau) \\ \text{length-$t$ path in } A^{-1} } }  (-1)^{t}  \cdot \cm_\psi \,.
	\]
\end{lemma}

\begin{proof}
	By~\cref{clm:linearized-property-vertical}, each term from the linearization is either $\tau$ or a properly concatenated shape. Hence it suffices for us to consider any such fixed shape $\tau$ and compute its coefficient.

Firstly, we observe that for any properly concatenated shape $\psi$, there are four terms that may give rise to this in the final linearized sum,\begin{enumerate}
	\item The properly concatenated shape itself;
	\item Intersection Shapes that reduce to $\psi$ via linearization of $\mathcal{T}$, call them $\mathcal{T}^{-1}(\psi)$;
	\item Intersection Shapes that reduce to $\psi$ via linearization of $\mathcal{B}$, call them $\mathcal{B}^{-1}(\psi)$;
	\item Intersection Shapes that reduce to $\psi$ via linearization of $\mathcal{TB}$, call them $\mathcal{TB}^{-1}(\psi)$.
\end{enumerate} 

For each term, the coefficient in addition to the coefficient of $c_\tau$ is given by the following,
\begin{enumerate}
	\item In the correction $\corr(\tau)$, $\mathcal{L}^*(w_\tau)$ receives a scaling of $  \frac{1}{1-\gam } $;
	\item There is an additional coefficient of $(-1)^t \cdot \frac{1}{1-\gamma}$ for each (intersection) term that uses a length-$t$ path from the $A^{-1}$ attachment.
\end{enumerate}
Next, we observe that for any fixed shape $\psi$, there are two shapes that linearize to $\psi$ via $\mathcal{T}$ and $\mathcal{B}$ respectively by adding a square vertex with the same pair of edges. The matching of the pair gives us the choice of $2$. Therefore, via~\cref{prop:linearization-well-behaved-vert}, we have \[ 
\cm_{\mathcal{T}^{-1}(\psi)} = 2\cdot \frac{m}{d^2} \cdot \cm_{\psi} +o_d(1) = \gam \cdot \cm_{\psi}+o_d(1);
\]
and similarly,
\[ 
\cm_{\mathcal{B}^{-1}(\psi)} = \gam \cdot \cm_{\psi}+o_d(1)\,.
\]
A consecutive application of $\mathcal{T}$ and $\mathcal{B}$ gives us \[ 
\cm_{\mathcal{TB}^{-1}(\psi)} = \gam^2 \cdot \cm_{\psi}+o_d(1)\,.
\]
For $\psi$ that uses a length-$t\geq 0$ path in the $A^{-1}$-path, shapes that linearize to $\psi$ via $\mathcal{T}$ and $\mathcal{B}$ have length $t+1$ while shapes via $\mathcal{TB}$ have length $t+2$. Write \[ 
\frac{1}{1-\gam} \cdot  \mathcal{L}^*(w) = \sum_{\substack{\psi \\ \text{shape from concatenation} \\ \text{possibly intersected} }} c(\psi) \cdot \cm_\psi 
\]
by expanding in the graph matrix decomposition without processing (i.e., linearizing) intersection terms. Let $c(\psi)$ be the coefficients in the raw decomposition, we have \[ 
c(\psi) = c(\mathcal{TB}^{-1}(\psi)) = -1\cdot c(\mathcal{B}^{-1}(\psi)) = -1\cdot c(\mathcal{T}^{-1}(\psi))  \,,
\]
by alternating the sign in $A^{-1}$, and additionally, for $\psi$ that uses a length-$t$ path in $A^{-1}$, we have\[
c(\psi) = (-1)^t\cdot \frac{1}{(1-\gam)^2}
\]
with one factor of $1/1-\gam$ from the normalization coefficient in front of $\calL^*(w)$, and the other factor from $A^{-1}$. 

Finally, grouping all the terms that produce $\psi$ (with length-$t$ in $A^{-1}$ path) after linearization, the total contribution via a proper shape $\psi$ in $\frac{1}{1-\gamma} \cdot \mathcal{L}^{*}(w_\tau)$ is  
\begin{align*}
	&c(\psi) \cdot \cm_\psi + c(\mathcal{T}^{-1}(\psi)) \cdot \cm_{\mathcal{B}^{-1}(\psi)} + c(\mathcal{B}^{-1}(\psi)) \cdot \cm_{\mathcal{B}^{-1}(\psi)} +  c(\mathcal{TB}^{-1}(\psi))\cdot \cm_{\mathcal{TB}^{-1}(\psi)}\\
	&= c(\psi) \cdot \cm_\psi (1- 2\gam +\gam^2) +o_d(1)\\
	&= (-1)^t \cdot \frac{1}{(1-\gam)^2} \cdot (1-\gam)^2 \cdot \cm_\psi+ o_d(1)\\
	&= (-1)^t \cdot \cm_\psi +o_d(1)\,. 
\end{align*}
This proves the second term in the desired equality. For the first term of $\cm_\tau$, apply the above argument to intersection shapes that linearize to $\tau$, i.e., $\calT^{-1}(\tau)$ and $\mathcal{TB}^{-1}(\tau)$, \[
c(\calT^{-1}(\tau)) \cdot \cm_{\calT^{-1}(\tau)} + c(\mathcal{TB}^{-1}(\tau)) \cdot \cm_{\mathcal{TB}^{-1}(\tau)} = \frac{1}{(1-\gam)^2}  (\gam-\gam^2) = \frac{\gam}{1-\gam}\,.
 \]
This completes the proof.% to our lemma.
\end{proof}

Crucially, this shows that the $\frac{1}{1-\gamma}$ coefficients are precisely offset in the weighted linear combination when we combine the properly concatenated shape with the linearization of well-behaved intersection shapes.  Next, for each $\cm_\psi$, we obtain its variance bound by the following,
 \begin{proposition}[Variance of Vertical Attachment] \label{prop:variance-vertical-concate}  For $\psi\in \mathsf{Proper Shapes}(\eta_\tau)$ with a length-$t$ path from $A^{-1}$, we have 
 	\[
 	\Var( \cm_\psi ) = \gam^{1+t} \cdot \Var(\tau)\,. 
 	 \]
 \end{proposition}
 
 
 
 \jnote{this might require a proof}
 
 Since each shape in the collection is backbone-dangling, our main theorem implies that they are freely independent, and therefore the variance of these matrices also adds linearly. 
 Therefore, ignoring truncation for $A^{-1}$, we obtain \begin{align*}
	\Var(\mathsf{Proper}(\eta_\tau) )&= \sum_{\substack{\psi \in \mathsf{Proper Shapes}(\eta_\tau) \\ \text{length-$t$ path in } A^{-1} } }   \Var(\cm_\psi) 
	= \Var(\tau)\cdot   \sum_{t\geq 0} \gamma^{t+1}  
	=\frac{\gam }{1-\gamma} \cdot \Var(\tau) \,.
 \end{align*} 

%

%\begin{proposition}
%For every \(\psi\in\mathsf{ProperShapes}(\eta_\tau)\) whose
%\(A^{-1}\)-path has length \(t\), we have
%\[
%    \Var(\cm_\psi)
%    =
%    \bigl(1+o_d(1)\bigr)\gam^{t+1}\Var(\tau).
%\]
%\end{proposition}
%
\begin{proof}[Proof of~\cref{prop:variance-vertical-concate}]
We compare the two-copy diagram defining
\[
    \Var(\cm_\psi)
    =
    \frac{1}{d}\E\Tr(\cm_\psi\cm_\psi^\top)
\]
with the corresponding diagram for \(\Var(\tau)\).

The deviation attachment adds one square vertex and two
\(h_1\)-edges connecting it to the boundary vertices of \(\tau\).
There are two possible matchings of these edges between the two
copies, so this attachment contributes
\[
    \bigl(1+o_d(1)\bigr)\,
    2m\cdot\frac{1}{d^2}
    =
    \bigl(1+o_d(1)\bigr)\gam.
\]

By definition of \(\mathsf{ProperShapes}(\eta_\tau)\), every
nontrivial gadget on the \(A^{-1}\)-path is an \(\alpha\)-gadget.
Each such gadget adds one square vertex, two circle vertices, and four
\(h_1\)-edges. Its contribution to the two-copy diagram is therefore
\[
    \bigl(1+o_d(1)\bigr)\,
    2m d^2\cdot\frac{1}{d^4}
    =
    \bigl(1+o_d(1)\bigr)\gam.
\]
The \(t\) gadgets on the path thus contribute
\(\bigl(1+o_d(1)\bigr)\gam^t\).

Finally, the outer backbone adds two circle boundary labels \(a\neq b\)
and two \(h_1\)-edges. Its contribution is
\[
    \sum_{a\neq b}
    \E\!\left[v_s[a]^2v_s[b]^2\right]
    =
    d(d-1)\cdot\frac{1}{d^2}
    =
    1+o_d(1).
\]
Multiplying these contributions gives
\[
    \Var(\cm_\psi)
    =
    \bigl(1+o_d(1)\bigr)
    \gam\cdot\gam^t\Var(\tau)
    =
    \bigl(1+o_d(1)\bigr)\gam^{t+1}\Var(\tau).
\]
The error accounts for the global injectivity restrictions and is
uniform over the allowed truncated path lengths. The alternating sign
\((-1)^t\) of the \(A^{-1}\)-coefficient does not affect the variance.
\end{proof}

%\begin{proof}[Proof to~\cref{clm:proper-decomp-deviation-attach}]
%Firstly, write \[ 
%\mathcal{L}^*(w) = \sum_{\substack{\psi \\ \text{shape from concatenation} \\ \text{possibly intersected} }} c(\psi) \cdot \cm_\psi 
%\]
%by expanding in the graph matrix decomposition without processing (i.e., linearizing) intersection terms. Let $c(\psi)$ be the coefficients in the raw decomposition.
%
%Next, recall that \[ 
%A^{-1} = \frac{1}{1-\gamma} \sum_{t\geq 0}  (-1)^t\q_t(A)\,.
%\]
%Thus, we have \[ 
%c(\psi) = \frac{1}{1-\gam } (-1)^t \numberthis \label{eq:raw-coef}
%\]	 
%for any $\psi$ that uses $t$-way concatenation from $A^{-1}$-term. Moreover, for any properly concatenated shape $\tau_R$, and $\tau_{I_a(R)}$, $\tau_{I_b(R)}$ the two well-behaved intersection shapes that reduce to $\tau_R$, we have \[ 
% c(\tau_{I_a(R)})= c(\tau_{I_b(R)}) = -1\cdot c(\tau_R)
%\]
%by the alternating sign. By symmetry, we define $c(\tau_{I(R)}) \coloneqq c(\tau_{I_a(R)})= c(\tau_{I_b(R)})$.
%
% Restricting our attention to the shapes in our definition for  $\mathsf{Proper}$, we group together each proper-shape with the well-behaved intersection shapes that reduce to that, and we have \begin{align*}
% \mathsf{Proper}(\eta_\tau) &\coloneqq \sum_{\tau_R: \text{proper}} \left( c(\tau_R)    \cm_{\tau_R}    + c(\tau_{I_a(R)}) \cm_{\tau_{I_a(R)}} +  c(\tau_{I_b(R)}) \cm_{\tau_{I_b(R)}}      \right)\\
% &= \sum_{\tau_R:\text{proper} } \left(c(\tau_R)  \cdot   \cm_{\tau_R}    +  c(\tau_{I(R)}) \cdot \cm_{\tau_I(R)}  \right)
%% &= \sum_{\tau_R:\text{proper} } \left( c(\tau_R)  + c(\tau_R)   \right) \cm_{\tau_R} 
%\\ &= \sum_{\tau_R:\text{proper} } \left( c(\tau_R)  -  \gam\cdot  c(\tau_R)     \right) \cm_{\tau_R} 
% \qquad \text{ (by~\cref{prop:linearization-well-behaved} and ignore $o_d(1)$ terms) }\\
% &=  (1-\gam) \cdot \sum_{\tau_R:\text{proper} }  c(\tau_R)\cdot \cm_{\tau_R}\\
% &=\sum_{\substack{\tau_R:\text{proper} \\ \text{length-$t$ in }A^{-1}  }} (-1)^t \cdot \cm_{\tau_R}  \qquad \text{(by~\cref{eq:raw-coef} ) }
%\end{align*}
%This proves our claim.
%\end{proof}
%
%
%
%\paragraph{Error Terms for Vertical Concatenation}
%We now shift gear to considering the error terms, and bound the error term by the following lemma.
%\begin{lemma}[Error Terms for Vertical Concatenation]
%	Recall that we define 
%	\[ 
%	\mathcal{L}^*(w)(w_\tau) = \mathsf{Proper}(\eta_\tau) + \mathsf{Error}(\eta_\tau)\,,
%	\]
%	we have \[ 
%	 \|\mathsf{Error}(\eta_\tau)\|_{sp}=o_d(1)\,.
%	\]
%\end{lemma}
%
% There are primarily two sources of error terms:\begin{enumerate}
%	\item Non well-behaved intersections in the concatenation of $A^{-1}$;
%	\item Terms from rank-$2$ components in $M^{-1}$.
%\end{enumerate}
%For the error term analysis, we will again apply block-value bound throughout.
%For starters, we consider the error terms from the concatenation of $A^{-1}$. 
%
%\paragraph{Pre-processing Intersection Shapes} Since the intersection shapes may be consist of a well-behaved intersection around the first square vertex + other vertex intersections further down the dangling path, we first process away the well-behaved intersection term by removing the initial square vertex and its incident edges. With this pre-processing, we may observe that there the there is no well-behaved intersection around the initial square vertex on the backbone path.
%
% This further allows to deduce the following claim that shows either the backbone-path is unaffected in the concatenation operation, or we can identify a intersection property that can later be used to deduce a block-value gap.
%
%
%%Our crucial observation is that as we process away well-behaved intersection around the first square vertex, we may assume that the separator vertex is at either boundary.
%
%\begin{claim}[Pe-processing Intersection Property] \label{clm:preprocessing-int-property}
%	There is at least one separator vertex along the backbone-path, i.e., from the two boundary circle vertices or the initial square vertex, or we have a mismatching intersection around the initial square vertex- it is incident to $3$ distinct vertices, with one pair of edges collapsing.
%\end{claim}
%\begin{proof} Note that this is immediate if there is no vertex intersection that collides with vertices on the backbone path.
%
%	On the other hand, by our preprocessing operation, there is no well-behaved intersection around the initial square vertex. The crucial observation is that the square vertex cannot be incident to only two distinct vertices.
%	
%	Suppose it does, the pre-processing operation would have applied and give us a contradiction. 	
%	Finally, if it is incident to $3$ distinct vertices, we have a mismatching intersection around the initial square vertex as described.
%\end{proof}
%
%\begin{claim}
%	In a pre-processed intersection shape $\tau_I$ from $\eta_\tau$, there is no "isolated" vertex: i.e., no vertex can be making a first and final appearance within the same block-step.
%\end{claim}
%\begin{proof}
%    This is straightforward for boundary vertices by definition. Next, since $\eta_\tau$ is injective, no edge can appear twice within $\eta_\tau$, while some edge may intersect with the edges on the backbone path. However, this may only appear for the edges around the initial square vertex. For the square vertex, unless it is incident to two distinct vertices exactly which is removed by our pre-processing operation, the square vertex cannot be making its first and final appearance in this step.
%\end{proof}
%
%
%
% \begin{lemma}
%	Let $\eta_i$ be the deviation vector at the $i$-th iteration that is a concatenation of backbone-dangling shapes, let \[
%	M_W \coloneqq  \frac{1}{s^2-ru}\, A^{-1} \left( u\,\frac{\1_m\1_m^T}{d} - s\,\frac{\eta\1_m^T+\1_m\eta^T}{d} + r\,\frac{ \eta\eta^T}{d} \right) A^{-1}
%	 \,,\] then for any $i\geq 1$, we have \[ 
%	 \eta^\top \cdot M_W \cdot \eta_i =  o_d(1)\,.	\]
%	 Quantitatively, we have \[ 
%	  \eta^\top \cdot M_W \cdot \eta_i =  \tilde{O}(1/\sqrt{d})\,.
%	 \]
%\end{lemma}
%
%
%
%
%
%
%It is useful to recall the following scalar bounds.
%\HyperParameterBounds*
%As a result, we also have\[
%\frac{1}{s^2 - ru } = \Theta \left(\frac{d}{m}\right)\,.
% \]
%Additionally, we split $M_W$ into three components depending on the inner term, 
%\[
%	M_W \coloneqq  \frac{1}{s^2-ru}\, A^{-1} \left( \underbrace{ u\,\frac{\1_m\1_m^T}{d}}_{M_J} -  \underbrace{s\,\frac{\eta\1_m^T+\1_m\eta^T}{d}}_{M_C} + \underbrace{r\,\frac{ \eta\eta^T}{d} }_{M_\eta} \right) A^{-1}
%	 \,,\]
%  Intuitively, this comes from the following observation:
% \begin{enumerate}
%  	\item $A^{-1}$ is an injective path consist of $\al$ and $\beta$-gadgets that ends with a square-vertex;
%  	\item Combining the scalar with the normalization of $\frac{1}{s^2-ru}$, each term comes with the following coefficient (up to $\Theta(1)$ factors), \begin{itemize}
%  		\item $M_J$ comes with a coefficient of $ \frac{d}{m} \cdot u \approx  \frac{d}{m}  $;
%  		\item $M_C$ comes with a coefficient of $\frac{d}{m}  $;
%  		\item $M_\eta$ comes with a coefficient of $\frac{d }{m}  \cdot \frac{m}{d} \cdot \frac{1}{1-\gamma}  = \Theta(1)\,.$
%  	\end{itemize}
%  	\item The only dominant term comes from $M_\eta$ which gives  \[A^{-1} \cdot  \frac{1}{s^2-ru}\cdot  r \cdot \frac{\eta\eta^\top}{d} A^{-1} \approx \left( A^{-1}\eta  \right)\cdot  \left( \frac{1}{d}\eta^TA^{-1} \right)\,; \]
%  	\item Notice we group out $A^{-1}\eta$ as that is indeed the form of the main dangling term for the base iteration), and therefore we expect the first term continues to be a dominant term (i.e., there is no inverse-polynomial gap in $d$);
%  	\item However, for the second term, we have for $\eta_i\neq \eta$, we have \[ 
%  	\left( \frac{1}{d} \cdot  \eta^TA^{-1} \cdot \eta_i \right )\ 
%  	\]
%  	We claim this is $o_d(1)$ norm for $\eta_i\neq \eta$. 
%  	\begin{itemize}
%  		\item For any $\eta_i \neq \eta$, $\eta^T A^{-1} \eta_i$ is of mean-$0$ (i.e., there is no square term). To see this, $\eta_i$ is a concatenation of $t>1$ $\al$ or $\beta$ gadgets by design of correction scheme only applied to $P_{t}$ for $t\geq 2$.
%  		\item To see why it is $o_d(1)$, consider the following charging process that traverses the expansion for $\eta^T A^{-1} \eta_i$ from the initial circle vertex corresponding to $\eta$ attachment.. Our goal is to show that a factor of $\tilde{O}(1/\sqrt{d})$ from the normalization factor of $\frac{1}{d}$ is sufficient for the circle vertex. In contrast, a full factor of $1/d$ (aka the full normalization factor) needs to be assigned to the circle vertex for $\eta^\top A^{-1} \eta$.
%  	\end{itemize}
% \end{enumerate} 
%
%\begin{proof} We now formally prove this via block-value bounds. Without loss of generality, we may assume $\eta_i$ does not contain any floating component from $M_W$, as our proof will in turn show that any floating component comes with a block-value gap of $\tilde{O}(1/\sqrt{d})$.  For starters, we apply graph-matrix decomposition:\begin{enumerate}
%	\item Each term is a backbone-dangling shape with an $A^{-1}$ dangling path;
%	\item The first $A^{-1}$ path may terminate in either a square vertex without $\eta$-attachment, or terminates with a square $h_2$-attachment;
%	\item For the terms that do not terminate with an $\eta$-attachment, we need to identify an additional $O(1/\sqrt{d})$ gap (recall that one of the $1/\sqrt{d}$ factor from the final $h_2$-attachment is assigned to the initial square vertex);
%	\item We then apply block-value bound for the subsequent terms: starting from a circle vertex in the case of $\eta^\top A^{-1}$ terms, and a square vertex in the case of $\1_m A^{-1}$ terms.
%	\item Each gadget in an $A^{-1}$-path contributes $O(1)$ block-value.
%\end{enumerate}
%\paragraph{Analysis for $M_\eta $}	
%This is the dominant term where we use $\eta_i\neq \eta$ for $i\geq 1$. Note that this is a term that contains an $h_2$-attachment for the backbone-dangling path, and the coefficient for this term is  $1/d$.  We apply the following charging argument in which we split $1/d$ into two copies evenly:
%\begin{enumerate}
%	\item The backbone-dangling path contains an $h_2$-attachment, hence it suffices for us to apply the charging argument to the backbone-dangling path and bound the component to have a block-value $O(1)$. Next, it suffices for us to bound the block-value corresponding to the floating component.
%	\item Assign $\tilde{O}(\frac{1}{\sqrt{d}}$) value to the initial circle vertex $s$ that is part of the $h_2$ attachment in the $\eta^\top A^{-1} \eta_i $-path. This leaves us an excess of $\tilde{O}( \frac{1}{\sqrt{d}})$ from the scalar coefficient;
%	\item Our key observation is that the initial-circle vertex $s$ cannot be making both of its first and final appearances in a $\eta^\top A^{-1} \cdot \eta_i$ component for $i\neq 0$, and the remaining vertex factors can be charged by the edge-factors in the component, hence leaving us with an excess of $\frac{1}{\sqrt{d}}$.
%\end{enumerate}
%Towards that end, we use the ideas adapted from prop. 4.14 in \cite{HKPX23}: we split each $h_2$ edge as two edge-copy such that each edge-copy carries a factor of $O(1/\sqrt{d})$, and  we would like to assign edge-copy to vertices in a local traversal process. The up-shot of the local traversal process is to assign factors to vertices locally that upper bounds their contribution in a global walk:
%\begin{enumerate}
%	\item A vertex appearance for the first time in the walk may correspond to its first/final appearance in the walk, in this case we will assign a local bound given by half of its vertex weight, matching the global bound;
%%	\item If a vertex makes a first and last appearance in the same local traversal, it is assigned a factor of full vertex weight, matching the global bound;
%	\item Any mismatch of first/last in a local-traversal with a middle appearance in the global walk corresponds to a $\tilde{O}(1/\sqrt{d})$ gap.
%\end{enumerate}
%
%
%\jnote{likely need to verify we never pass through a diagonal term inside}
% \end{proof}
 
 
%\paragraph{Vertical + Horizontal Charging Scheme} 
%To prove our norm bound for each backbone-dangling shape, we consider the following assignment scheme adapted from the proof of prop. 4.14 in \cite{HKPX23}. 
%To describe our scheme, we first make the following definitions for convenience of our discussion. 
%
%
%
%For each vertical attachment to an injective $A^{-1}$ path, we have two circle vertices (corresponding to the shape generated from Chebysehev polynomials) attached to the square vertex (corresponding to evaluation of affine constraint deviations). We will single out these two edges and call them \emph{anchor} edges for the corresponding vertical attachment, and call the two circle vertices \emph{anchor} vertices.
%
%
%
%%For starters, we will consider edge-assignment for non-anchor edges. 
%
%
%\begin{mdframed} \textbf{Top-Down and Horizontal Edge-Copy Charging Primitive.} \begin{enumerate}  
%\item  \textbf{Edges from an $A^{-1}$ path}: Traverse vertically from top to bottom, and assign each edge-copy to the vertex it leads to unless this is a \emph{reverse-charging edge} as in~\cref{def:reverse-charging}, in which case we assign to the source vertex.
%\item  \textbf{Edges in Horizontal Path}: there are two circle \emph{anchor vertices} corresponding to the two endpoints of the path. Traverse from the left anchor vertex to the right anchor vertex, and assign each edge-copy to the vertex it leads to, unless it is a \emph{reverse-charging edge}, in which case we assign to the source vertex.
%\item \textbf{Anchor Edges}: these are the edges between $A^{-1}$-path and horizontal path. \begin{itemize}
%	\item Assign both edge-copies to the source square vertex if this is a \emph{reverse-charging edge} as in \cref{def:reverse-charging}, i.e., the underlying edge-pair explore the source square vertex earlier on in the local traversal process;
%	\item If not an $h_2$ edge, assign both edge-copies to the circle vertices they lead to;
%	\item Otherwise, this is an $h_2$ edge, assign one edge-copy to the incident circle vertex, and we treat the other edge-copy as an \textbf{attachment-reserve}.
%\end{itemize}
%\end{enumerate}
%\end{mdframed}
 

\begin{restatable}[Local Charging]{proposition}{propLocalTraversal}
\label{prop:local-traversal}
We can assign each edge-copy to at most one vertex so that the following properties hold:
\begin{enumerate}
    \item For each circle vertex, one edge-copy is assigned to each of its global first and global last appearances. In particular, it is assigned one edge-copy if it makes exactly one of these appearances, and two edge-copies if it makes both.

    \item For each square vertex on the dangling path, two edge-copies are assigned to each of its global first and global last appearances. In particular, it is assigned two edge-copies if it makes exactly one of these appearances, and four edge-copies if it makes both.

    \item Whenever a square vertex makes a global middle appearance but its local first appearance within the current block, it is assigned at least one edge-copy.

    \item No edge-copy corresponding to a middle-appearance step is assigned to the global first or global last appearance of any vertex.
\end{enumerate}
\end{restatable}



 This is a generalization of prop. 4.14 in~\cite{HKPX23} in that we additionally incorporate horizontal paths of length $>1$. Our charging and traversal scheme is also an extension of that from the prior work. We defer the full analysis to the appendix. 

\paragraph{Analysis of Base Matrix}
Finally, we extend the above to bound the variance of our initialization at $Q_0$.
\begin{lemma}
	$Q_0$ admits a decomposition into proper backbone-dangling shapes \[
	Q_0 = \sum_{\tau\in \calB(Q_0):\text{proper}} c(\tau) \cdot \cm_\tau  + o_d(1) \,,
	\]
	such that \[ 
	\Var(Q_0) = C_F^2 \cdot \frac{\gamma}{1-\gam}\,.
	\]
\end{lemma}

The proof is in two steps. Firstly, by prop~4.1 of \cite{HKPX23}, we can rewrite \[
M^{-1}\eta = \frac{r+s}{s^2-ru} \cdot A^{-1}\eta - \frac{u+s}{s^2-ru} A^{-1} \1_m\,.
 \]
 The key is to observe that the second term is an $o_d(1)$ error term, while we can decompose the first term via proper backbone-dangling shapes with the desired variance bound. Once rearranged as above, we apply similar ideas to the above calculation for the subsequent levels with one minor distinction: since the final attachment is now an $h_2$-gadget instead of two $h_1$ anchors to a horizontal level, the definition of bottom well-behaved intersection $\mathcal{B}$ needs to be adjusted for $Q_0$.
 \begin{definition}[Well-behaved  Intersection for $Q_0$] Let $s_\eta$ be the final square-attachment vertex at the end of the $A^{-1}$ dangling path. An intersection is a well-behaved bottom intersection if
 	  $s_{\eta}$ is incident to the same $h_2$ edge in $A^{-1}$ and in the $h_2$-attachment. In other words, the two vertices around $s_\eta$ intersect.

 	  Analogously, the modification applies to the $\mathcal{TB}$ intersection for $Q_0$.
 \end{definition}
 It should be noted that by the above definition, it is implicit that the final $A^{-1}$-gadget is a $\beta$-gadget, to be contrasted with the previous discussion for the subsequent correction terms. With this definition, we now consider terms in $\mathcal{L}(M^{-1}\eta)$.
 
 We declare the following terms as the main terms in the decomposition. \begin{definition}	 For the base iteration, we define $
	\mathsf{Proper} \mathcal{L}^*(M^{-1}\eta )$ to be the weighted linear combination of the following terms in the $3$-way concatenation in $\frac{r+s}{s^2-ru} \cdot A^{-1}\eta$ :
	\begin{enumerate}
		\item Proper concatenations such that each non-trivial gadget in $A^{-1}$-path is an $\al$-gadget;
		\item Linearization of well-behaved intersections such that each non-trivial gadget in $A^{-1}$-path is an $\al$-gadget without interior $\beta$-gadget.
	\end{enumerate}
	Let $\mathsf{ProperShape} \mathcal{L}^*(M^{-1}\eta )$ denote the collection of underlying proper shapes.
\end{definition}

Before proceeding, let's observe that among the properly concatenated terms, there is in fact a unique term $\psi$ that uses a length-$t$ path in $A^{-1}$ for any $t\geq 0$ since each gadget in the $A^{-1}$ path is an $\al$-gadget.  This allows us to call such a term $\psi(t)$.

\begin{claim}[Decomposition of Proper Shapes in $Q_0$]\label{clm:proper-q0-decomp}
	\[ 
	\mathsf{Proper} \mathcal{L}^*(M^{-1}\eta ) = \sum_{\substack{\psi\in \mathsf{ProperShape} \mathcal{L}^*(M^{-1}\eta )  \\ \text{length-$t$ in }A^{-1}} } (-1)^t \cdot \cm_\psi 
	\]
	Moreover, each shape $\psi\in \mathsf{ProperShape} \mathcal{L}^*(M^{-1}\eta )$ is a backbone-dangling shape
\end{claim}  
\begin{proof}
	This is analogous to \cref{lem:proper-decomp-deviation-attach} except the linearization of any well-behaved intersection is also a backbone-dangling shape as we no longer have a non-$h_2$ horizontal level in the attachment. Fix any proper shape $\psi$, again there are $3$ (intersection) shapes that can contribute via $\psi$ after linearization:
	\begin{enumerate}
		\item The properly concatenated shape $\psi$ itself;
		\item Intersection Shapes that reduce to $\psi$ via linearization of $\mathcal{B}$ or $\mathcal{T}$;
		\item Intersection Shapes that reduce to $\psi$ via linearization of $\mathcal{TB}$.
	\end{enumerate}
	Let's first consider the coefficient of $\psi$ in $\mathcal{L}^*(A^{-1}\eta )$.
	Combining with the $\frac{1}{1-\gamma}$ coefficient in $A^{-1}$ and the alternating sign in the length-$t$, we have the final coefficient of $\psi$ in $ \mathcal{L}^*(A^{-1}\eta )$ be \[ 
	\frac{1}{1-\gamma} \cdot (-1)^t (1 - 2\gam + \gam^2) = (1-\gam)\cdot  (-1)^t
	\]
	Next, recall that we have an additional coefficient of \[ 
	\frac{r+s}{s^2 -ru } = \frac{1}{1-\gamma}
	\]
	via \cref{lem:hyper-parameter-bounds} recalling that $s=O(1), r= \frac{m}{d}\frac{1}{1-\gamma}$ and $s = -(1-\gamma)$ up to $o_d(1)$ terms. Combining the above proves the desired claim.
		\end{proof} 
		
\begin{proof}[Proof of~\cref{lem:q0-graph-variance}]
	%\anote{The grammar of the first two sentences is awkward.}
    %We now prove the variance of untruncated $Q_0$, with the truncated version deferred to the subsequent section. %We also defer to the appendix to 
    We now analyze the variance of the untruncated $Q_0$, deferring the analysis of the truncated version to the next section. We show that the remaining terms not in $\mathsf{Proper}\mathcal{L}^{*}(M^{-1}\eta)$ are of negligible norm in the appendix. By definition of $\mathsf{Proper}\mathcal{L}^{*}(M^{-1}\eta)$, each term is a backbone dangling shape without any intersection.
	
	Observe that for proper shape $\psi$ with length-$t\geq 0$ in $A^{-1}$ path, we have \[ 
	\Var(\cm_\psi) = \gam^{1+t}\,.
	\]
	To see this, we start by observing that for $t=0$, we have \[ 
	\Var(\cm_{\psi(t=0)} ) = \frac{1}{d^2} \cdot d^2 \cdot \frac{2}{d^2} \cdot m = \gam \,.
	\]
	Each additional gadget within the $A^{-1}$ path is an $\al$ gadget (as the $\beta$-gadget is negligible under an injective $A^{-1}$-path), and any such gadget contributes a factor of $\gam$. Assuming each backbone-dangling shape is freely independent (established in the later section), and recall that $Q_0$ is the identity perturbation scaled with an additional factor of $C_F$, by the decomposition in \cref{clm:proper-q0-decomp} we have \[ 
	\Var(Q_0) = C_F^2 \cdot \sum_{\psi(t)} \Var(\cm_{\psi(t)})  = C_F^2 \gam \sum_{t=0}^\infty \gam^t = \frac{\gamma}{1-\gam }\cdot C_F^2\,.
	\]
	
	% \todo{ add the robust version with truncation put in}
	\end{proof}		
\subsection{Concrete Primal Iterative Process via Graph Matrices}\label{sec:concrete-primal-graph-mat}
In this section, we instantiate the ideal iterative procedure described with graph matrices to streamline our analysis. For simplicity of the discussion, we focus on the primal program while analogous changes apply to the dual program in a straightforward manner. See~\cref{app:concrete-dual-iteration} for details regarding the dual program. We now focus on the primal.

We incorporate the following changes into the iterative process:
\begin{enumerate}
	\item Instead of reasoning about $H_i$ directly and adding update $\corr(H_i)$, we will restrict our attention to the properly concatenated shapes arising therein, i.e., $\fp_j(Q_i) - \fp_j(Q_{i-1})$ for any matrix $Q_i$ and any $j\ge 2$. 
	\item Our update will address solely proper shapes from  $\fp_j(Q_i) - \fp_j(Q_{i-1})$, and we will treat the remaining terms as negligible error terms. We call any error in this process $\textsf{HorizonErr}$.	\item We additionally truncate $A^{-1}$ expansion at degree-$D$ for each dangling $M^{-1}$-path. We defer these error terms for the truncation section.
	\item The three-way concatenation to correct for the proper shapes in $\fp_j(Q_i) - \fp_j(Q_{i-1})$ may incur another collection of (non-well-behaved) intersection terms as we address in the previous section. We call any error in this process $\mathsf{VerticalErr}(i)$. We discard these terms in the iteration update $\Delta_i$;
	\item We additionally discard diagonal terms in the update $\Delta_i$. These are small terms for which an iterative update is unnecessary, and we discard them to ensure $Q$ is definitionally a linear combination of backbone-dangling shapes recalling that they must be off-diagonal terms. Analogously, a similar trimming process is also applied to the initialization of $Q_0$.
\end{enumerate}
This is summarized in the following process. For any $i\ge 0$, define the proper residuals (to be corrected) by
\[
    H_0^{\mathrm{prop}}
    \coloneqq
    \sum_{j=2}^{D}b_j\fp_j(Q_0),
\]
and, for \(i\geq1\),
\[
    H_i^{\mathrm{prop}}
    \coloneqq
    \sum_{j=2}^{D}b_j
    \bigl(\fp_j(Q_i)-\fp_j(Q_{i-1})\bigr).
\]

Let \(\corr^{\leq D}(H_i^{\mathrm{prop}})\) denote the correction obtained after truncating
each MP expansion in the dangling \(M^{-1}\)-paths at degree \(D\) in the three-way concatenation of \(
\calL^* M^{-1} \calL \left( H_i^{\mathrm{prop}} \right)
\), and
let
\[
    \Delta_i
    \coloneqq
    \operatorname{Proper}\!
    \left(\corr^{\leq D}(H_i^{\mathrm{prop}})\right),
\]
where \(\operatorname{Proper}(\cdot)\) retains the proper
backbone-dangling shapes after linearizing all well-behaved
intersections. Importantly, intersection shapes from the vertical three-way concatenation as well as diagonal terms are discarded.
 
 Finally, we also apply similar procedure to the initialization at $Q_0$ to remove negligible components therein. Let the resulting matrix be $\operatorname{Proper}(\calR)$. 
 \begin{mdframed}[linewidth=0.2pt]
\textbf{Summary of the Concrete Truncated Iterative Procedure.}

\medskip
\textbf{Initialization:}
Set
\[
    Q_0\coloneqq-C_F \cdot \operatorname{Proper}(\calR),
\]
so that \(I+Q_0/C_F\) satisfies the affine constraints.

\medskip
\textbf{Iterative update \(i\to i+1\):}
For \(0\leq i<t^\ast\), set
\[
    \Delta_i
    \coloneqq
    \corr_{\mathrm{prop}}^{\leq D}
    \bigl(H_i^{\mathrm{prop}}\bigr),
    \qquad
    Q_{i+1}\coloneqq Q_i+\Delta_i.
\]

\medskip
\textbf{Output:}
Let
\[
    Q\coloneqq Q_{t^\ast}.
\]
The final matrix is obtained from
\(\frac1{C_F}F^{\leq D}(Q)\) by removing the aggregate discrepancy
defined below and correcting the final proper residual.
\end{mdframed}
\paragraph{Affine invariant for the concrete iteration.}
To compare the concrete iteration with the ideal one, define
\[
    \mathsf{HorizonErr}(i)
    \coloneqq
    \sum_{j=2}^{D} b_j
    \bigl(P_j(Q_i)-\fp_j(Q_i)\bigr)
\]
and
\[
    \mathsf{UpdateErr}(i)
    \coloneqq
    \Delta_i-\corr\!\left(H_i^{\mathrm{prop}}\right) = \corr_{\mathrm{prop}}^{\leq D}
    \bigl(H_i^{\mathrm{prop}}\bigr) -\corr\!\left(H_i^{\mathrm{prop}}\right) .
\]
Here \(\mathsf{UpdateErr}(i)\) collects the MP-truncation error and the
discarded vertical intersection terms as well as diagonal terms. Define the cumulative discrepancy
\[
    \mathsf{IterErr}_i
    \coloneqq \mathsf{HorizonErr}(i)
+
        \sum_{s=0}^{i} \mathsf{UpdateErr}(s) \,.\]

\begin{proposition}[Invariant for the concrete iteration]
\label{prop:concrete-iteration-invariant}
For every \(i\geq0\),
\[
    \calL\!\left(
        \frac1{C_F}
        \bigl(F^{\leq D}(Q_i)-\mathsf{IterErr}_i\bigr)
    \right)
    -
    \1_m
    =
    \frac1{C_F}\calL\!\left(H_i^{\mathrm{prop}}\right).
\]
\end{proposition}

At this point, we observe that there are three sources of error terms: vertical, horizontal, and MP truncation. We defer the last type of truncation term to the subsequent section alongside other truncation details. The remainder of this section is then dedicated to formalizing and establishing error bounds for vertical (\cref{sec:vertical-details}) and horizontal (\cref{sec:chebyshev-Q}) concatenation that culminate in \cref{lem:error-vertical} and \cref{lem:semicircle-shape-concatenation}.




\subsection{Combinatorial Analysis of the Error Terms for Vertical Concatenation}%{Defer Combinatorial Analysis for Vertical Concatenation}
\label{sec:vertical-details}

%\paragraph{Error Terms for Vertical Concatenation}
We now shift gear to considering the error terms, and bound the error term by the following lemma.
\begin{lemma}[Error Terms for Vertical Concatenation] \label{lem:error-vertical}
	Recall that we define 
	\[ 
	\mathcal{L}^*(w_\tau) = \mathsf{Proper}(\eta_\tau) + \mathsf{Error}(\eta_\tau)\,,
	\]
	we have \[ 
	 \|\mathsf{Error}(\eta_\tau)\|_{sp}=o_d(1)\,.
	\]
\end{lemma}

 There are primarily three sources of error terms:\begin{enumerate}
	\item Non well-behaved intersections in the concatenation of $A^{-1}$ and $\eta_\tau$;
	\item Terms from rank-$2$ components in $M^{-1}$;
	\item Diagonal terms from the vertical concatenation.
\end{enumerate}
For the error term analysis, we will again apply block-value bound throughout.
For starters, we consider the error terms from the concatenation of $A^{-1}$. We give a charging scheme for $A^{-1}$ terms, and then show it also applies for the rank-$2$ components.

\paragraph{Pre-processing Well-behaved Intersections} Since the intersection shapes may consist of a well-behaved intersection with extra vertex intersections, we first process away the well-behaved intersection term by removing the initial square vertex and its incident edges. With this pre-processing, we may assume there is no well-behaved intersection around the initial and final square vertex of the $A^{-1}$-path.

 This further allows us to deduce the following claim that shows either the backbone-path is unaffected in the concatenation operation, or we can identify an intersection property that can later be used to deduce a block-value gap.


%Our crucial observation is that as we process away well-behaved intersection around the first square vertex, we may assume that the separator vertex is at either boundary.

%\begin{claim}[Pe-processing Intersection Property] \label{clm:preprocessing-int-property}
%	There is at least one separator vertex along the backbone-path, i.e., from the two boundary circle vertices or the initial square vertex, or we have a mismatching intersection around the initial square vertex- it is incident to $3$ distinct vertices, with one pair of edges collapsing.
%\end{claim}
%\begin{proof} Note that this is immediate if there is no vertex intersection that collides with vertices on the backbone path.
%
%	On the other hand, by our preprocessing operation, there is no well-behaved intersection around the initial square vertex. The crucial observation is that the square vertex cannot be incident to only two distinct vertices.
%	
%	Suppose it does, the pre-processing operation would have applied and give us a contradiction. 	
%	Finally, if it is incident to $3$ distinct vertices, we have a mismatching intersection around the initial square vertex as described.
%\end{proof}

\paragraph{Preprocessing the \(A^{-1}\)-expansion.}
Throughout this subsection, every occurrence of \(A^{-1}\) is replaced by its
truncated expansion in properly concatenated \(\alpha\)- and \(\beta\)-gadgets.
Terms containing an MP approximation error have already been shown to have
negligible spectral norm in~\cref{lem:mp-error-quant} and are absorbed into
\(\operatorname{Error}(\eta_\tau)\).

We also preprocess all well-behaved top, bottom, and top--bottom intersections
from~\cref{def:well-behaved-intersection-vertical}. Namely, whenever the
incident edge-pairs around \(s_{\mathrm{bp}}\) or \(s_\eta\) match, we apply
the linearization operation of~\cref{def:well-behaved-intersection-vertical}.
%If the intersection is isolated, its linearization is included in
%\(\operatorname{Proper}{\calL^*}(w_\tau)\). If additional intersections
%remain, we retain the linearized record as an error record.
Consequently, every error term considered below has at least one vertex
intersection that is not removed by the well-behaved linearization.

%\begin{claim}
%	In a pre-processed intersection shape $\tau_I$ from $\eta_\tau$, there is no "isolated" vertex: i.e., no vertex can be making a first and final appearance within the same block-step.
%\end{claim}
%\begin{proof}
%    This is straightforward for boundary vertices by definition. Next, since $\eta_\tau$ is injective, no edge can appear twice within $\eta_\tau$, while some edge may intersect with the edges on the backbone path. However, this may only appear for the edges around the initial square vertex. For the square vertex, unless it is incident to two distinct vertices exactly which is removed by our pre-processing operation, the square vertex cannot be making its first and final appearance in this step.
%\end{proof}

 

\propLocalTraversal*

We introduce the following definitions for our analysis. Recall from our norm bounds analysis that we bound the norm of a graph matrix by analyzing length-\(2q\) ``block-walks'' of the underlying shape and controlling the vertex and edge factors contributed by each ``block-step.'' To carry out this local analysis, we must distinguish between the \emph{local} and \emph{global} appearances of labeled vertices and edges. We remind the reader that a labeled square vertex is an element of \([m]\), a labeled circle vertex is an element of \([d]\), and a labeled edge is an element of \([m]\times[d]\).

\begin{definition}[Local and global appearances]
    Consider a fixed block-walk. An appearance of a labeled vertex within a given block is called a \emph{local appearance}, whereas an appearance of that vertex anywhere in the full block-walk is called a \emph{global appearance}.

    A labeled vertex or edge makes its \emph{global first} or \emph{global last appearance} if the corresponding occurrence is its first or last appearance in the entire walk. Similarly, it makes its \emph{local first} or \emph{local last appearance} if the corresponding occurrence is its first or last appearance within the current block.
\end{definition}

We also introduce a special class of steps that will be treated differently by our charging scheme. The terminology ``reverse-charging'' will become clear once the edge-copy charging rule is described below.

\begin{definition}[Reverse-charging step/edge (Def.4.10 in~\cite{HKPX23})]
\label{def:reverse-charging}
    Fix a block in a given block-walk. We call a step \(u\to v\) \emph{reverse-charging} if the following conditions hold:
    \begin{enumerate}
        \item the underlying edge is making its global last appearance in the walk;
        \item the global first appearance of the same underlying edge also occurs within the current block;
        \item that first appearance traverses the edge in the opposite direction, namely from \(v\) to \(u\).
    \end{enumerate}
\end{definition}
% 
% We recall the charging scheme as follows.
% \begin{mdframed} \textbf{Top-Down Traversal for Edge-Copy Charging} \begin{enumerate}  
%\item  \textbf{Edges from an $A^{-1}$ Path}: Traverse vertically from top to bottom, and assign each edge-copy to the vertex it leads to unless this is a \emph{reverse-charging edge} as in~\cref{def:reverse-charging}, in which case we assign to the source vertex.
%\item  \textbf{Edges from Chebyshev Concatenation}: \begin{itemize}
%	\item Instead of viewing a Chebyshev path as a horizontal path which is useful for illustration, it is more intuitive to visualize them as again a $V$-shape vertical attachment connecting the two anchor vertices ;
%	\item Traverse top-down from both anchor vertices, and assign each edge-copy to the vertex it leads to unless this is a \emph{reverse-charging edge},, in which case we assign to the source vertex.
%
%	\end{itemize} 
%\item \textbf{Anchor Edges}: these are the edges between $A^{-1}$-path and a horizontal path, or the edges that are part of the final $h_2$ attachment. \begin{itemize}
%	\item Assign both edge-copies to the source square vertex if this is a \emph{reverse-charging edge} as in \cref{def:reverse-charging}, i.e., the underlying edge-pair explore the source square vertex earlier on in the local traversal process;
%	\item Otherwise, we case on it is incident to $2$ $h_1$-edges or an $h_2$-edge:
%	 \item If we have $2$ $h_1$ edges, assign one edge-copy to the "left" circle vertex, and we treat the other edge-copy as an \textbf{attachment-reserve};
%	 \item If we have an $h_2$ edge, assign one edge-copy to the vertex, and designate other edge-copy as an \textbf{attachment-reserve}.
%\end{itemize}
%\end{enumerate}
%\end{mdframed}

We first describe a traversal process for a backbone-dangling shape and then define the charging scheme induced by this traversal.


%\begin{mdframed} \textbf{Top-Down Traversal for Edge-Copy Assignment}:
% Assign depth-label for the vertices in the our backbone-dangling shape as follows. This assignment process is oblivious of the vertex-labels.  \begin{enumerate}  
%	\item \textbf{Base Case}: the boundary vertices in $U\cup V$ are of depth-$0$, as well as the square-vertex on the backbone path;
%	\item \textbf{Traversal for Dangling $A^{-1}$-path:} view each $A^{-1}$-path attached to the square vertex in a top-down manner.  
%	\item   \textbf{Traversal for Chebyshev Concatenation}: \begin{itemize}
%	 \item Each horizontal-level is attached by two \emph{anchor edges} to the prior square vertex. For $\tau\neq \emptyset$, we have two anchor edges connecting two distinct circle vertices. For the special case of $\eta_0$, we consider the $h_2$ edge as consisting of two anchor-edges as well.
%	 \item Without loss of generality, we assign one anchor edge to the left anchor vertex;
%	 \item We traverse each horizontal path from left to right.
%	 	\end{itemize} 
%\end{enumerate}
%\textbf{Charging via Traversal:} 
%
%\begin{enumerate}
%	\item \textbf{Backbone Edges:} Assign both edges on the backbone path to the square vertex .
%	\item \textbf{Dangling Edges in any $A^{-1}$ Path:} Assign each edge to the vertex it leads to.
%	\item \textbf{Horizontal Edges:} for each horizontal attachment, for each square vertex within a horizontal-interval, assign both incident edges to the square vertex.
%	\item \textbf{Attachment Edges:} for each horizontal attachment, we reserve their assignment.
%\end{enumerate}
%\end{mdframed} 
%
%We make the following immediate observations:

%
\begin{mdframed} \textbf{Top-Down Traversal for Edge-Copy Charging}:
 Assign depth-label for the vertices in the our backbone-dangling shape as follows. This assignment process is oblivious of the vertex-labels.  \begin{enumerate}  
	\item \textbf{Base Case}: the boundary vertices in $U\cup V$ are of depth-$0$, as well as the square-vertex on the backbone path;
	\item \textbf{Traversal for Dangling $A^{-1}$-path:} view each $A^{-1}$-path attached to the square vertex in a top-down manner.  
	\item   \textbf{Traversal for Chebyshev Concatenation}: \begin{itemize}
	 \item Each horizontal-level is attached by two \emph{attachment edges} to the prior square vertex. For $\tau\neq \emptyset$, we have two attachment edges connecting two distinct circle vertices. For the special case of $\eta_0$, we consider the $h_2$ edge as consisting of two attachment-edges as well.
	 \item Without loss of generality, we assign one attachment edge to the left attachment vertex, and call the other edge \emph{attachment-reserve}.
	 \item We traverse each horizontal path from left to right.
	 	\end{itemize} 
\end{enumerate}
\textbf{Charging via Traversal:} 

\begin{enumerate}
	\item \textbf{Backbone Edges:} assign both edges to the square vertex $s_{bp}$;
	\item \textbf{Dangling Edges and Subsequent Horizontal Levels:} assign each edge to the vertex it leads to, unless it is a \emph{reverse-charging edge}, in which case we assign the edge to the source. \end{enumerate}
\end{mdframed} 


\begin{proof}  
We make the following observations that are immediate from the charging scheme:
\begin{displayquote}
	 Square vertices that make their first appearance within any $A^{-1}$-path, $s_{bp}$ the square vertex on the initial backbone path, and  any circle vertex outside the $U\cup V$ boundary have been assigned the required factors.
\end{displayquote}
Any such vertex is assigned 1 or $2$ edges for its first local appearance, and these edges then protect the potential subsequent final appearance of the vertices within the same block-step by ``reverse-charging''.

Putting these vertices aside, it suffices for us to focus on square-attachment vertices (the vertices at the end of some $A^{-1}$-path that additionally spawn an additional horizontal level or $h_2$-gadget), and circle vertices within the boundary.

\paragraph{Square-Attachment Vertex (other than $s_{bp}$):}
 For any vertex that makes its first appearance as a square-attachment vertex, we first observe that it cannot make both its first and final appearances within the same block-step. This follows by our $3$-way concatenation property: the deviation-attachment $\tau$ that induces the correction term is a proper shape. Hence, any square-attachment vertex would come from the $\tau$ piece of the horizontal attachment, and it must be intersected by some vertex from $s_{bp}$ or the interior of the $A^{-1}$ path. Hence, the vertex is not making its \emph{first} appearance as a square-attachment vertex.

\paragraph{Finding Attachment Reserve}
 That said, we still need to assign two edge-copies to such vertices. It receives one factor from the edge leading to it in the traversal on the horizontal path, and we will identify an additional edge from the reserve-edge.  Without loss of generality, we may assume the left attachment edge is not reverse-charging to the square vertex (since by our preprocessing of well-behaved intersection, both edges cannot be both reverse-charging). We next claim that we can find an unassigned edge as \textbf{attachment-reserve}: starting from the other attachment edge, it is not assigned unless it is a reverse charging for the circle vertex. In that case, we observe that the square vertex must have made a prior appearance, and therefore any of the two edges assigned in the most recent appearance can be used as \textbf{attachment reserve}.

%
%
%
%, for any vertex reached via edges in $A^{-1}$ path (including the final square vertex for any $A^{-1}$-path) has received the assigned factors: each circle vertex is reached via one edge, and each square vertices via two edges. Moreover, one (or two) extra edge-factors are assigned if such vertex makes both first and final global appearance in the same block-step by design of the reverse-charing.
%
% In addition, this also holds for any circle vertex along a horizontal path. Therefore, it suffices for us to consider any square vertex: any such square vertex has only ben assigned one edge in the traversal edge, while two is needed to offset its corresponding factor. In addition, two more factors would be needed if the square vertex makes first and final appearance within the same block-step.
% 
% We now observe that any such square vertex cannot make its first "local" appearance within any horizontal level (since $\eta_\tau$ is injective), and therefore must appear within the $A^{-1}$-path: if it is within the interior of a dangling $A^{-1}$-path, the prior argument applies to show appropriate factors are assigned; otherwise, it is assigned an edge from the backbone-path, again giving us four factors from two edges for its first and last appearance factors as desired.
% 
%
%Next, to assign an edge-copy for such a square vertex,
% notice that any such square vertex splits out a deviation-attachment via two attachment-edges. The case of $h_2$ attachment follows from the proof of prop. 4.14 in \cite{HKPX23}, and it suffices for us to consider non-trivial horizontal attachment.
%  Without loss of generality, we may assume the left attachment edge is not reverse-charging to the square vertex (since by our preprocessing of well-behaved intersection, both edges cannot be both reverse-charging). We next claim that we can find an unassigned edge as \textbf{attachment-reserve}: starting from the other attachment edge, it is not assigned unless it is a reverse charging for the circle vertex. In that case, we observe the square vertex must have made prior appearance, and therefore the two edges assigned to the most recent appearance can be used as \textbf{attachment reserve}.
 
\paragraph{Charging Boundary Vertices:}
We first notice that we can apply the above \emph{finding-reserve} idea to the initial square vertex $s_{bp}$ as well to identify a reserve edge. Next, apply the corollary 4.16 and its preceding proposition 4.15 in \cite{HKPX23}:
\begin{displayquote} either one boundary vertex requires a $\sqrt{d}$ factor (a single edge copy), or there is a vertex making a middle appearance outside $U\cup V$
\end{displayquote}
It is clear in the first case as we assign the edge-copy to the contributing vertex. In the latter case, their assigned factor from the local traversal process is unnecessary and can be combined with the \textbf{backbone-reserve} edge and assigned to both boundary vertices.
\end{proof}


%Applying the above to $A^{-1}$ terms gives us the desired bound. We defer the formal proof to the end of this section. We now extend this analysis to the rank-$2$ component in $M^{-1}$.

\paragraph{Extension to Diagonal Terms} Recall that in the formal iterative process, we additionally discard the diagonal terms at the end of each vertical concatenation instead of applying Chebyshev polynomials with correction iteratively. The analysis follows from the corresponding component in the proof for proposition 4.15 from~\cite{HKPX23}.

%\todo{in construction}
\paragraph{Extension to Woodbury Terms} To get started, it is useful to recall the following scalar bounds. \HyperParameterBounds*
As a result, we also have\[
\frac{1}{s^2 - ru } = \Theta \left(\frac{d}{m}\right)\,.
 \]
Additionally, we split $M_W$ into three components depending on the inner term, 
\[
	M_W \coloneqq  \frac{1}{s^2-ru}\, A^{-1} \left( \underbrace{ u\,\frac{\1_m\1_m^T}{d}}_{M_J} -  \underbrace{s\,\frac{\eta\1_m^T+\1_m\eta^T}{d}}_{M_C} + \underbrace{r\,\frac{ \eta\eta^T}{d} }_{M_\eta} \right) A^{-1}
	 \,,\]
Following the language in \cite{HKPX23}, each $M_W$ term can be viewed as either concatenating a permutation jump ($\1_m \1_m^\top$ or $\eta\cdot \1_m^\top)$ at the end of the first $A^{-1}$ path, or an extra $\eta$-gadget ($\eta\eta^\top$).
At a high level, we continue applying the traversal-charging scheme described above, albeit with the following caveats:
\begin{enumerate}
	\item For the $\1_m\1_m^T$ term as well as the $\1_m\eta^\top$ factor, the charging for the  first $A^{-1}$ path is essentially identical up to the final square-attachment vertex as it is not incident to an $h_2$-attachment, and thereby missing a factor of $1/\sqrt{d}$ from the corresponding reserve factor;
	\item We no longer necessarily have vertices making middle appearances (via non-trivial vertex-intersections), hence we need to identify a slack from other factors beyond just an application from the prior scheme.
\end{enumerate}

\paragraph{Analysis for Terms in $M_J$}
We make the following observations:
\begin{enumerate}
	\item Each term is scaled by a coefficient of \[ 
	\frac{1}{s^2-ru} \cdot u  \cdot \frac{1}{d} = O(\frac{d}{m}) \cdot O(1) \cdot \frac{1}{d} =\frac{1}{m} \,.  
	\]
	\item Each term can be described as follows: we attach an additional permutation-jump between square vertices (corresponding to $\1_m\1_m^T$), and from the destination square vertex, we start a new $A^{-1}$ dangling path along with deviation attachment of $\eta_\tau$ at the end;
	\item The permutation-jump requires an additional factor of $\sqrt{m}$ for the first local appearance, call the destination square vertex $s$.  
	\item Unless $s$ is the destination of some edge in the local traversal process, $s$ cannot make its first and final global appearances in the same block-step;
	\item For the remaining $A^{-1}$ path, traverse from the start vertex, and assign each edge to the vertex it leads to. 
	\item For the $\eta_\tau$ attachment to the final square vertex $s$ along the dangling $A^{-1}$ path, there are two cases:
	\begin{itemize}
		\item Assign both anchor edges to the square vertex if the final square vertex  $s$ is incident to $2$ distinct vertices (i.e., the pair of its neighbors in $A^{-1}$ path and in $\eta_\tau$ intersect in pairs). Again, we can apply partial linearization to remove the $s$ vertex and its incidental edges.
		\item Apply the horizontal traversal scheme: assign one edge to the circle vertex, and traverse horizontally, assign the edge to the vertex it leads to. Finally, we declare the other anchor-edge as reserve.
	\end{itemize} 
	 \end{enumerate}
Therefore, we have a total block-value of \[ 
O(1) \cdot \frac{1}{m} \cdot \sqrt{m} \cdot O(1) = o_d(1)\,.
\]
where the $O(1)$ factor comes from our local charging for the $A^{-1}$ backbone path as well as the floating $A^{-1}$-path and deviation attachment rooted at $s$.
%  Intuitively, this comes from the following observation:
% \begin{enumerate}
%  	\item $A^{-1}$ is an injective path consist of $\al$ and $\beta$-gadgets that ends with a square-vertex;
%  	\item Combining the scalar with the normalization of $\frac{1}{s^2-ru}$, each term comes with the following coefficient (up to $\Theta(1)$ factors), \begin{itemize}
%  		\item $M_J$ comes with a coefficient of $ \frac{d}{m} \cdot u \approx  \frac{d}{m}  $;
%  		\item $M_C$ comes with a coefficient of $\frac{d}{m}  $;
%  		\item $M_\eta$ comes with a coefficient of $\frac{d }{m}  \cdot \frac{m}{d} \cdot \frac{1}{1-\gamma}  = \Theta(1)\,.$
%  	\end{itemize}
%  	\item The only dominant term comes from $M_\eta$ which gives  \[A^{-1} \cdot  \frac{1}{s^2-ru}\cdot  r \cdot \frac{\eta\eta^\top}{d} A^{-1} \approx \left( A^{-1}\eta  \right)\cdot  \left( \frac{1}{d}\eta^TA^{-1} \right)\,; \]
%  	\item Notice we group out $A^{-1}\eta$ as that is indeed the form of the main dangling term for the base iteration), and therefore we expect the first term continues to be a dominant term (i.e., there is no inverse-polynomial gap in $d$);
%  	\item However, for the second term, we have for $\eta_i\neq \eta$, we have \[ 
%  	\left( \frac{1}{d} \cdot  \eta^TA^{-1} \cdot \eta_i \right )\ 
%  	\]
%  	We claim this is $o_d(1)$ norm for $\eta_i\neq \eta$. 
%  	\begin{itemize}
%  		\item For any $\eta_i \neq \eta$, $\eta^T A^{-1} \eta_i$ is of mean-$0$ (i.e., there is no square term). To see this, $\eta_i$ is a concatenation of $t>1$ $\al$ or $\beta$ gadgets by design of correction scheme only applied to $P_{t}$ for $t\geq 2$.
%  		\item To see why it is $o_d(1)$, consider the following charging process that traverses the expansion for $\eta^T A^{-1} \eta_i$ from the initial circle vertex corresponding to $\eta$ attachment.. Our goal is to show that a factor of $\tilde{O}(1/\sqrt{d})$ from the normalization factor of $\frac{1}{d}$ is sufficient for the circle vertex. In contrast, a full factor of $1/d$ (aka the full normalization factor) needs to be assigned to the circle vertex for $\eta^\top A^{-1} \eta$.
%  	\end{itemize}
% \end{enumerate} 
%
%\begin{proof} We now formally prove this via block-value bounds. Without loss of generality, we may assume $\eta_i$ does not contain any floating component from $M_W$, as our proof will in turn show that any floating component comes with a block-value gap of $\tilde{O}(1/\sqrt{d})$.  For starters, we apply graph-matrix decomposition:\begin{enumerate}
%	\item Each term is a backbone-dangling shape with an $A^{-1}$ dangling path;
%	\item The first $A^{-1}$ path may terminate in either a square vertex without $\eta$-attachment, or terminates with a square $h_2$-attachment;
%	\item For the terms that do not terminate with an $\eta$-attachment, we need to identify an additional $O(1/\sqrt{d})$ gap (recall that one of the $1/\sqrt{d}$ factor from the final $h_2$-attachment is assigned to the initial square vertex);
%	\item We then apply block-value bound for the subsequent terms: starting from a circle vertex in the case of $\eta^\top A^{-1}$ terms, and a square vertex in the case of $\1_m A^{-1}$ terms.
%	\item Each gadget in an $A^{-1}$-path contributes $O(1)$ block-value.
%\end{enumerate}
\paragraph{Analysis for Terms in $M_\eta $}	
This is the dominant term where we use $\eta_i\neq \eta$ for $i\geq 1$. Note that this is a term that contains an $h_2$-attachment for the backbone-dangling path, and the coefficient for this term is  $1/d$.  We apply the following charging argument in which we split $1/d$ into two copies evenly:
\begin{enumerate}
	\item The backbone-dangling path contains an $h_2$-attachment, hence it suffices for us to apply the charging argument to the backbone-dangling path and bound the component to have a block-value $O(1)$. Next, it suffices for us to bound the block-value corresponding to the floating component.
	\item Assign $\tilde{O}(\frac{1}{\sqrt{d}}$) value to the initial circle vertex $s$ that is part of the $h_2$ attachment in the $\eta^\top A^{-1} \eta_i $-path. This leaves us an excess of $\tilde{O}( \frac{1}{\sqrt{d}})$ from the scalar coefficient;
	\item Our key observation is that the initial-circle vertex $s$ cannot be making both of its first and final appearances in a $\eta^\top A^{-1} \cdot \eta_i$ component for $i\neq 0$, and the remaining vertex factors can be charged by the edge-factors in the component, hence leaving us with an excess of $\frac{1}{\sqrt{d}}$.
\end{enumerate}
%Towards that end, we use the ideas adapted from prop. 4.14 in \cite{HKPX23}: we split each $h_2$ edge as two edge-copy such that each edge-copy carries a factor of $O(1/\sqrt{d})$, and  we would like to assign edge-copy to vertices in a local traversal process. The up-shot of the local traversal process is to assign factors to vertices locally that upper bounds their contribution in a global walk:
%\begin{enumerate}
%	\item A vertex appearance for the first time in the walk may correspond to its first/final appearance in the walk, in this case we will assign a local bound given by half of its vertex weight, matching the global bound;
%%	\item If a vertex makes a first and last appearance in the same local traversal, it is assigned a factor of full vertex weight, matching the global bound;
%	\item Any mismatch of first/last in a local-traversal with a middle appearance in the global walk corresponds to a $\tilde{O}(1/\sqrt{d})$ gap.
%\end{enumerate}

\paragraph{Analysis for Terms in $M_C$}

%\paragraph{Analysis for Terms in \(M_C\).}
Recall that
\[
    M_C
    =
    -\frac{s}{d(s^2-ru)}
    A^{-1}
    \bigl(\eta\1_m^\top+\1_m\eta^\top\bigr)
    A^{-1},
    \qquad
    \left|\frac{s}{d(s^2-ru)}\right|
    =
    O\left(\frac1m\right).
\]
The case of $\eta = \eta_\emptyset$ is covered in the previous work of \cite{HKPX23}.
Thus, it suffices for us to consider \(\eta_i\neq\eta\), and we have
\[
    \calL^*(M_C\eta_i)
    =
    -\frac{s}{d(s^2-ru)}
    \left[
        \bigl(\1_m^\top A^{-1}\eta_i\bigr)
        \calL^*(A^{-1}\eta)
        +
        \bigl(\eta^\top A^{-1}\eta_i\bigr)
        \calL^*(A^{-1}\1_m)
    \right].
\]
We bound the two terms as follows:
\begin{enumerate}
    \item
    In the first term, \(\calL^*(A^{-1}\eta)\) is a
    backbone-dangling component with an \(h_2\)-attachment and has
    block-value \(O(1)\). The \(\1_m\)-endpoint in the floating
    component contributes at most \(\sqrt m\), while the remaining
    factors are charged as in the \(M_J\)-case. Hence this term has
    block-value
    \[
        O\left(\frac1m\right)\cdot\sqrt m\cdot O(1)
        =
        O\left(\frac1{\sqrt m}\right).
    \]

    \item
    In the second term, the \(\1_m\)-endpoint in
    \(\calL^*(A^{-1}\1_m)\) again contributes at most \(\sqrt m\).
    After paying this factor, the remaining normalization is
    \[
        O\left(\frac{\sqrt m}{m}\right)
        =
        O\left(\frac1d\right),
    \]
    since \(m=\Theta(d^2)\). The preceding \(M_\eta\)-analysis therefore
    applies directly to the floating component
    \(\eta^\top A^{-1}\eta_i\), leaving a block-value gap
    \(\widetilde O(d^{-1/2})\).
\end{enumerate}
Consequently,
\[
    B_q\bigl(\calL^*(M_C\eta_i)\bigr)
    \leq
    \widetilde O\left(
        \frac1{\sqrt m}+\frac1{\sqrt d}
    \right)
    =
    \widetilde O\left(\frac1{\sqrt d}\right)
    =
    o_d(1).
\]



%\input{alternative_charging.tex}

%
%
%\begin{proof}
%%	We first make the following observations: any circle vertex in the shape, and any square vertex that makes its first appearance outside of horizontal path, has received the corresponding required vertex factors. Their first local appearance is a destination of some non-reverse-charging edge, and their possibly final appearance receives the corresponding factors from reverse-charging edge.
%%
%%
%Similarly to the proof of prop.4.14 in~\cite{HKPX23}, we first observe that the following vertices have been assigned their corresponding required factors by the top-down traversal scheme:
%\begin{enumerate}
%	\item Any labeled-circle vertex that makes its first local appearance outside $U\cup V$: any such vertex is explored via being destination of some edge that has assigned to it the corresponding first-appearance factor. On the other hand, if the vertex additionally makes a subsequent "final" appearance in the block, the corresponding edge is reverse-charging and has its factor assigned to this vertex as well;
%	\item The same reason applies to any labeled square vertex that makes its first local appearance within some $M^{-1}$-path. In other words, this is any labeled square vertex that makes its first appearance at some square vertex in the shape that does not contain a \emph{deviation-attachment}. 
%\end{enumerate}
%
%Therefore, it remains for us to focus on the following two type of vertices. \begin{enumerate}
%	\item Circle vertices in $U\cup V$: there is no edge-copy assigned to any of them;
%	\item Square vertices that have deviation-attachment are missing half of its factor: It is only assigned one edge-copy (a factor of $O(1/\sqrt{d})$) for its first local appearance, while a factor of $\sqrt{m}\approx d$ is needed. The absence of an extra edge-copy is identical for the corresponding reverse-charging factors.
%\end{enumerate}
%
%We next establish the following:\begin{enumerate}
%	\item For circle vertices in$U\cup V$: even though there are $2$ vertices, only one of them may contribute first/last appearance factor of $\sqrt{d}$ in the global walk, and it suffices for us to assign $1$ single edge-copy via the \textbf{backbone-reserve};
%	\item For any square-attachment vertex, we can identify an extra edge-copy 
%\end{enumerate}
%
%	 
%\end{proof}

%Norm Bounds
\subsection{Norm Bounds for Linear Combination of Backbone-Dangling Shapes}

In this section, we prove norm bounds for general weighted linear combination of backbone-dangling shapes. We use $K$ as a placeholder for such matrix. Our main result is the following theorem for the spectral radius of $K$, which shows it obeys semicircular distribution up to the edge.

We do not explicitly state our result for the bulk of the spectrum, while noting that by similar analysis in \cite{PotechinXu2026Theta}, our results also establish semicircular spectrum by considering $O(1)$ moments of the trace.

%\begin{theorem}[Spectral Radius for Linear Combination of Backbone-Dangling Shapes]
%\label{thm:Q-norm}
%	For a symmetric linear combination of backbone-dangling shapes $K$ with ground-set $\calB(K)$ s.t. \[ 
%	K = \sum_{\tau \in\calB(M)} c_\tau \cdot \cm_{\tau}\,,
%	\] we have \[ 
%	\|K \|_{sp} \leq (2+o_d(1))  ,
%	\]
%%<<<<<<< HEAD
%	provided $\Var(K)=1$
%%=======
%%	provided $\Var(R)=1$,
%%>>>>>>> 258cc9d9278f3d3112b6597fe09243ecb5da8fa1
%	where we define \[
%	\Var(K) \coloneqq \frac{1}{d}\Tr(KK^\top) =   \sum_{\tau\in\calB(K)} c_\tau^2 \cdot \Var(\tau)\,.
%	 \]
%\end{theorem}

\begin{theorem}[Spectral Radius for Linear Combination of Backbone-Dangling Shapes]
\label{thm:Q-norm}
	For a symmetric linear combination of backbone-dangling shapes $K$ with ground-set $\calB(K)$ s.t. \[ 
	K = \sum_{\tau \in\calB(M)} c_\tau \cdot \cm_{\tau}\,,
	\] we have \[ 
	\|K \|_{sp} \leq (2+o_d(1))  ,
	\]
	provided $\Var(K)=1$ 
	
	Quantitatively, let $\|c(K)\| = \sum_{\tau\in B(K)} |c_\tau|$ and $D_V$ be the size bound for shapes in $K$,  with high probability, \[ 
	\|K \|_{sp} \leq 2 + B_q(\mathsf{NonIdeal}) 
	\]	for \[ 
	B_q(\mathsf{NonIdeal}) \coloneqq  O(1) \cdot 3^{|V(\tau)|} \cdot   c(\tau) \cdot    \frac{1}{\sqrt{n}} \cdot \|c(K)\|_1^{  2|V(\tau_i)|}  \cdot (q\cdot D_V)^{2}\,.
	\]
\end{theorem}
	

We follow the ideas for the analogous result in \cite{PX25} by examining the trace moment calculation for $K^{2q}$ directly. \begin{enumerate}
	\item  We first apply the insight from block-value assignment scheme: it suffices to focus on $F$- and $R$-block-steps (i.e. block-steps in which all by one circle vertex are making their first/final appearances in the wlak) , as all other contributions can be shown to be negligible.

%	\item To show we can matching $F/R$-steps with the exact edge-set (and consequently, the underlying shape), we show the following: \begin{itemize}
%%		\item  For each $F$-step $i$, the subsequent block-step $j>i$ that shares some edge with $i$, is either an $R$-step with the same set of edges, or it is a slack step.
%%		\item Similarly, for each $R$-step $j$, we
%	\end{itemize}
    \item \textbf{Ideal pairs.}
If an \(F\)-step using shape \(\tau\) is matched to an \(R\)-step with the same labeled
edge-set, we pick up a total value of 
\[
   (c_\tau)^2 \cdot  \Var(\tau) 
\]
Summing over \(\tau\) gives the ideal-pair value
\[
    B_q(\text{ideal})^2 \le \Var (K)
\]

    \item \textbf{Non-ideal steps.}
Every non-ideal step is charged to a block-value gap, i.e. to an additional
middle-appearance vertex.  Such a gap gives a factor \(d^{-1/2}\) after normalization.
The only remaining issue is that a single gap may be used to pay for several coefficient
mismatches. We use the gap-witness accounting below bounds this overhead quantitatively. 

\end{enumerate}


The key component of our argument is to identify a slack from any non-ideal step. Towards that end, we introduce some preliminaries.

\paragraph{Preliminaries for Trace Moment Calculation}

\begin{definition}[F/R/S-Step]
 For a block-step $t$ in a trace walk of backbone-correction shapes from the current boundary vertex $U_{t}=u$ to the next boundary vertex $V_{t} =v$, we classify it as following:
\begin{enumerate}
\item $F$-step: each vertex in $\tau_t$ is making a first appearance except the current boundary vertex $u$ (which by definition cannot).
\item   $R$-step: each vertex in $\tau_t$ is making a last appearance except the current boundary vertex $v$ (which by definition cannot). 
 \item Any other step is an $S$-step.
\end{enumerate}

\end{definition}
Additionally, we make a further classification of $F/R$-steps by highlighting steps that can be bundled into steps of the same edg-set.
\begin{definition}[Ideal $F/R$-Step]
For a block-step $t$ in a trace walk of backbone-correction shapes from the current boundary vertex $U_{t}=u$ to the next boundary vertex $V_{t} =v$, we classify it as following:
\begin{enumerate}
\item Ideal-$F$-step:  There is a subsequent step (possibly an $S$-step to be defined) with the same underlying edges-set, i.e. there is some step $t'>t$ s.t. $E(\tau_t) = E(\tau_{t'})$. 
\item   Ideal-$R$-step: There is a a prior step (possibly an $S$-step) with the same underlying edges-set, i.e. there is some step $t'<t$ s.t. $E(\tau_t) = E(\tau_{t'})$ (pick $t'$ to be the smallest one if there are multiple).
 \item Any other $F/R$ step is a \emph{non-ideal} F/R-step.
% \item 
\end{enumerate}
 \end{definition}

\begin{proposition}[Block-Value for Ideal Pairs]
	For any $\tau$, an ideal-$F/R$-pair of shape $\tau$ gives block-value under our block-value assignment scheme, \[ 
	B_q(\mathsf{Ideal}(\tau) ) \leq (c_\tau)^2 \cdot \Var(\tau)\,. 
	\]
\end{proposition}
\begin{proof}
	Firstly, suppose the $F$-step has the same start vertex as the destination vertex as the $R$ step, then the corresponding block-value bound is precisely given by $\Var(\tau)$ by definition. On the other hand, suppose not and we call this an $R'$-step, we observe that the block-value assignment is irrespective of the specific vertex label - it only depends on vertex appearance status. Therefore, the corresponding $R'$-step receives the same value as the $R$-step, and our proposition holds again. 
	
	The extra factor of $c_\tau$ is added as we picked up two factors of $c_\tau$ in these two corresponding blocks for the trace moment calculation of the weighted linear combination.
\end{proof}
Summing over all $\tau\in \calB(K)$ gives us the following desired bound.
\begin{corollary}
	\[ 
	B_q(\mathsf{Ideal}) = \sum_{\tau \in \calB(K)} B_q(\mathsf{Ideal}(\tau) ) = \sum_{\tau\in\calB(K)} (c_\tau)^2 \cdot \Var(\tau) = \Var(K)
	\]
\end{corollary}

Next, to facilitate our accounting of block-value gaps, we introduce the following quantity that helps us keep track of each block-step that incurs a gap in the corresponding block-value does not get its gap assigned to \emph{too many} other block-steps.
%\paragraph{Gap Witness Vertex}
\begin{definition}[Gap Witness-Vertex]
	For any block-step $\tau_i$, let $S$ be the set of separator vertices in $V(\tau_i)$, i.e. the vertices making a middle appearance in $\tau_i$, when $|S|>1$, we call each vertex $v\in S$ a gap witness-vertex for block-step $i$.
\end{definition}

	It is crucial to notice a block-step only incurs gap witness vertices if it contains more than $1$ vertex in the separator, corresponding to when we expect to incur a gap in the block-value.

\paragraph{Finding Slack for Mismatched Block-Steps}

We are now ready to introduce the core of our argument in matching $F,R$ steps- finding some block-steps with a gap for each non-ideal $R$-steps. Notice that we assume all vertices except one - the vertex in $V_t$ - are making their final appearance throughout this section, as otherwise such a block-step is either not an $R$-step, or locally comes with a block-value gap in itself. 


\begin{definition}[Vertex Weight]\label{def:vtx-wt} 
	We assign each circle vertex a weight of $1$ and a square vertex of weight $2$. For a set of vertices $S$, we define its weight as the sum of its vertex weights.
	
	Additionally, as an edge-case, we assign a circle vertex in $\beta$-gadget a weight of $2$ as opposed to $1$. 
\end{definition}

The crux of our analysis is the following combinatorial argument that helps us identify extra vertices making middle appearance (i.e. extra separator vertices) for two block-steps sharing any edge. 


\begin{claim}[Structural Property of Backbone Dangling Shapes] \label{clm: structural-clm}
	For any backbone dangling shape, and any $2$ edge-coloring, there is a weight $\geq 2$ separator where a separator is defined as the set of vertices incident to edges of both colors. Moreover, in the case the separator contains $2$ circle vertices, it does not contain circle vertex of any final $h_2$ attachment.
\end{claim}

\begin{proof}
	Since the reduced shape $\tau$ is a single connected component and not all edges have the same color, there exists at least one vertex incident to both edges. Moreover, consider the case that this is a circle vertex, as otherwise we have found a weight-$2$ separator.
	
	Next, we observe that any such vertex cannot be a boundary vertex in $U_\tau $ or $V_\tau $, nor a circle vertex in any final attachment gadget as well, as any such vertex is also degree-$1$ (crucially we view $h_2$ edge as a single edge in $\tau$). Therefore, we may start with the circle vertex promised, let it be $c$, and observe that it must be part of some $\al$ and $\beta$ gadget. 
	
	Suppose it is part of a $\beta$ gadget, then we are again done as this is a single circle vertex that has weight $2$ by our rule in~\cref{def:vtx-wt}. On the other hand, it is part of an $\al$-gadget, and it is straightforward to observe that there must be another vertex (possibly square) that is incident to edges of both color. To see this, note that we have two edge-disjoint paths at $c$ (i.e., a cycle) and the edges incident to $c$ are additionally of different colors.
 \end{proof}
 
We now use the above structural property to identify slack in block-steps.

\begin{proposition}[Simultaneous Appearance] \label{prop:simultaneous-appearance} Given an $L$-block-step at block-step $i$, the edge(s) in the current block-step must have appeared altogether (potentially with some other edge) in a previous $F$-step, or there is a local-gap (potentially from a previous step).

 In other words, unless there is a local-gap (potentially from a previous step), there is a block-step $t<i$ s.t. $E(\tau_i) \subseteq E(\tau_t)$, and moreover, $\tau_t$ is an $F$ block-step.
 
 Formally, at least one vertex making a final appearance at the current block-step $i$ is a block-witness vertex for some previous block-step, particularly step-$t$.
\end{proposition}


%


\begin{proof} \label{proof: simul-appearance}
Let $\tau_i$ and $\tau_t$ be the shape of the block-steps at each time. Note that we may have $\tau_i\neq \tau_t$, and moreover, $\tau_i, \tau_t$ comes with underlying edges labeled in $[n]$ as well.
 In the graph of $\tau_t$, color any labeled edge in $E(\tau_t)\cap E(\tau_i)$ that are making $F$ appearance in the latest (most recent) block-step $t$,  red, and color any other edge blue.

%
%	Let $E(\tau_i)$ be the edge-set of the current block-step in inspection. Color any edge in $E(\tau_i)$ that are makes $F$ appearance the latest (most recent) block-step $t$,  red, and color any other edge blue.
	
%	 We now claim that, unless all edges are red, any red edge-component intercepts some blue-edge component in at least two vertices. Suppose there is only one such vertex, call it $r$, and let $b$ be a neighbor of $r$ alongside blue-edge. We reach a contradiction as there is only one vertex-disjoint path between $r$ and $b$ if there is no other vertex in the red component incident to a blue edge. 

%	We claim that, unless all edges are red and the previous $F$-step is thereby identified, 
	Apply \cref{clm: structural-clm} on any red edge-component $C$ in $\tau_t$,  unless all edges are red and step-$t$ satisfies the simultaneous-appearance property for block-step $i$,  $C$ has at least two
distinct vertices that are incident to some blue edge otherwise.


% Suppose some red
%edge-component $C$ has exactly one such vertex $r$. Since not all edges are red, there exists
%a blue edge and hence some vertex $b$ outside $C$ that is incident to a blue edge.
%Then every path from any vertex in $C\setminus\{r\}$ to $b$ must pass through $r$, so $r$
%is a cut-vertex of $\tau_i$, contradicting $2$-vertex-connectivity.	 
Finally, notice any red-component vertex incident to some blue edge is a vertex making a middle appearance at block-step $t$, therefore at least two vertices are making middle appearance at time-step $t$. That is the block-step with a gap in block-value as desired. Importantly, observe the following,
\begin{enumerate}
	\item  Both blue vertices are gap witness vertices for block-step $t<i$ by construction (each middle appearance vertex is a block-value witness of the step when there is more than $1$ middle-appearance vertex) ;
	\item At least one of the blue vertices is making a final appearance at step-$i$ (since this is an $L$-step, there is at most one-vertex not making a final appearance);
	\item Therefore, we can conclude there is at least one final-appearance vertex in the current block that is a block-witness vertex for some previous step (in particular, block-step $t$).
\end{enumerate}
	 \end{proof}
%	 
%\begin{proposition}[Maximal Appearance] \label{prop:maximal-appearance}  Let $i$ be the current block-step, and let $t$ be the previous $F$-step identified by the last proposition, either we have $E(\tau_t)\subseteq E(\tau_i)$, or there is a local gap. 
%\end{proposition}
%\jnote{I think we only one gap of $1/\sqrt{n}$ here from the first partial closure - it is possible to incur one gap, while followed by $O(|V(\tau_i)|$ many steps with unmatched coefficients. We need to handle this by truncation which poses a bound on $O(|V(\tau_i)|) \leq c \log n $ for some small constant $c<1$ .}
%
%\begin{proof}
%	Suppose $E(\tau_t)\subseteq E(\tau_i)$ is not true, we have some extra edge that appears in the block-step $t$: color $E(\tau_i)\subseteq E(\tau_t)$ red in $\tau_t$, and the rest of the edges $E(\tau_t)\subseteq E(\tau_i) $   blue. Let $C$ be a red-component, apply \cref{lem:2-vtx-disjoint-paths}, there are at least two vertex-disjoint blue edges $e_1, e_2$ incident to $C$. Let their blue endpoints and red endpoints be $r_1\neq r_2$	 and $b_1\neq  b_2$.
%	
%	Notice all the blue edges incident to any red component must be closed before step-$i$ by our assumption. 
%	 We next give an inductive argument to find a (local) gap that can be assigned to step-$i$.
%	\begin{enumerate}
%		\item If $e_1$ and $e_2$ are closed (appearing for the last time) at the same block-step, this means we have a block-step containing both $b$ and $b'$. However, by assumption these two vertices are both making a middle appearance, hence the local gap.
%		\item Otherwise, if $e_1$ or $e_2$ are closed at a step in which $b$ or $b'$ is not the only middle-appearance vertex, i.e., closed in a block-step in which we can assign a gap to either $b$, or $b'$, this is also a local gap identified.
%		\item Otherwise, $e_1$ (and $e_2$) must each be closed in a block-step in which $b$ and $b'$ is the unique middle-appearance vertex, in this case, color both edges red and recurse on them and their incidental vertices.
%			\end{enumerate}
%	This process necessarily terminates as we can see each edge processed as being oriented towards the red-component, and consider the path between $r_1$ and $r_2$ across the blue component, there must be some vertex that have multiple out-going edges, which corresponds to some step that contains at least two middle-appearance vertices. Assign the most recent one to the current-step.
%\end{proof}
%

\begin{proposition}[Simultaneous Closure]\label{prop:simul-closure}
	Given an $F$-block-step with edge-set $E(\tau_i)$, unless all edges are simultaneously closed, there is a local gap at the first subsequent block-step $t>i$ that uses some edge in $E(\tau_i)$. 
	
	Formally, at least one vertex making first appearance at step-$i$ is a gap witness-vertex for some block-step at the walk, specifically at step-$t$.
	\end{proposition}
\begin{proof}
	Suppose not all of $E(\tau_i)$ are closed (i.e. appearing for the final time) in the subsequent block-step that closes some edge in $E(\tau_i)$. Consider the subsequent block-step $\tau_t$ that uses some edge in $E(\tau_i)$, our goal is to identify that there are at least two vertices in $\tau_t$ making a middle appearance, forming a local gap in the block-value - these two vertices will be our candidates for a gap witness-vertex for block-step $i$. Since at most one of them is not making a \emph{first} appearance at the current step $i$, this proves the desired.
	
	In $\tau_i$, color the edges that appear in $E(\tau_i)$ red, and the remaining edges blue. Apply \cref{clm: structural-clm} on the red-component, it has at least two distinct vertices incident to some blue edge in $\tau_i$. By assumption, any blue-edge is unclosed since each makes their first appearance at $\tau_i$, and remains unclosed at $\tau_t$ - hence, they are bound to make future appearance again. Finally, notice that the two incidental vertices identified also appear in $V(\tau_t)$, and since each is incident to some unclosed edge bound to appear, each is making a middle appearance. This completes our proof that $\tau_t$ is a slack step.%	
%	
%	 Call such a block-step an \emph{intermediate closure } w.r.t.  $E(\tau_i)$, as it partially closes some edge in  $E(\tau_i)$. Additionally, call a block-step a \emph{final closure}  w.r.t.  $E(\tau_i)$ if it closes the final edge in $E(\tau_i)$.
%	
%	It suffices for us to show each intermediate closure step has a gap locally, and moreover, by splitting a local gap, this implies a local gap for the final closure as well.
%	
	
\end{proof}


\paragraph{One Gap to Rule All the Mismatching Coefficients from Individual Block-Step }

Finally, we exploit the definition of gap-witness vertex to make our identified slack-step quantified.

\begin{observation}
	Each gap witness-vertex is only assigned to a block-step in which it appears for the first or final time by~\cref{prop:simultaneous-appearance} and~\cref{prop:simul-closure}. In other words, each gap witness vertex is in charge of at most $2$ block-steps for violations of simultaneous appearance or closure.
\end{observation}


\paragraph{Formal Bounds for Linear Combination}
%\subsection{Bounding Block-Value for Linear Combination $M$}

We now give a formal description of how coefficients are assigned to each block-step, and complete the proof for the block-value bound for linear combination $M$. We apply the following rule to assign coefficient from the weighted linear combination in the block-value assignment:
\begin{enumerate}
	\item For any ideal-$F/R$-block and any $S$-blcok-step, assign coefficient $c_\tau$ to the same block;
	\item For any non-ideal $F$-step of shape $\tau$, assign the coefficient $c_\tau$ to the block-step-$t$ identified via~\cref{prop:simul-closure};
	\item For any non-ideal $R$-step of shape $\tau$,\begin{itemize}
		\item  If its underlying edge-set violates the simultaneous-appearance property, assign the coefficient $c_\tau$ to block-step $t$ identified in~\cref{prop:simultaneous-appearance};
	\item Otherwise, its underlying edge-set comes from some $F$-step that violates  the ``simultaneous closure'' property, assign the coefficient $c_\tau$ to the block-step-$t$ identified via~\cref{prop:simul-closure}.
	\end{itemize}
\end{enumerate}




\begin{claim}
	For any $F$ block-step $i$ that violates simultaneous-closure property, its associated coefficient assigned from the block-value assignment scheme is at most  \[ 
	\sum_{\tau_1,...\tau_r: \text{assigned to step-}i } \prod_{\tau_i} |c(\tau_i)| \leq \left( \sum_{\tau \in \calB(K)}|c(\tau)| \right)^{|V(\tau_i)|}  =  \|c(K)\|_1^{  |V(\tau_i)|/2}
	\]
\end{claim}

\begin{corollary}Each non-ideal block-step gets assigned a coefficient factor of at most \[ 
\|c(K)\|_1^{  |V(\tau_i)|/2}\cdot \|c(K)\|_1^{|V(\tau_i)|} \leq \|c(K)\|_1^{2|V(\tau_i)|}
\,.\]
\end{corollary}
 	
 	
\paragraph{Bounds for Non-Ideal Steps}
Next, we verify the non-ideal shapes.
\begin{proposition} \label{prop:block0non-ideal-M}
	For any non-ideal step it gets assigned a bound at most \[ 
	B_q(\mathsf{NonIdeal})  \leq O(\frac{1}{\sqrt{d}}) \cdot  (3\|c(K)\|_1)^{ 3\cdot D_V} \cdot (q\cdot D_V)^2 \,, 	\]
	where $D_V \leq O(\log n) $ is the size-limit s.t. $|V(\tau)|\leq D_V$ for any $\tau\in \calB(K)$.
\end{proposition}
\begin{proof}
	For each non-ideal step of shape $\tau$ with weight$>1$ separator, we have a block-value of
	\begin{align*}
		B_q(\tau) &\leq \sum_{2\leq v\leq |V(\tau)|} 3^{|V(\tau)|} \cdot   c_\tau  \cdot \sqrt{n}^{|V(\tau)|-v }\cdot \|c(K)\|_1^{  2|V(\tau_i)|}  \cdot  (q\cdot D_V)^{v}\\
%		&=  3^{|V(\tau)|}  \cdot c(\tau) \cdot (\frac{1}{\sqrt{n}})^{|V(\tau)|-1} \cdot \sqrt{n}^{|V(\tau)|-2 }\cdot |\calB(M)|^{  |V(\tau_i)|/2}\\
%		&\leq \frac{1}{\sqrt{n}} \cdot  3^{D_V} \cdot |c(\tau)| \cdot |\calB(M)|^{  |V(\tau_i)|/2}
	\end{align*}
	where we note 
	\begin{enumerate}
		\item $3^{|V(\tau)|}$ is a trivial bound specifying vertex appearance of any vertex in a given $\tau$;
		\item $c_\tau = c(\tau) \cdot c_n(\tau) = c(\tau) \cdot (\frac{1}{\sqrt{n}})^{|V(\tau)|-1}$ is the coefficient for shape $\tau$ in $K$;
		\item Any non-slack vertex outside the separator gives factor $\gam<1$ under our edge-charging scheme; 
		\item Each slack block comes with at least an $\frac{1}{\sqrt{d}}$ for extra separator vertex;
		\item $(q\cdot D_V)^{v}$ is an upper bound for the factor restricted to vertices making middle appearances;
		\item $\|c(M)\|_1^{  2|V(\tau_i)|} $ is an upper bound for coefficient factor assigned to this step.
	\end{enumerate}
	Provided $D_V\leq O(\log n)$ ,the above series is dominated by $v=2$, giving us a bound of \[ 
	B_q(\tau) \leq O(1) \cdot 3^{|V(\tau)|} \cdot   c(\tau) \cdot    \frac{1}{\sqrt{n}} \cdot \|c(K)\|_1^{  2|V(\tau_i)|}  \cdot (q\cdot D_V)^{2}
	\]
	
	Summing over all possible shapes $\tau$ gives us
	\begin{align*}
		B_q(\mathsf{NonIdeal})  &\leq O(1) \cdot \sum_{\tau} \frac{1}{\sqrt{d}} \cdot  3^{D_V} \cdot |c(\tau)| \cdot (3\|c(K) \|_1)^{  |V(\tau_i)|2} (q \cdot D_V)^2\\
		&\leq O(\frac{1}{\sqrt{d}}) \cdot  (3\|c(K)\|_1)^{ 3\cdot D_V} \cdot (q\cdot D_V)^2 
	\end{align*}
	% Finally, this is $\frac{1}{n^\delta}$ for $D_V=O_\delta( \sqrt{\log n}  )$ for $|\calB(M)|\leq c^{D_V}$ for any constant $c$.
 \end{proof}
 
We are now ready to wrap up this section by proving our main theorem for the spectral radius of a weighted linear combination of backbone-dangling shapes.

\begin{proof}[Proof to~\cref{thm:Q-norm}]
	We wrap up the proof for block-value bound. For each step, there are two cases by our factor assignment scheme- ideal-$F$ and ideal-$R$, each is assigned a value of $ \sqrt{\Var(K)}.$  Any non-ideal step gives block-value captured by our auxiliary function $B_q(\text{Non-Ideal})$. Therefore, we have \[ 
	\E[\Tr(KK^\top)^{q}] \leq d \cdot (2\sqrt{\Var(K)} + B_q(\text{Non-Ideal}))^{2q}
	\]
 	 	The proof for norm bound then follows from \cref{claim:trace-to-norm-rough}.
\end{proof}
	


%




\subsection{Chebyshev Polynomial and Horizontal Concatenations}
Throughout this section, we use $K$ to denote an arbitrary linear combination of backbone-dangling shapes. 


\begin{theorem}[Orthogonal Polynomials and Shape Concatenation for Semicircular Distributions (Formal)]  
\label{thm:SC-concate}
For any linear combination of backbone-dangling shapes \[
K = \sum_{\calB(K)} c_\tau \cdot \cm_\tau \,,
 \]
	from a ground set $\calB(K)$ of size $|\calB(K)|$ and each shape has at most $D_V$ vertices. Then, for any $t>0$, we have \[ 
	P^{\mathsf{sc}}_t(K) = \fp_t(K) + O( \frac{\poly(t, D_V) }{\sqrt{d}} )
	\]
	Quantitatively, we have \[ 
	P^{\mathsf{sc}}_t(K) = \fp_t(K) + \text{SC-Error}(t)
	\]
	for \[
	\|\text{SC-Error}(t)\|_{sp} \leq O(2^{t}) \cdot B_q(\text{non-ideal})	\]
	with $B_q(\text{non-ideal})$ the auxiliary function defined from~\cref{thm:Q-norm}.
\end{theorem}

\paragraph{Key Ingredient for Cancellation}

The overall structure is identical to analysis of orthogonal polynomial approximation for $A$, except now we use the basis with respect to the semicircular distribution with variance-$1$ and the recurrence now follows by Chebyshev polynomial of the second kind.


The main idea is the following cancellation lemma that shows $K^2$ is either proper concatenation or the identity matrix which shall be canceled out by the Chebyshev recurrence relation.
\begin{lemma}[Backtracking Residue for SC Polynomials]  
\[ 
 K^2 = \sum_{\tau \in \calP_2(K) } c_\tau \cdot \cm_\tau + I_d + o_d(1)\,, 
\]
provided \[ 
\Var(M) = 1\,.
\]
\end{lemma}

To shed light on the above lemma, it is crucial for us to highlight a phenomenon unique to the cancellation for $Q$ that is rather distinctive from that of the prior analysis for the Marchenko-Pastur distribution. Recall that in the prior analysis, the cancellation is made possible by backtracking intersections; in particular, this happens for any two consecutive steps with the exact same edge set, as shown in \cref{prop:diamond-cancellation} and analogously in \cref{prop:half-diamond-cancellation}. These two criteria may be viewed as combinatorial interpretations of the dominant terms in a length-$2$ walk (e.g., $\Tr(AA^\top )$ for diamond cancellations). This is formally proven via~\cref{prop:backtracking-sc}.

% That said, the corresponding adaptation to the analysis of $Q$, and concretely, the dominant terms in $\Tr(QQ^\top )$ requires some minor adjustment.

% The reason for the adjustment is that each summand in $Q$, unlike those in $A$, is not necessarily a product of mean-$0$ random variables, and hence the dominant term for $2$-step may contain a mix of backtracking terms, as well as fresh random variables (that are not mean-$0$). Combinatorially, we appeal to the following characterization, which we call \emph{generalized backtracking terms}.

% \paragraph{Generalized Backtracking Terms from Shape Preprocessing}

% To start with, we adopt the following preprocessing operation to identify the component of each shape in $Q$ that requires backtracking (revisiting) in the dominant term. Coincidentally, we peel off the component that does not require revisiting,

% Recall in each backbone-dangling shape, each \(A^{-1}\)-path is locally
% injective: it is represented as a proper concatenation of \(\alpha\)- and
% \(\beta\)-gadgets. The final \(h_2\)-attachment, however, is not necessarily
% injective. To remove this discrepancy, we preprocess each shape by eliminating the
% intersection created by the final \(h_2\)-attachment at the end of each
% \(A^{-1}\)-path.
% %This creates a minor notational discrepancy when matching backtracking steps: in some cases, the dominant contribution does not literally come from matching edge sets, even though it contributes the same value and should be grouped into the same variance term.



% \begin{definition}[\(h_2\)-Terminal Intersection]
% Consider an \(A^{-1}\)-path whose terminal gadget is an \(M_\beta\)-gadget
% incident to a square vertex \(s\), where the path is oriented according to the
% natural direction of the \(A^{-1}\)-attachment. Suppose an \(h_2\)-attachment is
% added at the terminal square vertex \(s\). We say that this attachment creates
% an \emph{\(h_2\)-terminal intersection} if the circle vertex of the attached
% \(h_2\)-edge is identified with the circle vertex of the terminal
% \(M_\beta\)-gadget.
% \end{definition}

% With the notion of an \(h_2\)-terminal intersection in place, we now define an operation that peels off the final square vertex from the intersected circle vertex. 
% \begin{definition}[Peeling off an \(h_2\)-Terminal Intersection] Consider an \(A^{-1}\)-path whose terminal gadget is an \(M_\beta\)-gadget incident to a square vertex \(s\). Suppose that the final \(h_2\)-attachment at \(s\) creates an \(h_2\)-terminal intersection. We define the  \emph{reduced} shape to be the shape obtained by removing the terminal square vertex \(s\), together with the two \(h_2\)-edges incident to \(s\) and the intersected circle vertex. All other vertices, edges, and boundary labels are left unchanged. \end{definition}

% \begin{remark}
% 	For a given shape with multiple terminal intersection, the reduced shape has all the terminal intersections removed.
% \end{remark}

% \begin{claim}[Shape Reduction Preserves Coefficient] \label{clm:shape-reduction-coef}
% For a backbone-dangling shape $\tau$ with coefficient $c_\tau$ in the iterative construction $Q$, let its reduced shape be $\tau'$ with coefficient $c_\tau'$ in $Q$, they have the same coefficient \[ 
% c_\tau = c_{\tau'}
% \,.\]
% \end{claim}
% \begin{proof}
% 	The coefficient is oblivious of the specific gadgets used in any $A^{-1}$ path, and the peeling off operation is equivalent to removing a $\beta$ gadget at the end of any dangling $A^{-1}$ path. 
% 	\end{proof}

% \begin{claim}[Shape Reduction Preserves Connectivity]\label{clm: shape-reduction-conn}
% 	Let $\tau$ be a backbone-dangling shape, its reduced shape $\tau'$ is also a single connected component.
% \end{claim}
% \begin{proof}
% 	The proof is identical as above in noticing the removal process only removes the final square vertex.
% \end{proof}

% From the peeling-off operation, we define the key ingredient for backtracking cancellation for the analysis of $Q$ as follows.

% \begin{definition}[Generalized Backtracking Cancellation]
% 	 Consider two block-steps $\tau_F$ and $\tau_R$ of the underlying backbone-dangling shape $\tau$, and let $\tau_{\text{reduced}}$ be the reduced shape of $\tau$. We say they form a backtracking cancellation if \begin{enumerate}
% 	 	\item Any vertex in $V(\tau_{\text{reduced}})\setminus V(\tau)$ appears for the first time in $\tau_F$ and $\tau_R$. In other words, these are the vertices that are intended to receive first appearance at each dominant step;
% 	 	\item The remaining edges in $E(\tau)$ match as an edge-set in $\tau_F$ and $\tau_R$.
% 	 \end{enumerate} 
% \end{definition}

% \begin{proposition}[Cancellation via Generalized Backtracking]
% For $R = \sum_{\tau\in \calB(R)} c_\tau \cdot \cm_\tau,
% $ we have
% 	\[ 
% 	\sum_{\tau \in Q}  c_{\tau}^2 \cdot  \cm_{\bti(\tau) }   = \Var(R) \cdot I_d+ o_d(1) \,.
% 	\]
% \end{proposition}

% \begin{proof}
% 	\todo{to fill in: give an indirect analysis and show any term not captured by above is small in spec norm }
% \end{proof}



\subsection{Formal Analysis for Chebyshev Polynomials and Concatenations}
We now give a formal analysis for the error terms that arise when we switch from Chebyshev polynomials of $R$ to graph matrices of concatenated shapes in $R$.

\label{sec:chebyshev-Q}

Firstly, we formally define well-behaved intersections at the heart of the cancellation.
\begin{definition}[Well-behaved Backtracking Intersections for SC Polynomials]
	Given $\tau \in \calB(K)$ that are backbone-dangling shapes, we say the composition $\tau\cdot \tau'$ forms a backtracking intersection if \begin{enumerate}
		\item $U_{\tau} \equiv V_{\tau'} $: the left-boundary vertex of $U_\tau$ matches with the right-boundary vertex of $V_{\tau'}$ ;
		\item $E(\tau) \equiv E(\tau')$: the underlying edge-set of the two sets receive the same labels.  
	\end{enumerate}
	We denote the corresponding backtracking intersection pattern as $\bti(\tau)$.
\end{definition}
The crux of the above observation is the following proposition that shows backtracking intersections gives a copy of identity, enabling the key cancellation for Chebyshev polynomials. 
%\begin{remark}
%	In particular, it should be noted that given a shape $\tau \in \calB(K)$, the collection of shapes $\tau'$ that may form backtracking-intersection with $\tau$ is any shape with the same matching edge-set (without enforcing matching vertex-locations in the shape).  
%\end{remark}

\begin{proposition}[Backtracking Residual for SC Polynomials] \label{prop:backtracking-sc}
	Define \[ 
	\Delta_{\bti} \coloneqq  \sum _{\tau \in \calB(K)} (c_\tau)^2 \cdot  \cm_{\bti(\tau)} - I_d \,,	\]
	we have \[ 
	\|\Delta_{\bti}\|_{sp} = o_d(1)\,.
	\]
	provided $\Var(K)=1$.
\end{proposition} 
\begin{proof}[Proof Sketch]
For any shape $\tau$, we have \[ 
c_\tau^2 \cdot \cm_{\bti(\tau)}  = c_\tau^2 \cdot  \Var(\tau) \cdot I_d -o_d(1)
\]	
with the discrepancy from falling factorial of the vertices while we assign each a full factor of $d$ and $m$ and the fluctuation of second moment Gaussian random variables from their mean. 

Summing over all $\tau\in \calB(K)$ gives the desired.
 We defer the formal proof to \cref{app:def-Q}.
\end{proof}

\begin{definition}[Local-Collision Pieces for Chebyshev Polynomials (SC-LCP)]
\label{def:SC-LCP}
Fix a concatenation record
\[
    \tau_t\cdot \tau_{t-1}\cdot \dots \cdot \tau_1,
%    \qquad \tau_i\in\{\alpha,\beta\}.
\]
A consecutive sub-record
\[
    \Theta
    =
    \tau_b\cdot \tau_{b-1}\cdot \dots \cdot \tau_a
\]
is called a \emph{local-collision piece} for Chebyshev polynomials if it is an MP-LCP as in~\cref{def:local-collision-pieces}, except that each constituent shape can be any backbone-dangling shape as in~\cref{def:backbone-dangling-shape}.

\end{definition}

	It should be pointed out that this is the precise translation of \cref{def:local-collision-pieces} for MP-polynomials to the matrix setting via Chebyshev polynomials of the second kind.
	
	
\begin{proposition}[Structural Property for Error Terms (Chebyshev)]
\label{prop:chebyshev-error-term} 
[Analog of \cref{prop:error-term-A-equivalence}]
For $K$ a linear combination of backbone dangling shapes with variance $\Var(K)=1$, any error term in $\text{Error}(t) = P_t(K)- \fp_t(K) $
admits an LCP-decomposition.
%
%Equivalently, after the MP recurrence cancels all isolated backtracking
%pairs, every surviving error term is generated by terminal collisions. In
%particular, there is no unresolved isolated backtracking intersection between
%adjacent blocks in the LCP-decomposition.
\end{proposition}

\begin{proof}
	The proof is identical to that in \cref{prop:error-term-A-equivalence} except that we only need to consider the full backtracking intersection of the reduced shapes with the corresponding backtracking residual.
	
	We apply the induction from the proof of~\cref{prop:error-term-A-equivalence}. The base cases, the observation
that deleting the leftmost shape preserves an LCP-decomposition, and
the proper-attachment, non-backtracking-collision, and non-isolated
backtracking cases are unchanged after replacing the MP gadget alphabet
by \(\calB(M)\) and MP backtracking by generalized backtracking.

It remains only to modify the isolated-backtracking case, and this is the only difference from the proof of
\cref{prop:error-term-A-equivalence}. In particular, since the
Chebyshev recurrence has no analogue of the term
\(-\gamma P^{\mathrm{sc}}_t(M)\), there are no half-diamond or half-flat
cancellation cases. The induction therefore gives the claimed
SC-LCP-decomposition by the cancellation from~\cref{prop:backtracking-sc}. 
\end{proof}

\paragraph{Reduction to the block-value bound.}
It remains to apply the block-value bound below to the
SC-LCPs appearing in this decomposition.
%	\todo{a full spelled-out verification is likely needed} 



%\subsection{Error Bounds for $Q$}

\begin{lemma}
For a locally-collision-piece $\tau$ of length-$t$, we have the following block-value bound   
\[
B_q(\tau) \leq \prod_{i=1}^{t-1} 2^{t-1}  \cdot  B_q(\text{non-ideal}) \]
for $B_q(\text{non-ideal})$ defined from \cref{thm:Q-norm}.
\end{lemma}
\begin{proof}
	


Apply the block-value scheme to the non-terminal component $\tau_1\circ\dots \circ \tau_{t-1}$ by viewing $\tau$ as a specific walk of length-$t$ that arises in the trace moment calculation,  we he have, \[
B_q(\text{Non-Terminal} ) \leq  2^{t-1} \,.
 \]
% Summing over all choices of $\tau_i\in \calB(K)$ gives us the second inequality.
  
 Finally, we notice the the terminal-gadget is a slack block-step in the analysis for \cref{thm:Q-norm}, and it suffices for us to bound it by the auxiliary function of $B_q(\text{non-ideal})$. 
\end{proof}
 
We are now ready to wrap up the proof for the equivalence theorem of Chebyshev polynomials and graph matrices of concatenated shapes.
\begin{proof}[Proof to~\cref{thm:SC-concate}]
	This follows by the the block-value function for any LCP that gives an upper bound for the expected trace, which then yields a norm bound via~\cref{claim:trace-to-norm-rough}.
\end{proof}
 
 % 	Our main idea is to show that each block-step can be classified as an $F/L$ block-step defined as follows.
%%\begin{definition}[Multi-$\beta$ Dangling $A^{-1}$ Path and Shape] We call an $A^{-1}$ path, i.e., a concatenation of $\{\al,\beta\}$ a multi-$\beta$ path if the dangling $A^{-1}$-path contains at least two $M_\beta$ gadgets. 
%%Moreover, we consider a backbone-dangling shape a multi-$\beta$ shape if it contains a multi-$\beta$ dangling path. 
%%This will be considered as an $o_d(1)$ error term.
%%\end{definition}
%
%\begin{definition}[Vertex Weight]\label{def:vtx-wt} 
%	We assign each circle vertex a weight of $1$ and a square vertex of weight $2$. For a set of vertices $S$, we define its weight as the sum of its vertex weights.
%	
%	Additionally, as an edge-case, we assign a circle vertex in $\beta$-gadget a weight of $2$ as opposed to $1$. 
%\end{definition}
%
%
%
%\begin{definition}[$F/L$ Block-Step]
%Consider a block-step of a backbone-dangling shape $\tau$ within a walk. Provided $\tau$ is not a multi-$\beta$ dangling shape, we call this block-step an $F$ step if every vertex in $V(\tau)$ except the left circle boundary vertex $U_\tau$ makes its first appearance in the walk. Similarly, we call it an $L$ step if every vertex in $V(\tau)$ except the right circle boundary vertex $V_\tau$ makes its final appearance in the walk.
%
%When the context is clear, we simply refer to a block-step as a step.
%\end{definition}
%
%\begin{definition}[Block-Value Gap Step]
%	For any step that is not an $F/L$ step, including a block-step with underling shape a multi-$\beta$ dangling shape, we call it a block-value gap step, abbreviated as a gap step. 
%\end{definition}

%%We show that each \(F\)-step can be bundled with an \(L\)-step of the same
%underlying shape. Moreover, after the preprocessing described below, the two
%steps have identical edge sets. We will use the following two propositions to
%establish this matching.



%\paragraph{Finding Matching Steps for Reduced Shapes}

%\begin{proposition}[Simultaneous Closure or Block-Value Gap] \label{prop:sim-closure}
%	For a given $F$-step of reduced shape $\tau_i$ , all first-appearance vertices are closed (appearing for the final time) in a later step, or there is a step with block-value gap that can be locally identified.
%	
%	Formally,  there is a local gap at the first subsequent block-step $t>i$ that uses some edge in $E(\tau_i)$: at least one vertex-weight making first appearance at step-$i$ is a gap witness-vertex-weight for some block-step at the walk, specifically at step-$t$.
%\end{proposition}
%
%\begin{proof}
%		Suppose not all of $E(\tau_i)$ are closed (i.e., appearing for the final time) in the subsequent block-step that closes some edge in $E(\tau_i)$, concretely step-$t$. Our goal is to identify that there are at least a vertex-weight of $2$ in $\tau_t$ making a middle appearance, forming a local gap in the block-value - these two vertices will be our candidates for gap witness-vertex for block-step $i$. Since at most one of them is not making a \emph{first} appearance at the current step $i$, this proves the desired.
%		
%		In the graph $\tau_i$, color edges that make their subsequent appearance in $\tau_t$ (i.e., $E(\tau_i)\cap E(\tau_t)$ red, while the remaining edges in $E(\tau_i)\setminus E(\tau_t)$ green). Apply the following~\cref{clm: structural-clm} gives a set of vertices of weight-$2$ incident to both green and red edges. Finally, notice any vertex in the identified separator must be making a middle appearance in $\tau_t$, and that gives us the desired block-value gap from step-$t$.
%% 
%		
%			
%\end{proof}
%
%\begin{claim}[Structural Property of Backbone Dangling Shapes] \label{clm: structural-clm}
%	For any reduced backbone dangling shape, and any $2$ edge-coloring, there is a weight $\geq 2$ separator where a separator is defined as the set of vertices incident to edges of both colors. Moreover, in the case the separator contains $2$ circle vertices, it does not contain circle vertex of any final $h_2$ attachment.
%\end{claim}
%
%\begin{proof}
%	Since the reduced shape $\tau$ is a single connected component and not all edges have the same color, there exists at least one vertex incident to both edges. Moreover, consider the case that this is a circle vertex, as otherwise we have found a weight-$2$ separator.
%	
%	Next, we observe that any such vertex cannot be a boundary vertex in $U_\tau $ or $V_\tau $, nor a circle vertex in any final attachment gadget as well, as any such vertex is also degree-$1$ (crucially we view $h_2$ edge as a single edge in $\tau$). Therefore, we may start with the circle vertex promised, let it be $c$, and observe that it must be part of some $\al$ and $\beta$ gadget. 
%	
%	Suppose it is part of a $\beta$ gadget, then we are again done as this is a single circle vertex that has weight $2$ by our rule in~\cref{def:vtx-wt}. On the other hand, it is part of an $\al$-gadget, and it is straightforward to observe that there must be another vertex (possibly square) that is incident to edges of both color. To see this, note that we have two edge-disjoint paths at $c$ (i.e., a cycle) and the edges incident to $c$ are additionally of different colors.
% \end{proof}
% 
% We now complete the proof of the block-value bound for terminal gadget in an LCP for $M$.
%\jnote{make sure "reduced" is removed everywhere}
% \begin{proof}[Proof of \cref{lem:block-val-gap-Q}]
% 	% We apply a two-step argument:\begin{enumerate}
% % 		\item  A gap for the fresh component among the square vertices that get peeled off is  straightforward;
%
% % 		\item Apply the structure lemma to identify a gap in the reduced shape, i.e., the component that requires backtracking;
% % 		\item Any shape that we consider here falls into at least one category of the above by definition, and this guarantees a slack for any error term considered.
% % 	\end{enumerate}
%	% We focus on the second condition where we want to identify a gap from the reduced shape provided the final gadget $\tau_t$ does not form a backtracking intersection. Without loss of generality, we assume each vertex in the fresh component is "new" and hence condition-$1$ does not apply.
%    Suppose there is some vertex collision, let $\tau_i$ be the shape of some block-step that shares edges with $E(\tau_t)$, if any.
%	
%	
%	
%	\textbf{Case 1: $\tau_i$ is not well-defined.} In this case, the final gadget does not share any edge with prior gadgets. Any collision vertex is then a slack vertex in our block-value bound.
%	
%	From this point on, assume there is some $\tau_i$ that satisfies the above properties. Notice we cannot have $E(\tau_i)= E(\tau_t)$ as a proper concatenation in the $t-1$ gadgets imply $i=t-1$ and $\tau_t$ would then form a backtracking intersection. Thus, we either have edges in $E(\tau_i)\setminus E(\tau_t)$ or in $E(\tau_t)\setminus E(\tau_i)$. We now discuss by cases.
%
%	
%	\textbf{Case 2:$E(\tau_i)\setminus E(\tau_t) \neq \emptyset $ } In this case, not all edges in $E(\tau_i)$ are closed in $\tau_t$.  Apply \cref{clm: structural-clm} to $\tau_i$ by coloring edges $E(\tau_t)$ red and edges in $E(\tau_i)\setminus E(\tau_t)$ blue. The vertices given by \cref{clm: structural-clm} make middle appearance in $\tau_t$ and give extra slack as their weight is at least $2$.
%	 
%	 
%	\textbf{Case 3: $E(\tau_t)\setminus E(\tau_i) \neq \emptyset$.} Similarly to the above, except we apply \cref{clm: structural-clm} to $\tau_t$ by coloring edges $E(\tau_i)$ red and edges in $E(\tau_t)\setminus E(\tau_i)$ blue. The vertices given by \cref{clm: structural-clm} make middle appearance in $\tau_t$ and give extra  slack as their weight is at least $2$.
%	
%	
%%	 Suppose not all of $E(\tau_i)$ are closed (i.e., appearing for the final time) in $\tau_t$, we apply \cref{clm: structural-clm} with in $\tau_t$ by coloring edges $E(\tau_i)$ red and edges in $E(\tau_t)\setminus E(\tau_i)$ blue. 
%	
%	
%	
%	
%	
%	
% \end{proof}
 
% \paragraph{alternative version}
 
 

%
%\begin{definition}[Gadget incursion] \label{def:gadget-incursion}
%    Let $A,B$ be two gadgets along the dangling path.
%    We call it a gadget-incursion if there are unnecessary vertex intersections between $V(A)$ and $V(B)$ beyond the ``necessary intersection'': when $A,B$ are adjacent, $V_A = U_B$ is the necessary boundary intersection, any intersection in $V(A)\setminus V_A $ and $V(B)\setminus U_B$ is a gadget incursion; similarly, when $A, B$ are not adjacent,  any intersection in $V(A) $ and $V(B)$ is a gadget incursion.
%\end{definition}


%\input{Qnorm.tex}


%\section{Formal Analysis: Inner and Outer Truncation}
%
%\jnote{To get a P.D. matrix, we can consider instead $F_\delta (x) = \delta$ for $x\leq 0$ and bound the coefficient   }


\section{Formal Analysis: Inner and Outer Truncation}\label{sec:formal-truncation}

We now give an overview of our analysis by putting in the details to handle truncation terms. On a high level, we proceed in three steps:
\begin{enumerate}
	\item We first understand the effects of truncation on the function $F$: concretely, we show that the variance of the inner matrix after polynomial truncation and early termination of the iterative process is still $1-o_d(1)$ by our choice of parameters;
	\item Next, we observe that our construction is not necessarily bound %blinded 
    to the particular choice of $F$. In particular, we show that the above bound can be readily adjusted to an analogous positive function $F_\delta$ to be defined;
	\item Finally, we recall that the inner matrix, and therefore its variance, is a function of the threshold parameter $\gamma$. By adjusting the threshold $\gamma$ with $o_d(1)$ slack, we can ensure the inner matrix has variance-$1$ (for our spectrum estimates to apply) and this concludes the proof of our main theorem.
\end{enumerate}


We first bound the effects of truncation by the following.
%
%Our starting point is to note that our choice of non-negative function $F$ is not necessarily fixed - we may instead pick a positive function by softening the prior choice of $F$ to enable a tiny $\delta$ positive mass everywhere in the support. 
%


 
 \paragraph{Truncation of the Inner Matrix $Q$}
 We work with the following parameters for truncations to be chosen. 
\begin{definition}[Parameter $t^*$: Termination of the Iterative Process]
    For $t^*$ to be picked, we will let it denote the termination level for the iterative process. In other words, we will construct a sequence of inner matrices $Q_1, Q_2,\dots  $ up to $Q_{t^*}$. 
\end{definition}

\begin{definition}[Parameter $D$: Inner Truncation for Chebyshev Polynomials] \label{def:inner-trunc}
	We define $D$ to be the truncation parameter for $P(Q_i)$ at each level of $Q_i$. In other words, at each level of $Q_i$, we only consider $P_{j}(Q_i)$ for $j\leq D$. 
\end{definition}

We will follow the recursion until some $t^*$ to be chosen which gives %, and this yields 
us a matrix $Q_{t^*}$. Since we stop the process early, and given the truncation of $P_{j\leq D_j}$ at each level of $Q_i$,  the variance of $Q_{t^*}$ is not exactly $1$ but $1-o_{D,t^*}(1) $ by the following lemma. Crucially, the variance continues to go to $1$ by taking large enough truncation threshold of $t^*$ and $D$.
% \begin{restatable}[Truncated Variance]{lemma}{lemTruncatedVariance} \label{lem:truncated-var}
% There exists a function $\delta(t^*, D)$ such that the variance of terminating at $Q_{t^*}$ with degree-$D$ truncation at each level satisfies
% \[
% \sum_{\substack{\tau \in \calB(Q_{t^*})} } c^2(\tau) = 1 - \delta\,.
% \]
% moreover, it suffices for us to take \[\delta(t^*,D)\leq   O(\frac{1}{D^2} + \frac{1}{(t^*)^2})\]
% %\jnote{currently unproven}
% \end{restatable}



Let \(\varepsilon_{D,t^*}\) measure the variance loss caused by the two
truncations.
\begin{lemma}[Truncated Variance is still almost $1$ (Lemma~D.11 of \cite{PotechinXu2026Theta})
] 
\label{lem:truncated-var-damage}
Let $Q_{t^*,D}$ be the inner matrix with inner truncation parameter $D$ and termination at iteration level $t^*$; we have
\[
 \Var(Q_{t^*,D} ) = \frac{\gam}{1-\gamma}\left( 1 - \varepsilon_{D,t^*} \right).
\]
for \[ 
\varepsilon_{D,t^*} = O( \frac{1}{(t^*)^2} + \frac{1}{D^2})\,\,.
\]
Notice this quantity is $1-o_{D,t^*}(1)$ at the threshold $\gamma=1/2$.
\end{lemma}
 
 
 
\paragraph{Shifting from a Non-Negative Function to a Positive Function}

%\jnote{here's an argument that shows we can get an $\delta$-P.D. matrix by taking the threshold at $(\frac{1}{4}-O(\delta)) \cdot d^2 $. We can then take $\delta=o_d(1)$. }
Next, we note that our choice of non-negative function $F$ is not necessarily fixed --- we may instead pick a positive function by softening the prior choice of $F$ to enable a tiny $\delta$ positive mass everywhere in the support.
Consider the following positive function $\widetilde{F}_\delta$ defined as \[
   \widetilde{F}_\delta(x) \coloneqq  \max(2x, \delta)\,.
\]

Moreover, we consider the rescaled version $F_\delta$ such that $P_1(x)$ has coefficient $1$ in the Chebyshev polynomial expansion for $F_\delta$. Concretely, we take \[
 \widetilde{F}_\delta(x) = \widetilde{C}_F + \widetilde{a}_1 x + \sum_{t>1} \widetilde{a}_i \cdot P_t(x)
 \]
 and obtain the rescaled version \[ 
 F_{\delta} (x) = C_\delta + x + \sum_{t>1} b^{(\delta)}_t \cdot P_t(x)\,,
 \]
 where we take \(
 C_\delta \coloneqq \frac{\widetilde{C}_F}{\widetilde{a}_1}, b^{(\delta)}_t\coloneqq  \frac{\widetilde{a}_t}{\widetilde{a}_1}\,.
  \)
  Let $\calQ$ be the new inner matrix when we apply the iterative process with $F_\delta$ instead of $F$, with the same initialization condition $\calQ_0 \coloneqq C_\delta \cdot \sum_{i\in[m]} w_0[i] \cdot v_iv_i^T$, the appropriately scaled starting point from the identity perturbation, and $\calQ_{\infty} = \lim_i \calQ_i$.  Moreover, let $\calQ_{t^*,D}$ denote the analogous truncation at level $t^*$ with degree $D$ at each level for the $\calQ$ matrix.

 At this point, we remind the reader that our inner matrix $Q$ (as well as $\calQ$) in the end is a function of the threshold parameter $\gamma$. The next claim shows that by allowing a tiny, vanishing gap from the threshold of $\frac{1}{2}$ (which corresponds to the $m\sim d^2/4$  threshold for ellipsoid fitting), we can ensure the inner matrix has variance-$1$ and the resulting matrix has $\delta$-P.D. mass.
 \begin{claim}[Threshold Shift for P.D. Mass: Truncated Version]
\label{claim:pd-shift-threshold}
There is a choice of
\[
    \gamma_{\delta,D,t^*}
    =
    \frac12-\Theta(\delta)
\]
for which the truncated inner matrix has variance $1$, i.e., \( 
\Var(\calQ_{D,t^*}(\gamma) )= 1\,,
\)
provided \(
    \delta \gg \varepsilon_{D,t^*}.
\)

\end{claim} 
 For simplicity, we will now call the final inner matrix $\calQ^*$ with truncation fixed.
 
 \paragraph{Outer Truncation for $F$}
 With the inner matrix   $\calQ^*$, we proceed to analyze \[ 
 \frac{1}{C_\delta} \cdot  F_\delta(\calQ^*)\,.
 \]
To that end, we consider $F_\delta^{\leq D}$, the degree-$D$ truncation of $F_\delta$, and we have  \[
\frac{1}{C_\delta} \cdot  F^{\leq D}_\delta(\calQ^*) = I + Q^*/C_\delta + \sum^D_{j>2} \frac{b^{(\delta)}_j}{C_\delta} \cdot P_j(\calQ^*)\,.
 \] 
 
By standard results from polynomial approximation, we show that as long as the inner matrix has the anticipated spectrum, the inevitable $o_n(1)$ deviation in the spectral edge as well as the outer truncation for $F_\delta$ does not pose substantial damage to our positive mass of $\delta$ that we have started with.


\begin{claim}\label{clm: outer-truncation}
   For an inner matrix with spectrum  \[ 
 \mathsf{spec}(\calQ^*) \in [-2+ \eps_{\mathrm{sp}} , 2+ \eps_{\mathrm{sp}}  ]
 \]
 for some parameter $\eps_{\mathrm{sp}}$, we have
 	\[ 
 	F_{\delta}^{\leq D}(\calQ^*) \succeq \left( \delta - \tilde{O}\!\left(\frac{1}{D} + \eps_{\mathrm{sp}} \cdot \poly(D)\cdot \exp(D\cdot  \eps_{\mathrm{sp}})\right)\cdot  \right) I.
 	\]
\end{claim}
 
%  \paragraph{Wrapping Up}
% We are now ready to wrap up the proof of our main theorem.
% \begin{proof}[Proof of \Cref{thm:ellipsoid-thm}]
% From the construction of our positive map, we have identified an inner matrix $\calQ^*$ s.t. \[ 
% \frac{1}{C_\delta} F_{\delta}(\calQ^*) \succeq \frac{1}{C_\delta} \left( \delta - \tilde{O}\!\left(\frac{1}{D} + \eps_{spec(Q)} \cdot \poly(D)\cdot \exp(D\cdot  \eps_{spec(Q)})\right)\cdot  \right) I_d
% \,.\]
% Moreover, \cref{lem:early-termination-correction} shows that there exists a matrix $D_E$ s.t. \[ 
% \frac{1}{C_\delta} F_{\delta}(\calQ^*) +D_E
% \]
% satisfies the affine constraints at the corresponding threshold $\gam^*= \frac{1}{2}-\Theta(\delta)$ with \[ 
% \|D_E\|_{sp} \leq  O(\frac{1}{C_\delta}\cdot  \sqrt{\varepsilon_{D,t^*}})\,.
% \]
% Finally, picking $\delta \gg  O(\sqrt{\varepsilon_{D,t^*}}) $ gives us the PSD matrix that additionally satisfies each affine constraint. Unpacking the parameters, this corresponds to \[ 
% m = \frac{d^2}{2} \cdot \gamma = (1-\Theta(\delta)) \cdot \frac{d^2}{4}\,.
% \]
% Taking $\delta = o_d(1)$ gives us the desired theorem.
% \end{proof}

% \subsection{Proofs for Shifting to a Positive Function with Truncation}
%We prove the claim for the shift. Firstly, we bound the shift of the coefficients when we switch from $F$ to $F_\delta$, and then incorporate the truncation.
%%Recall first the untruncated variance calculation.
% Recall the variance calculation for ideal $F$ is \[ 
%S_{k+1} =    \frac{\gam }{1-\gam } \cdot ( S_0 +  \sum_{j=2}^{\infty}{{b_j^2} \cdot S_k^j} )\,,
%\] 
%and we use $S_0 = C_F^2$ and $\sum_{j=2}^{\infty}{b_j^2} = 1 - C_F^2$ to deduce $S_{\infty} = 1$ is a solution of this equation at $\gam=\frac{1}{2}$.
%Keeping $\gamma$ as a parameter, the variance for $F_\delta$ now becomes \[ 
%S^{(\delta)}_0 =  C_{\delta}^2\,,
%\] 
%with the recurrence  \[ 
%S^{(\delta)}_{k+1} = \frac{\gam}{1-\gam} \cdot  S^{(\delta)}_0 +\frac{\gamma}{1-\gam }  \sum_{j=2}^{\infty}{({b^{(\delta)}_j)^2} \cdot (S^{(\delta)}_k)^j}
%\]
%For this to converge to $1$ in the limit, it suffices for us to take \[
%\frac{\gam_\delta}{1-\gam_\delta} = \frac{1}{C_\delta^2 + \sum_{j\geq 2} (b^{(\delta)}_j)^2 }
% \]
% which further simplifies to \[ 
% \gam_\delta = \frac{1}{1+C_\delta^2 + \sum_{j\geq 2} (b^{(\delta)}_j)^2  }\,.
% \]
% As a sanity check, plugging the set of ideal coefficients recovers $\gamma=1/2$ from above.
%
%
%
%  For
%\[
%    F_\delta(x)
%    =
%    C_\delta+x+\sum_{j\ge2}b_j^{(\delta)}P_j(x),
%\]
%the variance mass entering the ideal fixed-point equation is now given by 
%\[
%    A_\delta
%    :=
%    C_\delta^2+\sum_{j\ge2}(b_j^{(\delta)})^2.
%\]
%
%
%\begin{claim}[Coefficient Mass Shift]
%\label{claim:A-delta-shift}
%\[
%    A_\delta
%    :=
%    C_\delta^2+\sum_{j\ge2}(b_j^{(\delta)})^2
%    =
%    1+\Theta(\delta).
%\]
%\end{claim}
%
%\begin{proof}
%Let
%\[
%    \Delta_\delta(x):=\widetilde F_\delta(x)-F(x).
%\]
%Since \(F(x)=2x_+\), we have
%\[
%    \Delta_\delta(x)
%    =
%    \begin{cases}
%        \delta, & x<0,\\
%        \delta-2x, & 0\le x<\delta/2,\\
%        0, & x\ge \delta/2.
%    \end{cases}
%\]
%In particular, \(|\Delta_\delta(x)|\le \delta\) for all \(x\in[-2,2]\).
%
%We first estimate the change in the \(P_1\)-coefficient. Since \(P_1(x)=x\),
%\[
%    \widetilde a_1-1
%    =
%    \langle \Delta_\delta,x\rangle_{L^2(\mu_{\mathrm{sc}})}.
%\]
%The contribution from \(x<0\) is
%\[
%    \delta\int_{-2}^0 x\,d\mu_{\mathrm{sc}}(x)
%    =
%    -\Theta(\delta),
%\]
%while the contribution from \(0\le x<\delta/2\) is \(O(\delta^3)\). Hence
%\[
%    \widetilde a_1=1-\Theta(\delta).
%\]
%
%Next, we estimate the \(L^2\)-norm of \(\widetilde F_\delta\). We have
%\[
%    \|\widetilde F_\delta\|_2^2
%    =
%    \|F\|_2^2
%    +2\langle F,\Delta_\delta\rangle
%    +\|\Delta_\delta\|_2^2.
%\]
%The inner product \(\langle F,\Delta_\delta\rangle\) is supported only on
%\([0,\delta/2]\), where \(F(x)=O(\delta)\) and
%\(\Delta_\delta(x)=O(\delta)\). Therefore
%\[
%    \langle F,\Delta_\delta\rangle=O(\delta^3).
%\]
%Also, since \(|\Delta_\delta|\le \delta\),
%\[
%    \|\Delta_\delta\|_2^2=O(\delta^2).
%\]
%Using \(\|F\|_{L^2(\mu_{\mathrm{sc}})}^2=2\), we get
%\[
%    \|\widetilde F_\delta\|_{L^2(\mu_{\mathrm{sc}})}^2
%    =
%    2+O(\delta^2).
%\]
%
%Now apply Parseval to
%\[
%    F_\delta=\frac{\widetilde F_\delta}{\widetilde a_1}.
%\]
%Since the coefficient of \(P_1\) in \(F_\delta\) is normalized to be \(1\),
%\[
%    C_\delta^2+1+\sum_{j\ge2}(b_j^{(\delta)})^2
%    =
%    \|F_\delta\|_{L^2(\mu_{\mathrm{sc}})}^2
%    =
%    \frac{\|\widetilde F_\delta\|_{L^2(\mu_{\mathrm{sc}})}^2}
%         {\widetilde a_1^2}.
%\]
%Using
%\[
%    \|\widetilde F_\delta\|_{L^2(\mu_{\mathrm{sc}})}^2=2+O(\delta^2),
%    \qquad
%    \widetilde a_1=1-\Theta(\delta),
%\]
%we obtain
%\[
%    C_\delta^2+1+\sum_{j\ge2}(b_j^{(\delta)})^2
%    =
%    \frac{2+O(\delta^2)}{(1-\Theta(\delta))^2}
%    =
%    2+\Theta(\delta).
%\]
%Subtracting the normalized \(P_1\)-coefficient contribution, which is \(1\), gives
%\[
%    A_\delta
%    =
%    C_\delta^2+\sum_{j\ge2}(b_j^{(\delta)})^2
%    =
%    1+\Theta(\delta).
%\]
%This proves the claim.
%\end{proof}
%
%
% \begin{claim}[Threshold Shift for P.D. Mass: Truncated Version]
%(Restatement of \cref{claim:pd-shift-threshold})
%There is a choice of
%\[
%    \gamma_{\delta,D,t^*}
%    =
%    \frac12-\Theta(\delta)
%\]
%for which the truncated inner matrix has variance \[ 
%\Var(\calQ_{D,t^*}(\gamma) )= 1\,,
%\]
%provided \[
%    \delta \gg \varepsilon_{D,t^*}.
%\]
%
%\end{claim} 
%
%\begin{proof}
%The preceding calculation shows that
%\[
%    A_\delta=1+\Theta(\delta).
%\]
%
%Now incorporate the truncation.  By the variance estimate for the truncated
%construction with the original function \(F\), at \(\gamma=1/2\) the truncated
%variance is
%\(
%   \frac{\gamma}{1-\gam} (1-\varepsilon_{T,D}).
%\)
%Equivalently, the truncation removes at most \(O(\varepsilon_{T,D})\) from the
%effective variance mass.  The same truncation estimate applies to \(F_\delta\)
%for sufficiently small \(\delta\), since the Chebyshev coefficients of
%\(F_\delta\) differ from those of \(F\) by an \(O(\delta)\) perturbation and the
%same tail/finite-iteration bounds apply.  Therefore the effective truncated
%variance mass is
%\[
%    A_{\delta,T,D}
%    =
%    A_\delta-O(\varepsilon_{T,D})
%    =
%    1+\Theta(\delta)-O(\varepsilon_{T,D}).
%\]
%If \(\delta\gg \varepsilon_{T,D}\), then
%\[
%    A_{\delta,T,D}=1+\Theta(\delta).
%\]
%
%We now choose \(\gamma_{\delta,T,D}\) so that variance \(1\) is a fixed point of
%the truncated variance recurrence.  This requires
%\[
%    \frac{\gamma_{\delta,T,D}}{1-\gamma_{\delta,T,D}}
%    \cdot A_{\delta,T,D}
%    =
%    1.
%\]
%Solving for \(\gamma_{\delta,T,D}\), we obtain
%\[
%    \gamma_{\delta,T,D}
%    =
%    \frac{1}{1+A_{\delta,T,D}}.
%\]
%Using \(A_{\delta,T,D}=1+\Theta(\delta)\), this gives
%\[
%    \gamma_{\delta,T,D}
%    =
%    \frac{1}{2+\Theta(\delta)}
%    =
%    \frac12-\Theta(\delta).
%\]
%This proves the claim.
%\end{proof} 
%
%


%   \begin{claim}[Threshold Shift for P.D. Mass (Restatement of~\cref{claim:pd-shift-threshold})]   Apply the iterative process with $F_\delta$ instead of $F$, we have the corresponding (untruncated) inner matrix   $Q_{\text{ideal},\delta}  $ converging to variance $1$  for some choice of $\gamma_\delta$ such that  \[
% 	\gamma_\delta = \frac{1}{2} - \Theta(\delta)\,.
% 	 \]
% \end{claim}
% 
%\begin{proof}
%	 Recall the variance calculation for ideal $F$ is \[ 
%S_{k+1} =    \frac{\gam }{1-\gam } \cdot ( S_0 +  \sum_{j=2}^{\infty}{{b_j^2} \cdot S_k^j} )\,,
%\] 
%and we use $S_0 = C_F^2$ and $\sum_{j=2}^{\infty}{b_j^2} = 1 - C_F^2$ to deduce $S_{\infty} = 1$ is a solution of this equation at $\gam=\frac{1}{2}$.
%Keeping $\gamma$ as a parameter, the variance for $F_\delta$ now becomes \[ 
%S^{(\delta)}_0 =  C_{\delta}^2\,,
%\] 
%with the recurrence  \[ 
%S^{(\delta)}_{k+1} = \frac{\gam}{1-\gam} \cdot  S^{(\delta)}_0 +\frac{\gamma}{1-\gam }  \sum_{j=2}^{\infty}{({b^{(\delta)}_j)^2} \cdot (S^{(\delta)}_k)^j}
%\]
%For this to converge to $1$ in the limit, it suffices for us to take \[
%\frac{\gam_\delta}{1-\gam_\delta} = \frac{1}{C_\delta^2 + \sum_{j\geq 2} (b^{(\delta)}_j)^2 }
% \]
% which further simplifies to \[ 
% \gam_\delta = \frac{1}{1+C_\delta^2 + \sum_{j\geq 2} (b^{(\delta)}_j)^2  }\,.
% \]
% As a sanity check., plugging the set of ideal coefficients recovers $\gamma=1/2$ from above.
%
%	
%%\end{proof}
%We compare \(\widetilde F_\delta\) to \(F\). Let
%\[
%	\Delta_\delta(x):=\widetilde F_\delta(x)-F(x).
%\]
%Since \(F(x)=2x_+\), we have
%\[
%	\Delta_\delta(x)
%	=
%	\begin{cases}
%	\delta, & x<0,\\
%	\delta-2x, & 0\le x<\delta/2,\\
%	0, & x\ge \delta/2.
%	\end{cases}
%\]
%In particular,
%\[
%	|\Delta_\delta(x)|\le \delta
%	\qquad
%	\text{for all }x\in[-2,2],
%\]
%and hence
%\[
%	\|\Delta_\delta\|_{L^2(\mu_{\mathrm{sc}})}=O(\delta).
%\]
%Thus the Chebyshev coefficient vector of \(\widetilde F_\delta\) differs from that of \(F\)
%by \(O(\delta)\) in \(\ell_2\). In particular,
%\[
%	\widetilde C_\delta=C_F+O(\delta),
%	\qquad
%	\widetilde a_1=1+O(\delta),
%\]
%and
%\[
%	\sum_{j\ge 2}(\widetilde a_j-b_j)^2=O(\delta^2).
%\]
%Since \(\widetilde a_1=1+O(\delta)\) is bounded away from \(0\) for all sufficiently small
%\(\delta\), rescaling by \(\widetilde a_1\) gives
%\[
%	C_\delta=C_F+O(\delta)
%\]
%and
%\[
%	\sum_{j\ge 2}(b^{(\delta)}_j)^2
%	=
%	\sum_{j\ge 2}b_j^2+O(\delta)
%	=
%	1-C_F^2+O(\delta).
%\]
%Therefore
%\[
%	C_\delta^2+\sum_{j\ge 2}(b^{(\delta)}_j)^2
%	=
%	1+O(\delta).
%\]
%
%Next we show that the perturbation has the correct sign. Observe that
%\[
%	\widetilde a_1-1
%	=
%	\langle \Delta_\delta,x\rangle_{L^2(\mu_{\mathrm{sc}})}.
%\]
%On the interval \(x<0\), we have \(\Delta_\delta(x)=\delta\), so the contribution is
%\[
%	\delta\int_{-2}^0 x\,d\mu_{\mathrm{sc}}(x)
%	=
%	-\Omega(\delta).
%\]
%On the interval \(0\le x<\delta/2\), we have \(0\le \Delta_\delta(x)\le \delta\), and hence
%the contribution is at most
%\[
%	\int_0^{\delta/2}\delta x\,d\mu_{\mathrm{sc}}(x)
%	=
%	O(\delta^3).
%\]
%It follows that
%\[
%	\widetilde a_1=1-\Omega(\delta).
%\]
%Together with \(\widetilde a_1=1+O(\delta)\), this gives
%\[
%	\widetilde a_1=1-\Theta(\delta).
%\]
%
%Also, since \(\|\Delta_\delta\|_{L^2(\mu_{\mathrm{sc}})}=O(\delta)\) and
%\(\Delta_\delta\) only changes \(F\) by a floor of size \(\delta\),
%\[
%	\|\widetilde F_\delta\|_{L^2(\mu_{\mathrm{sc}})}^2
%	=
%	\|F\|_{L^2(\mu_{\mathrm{sc}})}^2+O(\delta^2)
%	=
%	2+O(\delta^2).
%\]
%Applying Parseval to 
%\(
%	F_\delta=\frac{\widetilde F_\delta}{\widetilde a_1},
%\)
%we have
%\[
%	C_\delta^2+1+\sum_{j\ge 2}(b^{(\delta)}_j)^2
%	=
%	\frac{\|\widetilde F_\delta\|_{L^2(\mu_{\mathrm{sc}})}^2}{\widetilde a_1^2}.
%\]
%Using
%\[
%	\|\widetilde F_\delta\|_{L^2(\mu_{\mathrm{sc}})}^2=2+O(\delta^2)
%	\qquad\text{and}\qquad
%	\widetilde a_1=1-\Theta(\delta),
%\]
%we obtain
%\[
%	C_\delta^2+1+\sum_{j\ge 2}(b^{(\delta)}_j)^2
%	=
%	2+\Theta(\delta).
%\]
%%Equivalently,
%%\[
%%	C_\delta^2+\sum_{j\ge 2}(b^{(\delta)}_j)^2
%%	=
%%	1+\Theta(\delta).
%%\]
%%
%%By definition,
%%\[
%%	\frac{\gamma_\delta}{1-\gamma_\delta}
%%	=
%%	\frac{1}{C_\delta^2+\sum_{j\ge 2}(b^{(\delta)}_j)^2},
%Plugging in the above gives us
%\[
%	\gamma_\delta
%	=
%	\frac{1}{2+\Theta(\delta)}
%	=
%	\frac12-\Theta(\delta).
%\]
%\end{proof}

%\subsection{Polishing attempt}

\subsection{Shifting to a Positive Function with Truncation}

We now justify the threshold shift after replacing the nonnegative activation
\(F(x)=2x_+\) by a strictly positive perturbation in a 3-step argument.
\begin{enumerate}
	\item  First, we recall how the variance recurrence identifies the relevant
coefficient mass;
\item  Then, we show that this coefficient mass increases by
\(\Theta(\delta)\) after shifting to \(F_\delta\);
\item Finally, we incorporate the
finite iteration and inner-degree truncation errors and show that \(\gamma_{\delta,D,t^*}\) can be chosen so that the truncated inner matrix has variance \(1\);
\end{enumerate} 
%First, we recall how the variance recurrence identifies the relevant
%coefficient mass.  Second, we show that this coefficient mass increases by
%\(\Theta(\delta)\) after shifting to \(F_\delta\).  Third, we incorporate the
%finite iteration and inner-degree truncation errors and choose
%\(\gamma_{\delta,D,t^*}\) so that the truncated inner matrix has variance \(1\).

\paragraph{Step-1: The variance mass.}

Consider a function
\(
    G(x)=C_G+x+\sum_{j\ge2} a_j P_j(x),
\)
whose \(P_1(x)=x\) coefficient is normalized to be \(1\).  In the ideal,
untruncated iteration, if \(S_k\) denotes the variance of the inner matrix at
iteration \(k\), then the variance recursion has the form
\[
    S_{k+1}
    =
    \frac{\gamma}{1-\gamma}
    \left(
        C_G^2+\sum_{j\ge2} a_j^2 S_k^j
    \right).
\]
Therefore, for \(S=1\) to be a fixed point of the ideal recursion, it is
sufficient to take
\[
    \frac{\gamma}{1-\gamma}
    \left(
        C_G^2+\sum_{j\ge2}a_j^2
    \right)
    =
    1.
\]
This motivates the definition
\(
    A_G
    :=
    C_G^2+\sum_{j\ge2}a_j^2.
\)
%For such an activation \(G\), the threshold parameter that makes variance \(1\)
%a fixed point is therefore
%\(
%    \gamma_G
%    =
%    \frac{1}{1+A_G}.
%\)
As a sanity check, for the original function $F$,  Parseval gives  \(A_F=1\), and hence
\(
    \gamma_F=\frac{1}{1+A_F}=\frac12\,,
\)
 recovering the threshold in the ideal calculation.

\paragraph{Step-2: Coefficient mass shift under \(F_\delta\).}

We now replace \(F\) by a strictly positive perturbation
\[
    \widetilde F_\delta(x):=\max(2x,\delta).
\]
Write its Chebyshev expansion as
\(
    \widetilde F_\delta(x)
    =
    \widetilde C_F+\widetilde a_1x+\sum_{j\ge2}\widetilde a_jP_j(x).
\)
We then rescale by the \(P_1\)-coefficient and define
\(
    F_\delta(x)
    :=
    \frac{\widetilde F_\delta(x)}{\widetilde a_1}
    =
    C_\delta+x+\sum_{j\ge2}b_j^{(\delta)}P_j(x),
\)
where
\(
    C_\delta:=\frac{\widetilde C_F}{\widetilde a_1},
       b_j^{(\delta)}:=\frac{\widetilde a_j}{\widetilde a_1}.
\)
The relevant variance mass for \(F_\delta\) is
\[
    A_\delta
    :=
    C_\delta^2+\sum_{j\ge2}(b_j^{(\delta)})^2.
\]
We bound the change of $A_\delta$ via the following claim.
\begin{claim}[Coefficient Mass Shift]
\label{claim:A-delta-shift}
For sufficiently small \(\delta>0\),
\[
    A_\delta=1+\Theta(\delta).
\]
\end{claim}

%\begin{proof}
%Let
%\[
%    \Delta_\delta(x):=\widetilde F_\delta(x)-F(x).
%\]
%Since \(F(x)=2x_+\), we have
%\[
%    \Delta_\delta(x)
%    =
%    \begin{cases}
%        \delta, & x<0,\\
%        \delta-2x, & 0\le x<\delta/2,\\
%        0, & x\ge \delta/2.
%    \end{cases}
%\]
%In particular,
%\[
%    |\Delta_\delta(x)|\le \delta
%    \qquad
%    \text{for all }x\in[-2,2].
%\]
%
%We first estimate the shift in the \(P_1\)-coefficient.  Since \(P_1(x)=x\),
%\[
%    \widetilde a_1-1
%    =
%    \langle \Delta_\delta,x\rangle_{L^2(\mu_{\mathrm{sc}})}.
%\]
%On the interval \(x<0\), the contribution is
%\[
%    \delta\int_{-2}^0 x\,d\mu_{\mathrm{sc}}(x)
%    =
%    -\Theta(\delta).
%\]
%On the interval \(0\le x<\delta/2\), we have
%\(|\Delta_\delta(x)|\le\delta\) and \(x\le\delta/2\), so the contribution is
%\(O(\delta^3)\).  Therefore
%\[
%    \widetilde a_1=1-\Theta(\delta).
%\]
%
%Next we estimate the \(L^2\)-norm of \(\widetilde F_\delta\).  We have
%\[
%    \|\widetilde F_\delta\|_2^2
%    =
%    \|F\|_2^2
%    +2\langle F,\Delta_\delta\rangle
%    +\|\Delta_\delta\|_2^2.
%\]
%The product \(F\Delta_\delta\) is supported only on
%\([0,\delta/2]\), where both \(F(x)=O(\delta)\) and
%\(\Delta_\delta(x)=O(\delta)\).  Hence
%\[
%    \langle F,\Delta_\delta\rangle=O(\delta^3).
%\]
%Also, since \(|\Delta_\delta|\le \delta\),
%\[
%    \|\Delta_\delta\|_2^2=O(\delta^2).
%\]
%Using \(\|F\|_{L^2(\mu_{\mathrm{sc}})}^2=2\), we get
%\[
%    \|\widetilde F_\delta\|_{L^2(\mu_{\mathrm{sc}})}^2
%    =
%    2+O(\delta^2).
%\]
%
%Now apply Parseval to
%\(
%    F_\delta=\frac{\widetilde F_\delta}{\widetilde a_1}.
%\)
%Since the \(P_1\)-coefficient of \(F_\delta\) is normalized to be \(1\),
%\[
%    C_\delta^2+1+\sum_{j\ge2}(b_j^{(\delta)})^2
%    =
%    \|F_\delta\|_{L^2(\mu_{\mathrm{sc}})}^2
%    =
%    \frac{\|\widetilde F_\delta\|_{L^2(\mu_{\mathrm{sc}})}^2}
%         {\widetilde a_1^2}.
%\]
%Using
%\[
%    \|\widetilde F_\delta\|_{L^2(\mu_{\mathrm{sc}})}^2=2+O(\delta^2),
%    \qquad
%    \widetilde a_1=1-\Theta(\delta),
%\]
%we obtain
%\[
%    C_\delta^2+1+\sum_{j\ge2}(b_j^{(\delta)})^2
%    =
%    \frac{2+O(\delta^2)}{(1-\Theta(\delta))^2}
%    =
%    2+\Theta(\delta).
%\]
%Rearranging gives us \[
%    A_\delta
%    =
%    C_\delta^2+\sum_{j\ge2}(b_j^{(\delta)})^2
%    =
%    1+\Theta(\delta).
%\]
%%This proves the claim.
%\end{proof}

\paragraph{Step-3: Incorporating truncation.}

We now return to the truncated construction. Recall that \(
    \varepsilon_{D,t^*}
    :=O(
    \frac{1}{D^2}+\frac{1}{(t^*)^2})
\)
denotes the variance damage coming from the inner Chebyshev truncation degree
\(D\) and the finite iteration level \(t^*\). With these ingredients in hand, we are ready to wrap up the proof of the following claim. This shows that with the truncated construction for the inner matrix, we can pick a $\gamma$ (which ultimately governs the threshold of $m$ we obtain) to ensure the inner matrix still has variance-$1$, albeit at a $\delta$-cost away from the threshold. 


\begin{claim}[Threshold Shift for P.D. Mass: Truncated Version (Restatement of \cref{claim:pd-shift-threshold})]
%\label{claim:pd-shift-threshold}
There is a choice
\[
  \gamma_{\delta,D,t^*}
    =
    \frac12-\Theta(\delta)
\]
for which the truncated inner matrix produced by the \(F_\delta\)-based
construction has variance
\[
    \Var\!\left(\calQ_{D,t^*}(\gamma_{\delta,D,t^*})\right)=1
\]
provided 
\[
    \delta\gg \varepsilon_{D,t^*}.
\]
For convenience, we write $\gam_\delta \coloneqq   \gamma_{\delta,D,t^*}$ when the dependence is clear.
\end{claim}
As a reminder, we will eventually take $\delta=o_d(1)$, so this cost is essentially harmless for our result.
%The complete proofs to the claims in this section can be found in~\cref{sec:defer-details-truncation}.

\begin{proof}
By \cref{claim:A-delta-shift}, the untruncated variance mass associated with
\(F_\delta\) is
\[
    A_\delta=1+\Theta(\delta).
\]

Let \(A_{\delta,D,t^*}\) denote the effective variance mass retained after
truncating the \(F_\delta\)-based construction at inner degree \(D\) and
iteration depth \(t^*\).  \Cref{lem:truncated-var-damage} gives
\[
    A_{\delta,D,t^*}
    =
    A_\delta - O(\varepsilon_{D,t^*}).
\]
i.e., \[
    A_{\delta,D,t^*}
    =
    1+\Theta(\delta) - O(\varepsilon_{D,t^*}) = 1+\Theta(\delta)
\]
for \(\delta\gg \varepsilon_{D,t^*}\).
For any \(D,t^*\), the variance of the truncated inner matrix has the form
\[
    \Var\!\left(\calQ_{D,t^*}(\gamma)\right)
    =
    \frac{\gamma}{1-\gamma}\,A_{\delta,D,t^*}.
\]
Thus, to make the variance equal to \(1\), it suffices to choose
\(\gamma=\gamma_{\delta,D,t^*}\) such that
\[
    \gamma_{\delta,D,t^*}
    =
    \frac{1}{1+A_{\delta,D,t^*}}.
\]
Using \(A_{\delta,D,t^*}=1+\Theta(\delta)\), we conclude that
\[
    \gamma_{\delta,D,t^*}
    =
    \frac{1}{2+\Theta(\delta)}
    =
    \frac12-\Theta(\delta).
\]
%This proves the claim.
\end{proof}


%\begin{proof}
%By \cref{claim:A-delta-shift}, the untruncated variance mass associated with
%\(F_\delta\) is
%\[
%    A_\delta=1+\Theta(\delta).
%\]
%
%Let \(A_{\delta,D,t^*}\) denote the effective variance mass retained after
%truncating the \(F_\delta\)-based construction at inner degree \(D\) and
%iteration depth \(t^*\).  The truncated-variance estimate gives
%\[
%    A_{\delta,D,t^*}
%    =
%    A_\delta - O(\varepsilon_{D,t^*}).
%\]
%Equivalently,
%\[
%    A_{\delta,D,t^*}
%    =
%    1+\Theta(\delta) - O(\varepsilon_{D,t^*}) = 1+\Theta(\delta)
%\]
%for \(\delta\gg \varepsilon_{D,t^*}\).
%For fixed \(D,t^*\), the variance of the truncated inner matrix has the form
%\[
%    \Var\!\left(\calQ_{D,t^*}(\gamma)\right)
%    =
%    \frac{\gamma}{1-\gamma}\,A_{\delta,D,t^*}.
%\]
%Thus, to make the variance equal to \(1\), it suffices to choose
%\(\gamma=\gamma_{\delta,D,t^*}\) such that
%\[
%    \gamma_{\delta,D,t^*}
%    =
%    \frac{1}{1+A_{\delta,D,t^*}}.
%\]
%Using \(A_{\delta,D,t^*}=1+\Theta(\delta)\), we conclude that
%\[
%    \gamma_{\delta,D,t^*}
%    =
%    \frac{1}{2+\Theta(\delta)}
%    =
%    \frac12-\Theta(\delta).
%\]
%This proves the claim.
%\end{proof}

\subsection{Bounds on the Truncation Error for \texorpdfstring{$A^{-1}$}{1/A}}
To avoid infinite series in our construction of the ellipsoid matrix, we would also need to truncate $A$. That said, explicit control on the spectrum of $A$ in turn allows us to obtain a reasonable truncation. This is the focus of this section.

There are two sources of error as we transfer from an ideal-world proof to a real-world proof:\begin{enumerate}
	\item The truncation error incurred by truncating $A^{-1}$ at degree-$D$. 	\item The edge-deviation error: our chosen polynomial expansion is only valid for the interval of $[(1-\sqrt{\gam})^2, (1+\sqrt{\gam})^2] $ while our instantiated matrix $A$ may further have inevitable $\varepsilon_A= o_d(1)$ fluctuation at the edge. \end{enumerate}
It can be shown that the second type of error term dominates in our regime of interest; we show the following.
%\begin{restatable}[Truncation error for $x^{-1}$ within the MP support]{lemma}{mpinversetruncation} \label{lem:mp-inverse-truncation} Let \(D\) be the truncation parameter for \(A^{-1}\), and define the truncated polynomial as \[ (x^{-1})_{\le D} \coloneqq \frac{1}{1-\gamma} \sum_{t=0}^{D}(-1)^t \q_t(x). \] Then, uniformly for \[ x\in\bigl[(1-\sqrt{\gamma})^2,\,(1+\sqrt{\gamma})^2\bigr], \] we have \[ \left| \frac{1}{x}-(x^{-1})_{\le D} \right| \le O_\gamma\!\left((D+1)(\sqrt{\gamma})^{D+1}\right). \] \end{restatable}


\begin{restatable}[Truncation error under spectral-edge fluctuation]{lemma}{mpinverseedgefluctuation} \label{lem:mp-inverse-edge-thickening} Suppose \[ \operatorname{spec}(A) \subseteq \bigl[(1-\sqrt{\gam})^2-\varepsilon_A,\, (1+\sqrt{\gam})^2+\varepsilon_A\bigr] \] for some \(\varepsilon_A>0\), and assume \[ (1-\sqrt{\gam})^2-\varepsilon_A>0. \]  We have\[ \left\| A^{-1}-(A^{-1})_{\le D} \right\|_{sp} \le O_\gam\!\left(\poly(D) (\sqrt{\gam})^{D+1}\right) + O_\gam\!\left( \varepsilon_A \cdot  \poly(D) \cdot  e^{O_\gam(D\sqrt{\varepsilon_A})} \right). \] \end{restatable}
We prove this by standard techniques from scalar polynomial approximation in~\cref{sec:defer-details-truncation}.



\subsection{Bounds on Truncation Error for Inner and Outer Truncation}
Let $\calQ^*= \calQ_{D,t^*}(\gamma_{\delta,D,t^*})$ be the final inner matrix at the chosen threshold such that the inner matrix has variance-$1$ after inner truncation.
%\begin{lemma}
%	There exists a matrix $D_E$ such that \[ 
%	\frac{1}{C_F} \cdot  F_{\delta}^{\leq D}(Q^*)  + D_E
%	\]
%	satisfies the affine constraints exactly, i.e., $v_i^\top (\frac{1}{C_F} \cdot  F_{\delta}^{\leq D}(Q^*)  + D_E) v_i=1$ for all $i\in[m]$. Moreover,
%	 we have \[
%	 \| D_E \| \ll O(\delta)\,.	  \]
%\end{lemma}
\begin{restatable}[Early-termination correction]{lemma}{lemEarlyTerminationCorrection}
\label{lem:early-termination-correction}
There exists a matrix \(D_E\) such that
\(
    \frac{1}{C_\delta}F_\delta^{\le D}(\calQ^*)+D_E
\)
satisfies the affine constraints exactly; namely,
\[
    v_i^\top
    \left(
        \frac{1}{C_\delta}F_\delta^{\le D}(\calQ^*)+D_E
    \right)
    v_i
    =
    1
    \qquad \text{for all } i\in[m].
\]
Moreover,
\[
    \|D_E\|_{sp}
    \le
    O\!\left(\sqrt{\varepsilon_{D,t^*}}\right).
\]
%In particular, if \(D\) and \(t^*\) are chosen so that
%\(
%    \varepsilon_{D,t^*}=o(\delta^2),
%\)
%then
%\[
%    \|D_E\|_{sp}=o(\delta)\,.
%\]
\end{restatable}

\begin{proof}
	
Let's start by identifying where error terms can arise. On a high level, other than terms in the initial $Q_0$, each other shape in the final polynomial expansion is intended to pair with one other shape so that they group into a term in the kernel of the affine constraints, in other words, \[
 v_i^\top (c_\tau \cdot \cm_\tau + c_{\corr(\tau)} \cdot \cm_{\corr(\tau)}) v_i = 0
\,.\]
That said, notice that as we pick the polynomial expansion threshold to be $D$ for the inner truncation as we design the update for each iteration, as well as the final outer truncation for $F_{\leq D}(\calQ^*)$, there is no missing term due to the inner truncation.  Therefore, the only error term is incurred by the early termination as we terminate the iterative process at some finite level.

Next, we observe that, at the chosen $\gam^*$ s.t. the truncated inner matrix has variance-$1$, the untruncated inner matrix has corresponding variance \(
\frac{\gam^*}{1-\gam^*} \cdot A_\delta \,.
 \)
Recall that $\gam^*$ is chosen as \(
\gam^* = \frac{1}{1+ A_{\delta,D,t^*}} 
 \) with \( A_{\delta,D,t^*} = A_\delta - O(\varepsilon_{D,t^*}) \); we have

\[
    \frac{\gam^*}{1-\gam^*}\cdot A_\delta
    =
    \frac{A_\delta}{A_{\delta,D,t^*}}
    =
    \frac{A_\delta}{A_\delta-\Delta_{\delta,D,t^*}}
    =
    1+
    \frac{\Delta_{\delta,D,t^*}}
         {A_\delta-\Delta_{\delta,D,t^*}}
    =
    1+O(\varepsilon_{D,t^*})\,.
\]
Thus, this gives us a bound on the variance of the terms removed by early truncation, \[ 
\text{Early-Termination-Error} \leq O(\varepsilon_{D,t^*})
\]

Therefore, combining with our matrix norm bounds which apply as each missing correction term is a backbone-dangling term, and satisfies the size constraint $D_V$, \begin{align*}
	\|\sum_{\tau: \text{missing correction term}} c_{\corr(\tau)} \cdot \cm_{\tau} \|_{sp}  &\leq (2+o_d(1)) \cdot \Var( \sum_{\tau: \text{missing correction term}} c_{\corr(\tau)} \cdot \cm_{\tau}) \\
	&\leq (2+o_d(1))\cdot   \sqrt{\text{Early-Termination-Error}} \\
	&=  O(\sqrt{\varepsilon_{D,t^*}}) \,.
\end{align*} 
\end{proof}

% the shift from $F$ to $F_\delta$, and subsequently our choice of $\gamma$ such that the inner matrix $\calQ^*$ has variance-$1$. We make the following observation in the prior calculation: putting truncation aside, we have \[
%    A_\delta
%    =
%    C_\delta^2+\sum_{j\ge2}(b_j^{(\delta)})^2
%    =
%    1+\Theta(\delta)\,,
%\]
%and therefore the untruncaed matrix has variance \( 
%\frac{\gam}{1-\gam} \cdot A_\delta\,.
%\)
%However, as we apply inner truncation on the polynomial as well early termination of the iterative process, the truncated variance becomes \[ 
% A_{\delta,D,t^*}
%    =
%    A_\delta - O(\varepsilon_{D,t^*}).
%\]
%and we pick the final $\gam^*$ s.t. $\calQ^*$ the inner matrix has variance-$1$. Thus, the truncated mass by truncation 


%\subsection{Proof of Main Theorem}
%We are now ready to wrap up the proof to our main theorem.
%\begin{proof}[Proof to \Cref{thm:ellipsoid-thm}]
%From the construction of our positive map, we have identified an inner matrix $\calQ^*$ s.t. \[ 
%\frac{1}{C_\delta} F_{\delta}(\calQ^*) \succeq \frac{1}{C_\delta} \cdot \delta \cdot I_d
%\,.\]
%Moreover, \cref{lem:early-termination-correction} shows that there exists a matrix $D_E$ s.t. \[ 
%\frac{1}{C_\delta} F_{\delta}(\calQ^*) +D_E
%\]
%that satisfies the affine constraints at the corresponding threshold $\gam^*= \frac{1}{2}-\Theta(\delta)$ with \[ 
%\|D_E\|_{sp} \leq  O(\frac{1}{C_\delta}\cdot  \sqrt{\varepsilon_{D,t^*}})\,.
%\]
%Finally, picking $\delta \gg  O(\sqrt{\varepsilon_{D,t^*}}) $ gives us the PSD matrix that additionally satisfies each affine constraint. Unpacking the parameters, this corresponds to \[ 
%m = \frac{d^2}{2} \cdot \gamma = (1-\Theta(\delta)) \cdot \frac{d^2}{4}\,.
%\]
%Taking $\delta = o_d(1)$ gives us the desired theorem.
%\end{proof}

\subsection{Choice of Parameters and Proof of Main Theorems}
\label{sec:parameter-instantiation}

For convenience, we use the same degree \(D\) for the inner and outer Chebyshev
truncations and for the MP-polynomial expansion of \(A^{-1}\). Fix a
sufficiently small absolute constant \(a>0\), and set
\[
    D=t^\ast=\left\lceil(\log \log d)^a\right\rceil,
    \qquad
    q=\left\lceil(\log d)^4\right\rceil,
    \qquad
    \delta=D^{-1/2}.
\]

\paragraph{Verification of the norm-bound regime.}
%Let \(V_i\) be the maximum number of vertices in a shape appearing in
%\(Q_i\). A truncated MP path has \(O(D)\) vertices, while a degree-\(D\)
%Chebyshev term concatenates at most \(D\) shapes. Hence
%\[
%    V_0=O(D),
%    \qquad
%    V_{i+1}\leq DV_i+O(D),
%\]
%so every shape arising in the finite construction, including the final
%polynomial and terminal correction, has at most
%\[
%    D_V\leq D^{t^\ast+O(1)}=d^{o(1)}
%\]
%vertices. 
To get started, we verify the following size bound on the shapes that arise in the truncated process to ensure our main theorem for matrix norm bounds applies.
\begin{claim}[Size of shapes in the truncated construction]
\label{claim:shape-size-truncated-construction}
Suppose all Chebyshev and MP expansions are truncated at degree \(D\).
Then every shape appearing in \(Q_i\) has at most
\[
    D^{i+O(1)}
\]
vertices. Consequently, every shape in the final degree-\(D\) polynomial
or terminal correction has at most
\[
    D_V\leq D^{t^\ast+O(1)}.
\]
In particular, if \(D=t^\ast=(\log d)^a\) for any fixed \(a<1\), then
\(D_V=d^{o(1)}\).
\end{claim}

\begin{proof}
Let \(V_i\) be the maximum shape size in \(Q_i\). The initialization
contains only constantly many MP paths of length at most \(D\), so
\(V_0=O(D)\). Each update concatenates at most \(D\) previous shapes and
adds \(O(D)\) vertices through the correction, giving
\[
    V_{i+1}\leq DV_i+O(D).
\]
Thus \(V_i\leq D^{i+O(1)}\); the final polynomial and terminal correction
change the exponent only by \(O(1)\). Finally,
\[
    \log D_V=O(t^\ast\log D)=o(\log d),
\]
which implies \(D_V=d^{o(1)}\).
\end{proof}

\begin{claim}[Shape Count Bound] \label{clm:shape-count}
	There are at most $c^{D_V}$ backbone-dangling shapes of $D_V$ vertices for some constant $c>1$.
\end{claim}

\begin{proof}
Encode a shape by its depth-first construction. Each new vertex is
introduced in one of \(O(1)\) ways: through an
\(\alpha\)-, \(\beta\)-, or \(M_D\)-gadget, a deviation attachment,
or a horizontal-concatenation operation. The active attachment site
is implicit in the traversal, and orientations and opening/closing
symbols enlarge the alphabet only by a constant factor. Thus a
\(T\)-vertex shape has an injective \(O(T)\)-length encoding over a
fixed finite alphabet, giving at most \(C^T\) possibilities.
\end{proof}



Since each shape has coefficient at most $1$, we can therefore deduce a bound on $|c(Q^*)|$.
\begin{corollary}
	Let $|c(Q^*)| = \sum_{\tau \in \calB(Q^*)} c_\tau $, we have \[ 
	|c(Q^*)|  \leq O(C^{D^D}).
	\]
\end{corollary}

\begin{remark}
For every fixed \(\delta>0\), \(C^{D^D}\leq d^\delta\) provided
\[
    D\leq c_{C,\delta}
    \frac{\log\log d}{\log\log\log d}.
\]
\end{remark}


Since \(q=(\log d)^4\), we also have \(qD_V=d^{o(1)}\).
Thus the polynomial losses in the MP-LCP, SC-LCP, and
vertical-attachment bounds are absorbed by their \(d^{-\Omega(1)}\)
gains. Together with \cref{thm:Q-norm} and
\cref{cor:trace-to-norm-quant}, this gives
\[
    \eps_{\mathrm{sp}}
    =
    O\!\left(\frac1{\log^2 d}\right)+d^{-1/2+o(1)},
    \qquad
    \mathsf{spec}(Q^\ast)
    \subseteq
    [-2-\eps_{\mathrm{sp}},\,2+\eps_{\mathrm{sp}}].
\]
In particular, \(D\sqrt{\eps_{\mathrm{sp}}}=o(1)\).
 By \cref{lem:truncated-var-damage},
\(
    \varepsilon_{D,t^\ast}
    =
    O\!\left(\frac1{D^2}+\frac1{(t^\ast)^2}\right)
    =
    O(D^{-2})
    =
    O(\delta^4).
\)
Let \(C>0\) be a fixed constant dominating the polynomial dependence on
\(D\) in \cref{lem:mp-inverse-edge-thickening} and
\cref{claim:pd-shift-threshold}. Since \(D=(\log \log d)^a\) for some $a>0$  and
\(\eps_{\mathrm{sp}}=O(\log^{-2}d)+d^{-1/2+o(1)}\), we have
\[
    D\sqrt{\eps_{\mathrm{sp}}}=o(1),
    \qquad
    \eps_{\mathrm{sp}}D^C=o(D^{-2}).
\]
Consequently, \cref{lem:mp-inverse-edge-thickening} implies
that the degree-\(D\) truncation error for \(A^{-1}\) is \(o(D^{-2})\),
while \cref{clm: outer-truncation} gives an outer-truncation error of
\(o(\delta)\).

Finally, the early-termination correction from
\Cref{lem:early-termination-correction} satisfies
\[
    \|D_E\|_{\mathrm{sp}}
    =
    O\!\left(\sqrt{\varepsilon_{D,t^\ast}}\right)
    =
    O(D^{-1})
    =
    o(\delta).
\]
Thus the positive spectral floor \(\delta\) dominates every truncation
and termination error.

\paragraph{Wrapping up the proof of \cref{thm:ellipsoid-thm}.}
By~\cref{claim:A-delta-shift} and~\cref{claim:pd-shift-threshold},
we may choose
\(
    \gamma_{\delta,D,t^\ast}
    =
    \frac12-\Theta(\delta)
\)
so that the truncated inner matrix \(Q^\ast\) has variance \(1\).
The preceding spectral and outer-truncation bounds give
\[
    \frac1{C_\delta}F_\delta^{\leq D}(Q^\ast)
    \succeq
    \left(\frac{\delta}{C_\delta}-o(\delta)\right)I.
\]
Moreover, \Cref{lem:early-termination-correction} provides \(D_E\) such
that
\(
    \Lambda
    \coloneqq
    \frac1{C_\delta}F_\delta^{\leq D}(Q^\ast)+D_E
\)
satisfies the affine constraints exactly. Since
\(\|D_E\|_{\mathrm{sp}}=o(\delta)\), we have
\(
    \Lambda\succeq \Omega(\delta) \cdot I.
\)
Thus, this proves feasibility for
\[
    m
    = \left(1-\Theta(\delta)\right) \cdot   
    \frac{d^2}{4}.
\]
%
%\paragraph{Wrapping up the proof of \cref{thm:ellipsoid-refutation-thm}.}
%The dual correction differs from the primal correction only by the
%bounded factor
%\[
%    \operatorname{correct}_{\mathrm{dual}}(H)
%    =
%    -\frac{1-\gamma}{\gamma}
%    \operatorname{correct}_{\mathrm{primal}}(H).
%\]
%Hence the preceding shape-size, norm, and truncation estimates apply
%unchanged to the dual construction.
%
%Take
%\[
%    \gamma_+=\frac12+\Theta(\delta),
%    \qquad
%    \rho\coloneqq\frac{1-\gamma_+}{\gamma_+},
%\]
%and choose \(c_R\) according to the dual variance calculation so that
%the truncated inner matrix \(W^\ast\) has variance \(1\). By
%\cref{sec:refutation}, 
%\(
%    c_R-\rho=\Theta(\delta).
%\)
%Combining the above, we can conclude that
%\(\Lambda_R\) satisfies
%\(
%    \Lambda_R\succeq0\) and \(    \Lambda_R\in \calS=  \operatorname{Im}(\mathcal L^\ast).
%\)
%
%Finally, the variance and correlation estimates from section~3 give
%\[
%\begin{aligned}
%    \langle\Lambda_R,I-\calR\rangle
%    &=
%    d-c_R\langle \calR,\calR\rangle+o(\delta d)\\
%    &=
%    d\left(1-\frac{c_R}{\rho} + o(\delta) \right)\\
%    &=
%    -\Theta(\delta d)<0.
%\end{aligned}
%\]
%Thus \(\Lambda_R\) is a valid dual refutation witness, proving
%infeasibility for
%\(
%    m
%    = \left( 1+ \Theta(\delta) \right) \cdot d^2/4.
%\)
%Finally, since
%\[
%    \delta=D^{-1/2}=(\log d)^{-a/2+o(1)}=o_d(1),
%\]
%both of our theorems hold with vanishing slack $o_d(1)$ as desired.
% 
% \paragraph{alt version}
 
 \paragraph{Wrapping up the proof of
\cref{thm:ellipsoid-refutation-thm}.}
The dual correction differs from the primal correction only by the
bounded factor
\[
    \operatorname{correct}_{\mathrm{dual}}(H)
    =
    -\frac{1-\gamma}{\gamma}
    \operatorname{correct}_{\mathrm{primal}}(H).
\]
Hence the preceding shape-size, norm, and truncation estimates apply
unchanged to the dual construction.

Take
\[
    \gamma_+
    =
    \frac12+\Theta(\delta),
    \qquad
    \rho
    \coloneqq
    \frac{1-\gamma_+}{\gamma_+},
\]
and choose \(c_R\) according to the dual variance calculation so that
the truncated inner matrix \(W^\ast\) has variance \(1\). By
\cref{sec:refutation},
\[
    c_R-\rho=\Theta(\delta).
\]
Combining the above, we conclude that $\Lambda_R\succeq 0 $ and $\Lambda_R \in \calS=\operatorname{Im}(\calL^*)$.


Next, we verify the objective value. Finally, write
\[
    \Lambda_R
    =
    I-\Pi^\perp I+c_R\calR+Z.
\]
Since \(\calR\in\calS\),
\[
    \langle\Pi^\perp I,\calR\rangle=0,
    \qquad
    \langle\Pi^\perp I,I\rangle
    =
    \|\Pi^\perp I\|_F^2.
\]
Therefore,
\begin{align*}
    \langle\Lambda_R,I-\calR\rangle
    &=
    d-\|\Pi^\perp I\|_F^2
    +(c_R-1)\Tr(\calR)
    -c_R\langle\calR,\calR\rangle \\
    &\qquad
    +\Tr(Z)-\langle Z,\calR\rangle.
\end{align*}
By the quantitative variance estimates from
\cref{sec:refutation} and the sub-sequential error bounds from \cref{lem:hyper-parameter-bounds},
\[
    \langle\calR,\calR\rangle
    =
    \frac d\rho+\widetilde O(\sqrt d).
\]

\begin{claim}[Auxiliary bounds] \label{clm:refutation-aux-bounds}
	We have the following bounds, \begin{enumerate}
	\item $ \|\Pi^\perp I\|_F^2 = \tilde{O}(\sqrt{d})$;
	\item  $ |\Tr(\calR)|\leq d \cdot B_q(\text{non-ideal}) = o(\delta d)  $, and similarly $ |\Tr(Z)| \leq d \cdot  B_q(\text{non-ideal})$ for the auxiliary function from~\cref{thm:Q-norm};
	\item $ \langle Z,\calR \rangle = o(\delta d) $ .
\end{enumerate}
\end{claim}
Combining the above gives us \[
    \|\Pi^\perp I\|_F^2
    +
    |c_R-1|\,|\Tr(\calR)|
    +
    |\Tr(Z)|
    +
    |\langle Z,\calR\rangle|
    =
    \widetilde O(\sqrt d)+o(\delta d) = o(\delta d)\,,
\]
since $\delta = 1/\poly(\log \log d)$.
Therefore, we obtain
\begin{align*}
    \langle\Lambda_R,I-\calR\rangle
    &=
    d-c_R\frac d\rho+o(\delta d)\\
    &=
    -\frac{c_R-\rho}{\rho}\,d+o(\delta d)\\
    &=
    -\Theta(\delta d)<0.
\end{align*}

Finally, since
\[
    \delta=D^{-1/2}= 1/\poly(\log \log d)= o_d(1),
\]
both of our theorems hold with vanishing slack $o_d(1)$ as desired.

\clearpage
\newpage
\bibliographystyle{alpha}
\bibliography{bib}
\clearpage
\newpage
\appendix
%\input{scalar.tex}

\clearpage
\newpage


\section{Deferred Details for the Iterative Construction} \label{sec:defer-construction}

\subsection{Concrete Dual Iterative Process via Graph Matrices}
\label{app:concrete-dual-iteration}

We modify the ideal dual iteration exactly as in the primal program in~\cref{sec:concrete-primal-graph-mat}: all Chebyshev and MP expansions
are truncated at degree \(D\), only proper horizontal concatenations are
corrected, and only the proper shapes produced by the vertical
concatenation are retained. Since
\[
    \corr_{\mathrm{dual}}(H)
    =
    -\frac{1-\gamma}{\gamma}\,
    \corr_{\mathrm{primal}}(H),
\]
the shape decomposition and error bounds are unchanged up to a bounded
scalar factor.

Define
\[
    H_0^{\mathrm{prop}}
    \coloneqq
    \sum_{j=2}^{D}b_j\fp_j(W_0),
\]
and, for \(i\geq1\),
\[
    H_i^{\mathrm{prop}}
    \coloneqq
    \sum_{j=2}^{D}b_j
    \bigl(\fp_j(W_i)-\fp_j(W_{i-1})\bigr).
\]
Let
\[
    \corr_{\mathrm{dual}}(H)
    \coloneqq
    \frac1\gamma
    \left(
        \calL^\ast M^{-1}\calL(H)-\gamma H
    \right),
\]
and let
\(\corr_{\mathrm{dual,prop}}^{\leq D}(H)\) denote the proper part of
its degree-\(D\) MP expansion, after linearizing all well-behaved
vertical intersections.

\begin{mdframed}[linewidth=0.2pt]
\textbf{Summary of the Concrete Truncated Dual Procedure.}

\medskip
\textbf{Initialization:}
Set
\[
    W_0
    \coloneqq
    C_Fc_R\operatorname{Proper}(\calR).
\]

\medskip
\textbf{Iterative update \(i\to i+1\):}
For \(0\leq i<t^\ast\), set
\[
    \Delta_i^{\mathrm{dual}}
    \coloneqq
    \corr_{\mathrm{dual,prop}}^{\leq D}
    \bigl(H_i^{\mathrm{prop}}\bigr),
    \qquad
    W_{i+1}\coloneqq W_i+\Delta_i^{\mathrm{dual}}.
\]

\medskip
\textbf{Output:}
Let \(W\coloneqq W_{t^\ast}\). The final dual witness is obtained from
\[
    \frac1{C_F}F^{\leq D}(W)-\Pi^\perp I
\]
by removing the aggregate discrepancy below and correcting the final
proper residual.
\end{mdframed}
To account for the change, define
\[
    \mathsf{HorizonErr}^{\mathrm{dual}}(i)
    \coloneqq
    \sum_{j=2}^{D}b_j
    \bigl(P_j(W_i)-\fp_j(W_i)\bigr),
\]
\[
    \mathsf{UpdateErr}^{\mathrm{dual}}(i)
    \coloneqq
    \Delta_i^{\mathrm{dual}}
    -
    \corr_{\mathrm{dual}}(H_i^{\mathrm{prop}}),
\]
and
\[
    \mathsf{IterErr}^{\mathrm{dual}}_i
    \coloneqq
    \mathsf{HorizonErr}^{\mathrm{dual}}(i)
    +
    \sum_{s=0}^{i-1}
    \mathsf{UpdateErr}^{\mathrm{dual}}(s).
\]
Here the update error collects the MP-truncation and discarded vertical
intersection terms.

\begin{proposition}[Invariant for the concrete dual iteration]
\label{prop:concrete-dual-iteration-invariant}
For every \(i\geq0\),
\[
    \Pi^\perp
    \left(
        \frac1{C_F}
        \bigl(
            F^{\leq D}(W_i)
            -
            \mathsf{IterErr}^{\mathrm{dual}}_i
        \bigr)
        -
        I
    \right)
    =
    \frac1{C_F}\Pi^\perp H_i^{\mathrm{prop}}.
\]
\end{proposition}
%
%The proof is identical to the primal case, using
%\[
%    H+\corr_{\mathrm{dual}}(H)
%    =
%    \frac1\gamma
%    \calL^\ast M^{-1}\calL(H)
%    \in\calS.
%\]
%
%At the terminal level, set
%\[
%    D_E^{\mathrm{dual}}
%    \coloneqq
%    \frac1{C_F}
%    \corr_{\mathrm{dual}}
%    \bigl(H_{t^\ast}^{\mathrm{prop}}\bigr)
%\]
%and define
%\[
%    \Lambda_R
%    \coloneqq
%    \frac1{C_F}
%    \left(
%        F^{\leq D}(W_{t^\ast})
%        -
%        \mathsf{IterErr}^{\mathrm{dual}}_{t^\ast}
%    \right)
%    +
%    D_E^{\mathrm{dual}}
%    -
%    \Pi^\perp I.
%\]
%Since
%\[
%    \Pi^\perp\corr_{\mathrm{dual}}(H)
%    =
%    -\Pi^\perp H,
%\]
%\Cref{prop:concrete-dual-iteration-invariant} implies
%\[
%    \Lambda_R\in\calS.
%\]
%
%As in the primal case, the horizontal and vertical errors are controlled
%by \Cref{lem:semicircle-shape-concatenation,lem:error-vertical}, while
%the MP-truncation error is handled in Section~6. Consequently,
%\[
%    \bigl\|
%        \mathsf{IterErr}^{\mathrm{dual}}_{t^\ast}
%    \bigr\|_{\mathrm{sp}}
%    =
%    o(\delta).
%\]


\begin{comment}
	

\subsection{Formal Truncated Iterative Procedures}
\label{app:formal-truncated-iteration}

The finite process used in the proof differs from the ideal iteration in
three ways: all Chebyshev expansions are truncated at degree \(D\); only
the proper-concatenation part \(\fp_j\) is used in each update; and the
omitted intersection terms are subtracted from the final matrix. In
addition, the MP expansion used to analyze each occurrence of \(A^{-1}\)
is truncated at the same degree \(D\), with the resulting tail controlled
separately by Lemma~6.8.

We state the construction for \(F_\delta\). For any matrix \(X\), define
the aggregate \emph{intersection remainder}
\[
    E_{\mathrm{int}}(X)
    \coloneqq
    \sum_{j=2}^{D} b_j^{(\delta)}
    \bigl(P_j(X)-\fp_j(X)\bigr),
\]
where \(\fp_j(X)\) denotes the sum of the \(j\)-fold proper
concatenations in the graph-matrix expansion of \(P_j(X)\). Thus
\begin{equation}
\label{eq:proper-intersection-decomposition}
    F_\delta^{\leq D}(X)-E_{\mathrm{int}}(X)
    =
    C_\delta I+X+
    \sum_{j=2}^{D}b_j^{(\delta)}\fp_j(X).
\end{equation}

\begin{mdframed}[linewidth=0.2pt]
\textbf{Formal Truncated Primal Iteration.}

\medskip
\textbf{Initialization:}
Set
\[
    Q_0\coloneqq -C_\delta\calR,
\]
so that \(I+Q_0/C_\delta\) satisfies the affine constraints.

\medskip
\textbf{Iterative update \(i\to i+1\):}
Define
\[
    H_0
    \coloneqq
    \sum_{j=2}^{D}b_j^{(\delta)}\fp_j(Q_0),
\]
and, for \(i\geq1\),
\[
    H_i
    \coloneqq
    \sum_{j=2}^{D}b_j^{(\delta)}
    \bigl(\fp_j(Q_i)-\fp_j(Q_{i-1})\bigr).
\]
For \(0\leq i<t^\ast\), set
\[
    Q_{i+1}
    \coloneqq
    Q_i+\corr_{\mathrm{primal}}(H_i),
\]
where
\[
    \corr_{\mathrm{primal}}(H)
    \coloneqq
    \frac{1}{1-\gamma}
    \left(
        \calL^\ast M^{-1}\calL(-H)+\gamma H
    \right).
\]

\medskip
\textbf{Output:}
Define the terminal correction
\[
    D_E
    \coloneqq
    \frac1{C_\delta}
    \corr_{\mathrm{primal}}(H_{t^\ast})
\]
and output
\[
    \Lambda
    \coloneqq
    \frac1{C_\delta}
    \left(
        F_\delta^{\leq D}(Q_{t^\ast})
        -
        E_{\mathrm{int}}(Q_{t^\ast})
    \right)
    +D_E.
\]
\end{mdframed}

The affine constraints remain exact under these modifications. Indeed,
using
\[
    \calL\!\left(\corr_{\mathrm{primal}}(H)\right)
    =
    -\calL(H),
\]
the same telescoping argument as in the ideal process gives
\begin{equation}
\label{eq:formal-truncated-primal-invariant}
    \calL\!\left(
        \frac1{C_\delta}
        \left(
            F_\delta^{\leq D}(Q_i)
            -
            E_{\mathrm{int}}(Q_i)
        \right)
    \right)
    -
    \1_m
    =
    \frac1{C_\delta}\calL(H_i).
\end{equation}
The terminal correction therefore yields
\[
    \calL(\Lambda)=\1_m.
\]

It remains to account for the terms removed in the construction. By
\Cref{claim:shape-size-truncated-construction}, every shape in
\(E_{\mathrm{int}}(Q_{t^\ast})\) has \(d^{o(1)}\) vertices. The SC-LCP
bounds from Section~5.4 consequently give
\[
    \left\|
        E_{\mathrm{int}}(Q_{t^\ast})
    \right\|_{\mathrm{sp}}
    =
    d^{-1/2+o(1)}
    =
    o(\delta).
\]
The degree-\(D\) MP tail used in the analysis of the correction terms is
controlled by Lemma~6.8, while Lemma~6.9 gives
\[
    \|D_E\|_{\mathrm{sp}}
    =
    O\!\left(\sqrt{\varepsilon_{D,t^\ast}}\right)
    =
    o(\delta).
\]
Thus subtracting the intersection remainder contributes only one
additional \(o(\delta)\) term to the final minimum-eigenvalue estimate.

\paragraph{Formal truncated dual iteration.}
The dual process is modified in exactly the same manner. Set
\[
    W_0\coloneqq C_\delta c_R\calR,
\]
and define
\[
    H_0^{\mathrm{dual}}
    \coloneqq
    \sum_{j=2}^{D}b_j^{(\delta)}\fp_j(W_0),
\]
and, for \(i\geq1\),
\[
    H_i^{\mathrm{dual}}
    \coloneqq
    \sum_{j=2}^{D}b_j^{(\delta)}
    \bigl(\fp_j(W_i)-\fp_j(W_{i-1})\bigr).
\]
For \(0\leq i<t^\ast\), update
\[
    W_{i+1}
    \coloneqq
    W_i+
    \corr_{\mathrm{dual}}(H_i^{\mathrm{dual}}),
\]
where
\[
    \corr_{\mathrm{dual}}(H)
    \coloneqq
    \frac1\gamma
    \left(
        \calL^\ast M^{-1}\calL(H)-\gamma H
    \right).
\]
The final dual witness is
\[
    \Lambda_R
    \coloneqq
    \frac1{C_\delta}
    \left(
        F_\delta^{\leq D}(W_{t^\ast})
        -
        E_{\mathrm{int}}(W_{t^\ast})
        +
        \corr_{\mathrm{dual}}
        (H_{t^\ast}^{\mathrm{dual}})
    \right)
    -
    \Pi^\perp I.
\]

Since
\[
    \Pi^\perp\corr_{\mathrm{dual}}(H)
    =
    -\Pi^\perp H,
\]
the dual analogue of
\eqref{eq:formal-truncated-primal-invariant} shows that
\[
    \Lambda_R\in\operatorname{Im}(\calL^\ast)
\]
exactly. Moreover, the same SC-LCP estimate gives
\[
    \left\|
        E_{\mathrm{int}}(W_{t^\ast})
    \right\|_{\mathrm{sp}}
    =
    o(\delta).
\]
Hence the shape-size, MP-truncation, intersection, and early-termination
bounds established for the primal construction apply unchanged to the
dual construction.
\end{comment}

\begin{comment}
\subsection{Summary of the Iterative Construction via Graph Matrices}

%\subsection{Iterative Construction Made Concrete via Graph Matrices}
\label{sec:iterative-construction}
%This correspondence lets us identify which terms have already been corrected
%and which terms are new. 
In this section, we give a full description of the iterative process, which will in turn identify the collection of graph matrices that arise throughout our analysis. Throughout this section, we restrict our attention to proper shapes.

For an inner matrix $Q_i$, we first consider what arises when we apply $F$, more concretely, in the polynomial expansion of $\frac{1}{C_F} F(Q_i)$, which is \[ 
\frac{1}{C_F} F(Q_i) = I + Q_i/C_F  + \sum_{j\geq 2} \frac{b_j}{C_F} P_j(Q_i)
\,.\]
Let's zoom into the higher-degree term. For each \(j\ge2\), define the collection of  shapes in the $j$-way concatenation of shapes in $Q_i$ as
\[
    \calP_j(Q_i)
    :=
    \left\{
        \tau_1\circ\cdots\circ\tau_j:
        \tau_t\in\calB(Q_i)\ \forall t\in[j]
    \right\}.
\]
Moreover, depending on whether the concatenation uses any shape from the most recent update, specifically in $\calB(Q_i)\setminus \calB(Q_{i-1})$,  we can partition the collection of $j$-way concatenation shapes as \[
    \textsf{Old}\text{-}\calP_j(Q_i)
    =
    \left\{
        \beta=\tau_1\circ\cdots\circ\tau_j\in\calP_j(Q_i):
        \tau_t\in\calB(Q_{i-1})\ \forall t\in[j]
    \right\},
\]
and
\[
    \textsf{New}\text{-}\calP_j(Q_i)
    =
    \left\{
        \beta=\tau_1\circ\cdots\circ\tau_j\in\calP_j(Q_i):
        \exists t\in[j]\ \text{such that }\
        \tau_t\in\calB(Q_i)\setminus\calB(Q_{i-1})
    \right\}.
\]
%Here and below, the notation also denotes the corresponding weightedgraph-matrix sums.
This partition then allows us to make explicit the following invariant that we would like to maintain across the iterations: in particular, the ``new'' terms are the only source of deviation with respect to the affine constraints.
\begin{proposition}[Iteration Invariant]
\label{prop:iterative-invar}
For any \(i\ge1\), consider the inner matrix \(
Q_i = \sum_{\tau} c_\tau \cdot \cm_\tau 
\),
 we have
\[
\frac{1}{C_F}F(  Q_i)
=
I+Q_i/C_F + \frac{1}{C_F}\cdot 
\sum_{j\ge2}
\left(
    \sum_{\tau\in\textsf{Old}\text{-}\calP_j(Q_i)}
    c_\tau\cm_\tau
    +
    \sum_{\tau\in\textsf{New}\text{-}\calP_j(Q_i)}
    c_\tau\cm_\tau
\right).
\]
Moreover,
\[
    I+Q_i/C_F + \frac{1}{C_F}\cdot 
    \sum_{j\ge2}
    \sum_{\tau\in\textsf{Old}\text{-}\calP_j(Q_i)}
    c_\tau\cm_\tau
\]
satisfies the affine constraints.
\end{proposition}

\begin{remark}
The same statement applies to the base level \(i=0\) if the old part is treated
as empty.
\end{remark}

We defer the proof of \cref{prop:iterative-invar} to
\cref{sec:defer-construction}.  The invariant shows that, at iteration \(i\),
the only remaining source of affine-constraint error is the new part
\(
    \sum_{j\ge2}
    \sum_{\tau\in\textsf{New}\text{-}\calP_j(Q_i)}
    c_\tau\cm_\tau.
\)
This is therefore the correction target for the next iteration. One lingering question
remains: how do we correct these affine deviations?

To answer this, we again draw inspiration from the identity-perturbation
construction. We seek a correction matrix of the form
\(
    \sum_{t\in[m]} w[t]\,v_tv_t^\top
\)
whose constraint values cancel the current deviation. Concretely, if
\[
    \eta_{\mathrm{new}}[i]
    :=
    v_i^\top
    \left(
        \sum_{\tau\in \mathsf{New}\text{-}\calP_j(Q_i)}
        c_\tau \cm_\tau
    \right)
    v_i
\]
denotes the affine deviation to be corrected, then we want
\[
    v_i^\top
    \left(
        \sum_{t\in[m]} w[t]\,v_tv_t^\top
    \right)
    v_i
    =
    -\eta_{\mathrm{new}}[i],
    \qquad \forall i\in[m].
\]
Equivalently, this reduces to solving the linear system
\(
    Mw=-\eta_{\mathrm{new}}.
\)
Thus, just as in the identity-perturbation construction, the desired correction
is obtained by applying \(M^{-1}\) to the current affine deviation, and thereby yielding an update to our inner matrix.


Importantly, due to correlated randomness between the rank-$1$ forms and the deviation vector, we carefully choose our correction update as follows for any $\tau \in \mathsf{New-}P_j(Q_i)$:
 \[ 
    \corr(\tau) \coloneqq -  \frac{1}{1-\gam } \left( \mathcal{L}^*(M^{-1} \eta_{\tau}) - \gam \cdot \tau  \right)\,.
    \]
    \jnote{this needs to be highlighted}

%In other words, each iteration corrects the newly produced affine deviation by
%solving the same constraint matrix \(M\), and then converts the resulting
%coefficient vector \(w\) into the rank-one correction
%\[
%    \sum_{t\in[m]} w[t]\,v_tv_t^\top .
%\]

\paragraph{Concrete Iterative construction.}
We summarize our iterative process as follows.
\begin{mdframed}[linewidth=0.2pt]
\textbf{Iterative update \(i\to i+1\).}
\begin{enumerate}
    \item For each \(j\ge2\), form the set of \(j\)-fold concatenations
    \[
        \calP_j(Q_i)
        =
        \left\{
            \tau_1\circ\cdots\circ\tau_j:
            \tau_t\in\calB(Q_i)\ \forall t\in[j]
        \right\}.
    \]

    \item Split \(\calP_j(Q_i)\) into old and new terms:
    \[
        \textsf{Old}\text{-}\calP_j(Q_i)
        =
        \left\{
            \tau_1\circ\cdots\circ\tau_j:
            \tau_t\in\calB(Q_{i-1})\ \forall t\in[j]
        \right\},
    \]
    and
    \[
        \textsf{New}\text{-}\calP_j(Q_i)
        =
        \left\{
            \tau_1\circ\cdots\circ\tau_j:
            \exists t\in[j]\ \text{such that }\
            \tau_t\in\calB(Q_i)\setminus\calB(Q_{i-1})
        \right\}.
    \]

    \item For each new concatenation
    \[
        \beta=\tau_1\circ\cdots\circ\tau_j
        \in \textsf{New}\text{-}\calP_j(Q_i),
    \]
    its coefficient is
    \[
        c_\beta
        =
        b_j \cdot \prod_{t=1}^j c_{\tau_t}.
    \]
    in the polynomial expansion of $ F(Q_i)$.
    This term appears in the spectral expansion of \(F(Q_i)\), but has not yet
    been corrected inside \(Q_i\).

    \item Define the new unscaled affine residual by
    \[
        \eta_{i+1}[t]
        =
        v_t^\top
        \left(
            \sum_{j\ge2}
            \sum_{\tau\in\textsf{New}\text{-}\calP_j(Q_i)}
            c_\tau\cm_\tau
        \right)
        v_t,
        \qquad t\in[m].
    \]
    Solve
    \[
        Mw_{i+1}=-\eta_{i+1}.
    \]
    Equivalently, one may decompose
    \[
        w_{i+1}
        =
        \sum_{j\ge2}
        \sum_{\tau\in\textsf{New}\text{-}\calP_j(Q_i)}
        w_\tau,
    \]
    where \(w_\tau\) is the correction vector associated with the new shape
    \(\tau\). We will set its corresponding correction term as \[ 
    \corr(\tau) \coloneqq -  \frac{1}{1-\gam } \left( \mathcal{L}^*(M^{-1} \eta_{\tau}) - \gam \cdot \cm_\tau  \right)\,.
    \]

    \item Update the inner matrix by
    \[
        Q_{i+1}-Q_i
        =
        \sum_{t\in[m]}w_{i+1}[t]\,v_t v_t^\top.
    \]
    Equivalently,
    \[
        Q_{i+1}-Q_i
        =
        \sum_{t\in[m]}
        \sum_{j\ge2}
        \sum_{\tau\in\textsf{New}\text{-}\calP_j(Q_i)}
        w_\tau[t]\,v_t v_t^\top.
    \]
    In graph-matrix language, the ground set evolves by adding the correction
    shapes associated with the newly appearing concatenations:
    \[
        \calB(Q_{i+1})
        =
        \calB(Q_i)
        \cup
        \left\{
            \text{correction shapes for }
            \tau:
            \tau\in\textsf{New}\text{-}\calP_j(Q_i),\ j\ge2
        \right\}.
    \]
\end{enumerate}
\end{mdframed}

\begin{remark}
In the formal analysis, we use an inner truncation: for each \(Q_i\), we only
retain \(P_j(Q_i)\) for \(j\le D\).
\end{remark}

For the final construction, we take \(Q=Q_t\) for a truncation level \(t\), and
then handle the remaining error terms directly.  This is addressed in~\cref{sec:formal-truncation}.  For the purposes of the overview, however, it is
useful to think of \(Q\) as the limit of the iterates \(Q_i\) as \(i\to\infty\). 

%
%
%\begin{proof}[Proof to \cref{prop:iterative-invar}]
%\begin{proof}
%	Let
%	\(
%		L(X)
%		:=
%		\bigl(v_s^\top Xv_s\bigr)_{s\in[m]}
%	\)
%	denote the affine-constraint map. Thus a matrix \(X\) satisfies the affine constraints if
%	\(
%		 L(X)=\mathbf 1_m,
%	\)
%	and a correction term \(Y\) preserves the affine constraints if
%	\(
%		 L(Y)=0.
%	\)
%
%	We prove the second statement by induction on the iteration level. Define
%	\[
%		R_i
%		:=
%		I+Q_i/C_F+
%		\sum_{j\ge 2}
%		\sum_{\tau\in \textsf{Old}\text{-}\calP_j(Q_i)}
%		c_\tau \cm_\tau .
%	\]
%	be the restriction of $\frac{1}{C_F} F(Q_i)$ to the degree-$1$, and the old components from $F(Q_{i-1})$
%	We will show that
%	\[
%	 L(R_i)=\mathbf 1_m
%	\]
%	for every \(i\ge 0\).
%
%	For the base case, the old part is empty, so
%	\[
%		R_0=I+Q_0/C_F.
%	\]
%	By initialization, \(Q_0\) is chosen so that
%	\[
%		v_s^\top (I+Q_0/C_F)v_s=1
%		\qquad \forall s\in[m],
%	\]
%	or equivalently \(\mathcal A(R_0)=\mathbf 1_m\). Now assume inductively that
%	\(
%		L(R_i)=\mathbf 1_m.
%	\)
%	We prove that \(\mathcal A(R_{i+1})=\mathbf 1_m\). By the update rule, every new concatenation
%	\[
%		\tau\in \textsf{New}\text{-}\calP_j(Q_i)
%	\]
%	is corrected at the next level by adding its corresponding correction shape \(\bar\tau\) to
%	\(Q_{i+1}\) with the same coefficient \(c_\tau\). Thus, in graph-matrix notation,
%	\[
%		Q_{i+1}
%		=
%		Q_i+
%		\sum_{j\ge 2}
%		\sum_{\tau\in \textsf{New}\text{-}\calP_j(Q_i)}
%		c_\tau \cm_{\corr(\tau)}.
%	\]
%	Moreover, the old concatenations at level \(i+1\) are exactly the concatenations that use
%	only shapes already present in \(Q_i\). Hence they decompose as
%	\[
%		\textsf{Old}\text{-}\calP_j(Q_{i+1})
%		=
%		\textsf{Old}\text{-}\calP_j(Q_i)
%		\sqcup
%		\textsf{New}\text{-}\calP_j(Q_i).
%	\]
%	Therefore
%	\[
%	\begin{aligned}
%		R_{i+1}
%		&=
%		I+Q_{i+1}/C_F
%		+
%		\sum_{j\ge 2}
%		\sum_{\tau\in \textsf{Old}\text{-}\calP_j(Q_{i+1})}
%		c_\tau\cm_\tau  \\
%		&=
%		R_i+
%		\sum_{j\ge 2}
%		\sum_{\tau\in \textsf{New}\text{-}\calP_j(Q_i)}
%		c_\tau \cm_\tau+ c_{\corr(\tau)} \cdot \cm_{\corr(\tau)}.
%	\end{aligned}
%	\]
%\jnote{requires examination on the coef.}
%	By construction of the correction shape \(\corr(\tau)\) (with the corresponding coefficient) obtained by solving for the inverse of the linear system, the pair
%	\[
%		c_\tau \cm_\tau+c_{\corr(\tau)}\cm_{\corr(\tau)}
%	\]
%	lies in the kernel of the affine-constraint map, namely
%	\[
%		 L(c_\tau  \cdot \cm_\tau+ c_{\corr(\tau)}  \cdot \cm_{\corr(\tau)})=0.
%	\]
%	Thus each added term
%	\(
%		c_\tau \cm_\tau+ c_{\corr(\tau)} \cm_{\corr(\tau)}
%	\)
%	groups to some element in the kernel of $L$ and thereby preserves the affine constraints. Applying \(L\) gives
%	\[
%	\begin{aligned}
%		  L(R_{i+1})
%		&=
%		  L(R_i)
%		+
%		\sum_{j\ge 2}
%		\sum_{\tau\in \textsf{New}\text{-}\calP_j(Q_i)}
%		(c_\tau \cm_\tau+ c_{\corr(\tau)} \cm_{\corr(\tau)})
% \\
%		&=
%		\mathbf 1_m.
%	\end{aligned}
%	\]
%	This proves the induction step.
%\jnote{elaborate on the choice of coefficient?}
%	Therefore \( L(R_i)=\mathbf 1_m\) for every \(i\ge0\), and in particular the claimed
%	affine-constraint invariant holds for every \(i\ge1\).
%\end{proof} \end{proof}

\begin{proposition}[No reappearance of old shapes in later updates]
\label{prop:no-reappearance-old-shapes}
Suppose a shape \(\sigma\) is first added to the ground set \(\calB(Q_i)\) at
iteration level \(i\). Then \(\sigma\) does not appear in any subsequent update:
for every \(j\ge i\),
\[
    \sigma \notin \calB(Q_{j+1})\setminus \calB(Q_j).
\]
Consequently, the coefficient of \(\sigma\) remains unchanged throughout all
later iterations.
\end{proposition}

\begin{proof}
It suffices to keep track of the deviation term that creates each added shape.
Indeed, every shape added in the update \(Q_{j+1}-Q_j\) is the correction shape
associated with some new deviation term in
\[
    \textsf{New}\text{-}P_\ell(Q_j)
    =
    P_\ell(Q_j)\setminus P_\ell(Q_{j-1}),
    \qquad \ell\ge 2.
\]
By definition, such a new deviation term contains at least one shape from the
most recent level of \(
    \calB(Q_j)\setminus \calB(Q_{j-1}).
\)
Consequently, the correction shape added to \(Q_{j+1}\) contains a newly created
dangling path that was not present before level \(j\).

Thus, once a shape is added to the inner matrix at level $i$, any subsequent addition to the inner matrix contains a dangling path unseen at level-$i$, hence once \(\sigma\) is added, it never appears in any later update. Since the
iterative construction only appends new shapes and never modifies existing
coefficients, the coefficient of \(\sigma\) remains unchanged in all subsequent
iterations.
\end{proof}
\end{comment}

\begin{proposition}[No reappearance of old shapes in later updates]
\label{prop:no-reappearance-old-shapes}
Suppose a shape \(\sigma\) is first added to the ground set \(\calB(Q_i)\) at
iteration level \(i\). Then \(\sigma\) does not appear in any subsequent update:
for every \(j\ge i\),
\[
    \sigma \notin \calB(Q_{j+1})\setminus \calB(Q_j).
\]
Consequently, the coefficient of \(\sigma\) remains unchanged throughout all
later iterations.
\end{proposition}

\begin{proof}
It suffices to keep track of the deviation term that creates each added shape.
Indeed, every shape added in the update \(Q_{j+1}-Q_j\) is the correction shape
associated with some new deviation term in
\[
    \textsf{New}\text{-}P_\ell(Q_j)
    =
    P_\ell(Q_j)\setminus P_\ell(Q_{j-1}),
    \qquad \ell\ge 2.
\]
By definition, such a new deviation term contains at least one shape from the
most recent level of \(
    \calB(Q_j)\setminus \calB(Q_{j-1}).
\)
Consequently, the correction shape added to \(Q_{j+1}\) contains a newly created
dangling path that was not present before level \(j\).

Thus, once a shape is added to the inner matrix at level $i$, any subsequent addition to the inner matrix contains a dangling path unseen at level-$i$; hence, once \(\sigma\) is added, it never appears in any later update. Since the
iterative construction only appends new shapes and never modifies existing
coefficients, the coefficient of \(\sigma\) remains unchanged in all subsequent
iterations.
\end{proof}



\subsection{Deferred Details for $A^{-1}$ Expansion}
\begin{claim}[Restatement of \cref{clm:inverse-expansion}]
For \(x\) in the Marchenko--Pastur support and \(\gamma<1\),
\[
    \frac{1}{x}
    =
    \frac{1}{1-\gamma}
    \sum_{t\geq 0} (-1)^t \q_t(x).
\]
\end{claim}

\begin{proof}
Let
\(
    G(z,x)
    :=
    \sum_{t\geq 0} \q_t(x) z^t
\)
be the generating function of the polynomials. Using the recurrence above,
we have\[
    G(z,x)
    =
    \frac{1+\gamma z}
    {1-\bigl(x-(1+\gamma)\bigr)z+\gamma z^2}.
\]
Evaluating at \(z=-1\) gives
\[
    \sum_{t\geq 0} (-1)^t \q_t(x)
    =
    G(-1,x)
    =
    \frac{1-\gamma}{x}.
\]
Rearranging proves the claim.
\end{proof}


\subsection{Diagram Illustrations of the Iterative Scheme}

\begin{figure}[h]
    \centering

    \begin{subfigure}[b]{0.6\textwidth}
        \includegraphics[width=\textwidth]{diagrams/old_3.pdf}
%        \caption{Shape $\tau \in P_2(Q_0)$}
%        \label{fig:p2(Q)}
    \end{subfigure}

\begin{subfigure}[b]{0.6\textwidth}
        \includegraphics[width=\textwidth]{diagrams/old_3p.pdf}
%        \caption{Shape $\tau \in P_2(Q_0)$}
%        \label{fig:p2(Q)}
    \end{subfigure}

    %	    \vspace{1em}
%
%	\cenring 
%        \begin{subfigure}[b]{0.48\textwidth}
%        \includegraphics[width=0.94\textwidth]{diagrams/old-3.pdf}
%        \caption{Backbone-dangling shape in $ Q_1\setminus Q_0$}
%            \end{subfigure}
%
    % \captionsetup{width=.9\linewidth}
    % \caption{Examples of graph matrices in the decomposition of $Q$.
    % Recall that square vertices take labels in $[m]$ and circle vertices take labels in $[d]$.
    % Unlabeled edges have Hermite index $1$.}
    % \label{fig:examples}
\end{figure}

\section{Deferred Deviation Bounds}
\label{sec:def-scalar}
\subsection{Calculations for Scalar Concentration}

\begin{proof}[Proof of~\cref{clm:base-scalar-variance}]
	
Recall that $M^{-1}$, as well as $M$, has two components. We bound the quadratic form of $\eta$ with respect to each component separately as follows.
% We record the following scalar concentration estimates, deferring their proofs to the appendix. For the variance calculation, we state only their first-moment versions, although the same arguments also yield \(1+o_d(1)\) concentration around the corresponding means with high probability.

 \begin{restatable}{claim}{EtaAInvEtaConcentration} 
\label{clm:concentration-eta-A-eta}
We have
\[
    \frac{1}{d}\cdot \E\!\left[\eta^\top A^{-1}\eta\right]
    =
    \gamma + o_d(1).
\]
\end{restatable}
% 	 and analogously for the second term from Woodbury expansion, \[ 
%\frac{1}{d} \cdot  \E\bigg[\eta^T \left( \frac{1}{s^2-ru}\cdot A^{-1} \cdot (u\cdot \frac{J_m}{d} -s \frac{\eta\cdot \1_m^T+ 1_m^T \cdot \eta }{d}
% + r \cdot \frac{\eta \cdot \eta^T }{d}) \cdot A^{-1}\right) \eta  \bigg ] = \frac{\gam^2}{1-\gamma}\,. 
%\]
 
\begin{restatable}{claim}{WoodburyCorrectionConcentration} From the rank-$2$ update in Woodbury, 
\label{clm:concentration-woodbury}
we have
\[
\frac{1}{d} \cdot
\E\!\left[
\eta^\top
\left(
\frac{1}{s^2-ru}\,
A^{-1}
\left(
u\cdot \frac{J_m}{d}
-
s\cdot \frac{\eta \mathbf 1_m^\top+\mathbf 1_m \eta^\top}{d}
+
r\cdot \frac{\eta\eta^\top}{d}
\right)
A^{-1}
\right)
\eta
\right]
=
\frac{\gamma^2}{1-\gamma}.
\]

\end{restatable}
The proof of the claim follows by combining the above two claims.

\end{proof}

\HyperParameterBounds*
\begin{proof}
	The dominant term is already handled, and we focus on the deviations terms here.
	
	\paragraph{Analysis for $r$}
		Recall that $r = \frac{1}{d}  \cdot \1_m^T \cdot A^{-1} \cdot \1_m  $, and  \[ 
 	A^{-1} = \frac{1}{1-\gamma} \sum_{t\geq 0} (-1)^t \cdot \q_t(A)  = \frac{1}{1-\gamma} \sum_{t\geq 0} (-1)^t \cdot \left(\fp_t(A)+o_d(1) \right)
 	\]
 	where the second equality ignores the $o_d(1)$ error term in passing from $\q_t(A)$ to graph matrices of concatenated shapes.
	 
	 	Expanding the \(2q\)-th power, each summand is described in graph-matrix language by a closed walk of \(2q\) floating components, where each component has the form \[
    \Box
    \;-\;
    \bigl(A^{-1}\text{ path via concatenation}\bigr)
    \;-\;
    \Box .
\] Here the two boundary square vertices come from the summation of \(\mathbf 1_m\), and the global normalization contributes a factor \(1/d\) for each component.
	  
	  It suffices for us to bound \[
	 \E[(r - \frac{m}{d})^{2q} ] = \E[(\frac{1}{d} \sum_{t>0} \1_m^\top \q_t(A) \cdot \1_m)^{2q}]  
	 \,,\]  
	 Next, we claim that for any $t$, the block-value is at most \(O(\frac{\sqrt{2q}}{d}\cdot  \sqrt{m} \cdot \sqrt{\gamma}^t ) \):\begin{enumerate}
	 	\item The $\frac{1}{d}$-factor comes from the normalization factor in $r$;
	 	\item For any $t>0$, each edge needs to appear at least twice in the calculation, similar to the usual trace method calculation. Hence, each vertex also makes at least two appearances;
	 	\item The first square vertex needs to appear at least twice, hence we assign a factor of $\sqrt{m} \sqrt{2q}$ for each of its appearance;
	 	\item The remaining factors are identical as those in the calculation for the spectral norm for $\fp_t(A)$, in particular, each $\al$ gadget gives a factor of $\sqrt{\gamma}$ and $\beta$ gadget, as well as error terms from approximation from shape concatenation gives $o_d(1)$ factor.
	 \end{enumerate} 
	 Setting \(q=\polylog(d)\) and recalling  \(m=\Theta(d^2)\), this bound is \[ O_\gamma\!\left( \frac{\sqrt{2q}}{d}\sqrt m \right) = O_\gamma(\sqrt q) = o_d\!\left(\frac{m}{d}\right), \] because \(m/d=\Theta(d)\). By Markov's inequality, the same estimate holds with high probability: \[ \frac{1}{d}\sum_{t\ge 1}\mathbf 1_m^\top P_t(A)\mathbf 1_m = o_d\!\left(\frac{m}{d}\right). \] Combining this with the \(P_0(A)\) contribution, we conclude that \[ r = \frac{1}{1-\gamma}\cdot \frac{m}{d} + o_d\!\left(\frac{m}{d}\right) = (1+o_d(1))\frac{m}{d}\cdot \frac{1}{1-\gamma}. \]
	 
	 
	\paragraph{Analysis for $s$} 
	 Recall that $s = 1+ \frac{1}{d} \cdot \eta^T A^{-1} \1_m$; it suffices for us to bound the $2q$-th moments of $\frac{1}{d} \eta^T A^{-1} \1_m$. In graph matrix language, this is a length-$2q$ walk across floating components where each component has the form\[
    \circ
    \xleftrightarrow{\;h_2\;}
    \Box
    \;-\;
    \bigl(A^{-1}\text{ path via concatenation}\bigr)\;-\;\Box
        \,.
\]
We bound the factors as follows,\begin{enumerate}
	\item The starting circle vertex gets a factor of $\sqrt{d\cdot (2q) }$, and we have a normalization factor of $\frac{1}{d}$;
	\item The $h_2$ edge gives a factor of $\frac{1}{d}$, and we combine that with the factor of the square vertex, that is \[ 
	\frac{\sqrt{2} }{d} \cdot \sqrt{m} =O(1) \,. 
	\]
	\item The subsequent factors are identical as the previous calculation, that we have a factor of $(1+o_d(1))\sqrt{\gamma}$ for each $\cm_\al$ gadget from the concatenation as well as its error term.
\end{enumerate}
For any length-$t\geq 0$ concatenation, this combines to a factor of \[ 
\frac{1}{d } \cdot  \sqrt{d\cdot (2q) }  \cdot O(1) \cdot \sqrt{\gam}^t\,.
\]
Since $\gam<1$, setting $q=\Theta(\polylog(d))$ gives us the desired bound.

\paragraph{Analysis for $u$}
For $t=0$, we have \[
\frac{1}{d} \cdot \eta^\top \q_0(A) \eta  = (1+o_d(1)) \cdot \frac{2md }{d^2} = (1+o_d(1)) \gam \,.
\]
 An analogous calculation applies for the higher moments.
   	
  	For $t=1$, we have \[
  	\q_1(A) = \fp_1(A) = \cm_\al + \cm_\beta \,.
  	 \]
  	 The crucial observation is $\eta^T \cm_\al \eta = o_d(1)$  while the second term gives $\gamma^2$. We first show \[ 
  	 \frac{1}{d } \eta^\top \cdot \cm_\beta \cdot \eta =(1+o_d(1)) \cdot \gamma^2 \,. 
  	 \]  	 w.h.p.
  	 In graph matrix language, each term is described \[ 
  	 \circ \xleftrightarrow{\;h_2\;} (\Box \xleftrightarrow{\;h_2\;} \circ  \xleftrightarrow{\;h_2\;} \Box) \xleftrightarrow{\;h_2\;}\circ 
  	 \,,\] and the dominant term is given by colliding all three circle vertices. Formally, this gives a contribution of \[
  	 \frac{1}{d }\cdot d \cdot  \frac{2}{d^2}\cdot m \cdot \frac{2}{d^2}\cdot (m-1)  =  (1+o_d(1))\cdot \gam^2\,.
  	  \]
  	 as we have $m(m-1)$ choices in total for the two square vertices.
  	 
  	 Next, to see that the other terms give an $o_d(1)$ contribution, we note that by injectivity of $\fp_t(A)$, each other term has mean-$0$. In other words, each edge needs to appear twice, and it cannot happen in a single floating component; thus we have \begin{enumerate}
  	 	\item A factor of $\frac{1}{d}$ from the normalization, and a factor of $\sqrt{d (2q)}$ for the starting circle vertex;
  	 	\item Each subsequent factor can be bounded by similar calculation for $A^{-1}$-path (as well as its error terms) to be at most $O(1)$.
  	 \end{enumerate}
  	 Finally, setting $q=\Theta(\polylog d)$ again gives us a bound of $O(\frac{\poly\log d}{\sqrt{d}}) = o_d(1)$ for each deviation term.   	 Rearranging gives us the desired bound for $u$ as we have \begin{align*}
  	 	 u = -1 + \frac{1}{d}\cdot \eta^\top A^{-1}\eta  &= -1+ \frac{1}{1-\gamma} (\eta^\top \q_0(A) \eta - \eta^\top \q_1(A) \eta  +o_d(1) ) \\&= -1 +\frac{1}{1-\gam }(\gam -\gam^2)\\&=-1+\gam +o_d(1)\,.  	    	 \end{align*}
	\end{proof}
	
	
Next, we complete the proofs for the variance calculation for the base inner matrix.
\EtaAInvEtaConcentration*
\WoodburyCorrectionConcentration*

\begin{proof}
	For starters, \cref{clm:concentration-eta-A-eta} follows from the calculation above where we give concentration of $u$. Next, for the second term, it is crucial to recall the following parameters from our previous bound, namely,\[ 
	 r= (1+o_d(1)) \cdot  \frac{m}{d} \cdot \frac{1}{1-\gamma} \,,
	\]
	\[
	s = 1+ o_d(1)\,,
	 \]
	 and \[
	 u = -1 +\gam +o_d(1) \,.
	  \]
	There are three terms that we need to bound here, with our focus on \[ 
	\frac{r}{s^2-ru} \cdot A^{-1} \cdot (\frac{\eta\cdot \eta^T}{d}) A^{-1}\,.
	\] 
%	By our concentration of $u$ (i.e., concentration for $\eta^\top A^{-1} \eta$), we have 
Next, we verify that we have
\[ 
	\frac{1}{d}\E [\eta^T ( A^{-1} \cdot (\frac{\eta\cdot \eta^T}{d}) A^{-1}) \eta ] =    \gam^2 +o_d(1)\,.
	\]
	since we have shown concentration of $\frac{1}{d}\cdot \eta^\top A^{-1} \eta =  \gamma +o_d(1) $ in the analysis for $u$.
	Moreover, by our previous bounds, we have
	\[
	\frac{r}{s^2-ru} =  \frac{1}{1/r-u} =  (1+o_d(1)) \cdot \frac{1}{1-\gamma}\,.
	 \]
	 
	 Therefore, \[ 
	 \frac{1}{d}\E\left[\eta^T(  \cdot \frac{r}{s^2-ru} \cdot  A^{-1} \cdot (\frac{\eta\cdot \eta^T}{d}) A^{-1} ) \eta   \right]  = \frac{\gamma^2}{1-\gamma} +o_d(1) \,.
	 \]
	We complete the proof by controlling the remaining two terms. Both are
\(o_d(1)\) with high probability, using the scalar concentration bounds above;
their expectations can be checked directly in the same manner.

For the \(u\cdot J_m\) term, we rewrite its contribution as
\[
    u \cdot
    \frac{\eta^\top A^{-1}\mathbf 1_m}{d}
    \cdot
    \frac{\mathbf 1_m^\top A^{-1}\eta}{d}
    =
    u\cdot (s-1)^2.
\]
Since \(s=1+o_d(1)\) and \(u=O(1)\) with high probability, this term is
\[
    u\cdot (s-1)^2=o_d(1).
\]

For the \(s\cdot(\eta\mathbf 1_m^\top+\mathbf 1_m\eta^\top)\) term, its contribution is
\[
    2s\cdot
    \frac{\eta^\top A^{-1}\eta}{d}
    \cdot
    \frac{\mathbf 1_m^\top A^{-1}\eta}{d}
    =
    2s\cdot
    \frac{\eta^\top A^{-1}\eta}{d}
    \cdot (s-1).
\]
Using \(s=1+o_d(1)\) and
\[
    \frac{\eta^\top A^{-1}\eta}{d}=O(\gamma)
\]
with high probability, we obtain
\[
    2s\cdot
    \frac{\eta^\top A^{-1}\eta}{d}
    \cdot (s-1)
    =
    O(\gamma)\cdot o_d(1)
    =
    o_d(1).
\]
Thus both remaining contributions are negligible.\end{proof}
	

\begin{comment}
	
\subsection{Deviations in Variance Calculation}
\CorrectionVarianceDeviationShape*
\begin{lemma} \jnote{alt version}
	For any deviation vector $\eta_\tau$ associated with a non-trivial shape $\tau$ such that
$\eta_\tau \neq \eta_0$, we have
\[
    \E\!\left[\eta_\tau^\top M^{-1}\eta_\tau\right]
    =   
    \frac{1}{1-\gamma}\,
    \E\!\left[\|\eta_\tau\|_2^2\right] + O(1) \,.
\]
Equivalently, \[ 
\Var(\corr(\tau)) =   \frac{1}{1-\gamma} \cdot \Var(\eta_\tau) +O(1) \,.
\]
\jnote{notice an error bound of $O(d)$ is interesting here as $\eta_0$ would have deviation $m/d$ by itself, so we are making polynomial improvement}
\end{lemma}

\begin{proof}
	For starters, consider the identity term from $A^{-1}$ which comes with a coefficient of $\frac{1}{1-\gamma}$, which gives us the first component of \(\frac{1}{1-\gamma} \E[\|\eta_\tau\|_2^2]\). It suffices for us to show the remaining terms are $o_d(1)$. Again, we apply our block value function.
	
	We note that the terms in $\eta_\tau^\top M^{-1} \eta_\tau$ admit the following description in diagram language,
	\begin{enumerate}
		\item Starting with a square vertex $s$, we attach to it a shape of $\eta_\tau$ (call this left-attachment);
		\item Next, we attach to the same square vertex the diagram corresponding to $M^{-1}$, that is either a (scaled) non-trivial $A^{-1}$-path from concatenation of $\al$ or $\beta$ shapes, or the corresponding shape from the low-rank component from Woodbury expansion. We call this the backbone-attachment, and the backbone-attachment ends with a square vertex $s'$ (with possibly $s'=s$).
		\item Finally, we attach to $s'$ a shape of $\eta_\tau$, and call this the right-attachment.
%		\item We do not enforce injectivity condition across any piece of left/backbone/right attachment.
	\end{enumerate}
	Next, we adopt the following strategy to bound the block-value of each term,
	\begin{enumerate}
		\item We start from any circle vertex on the latest backbone-path for the left $\eta$-attachment, and assign the first circle vertex a factor of $\sqrt{d}$ if it is making first/final appearance.
		\item Next, we apply the assignment scheme to the left-piece corresponding to $\eta$ attachment, and bound its block-value by $O(1)$. This process terminates at the $s$-vertex (which is the right boundary for the left-piece).
		\item From the $s$-vertex, we continue to apply block-value scheme and bound the block-value of the backbone-path by $O(1)$, with the process terminating at $s'$  (which is the right boundary for the backbone-piece)..
		\item Apply block-value bound on the right $\eta$-attachment from $s'$, and its value can be bounded by $O(\frac{1}{\sqrt{d}})$. To see this, note that all but the final gadget has $(1+o_d(1)) \sqrt{\gam}$ bound, while the final gadget is a $h_2$ attachment by our structural property. 
		\item Combining the above gives an error bound of $O(1)$ .
	\end{enumerate}	
	
	
\end{proof}
\end{comment}

\subsection{MP Derivations}\label{sec:MP-def-proof}

\begin{lemma}[Full Verification for MP Cancellation]
	\[\cm_{\mathsf{HalfDiamondInt}} =\gam \cdot M_\al + \Delta_{\mathsf{HalfDiamondInt}} \]
where we have \[\|\Delta_{\mathsf{HalfDiamondInt}}\|_{sp} =o_d(1)\,.\]
\end{lemma}

\begin{proof}
	We formally decompose $\Delta_{\mathsf{HalfDiamondInt}} $ into graph matrix basis. Recall that
\[
    M_\alpha[i,j]
    =
    \sum_{a\neq b\in[d]}
    v_i[a]v_i[b]v_j[a]v_j[b]
    =
    2\sum_{a<b\in[d]}v_i[a]v_i[b]v_j[a]v_j[b]
\]
for \(i\neq j\). For such \(i,j\), write
\[
    T_{ab}^{ij}
    :=
    v_i[a]v_i[b]v_j[a]v_j[b].
\]
Then
\[
\begin{aligned}
    \cm_{\mathsf{HalfDiamondInt}}[i,j]
    &=
    2\sum_{a<b\in[d]} T_{ab}^{ij}
    \left(
        \sum_{k\notin\{i,j\}}2v_k[a]^2v_k[b]^2
    \right)                                                     \\
    &=
    2\sum_{a<b\in[d]} T_{ab}^{ij}
    \left(
        \frac{2m}{d^2}
    \right)
    +
    2\sum_{a<b\in[d]} T_{ab}^{ij}
    \left[
        \sum_{k\notin\{i,j\}}2v_k[a]^2v_k[b]^2
        -
        \frac{2m}{d^2}
    \right]                                                     \\
    &=
    \gamma M_\alpha[i,j]
    +
    \Delta_{\mathsf{HalfDiamondInt}}[i,j],
\end{aligned}
\]
where we used \(\gamma=2m/d^2\), and where
\[
    \Delta_{\mathsf{HalfDiamondInt}}[i,j]
    :=
    2\sum_{a<b\in[d]} T_{ab}^{ij}
    \left[
        \sum_{k\notin\{i,j\}}2v_k[a]^2v_k[b]^2
        -
        \frac{2m}{d^2}
    \right].
\]
It remains to show that
\[
    \|\Delta_{\mathsf{HalfDiamondInt}}\|_{\mathrm{sp}}
    =
    o_d(1).
\]

We first rewrite the centered inner term. Since
\[
    v_k[a]^2v_k[b]^2-\frac1{d^2}
    =
    \frac1d\,h_2(v_k[a])
    +
    \frac1d\,h_2(v_k[b])
    +
    h_2(v_k[a])h_2(v_k[b]),
\]
we have
\[
\begin{aligned}
    \sum_{k\notin\{i,j\}}2v_k[a]^2v_k[b]^2-\frac{2m}{d^2}
    &=
    2\sum_{k\notin\{i,j\}}
    \left(
        v_k[a]^2v_k[b]^2-\frac1{d^2}
    \right)
    -
    \frac{4}{d^2}                                               \\
    &=
    \frac{2}{d}
    \sum_{k\notin\{i,j\}}
    \bigl(h_2(v_k[a])+h_2(v_k[b])\bigr)                         \\
    &\qquad
    +
    2\sum_{k\notin\{i,j\}}
    h_2(v_k[a])h_2(v_k[b])
    -
    \frac{4}{d^2}.
\end{aligned}
\]
Therefore
\[
    \Delta_{\mathsf{HalfDiamondInt}}
    =
    \Delta_1+\Delta_2+\Delta_3,
\]
where
\[
\begin{aligned}
    \Delta_1[i,j]
    &:=
    \frac{4}{d}
    \sum_{a<b}
    \sum_{k\notin\{i,j\}}
    T_{ab}^{ij}
    \bigl(h_2(v_k[a])+h_2(v_k[b])\bigr),                       \\
    \Delta_2[i,j]
    &:=
    4
    \sum_{a<b}
    \sum_{k\notin\{i,j\}}
    T_{ab}^{ij}
    h_2(v_k[a])h_2(v_k[b]),                                    \\
    \Delta_3[i,j]
    &:=
    -\frac{8}{d^2}
    \sum_{a<b}T_{ab}^{ij}
    =
    -\frac{4}{d^2}M_\alpha[i,j].
\end{aligned}
\]

For \(\Delta_1\), each summand is a graph matrix obtained from the \(M_\alpha\) shape by adding
one extra square vertex \(k\) and one \(h_2\)-edge from \(k\) to one of the two circle vertices
\(a,b\). In the block-value calculation, the four \(h_1\)-edges from the \(M_\alpha\) part contribute
a factor \(d^{-2}\), the extra \(h_2\)-edge contributes a factor \(d^{-1}\), while the additional internal
square vertex and the two circle vertices contribute at most
\[
    \sqrt m\cdot d.
\]
Thus, before the explicit prefactor \(1/d\), the block value is at most
\[
    \widetilde O\!\left(
        \sqrt m\cdot d\cdot \frac1{d^2}\cdot \frac1d
    \right)
    =
    \widetilde O\!\left(\frac{\sqrt m}{d^2}\right).
\]
Including the prefactor \(1/d\), and using \(m=\Theta(d^2)\), we obtain
\[
    B_q(\Delta_1)
    \le
    \widetilde O\!\left(\frac{\sqrt m}{d^3}\right)
    =
    \widetilde O\!\left(\frac1{d^2}\right)
    =
    o_d(1).
\]
By the trace moment bound for graph matrices, this implies
\(
    \|\Delta_1\|_{\mathrm{sp}}=o_d(1)
\)
with high probability.

Similarly, \(\Delta_2\) is a graph matrix obtained from the \(M_\alpha\) shape by adding one
extra square vertex \(k\) and two \(h_2\)-edges from \(k\) to the two circle vertices \(a,b\). The
four \(h_1\)-edges again contribute \(d^{-2}\), and the two \(h_2\)-edges contribute another
\(d^{-2}\). The vertex factor is again at most \(\sqrt m\cdot d\). Hence
\[
    B_q(\Delta_2)
    \le
    \widetilde O\!\left(
        \sqrt m\cdot d\cdot \frac1{d^2}\cdot \frac1{d^2}
    \right)
    =
    \widetilde O\!\left(\frac{\sqrt m}{d^3}\right)
    =
    \widetilde O\!\left(\frac1{d^2}\right)
    =
    o_d(1).
\]
Therefore
\[
    \|\Delta_2\|_{\mathrm{sp}}=o_d(1)
\]
with high probability.

Combining the three bounds gives
\[
    \|\Delta_{\mathsf{HalfDiamondInt}}\|_{\mathrm{sp}}
    \le
    \|\Delta_1\|_{\mathrm{sp}}
    +
    \|\Delta_2\|_{\mathrm{sp}}
    +
    \|\Delta_3\|_{\mathrm{sp}}
    =
    o_d(1).
\]
\end{proof}

\begin{claim}[Half-flat intersection]
\label{clm:half-flat-int}
We have
\[
    \cm_{\mathsf{HalfFlatInt}}
    =
    \gamma \cm_\beta
    +
    \Delta_{\mathsf{HalfFlatInt}},
\]
where
\[
    \|\Delta_{\mathsf{HalfFlatInt}}\|_{\mathrm{sp}}
    =
    o_d(1).
\]
Equivalently,
\[
    \cm_{\mathsf{HalfFlatInt}}
    =
    \gamma \cm_\beta
    +
    o_d(1)
\]
in spectral norm.
\end{claim}

\begin{proof}
For \(i\neq j\), recall that
\[
    \cm_\beta[i,j]
    =
    \sum_{a\in[d]} h_2(v_i[a])h_2(v_j[a]).
\]
The half-flat intersection has entries
\[
    \cm_{\mathsf{HalfFlatInt}}[i,j]
    =
    \sum_{a\in[d]} h_2(v_i[a])h_2(v_j[a])
    \left(
        \sum_{k\notin\{i,j\}} \frac{2}{d}v_k[a]^2
    \right).
\]
Since \(\E[v_k[a]^2]=1/d\), the deterministic contribution of the inner factor is
\[
    \frac{2(m-2)}{d^2}
    =
    \gamma+o_d(1).
\]
Thus
\[
    \cm_{\mathsf{HalfFlatInt}}[i,j]
    =
    \gamma \cm_\beta[i,j]
    +
    \Delta_{\mathsf{HalfFlatInt}}[i,j],
\]
where
\[
    \Delta_{\mathsf{HalfFlatInt}}[i,j]
    =
    \sum_{a\in[d]} h_2(v_i[a])h_2(v_j[a])
    \left[
        \sum_{k\notin\{i,j\}} \frac{2}{d}v_k[a]^2
        -
        \frac{2m}{d^2}
    \right].
\]
Using \(v_k[a]^2-1/d=h_2(v_k[a])\), we rewrite
\[
    \sum_{k\notin\{i,j\}} \frac{2}{d}v_k[a]^2
    -
    \frac{2m}{d^2}
    =
    \frac{2}{d}
    \sum_{k\notin\{i,j\}} h_2(v_k[a])
    -
    \frac{4}{d^2}.
\]
Hence \(\Delta_{\mathsf{HalfFlatInt}}\) is the sum of two terms: a centered graph matrix obtained
from the \(\cm_\beta\) shape by adding one extra square vertex and one extra centered \(h_2\)-edge,
with prefactor \(1/d\), and a deterministic correction equal to \(O(d^{-2})\cm_\beta\).

The first term has block value \(o_d(1)\), since the additional centered \(h_2\)-edge together with
the prefactor \(1/d\) gives an extra negative power of \(d\). The second term has norm \(o_d(1)\)
because \(\|\cm_\beta\|_{\mathrm{sp}}=O_d(1)\). \end{proof}

\begin{proposition}[Diamond Cancellation]
%\label{prop:formal-diamond-cancellation}
We have
\[
    \cm_{\mathsf{DiamondInt}}
    =
    \gamma I_m
    +
    \Delta_{\mathsf{DiamondInt}},
\]
where
\[
    \|\Delta_{\mathsf{DiamondInt}}\|_{\mathrm{sp}}
    =
    o_d(1).
\]
Equivalently,
\[
    \cm_{\mathsf{DiamondInt}}
    =
    \gamma I_m
    +
    o_d(1)
\]
in spectral norm.
\end{proposition}

\begin{proof}
A full-diamond intersection identifies the two boundary square vertices,
so \(\cm_{\mathsf{DiamondInt}}\) is diagonal. For \(i\in[m]\),
\[
\cm_{\mathsf{DiamondInt}}[i,i]
=
4\sum_{a<b}
v_i[a]^2v_i[b]^2
\sum_{k\neq i}v_k[a]^2v_k[b]^2.
\]
Set
\[
\rho_i
\coloneqq
2\sum_{a<b}v_i[a]^2v_i[b]^2
=
\|v_i\|_2^4-\sum_a v_i[a]^4
\]
and
\[
S_{ab}^{(-i)}
\coloneqq
2\sum_{k\neq i}v_k[a]^2v_k[b]^2.
\]
Then
\[
\cm_{\mathsf{DiamondInt}}[i,i]
=
2\sum_{a<b}v_i[a]^2v_i[b]^2S_{ab}^{(-i)}.
\]

With high probability,
\[
\max_i|\rho_i-1|=O(\frac{\poly\log d}{\sqrt{d}})= o_d(1),
\qquad
\max_{i,a<b}|S_{ab}^{(-i)}-\gamma|=O(\frac{\poly\log d}{\sqrt{d}})= o_d(1),
\]
where
\[
\E S_{ab}^{(-i)}
=
\frac{2(m-1)}{d^2}
=
\gamma+o_d(1).
\]
Consequently, for any \(i\),
\begin{align*}
\left|
\cm_{\mathsf{DiamondInt}}[i,i]-\gamma
\right|
&\leq
\gamma|\rho_i-1|
+
\rho_i\max_{a<b}|S_{ab}^{(-i)}-\gamma|  \\
&=o_d(1).
\end{align*}

Thus, writing
\[
\Delta_{\mathsf{DiamondInt}}
\coloneqq
\cm_{\mathsf{DiamondInt}}-\gamma I_m,
\]
we have
\[
\|\Delta_{\mathsf{DiamondInt}}\|_{\mathrm{sp}}
=
\max_i
\left|
\cm_{\mathsf{DiamondInt}}[i,i]-\gamma
\right|
=
o_d(1),
\]
which proves the claim.
\end{proof}

\begin{claim}[Proposition B.2 of~\cite{HKPX23}]
For $\cm_D$ defined in~\cref{eq:M_D_def} from the decomposition of $M$, we have
	\[ 
	\|\cm_D\|_{sp} =\tilde{O}(\frac{1}{\sqrt{d}} )= o_d(1)\,.
	\]
	with high probability.
\end{claim}

 \begin{claim}[Coefficient bound in the MP basis (Restatement of~\cref{clm:MP-coef-bnd})]
%\label{clm:MP-coef-bnd} 
  Assume \(0<\gamma\le1\). For every \(k\ge0\) and every \(t\ge0\), \[ |c_t(k)| \le \gamma^{-t/2}(2\sqrt{\gamma})^k. \] 
 In particular, for \(k=2q\), \[ |c_t(2q)| \le \gamma^{-t/2}(2\sqrt{\gamma})^{2q}. \] \end{claim} 

\begin{proof}[Proof of~\cref{clm:MP-coef-bnd}]
  Define \( B_k:=\max_{t\ge0}\gamma^{t/2}|c_t(k)|. \) We prove that \[ B_{k+1}\le 2\sqrt{\gamma}\,B_k. \] First consider \(t\ge2\). Using the recurrence, \[ \begin{aligned} \gamma^{t/2}|c_t(k+1)| &\le \gamma^{t/2}|c_{t-1}(k)| + \gamma^{t/2}\gamma |c_{t+1}(k)| \\ &= \sqrt{\gamma}\,\gamma^{(t-1)/2}|c_{t-1}(k)| + \sqrt{\gamma}\,\gamma^{(t+1)/2}|c_{t+1}(k)| \\ &\le 2\sqrt{\gamma}\,B_k. \end{aligned} \] For \(t=1\), we have \[ \begin{aligned} \gamma^{1/2}|c_1(k+1)| &\le \sqrt{\gamma}|c_0(k)| + \sqrt{\gamma}\gamma |c_2(k)| \\ &\le 2\sqrt{\gamma}\,B_k. \end{aligned} \] Finally, for \(t=0\), we use the boundary recurrence: \[ \begin{aligned} |c_0(k+1)| &\le \gamma |c_1(k)|+\gamma |c_0(k)| \\ &= \sqrt{\gamma}\,\gamma^{1/2}|c_1(k)| + \gamma |c_0(k)|. \end{aligned} \] Since \(0<\gamma\le1\), we have \(\gamma\le \sqrt{\gamma}\), and therefore \[ |c_0(k+1)| \le 2\sqrt{\gamma}\,B_k. \] 
Combining the three cases gives \[ B_{k+1}\le 2\sqrt{\gamma}\,B_k. \] Since \(B_0=1\), induction yields \[ B_k\le (2\sqrt{\gamma})^k. \] Equivalently, \[ |c_t(k)| \le \gamma^{-t/2}(2\sqrt{\gamma})^k. \] This proves the claim. \end{proof}


\subsection{Concentration for Matrix Norm Bounds}
\begin{claim}[Trace Moment Bound Implies Spectral Norm Control (Multiplicative)] \label{claim:trace-to-norm-rough}
Let $M\in \mathbb{R}^{n\times n}$ be a random matrix. Suppose that for some integer $q\ge 1$,
\[
\E\!\left[\Tr\!\big((MM^\top)^q\big)\right]
\le
n\left(B_q(M)\right)^{2q}.
\]
Then for any $\eta>0$,
\[
\Pr\!\left(
\|M\|_{\mathrm{sp}} \ge e^{\eta}\, B_q(M)
\right)
\le
n\, e^{-2q\eta}.
\]
In particular, if $q=\Omega(\log^2 n)$, then
\[
\|M\|_{\mathrm{sp}}
\le
(1+o_n(1))\, B_q(M)
\]
with probability $1-o_n(1)$.
\end{claim}
%
%\begin{proof}
%Since $MM^\top \succeq 0$, we have
%\[
%\|M\|_{\mathrm{sp}}^{2q}
%=
%\|MM^\top\|_{\mathrm{sp}}^q
%\le
%\Tr\!\big((MM^\top)^q\big).
%\]
%Therefore, for any $t>0$,
%\[
%\Pr\!\left(
%\|M\|_{\mathrm{sp}} \ge t\, B_q(M)
%\right)
%\le
%\Pr\!\left(
%\Tr\!\big((MM^\top)^q\big)
%\ge
%t^{2q} B_q(M)^{2q}
%\right).
%\]
%By Markov's inequality,
%\[
%\Pr\!\left(
%\|M\|_{\mathrm{sp}} \ge t\, B_q(M)
%\right)
%\le
%\frac{
%\E[\Tr((MM^\top)^q)]
%}{
%t^{2q} B_q(M)^{2q}
%}
%\le
%\frac{n}{t^{2q}}.
%\]
%Setting $t = e^{\eta}$ yields
%\[
%\Pr\!\left(
%\|M\|_{\mathrm{sp}} \ge e^{\eta}\, B_q(M)
%\right)
%\le
%n\, e^{-2q\eta}.
%\]
%
%Finally, take $\eta = \frac{\log n}{q}$. Then
%\[
%\Pr\!\left(
%\|M\|_{\mathrm{sp}} \ge n^{1/q} B_q(M)
%\right)
%\le
%n^{-1},
%\]
%and
%\[
%n^{1/q}
%=
%\exp\!\left(\frac{\log n}{q}\right)
%=
%1 + O\!\left(\frac{\log n}{q}\right)
%=
%1 + o_n(1),
%\]
%provided $q=\Omega(\log^2 n)$. This completes the proof.
%\end{proof}
%
We also make note of the following version that gives an explicit control for the additive deviation.
%\begin{claim}
%	With high probability, \[ 
%	\|M\|_{sp} \leq 
%	\]
%\end{claim}

% \subsection{Stronger Edge Control if needed}
\begin{corollary}[Quantitative version of \Cref{claim:trace-to-norm-rough}]
\label{cor:trace-to-norm-quant}
Under the assumptions of \Cref{claim:trace-to-norm-rough}, if $q=\log^4 n$, then
\[
\|M\|_{\mathrm{sp}}
\le
\left(1+\frac{1}{\log^2 n}\right) B_q(M)
\]
with probability $1-o_n(1)$.
\end{corollary}

\subsection{Auxiliary Bounds for Refutation}

\begin{proof}[Proof to~\cref{clm:refutation-aux-bounds}]
For convenience, write
\(
    \mathsf B_q\coloneqq B_q(\mathrm{non\text{-}ideal}).
\)
By the parameter verification above,
\[
    \mathsf B_q=d^{-1/2+o(1)}=o(\delta).
\]

Since \(\frac d m\calL^*(\1_m)\in\calS\), the variational
characterization of the orthogonal projection gives
\[
    \|\Pi^\perp I\|_F^2
    \leq
    \left\|I-\frac d m\sum_{i=1}^m v_iv_i^\top\right\|_F^2.
\]
Moreover,
\[
    \E\left\|I-\frac d m\sum_{i=1}^m v_iv_i^\top\right\|_F^2
    =
    \frac{d^2}{m}
    \E\left\|vv^\top-\frac1dI\right\|_F^2
    =
    O(1),
\]
The desired bound then follows by Markov.

Next, every diagonal term on $\calR$ and thereby $Z$ is treated as an error term in our combinatorial analysis in~\cref{lem:error-vertical}, consequently gives
\[
    |\Tr(\calR)|,\ |\Tr(Z)|
    \leq d\,\mathsf B_q
    =
    o(\delta d).
\]

Finally, consider
\[
    \langle Z,\calR\rangle=\Tr(Z\calR).
\]
An ideal closure would require a proper summand of \(Z\) to coincide
with a level-zero summand of \(\calR\). After subtracting
\(c_R\calR\), this cannot occur: horizontal terms contain at least two
backbone pieces, while later correction shapes do not reappear. Thus every mixed closure is non-ideal.
The block-value bound gives \(O(d\mathsf B_q)\), while the aggregate
truncation and termination errors contribute \(o(\delta d)\) by their
\(o(\delta)\) spectral-norm bounds and
\(\|\calR\|_{\mathrm{sp}}=O(1)\). Therefore,
\[
    |\langle Z,\calR\rangle|
    \leq O(d\mathsf B_q)+o(\delta d)
    =
    o(\delta d).
\]
\end{proof}




\section{Deferred Details for Semicircle Polynomials}
\label{app:def-Q}
\begin{proposition}[Backtracking Residual for Semicircle Polynomials]
\label{prop:sc-backtracking-residual}
Define \[ 
	\Delta_{\bti} \coloneqq  \sum _{\tau \in \calB(K)} (c_\tau)^2 \cdot  \cm_{\bti(\tau)} - I_d \,,	\]
	we have \[ 
	\|\Delta_{\bti}\|_{sp} = o_d(1)\,.
	\]
	provided $\Var(K)=1$. Suppose each shape has at most $D_V$ vertices. Then, with high probability,
\[
    \|\Delta_{\bti}\|_{\mathrm{sp}}
    \leq
    \frac{\poly(q,D_V)}{\sqrt d}
    +
    O\left(
        \frac{D_V^2}{d}
        +
        \frac{D_V^2}{m}
    \right).
\]
In particular, if \(m=\Theta(d^2)\) and \(q,D_V=\polylog(d)\), then
\[
    \|\Delta_{\bti}\|_{\mathrm{sp}}
    \leq
    \frac{\polylog(d)}{\sqrt d}
    =
    o_d(1).
\]
\end{proposition}

\begin{proof}
Write
\[
    \Delta_{\bti}=Z_{\bti}+B_{\bti},
\]
where
\[
    Z_{\bti}
    \coloneqq
    \sum_{\tau}
    c(\tau)^2
    \left(
        \cm_{\bti(\tau)}
        -
        \E\cm_{\bti(\tau)}
    \right)
\]
is the centered fluctuation and
\[
    B_{\bti}
    \coloneqq
    \sum_{\tau}
    c(\tau)^2\E\cm_{\bti(\tau)}-I_d
\]
is the injectivity bias.

For a fixed \(\tau\), the backtracking intersection gives
\[
    \E\cm_{\bti(\tau)}
    =
    \vartheta_\tau\Var(\tau)I_d,
\]
where
\[
    \vartheta_\tau
    =
    \frac{(d-1)_{a_\tau}}{d^{a_\tau}}
    \frac{(m)_{b_\tau}}{m^{b_\tau}},
    \qquad
    (x)_r=x(x-1)\cdots(x-r+1).
\]
Indeed, after fixing the boundary circle label, injectivity leaves
\((d-1)_{a_\tau}\) choices for the free circle labels and
\((m)_{b_\tau}\) choices for the square labels. Moreover,
\[
    0\leq 1-\vartheta_\tau
    \leq
    \frac{a_\tau(a_\tau+1)}{2d}
    +
    \frac{b_\tau(b_\tau-1)}{2m}
    =
    O\left(
        \frac{D_V^2}{d}
        +
        \frac{D_V^2}{m}
    \right).
\]
Using
\[
    \sum_{\tau}c(\tau)^2\Var(\tau)=1,
\]
we obtain
\[
    \|B_{\bti}\|_{\mathrm{sp}}
    =
    O\left(
        \frac{D_V^2}{d}
        +
        \frac{D_V^2}{m}
    \right).
\]

Finally, each summand in \(Z_{\bti}\) is a centered well-behaved
backtracking residual. From the subsequent~\cref{lem:centered-sc-backtracking}, \[
    \|Z_{\bti}\|_{\mathrm{sp}}
    \leq
    \frac{\poly(q,D_V)}{\sqrt d}
\]
with high probability. Combining the two estimates proves the claim.
\end{proof}

\begin{lemma}[Centered Backtracking Residual]
\label{lem:centered-sc-backtracking}
Assume \(m=\Theta(d^2)\) and every
\(\tau\in\calB(K)\) has at most \(D_V\) vertices
\[
    \|Z_{\bti}\|_{\mathrm{sp}}
    \leq
    \frac{\poly(q,D_V)}{\sqrt d}.
\]
\end{lemma}

\begin{proof}
For each \(\tau\), Hermite-expand the matched edge-products in the
well-behaved backtracking intersection and remove the all-constant
component. This gives a graph-matrix decomposition
\[
    \cm_{\bti(\tau)}
    -
    \E\cm_{\bti(\tau)}
    =
    \sum_{\rho\in\calR_{\bti}(\tau)}
    a_{\tau,\rho}\cm_\rho,
\]
where every residual shape \(\rho\) retains the decorated record
\(\tau\cdot\tau^\top\) and contains at least one centered,
positive-degree edge factor.

Consider a nonzero trace-walk contribution containing such a residual
block. Its centered edge factor must share its labeled input edge with
another block; otherwise the expectation over that input is zero.
The two matched copies of this edge already occur consecutively inside
the record \(\tau\cdot\tau^\top\). Hence one of its nonboundary
endpoints appears twice in the present block and at least once outside
it, so one of its appearances in the present block is a global middle
appearance.

Thus every residual block has an additional middle appearance relative
to the dominant well-behaved backtracking block. By the existing
factor-assignment bound, this gives a multiplicative gap
\[
    \frac{\poly(q,D_V)}
         {\sqrt{\wt(v)}}
    \leq
    \frac{\poly(q,D_V)}{\sqrt d}.
\]
Therefore,
\[
    B_q\!\left(
        \cm_{\bti(\tau)}
        -
        \E\cm_{\bti(\tau)}
    \right)
    \leq
    \frac{\poly(q,D_V)}{\sqrt d}\Var(\tau).
\]
Applying the linear-combination block-value bound and using
\(\sum_\tau c(\tau)^2\Var(\tau)=1\), we obtain
\[
    B_q(Z_{\bti})
    \leq
    \frac{\poly(q,D_V)}{\sqrt d}.
\]
The spectral-norm bound follows from the block-value-to-norm
implication.
\end{proof}
\section{Deferred Details for Truncation}
\label{sec:defer-details-truncation}

\subsection{Outer-Inner Truncation of the Iterative Process}
\begin{claim}[Coefficient Mass Shift (Restatement of \cref{claim:A-delta-shift})]
%\label{claim:A-delta-shift}
For sufficiently small \(\delta>0\),
\[
    A_\delta=1+\Theta(\delta).
\]
\end{claim}

\begin{proof}
Let
\[
    \Delta_\delta(x):=\widetilde F_\delta(x)-F(x).
\]
Since \(F(x)=2x_+\), we have
\[
    \Delta_\delta(x)
    =
    \begin{cases}
        \delta, & x<0,\\
        \delta-2x, & 0\le x<\delta/2,\\
        0, & x\ge \delta/2.
    \end{cases}
\]
In particular,
\[
    |\Delta_\delta(x)|\le \delta
    \qquad
    \text{for all }x\in[-2,2].
\]

We first estimate the shift in the \(P_1\)-coefficient.  Since \(P_1(x)=x\),
\[
    \widetilde a_1-1
    =
    \langle \Delta_\delta,x\rangle_{L^2(\mu_{\mathrm{sc}})}.
\]
On the interval \(x<0\), the contribution is
\[
    \delta\int_{-2}^0 x\,d\mu_{\mathrm{sc}}(x)
    =
    -\Theta(\delta).
\]
On the interval \(0\le x<\delta/2\), we have
\(|\Delta_\delta(x)|\le\delta\) and \(x\le\delta/2\), so the contribution is
\(O(\delta^3)\).  Therefore
\[
    \widetilde a_1=1-\Theta(\delta).
\]

Next we estimate the \(L^2\)-norm of \(\widetilde F_\delta\).  We have
\[
    \|\widetilde F_\delta\|_2^2
    =
    \|F\|_2^2
    +2\langle F,\Delta_\delta\rangle
    +\|\Delta_\delta\|_2^2.
\]
The product \(F\Delta_\delta\) is supported only on
\([0,\delta/2]\), where both \(F(x)=O(\delta)\) and
\(\Delta_\delta(x)=O(\delta)\).  Hence
\[
    \langle F,\Delta_\delta\rangle=O(\delta^3).
\]
Also, since \(|\Delta_\delta|\le \delta\),
\[
    \|\Delta_\delta\|_2^2=O(\delta^2).
\]
Using \(\|F\|_{L^2(\mu_{\mathrm{sc}})}^2=2\), we get
\[
    \|\widetilde F_\delta\|_{L^2(\mu_{\mathrm{sc}})}^2
    =
    2+O(\delta^2).
\]

Now apply Parseval to
\(
    F_\delta=\frac{\widetilde F_\delta}{\widetilde a_1}.
\)
Since the \(P_1\)-coefficient of \(F_\delta\) is normalized to be \(1\),
\[
    C_\delta^2+1+\sum_{j\ge2}(b_j^{(\delta)})^2
    =
    \|F_\delta\|_{L^2(\mu_{\mathrm{sc}})}^2
    =
    \frac{\|\widetilde F_\delta\|_{L^2(\mu_{\mathrm{sc}})}^2}
         {\widetilde a_1^2}.
\]
Using
\[
    \|\widetilde F_\delta\|_{L^2(\mu_{\mathrm{sc}})}^2=2+O(\delta^2),
    \qquad
    \widetilde a_1=1-\Theta(\delta),
\]
we obtain
\[
    C_\delta^2+1+\sum_{j\ge2}(b_j^{(\delta)})^2
    =
    \frac{2+O(\delta^2)}{(1-\Theta(\delta))^2}
    =
    2+\Theta(\delta).
\]
Rearranging gives us \[
    A_\delta
    =
    C_\delta^2+\sum_{j\ge2}(b_j^{(\delta)})^2
    =
    1+\Theta(\delta).
\]
%This proves the claim.
\end{proof}

%
%\begin{claim}[Threshold Shift for P.D. Mass: Truncated Version (Restatement of \cref{claim:pd-shift-threshold})]
%%\label{claim:pd-shift-threshold}
%There is a choice
%\[
%  \gamma_{\delta,D,t^*}
%    =
%    \frac12-\Theta(\delta)
%\]
%for which the truncated inner matrix produced by the \(F_\delta\)-based
%construction has variance
%\[
%    \Var\!\left(\calQ_{D,t^*}(\gamma_{\delta,D,t^*})\right)=1
%\]
%provided 
%\[
%    \delta\gg \varepsilon_{D,t^*}.
%\]
%For convenience, we write $\gam_\delta \coloneqq   \gamma_{\delta,D,t^*}$ when the dependence is clear.
%\end{claim}

%\begin{proof}
%By \cref{claim:A-delta-shift}, the untruncated variance mass associated with
%\(F_\delta\) is
%\[
%    A_\delta=1+\Theta(\delta).
%\]
%
%Let \(A_{\delta,D,t^*}\) denote the effective variance mass retained after
%truncating the \(F_\delta\)-based construction at inner degree \(D\) and
%iteration depth \(t^*\).  The truncated-variance estimate gives
%\[
%    A_{\delta,D,t^*}
%    =
%    A_\delta - O(\varepsilon_{D,t^*}).
%\]
%Equivalently,
%\[
%    A_{\delta,D,t^*}
%    =
%    1+\Theta(\delta) - O(\varepsilon_{D,t^*}) = 1+\Theta(\delta)
%\]
%for \(\delta\gg \varepsilon_{D,t^*}\).
%For fixed \(D,t^*\), the variance of the truncated inner matrix has the form
%\[
%    \Var\!\left(\calQ_{D,t^*}(\gamma)\right)
%    =
%    \frac{\gamma}{1-\gamma}\,A_{\delta,D,t^*}.
%\]
%Thus, to make the variance equal to \(1\), it suffices to choose
%\(\gamma=\gamma_{\delta,D,t^*}\) such that
%\[
%    \gamma_{\delta,D,t^*}
%    =
%    \frac{1}{1+A_{\delta,D,t^*}}.
%\]
%Using \(A_{\delta,D,t^*}=1+\Theta(\delta)\), we conclude that
%\[
%    \gamma_{\delta,D,t^*}
%    
%    \frac{1}{2+\Theta(\delta)}
%    =
%    \frac12-\Theta(\delta).
%\]
%This proves the claim.
%\end{proof}

\subsection{Truncation for \texorpdfstring{$A^{-1}$}{1/A}}
\begin{restatable}[Truncation error for $x^{-1}$ within the MP support]{lemma}{mpinversetruncation} \label{lem:mp-inverse-truncation} Let \(D\) be the truncation parameter for \(A^{-1}\), and define the truncated polynomial as \[ (x^{-1})_{\le D} \coloneqq \frac{1}{1-\gamma} \sum_{t=0}^{D}(-1)^t \q_t(x). \] Then, uniformly for \[ x\in\bigl[(1-\sqrt{\gamma})^2,\,(1+\sqrt{\gamma})^2\bigr], \] we have \[ \left| \frac{1}{x}-(x^{-1})_{\le D} \right| \le O_\gamma\!\left((D+1)(\sqrt{\gamma})^{D+1}\right). \] \end{restatable}\begin{proof}
	
For \(x\in [(1-\sqrt{\gamma})^2, (1+\sqrt{\gam} )^2 ]\), we may parameterize \[ x=1+\gamma+2\sqrt{\gamma}\cos\theta \] for some \(\theta\in[0,\pi]\). We observe that \[ \q_t(x) = \gamma^{t/2} \left( U_t(\cos\theta)+\sqrt{\gamma}\,U_{t-1}(\cos\theta) \right), \] where \(U_t\) is the Chebyshev polynomial of the second kind and \(U_{-1}:=0\). 

Since \( |U_t(\cos\theta)|\le t+1 \) for \(\theta\in[0,\pi]\), we get the uniform bound \[ |\q_t(x)| \le \gamma^{t/2}\bigl((t+1)+\sqrt{\gamma}\,t\bigr) \le 2(t+1)(\sqrt{\gamma})^t. \] Because \(\sqrt{\gamma}<1\), the series \( \sum_{t\ge0}(-1)^t q_t(x) \) converges absolutely and uniformly on \([(1-\sqrt{\gamma})^2, (1+\sqrt{\gam} )^2 ]\). 
Next, we have \[ \frac1x-(x^{-1})_{\le D} = \frac{1}{1-\gamma} \sum_{t>D}(-1)^t \q_t(x). \] Using the bound on \(\q_t(x)\), \[ \left| \frac1x-(x^{-1})_{\le D} \right| \le \frac{2}{1-\gamma} \sum_{t>D}(t+1)(\sqrt{\gamma})^t. \] Letting \(r=\sqrt{\gamma}<1\), we have \[ \sum_{t>D}(t+1)r^t = r^{D+1} \left( \frac{D+2}{1-r} + \frac{r}{(1-r)^2} \right). \] Therefore \[ \left| \frac1x-(x^{-1})_{\le D} \right| \le \frac{2r^{D+1}}{1-\gamma} \left( \frac{D+2}{1-r} + \frac{r}{(1-r)^2} \right) = O_\gamma\!\left((D+1)(\sqrt{\gamma})^{D+1}\right), \] as desired. \end{proof}

\begin{lemma}[Deviation of MP polynomials at the edge]
\label{lem:mp-edge-deviation-detailed}
Let
\[
    G_D(x)=\sum_{t=0}^D c_t\q_t(x),
    \qquad
    B_D:=\sum_{t=0}^D |c_t|.
\]
Then, for every sufficiently small \(\varepsilon\ge0\),
\[
    |G_D(b_\gam+\varepsilon)-G_D(b_\gam)|
    \le
    O_\gam\!\left(
        B_D\,\varepsilon (D+1)^3 e^{O_\gam(D\sqrt{\varepsilon})}
    \right),
\]
and similarly
\[
    |G_D(a_\gam-\varepsilon)-G_D(a_\gam)|
    \le
    O_\gam\!\left(
        B_D\,\varepsilon (D+1)^3 e^{O_\gam(D\sqrt{\varepsilon})}
    \right),
\]
where
\[
    a_\gam=(1-\sqrt{\gam})^2,\qquad
    b_\gam=(1+\sqrt{\gam})^2.
\]
\end{lemma}

\begin{proof}
We prove the upper edge. Write
\[
    b_\gam+\varepsilon
    =
    1+\gam+2\sqrt{\gam}\cosh s,
    \qquad
    s=\arccosh\!\left(1+\frac{\varepsilon}{2\sqrt{\gam}}\right)
    =
    O_\gam(\sqrt{\varepsilon}).
\]
Using
\[
    \q_t(1+\gam+2\sqrt{\gam}y)
    =
    \gam^{t/2}\bigl(U_t(y)+\sqrt{\gam}\,U_{t-1}(y)\bigr),
\]
together with
\[
    |U_j(\cosh s)-U_j(1)|
    \le
    O\!\left((j+1)^3s^2 e^{js}\right),
\]
we get
\[
    |\q_t(b_\gam+\varepsilon)-\q_t(b_\gam)|
    \le
    O_\gam\!\left(
        (t+1)^3\varepsilon e^{O_\gam(t\sqrt{\varepsilon})}
    \right).
\]
Therefore
\[
\begin{aligned}
    |G_D(b_\gam+\varepsilon)-G_D(b_\gam)|
    &\le
    \sum_{t=0}^D |c_t|\,
    |\q_t(b_\gam+\varepsilon)-\q_t(b_\gam)| \\
    &\le
    O_\gam\!\left(
        B_D\,\varepsilon (D+1)^3 e^{O_\gam(D\sqrt{\varepsilon})}
    \right).
\end{aligned}
\]
The lower edge follows identically from \(U_j(-x)=(-1)^jU_j(x)\).
\end{proof}

\mpinverseedgefluctuation*
\begin{proof}
Let
\[
    p_D(x):=(x^{-1})_{\le D}.
\]
On the exact MP support, via~\cref{lem:mp-inverse-truncation}, we have
\[
    \sup_{x\in[a_\gam,b_\gam]}
    \left|
        \frac1x-p_D(x)
    \right|
    \le
    O_\gam\!\left((D+1)(\sqrt{\gam})^{D+1}\right).
\]
Now take \(\lambda\in[a_\gam-\varepsilon_A,b_\gam+\varepsilon_A]\). If \(\lambda\in[a_\gam,b_\gam]\), this already gives the claim. Otherwise, \(\lambda\) lies within distance at most \(\varepsilon_A\) of one of the two edges. For example, if \(\lambda=b_\gam+\delta\), \(0\le\delta\le\varepsilon_A\), then
\[
\begin{aligned}
    \left|\frac1\lambda-p_D(\lambda)\right|
    &\le
    \left|\frac1\lambda-\frac1{b_\gam}\right|
    +
    \left|\frac1{b_\gam}-p_D(b_\gam)\right|
    +
    |p_D(b_\gam)-p_D(\lambda)| \\
    &\le
    O_\gam\!\left((D+1)(\sqrt{\gam})^{D+1}\right)
    +
    O_\gam\!\left(
        \delta (D+1)^4 e^{O_\gam(D\sqrt{\delta})}
    \right),
\end{aligned}
\]
where the last term uses \cref{lem:mp-edge-deviation-detailed} with
\[
    B_D=\frac{1}{1-\gam}\sum_{t=0}^D 1=O_\gam(D+1).
\]
The lower edge is the same, using \(a_\gam-\varepsilon_A>0\). Taking the supremum over eigenvalues of \(A\) gives the spectral-norm bound.
\end{proof}




%\section{Loose lemmas}
%\jnote{comment this out?}
%Define the projection $F_\epsilon(X)=\max(X,\epsilon)$ to the cone of matrices $\{X:X\succcurlyeq \epsilon I
%\}$ and the deviation $\textup{dev}(X)=\left(\left<Xv_i, v_i\right>-1\right)_{i\in [m]}$ which measures how far $X$ is from satisfying the linear constraints. The following lemma gives a sufficient condition 
%\jnote{need to think thru how to get such a starting point from the current solution...not sure about that myself - the plan i had was to add $\eps I$ to the solution of $F(Q)$, but that naturally gives deviation $O(\eps)$ as opposed to $o_d(1)$ independent of $\eps$. That's no good :p}
%\begin{lemma}
%    If $X\succcurlyeq \left(\epsilon-o\left(\frac 1{d^2}\right)\right)I$ and $|\textup{dev}(X)|=o\left(\frac 1{d^2}\right)$ then, with high probability as $d\to\infty$, $\pi\circ F_\epsilon (X)$ is a solution to the ellipsoid fitting problem, that is, it is PSD and satisfies the linear constraints.
%\end{lemma}

%\begin{proof}
%The image of $\pi$ is the set of elements that satisfies the linear constraints, therefore it suffices to prove that $\pi\circ F_\epsilon (X)$ is PSD. Recall that
%    \begin{align*}
%        \pi\circ F_\epsilon (X)&=F_\epsilon(X)+\sum_{i=1}^m \delta_i(F_\epsilon(X))v_iv_i^t
%    \end{align*}
%    where $\delta=M^{-1}\textup{dev}(X)$. So it is sufficient to prove that the above sum is bounded in operator norm by $\epsilon$. Towards that, write using the Cauchy-Schwarz inequality
%    \begin{align*}
%        \left\rVert\sum_{i=1}^m \delta_i(F_\epsilon(X))v_iv_i^t\right\rVert_{op}
 %       &\leq |\delta|\sqrt{\left\rVert\sum_{i=1}^m v_iv_i^t\right\rVert_{op}^2}\\
  %      &\leq \lVert M^{-1}\rVert_{op}|\textup{dev}(F_\epsilon(X))|\sqrt{\sum_{i=1}^m |v_i|^4}.
   % \end{align*}
    %Now we can bound $\lVert M^{-1}\rVert_{op}<(1-\sqrt{2\alpha})^{-2}+o(1)$ and an elementary bound gives $\sum_{i=1}^m |v_i|^4<m+o(1)$ with high probability. Therefore
%    \begin{align*}
%        \left\rVert\sum_{i=1}^m \delta_i(F_\epsilon(X))v_iv_i^t\right\rVert_{op}
%        &\leq O(d)|\textup{dev}(F_\epsilon(X))|.
%    \end{align*}
%    Furthermore, 
%    \begin{align*}
%        \textup{dev}(F_\epsilon(X))&=\left<(F_\epsilon(X)-X)v_i, v_i\right>)_{i\in [m]}+\textup{dev}(X)\\
%        |\textup{dev}(F_\epsilon(X))|
%        &\leq |\left<(F_\epsilon(X)-X)v_i, v_i\right>)_{i\in [m]}|+|\textup{dev}(X)|\\
%        &\leq (\lVert F_\epsilon(X)-X)\rVert_{op}|v_i|^2)_{i\in [m]}|+|\textup{dev}(X)|\\
%        &\leq \lVert F_\epsilon(X)-X\rVert_{op}\sqrt{\sum_{i=1}^m |v_i|^4}+o\left(\frac{1}{d^2}\right).
%    \end{align*}
%    Now the hypothesis $X\succcurlyeq \left(\epsilon-o\left(\frac 1{d^2}\right)\right)I$ gives $\lVert F_\epsilon(X)-X\rVert_{op}=o\left(\frac{1}{d^2}\right)$ so that
%    \begin{align*}
%        |\textup{dev}(F_\epsilon(X))|&=o(1)
%    \end{align*}
%\end{proof}

%\input{defer_charging.tex}
% {\color{red}
% \section{Definition recap}
% $d$ is the dimension of the underlying space, $m=m(d)=\frac\lambda 2 d^2+o(d^2)$, and $v_1, \dots, v_m\sim \mathcal N(0,I_d)$ and $\lambda<\frac 12$.
% \begin{align*}
%     \mathcal L:\textup{Sym}_d&\to \mathbb R^m\\
%     M&\mapsto (\left<Mv_i,v_i\right>)_{i\in [m]}\\
%     \mathcal L^*:\mathbb R^m&\to \textup{Sym}_d\\
%     w&\mapsto \sum_{i=1}^m w_i v_iv_i^\top
% \end{align*}
% The function $F$ is defined as the basis invariant function that acts as $F(x)=\max(2x,0)$ on the spectrum. Its expansion
% $$F(x)=C_F+x+\sum_{j=2}^\infty c_jP_j(x)$$
% where $P_j(x)=U_j(x/2)$ and $U_j$ is a Chebyshev polynomial of the second kind and $C_F=\frac{8}{3\pi}$ and $c_{2j+1}=0$ and
% $$c_{2j}=\frac{8(-1)^{j+1}}{\pi(2j-1)(2j+3)}.$$
% The affine space of matrices satisfying the linear constraints is $\textup{LC}=\cap_{i\in [m]} \{M\in\textup{Sym}_d:\left<Mv_i, v_i\right>=1\}$ and $\mathcal S=\textup{span}(v_iv_i^T)_{i\in [m]}$. 

% Define orthogonal projections $\pi$, $\Pi$, and $\Pi^\perp$ to $\textup{LC}$, $\mathcal S$, and $\mathcal S^\perp$ respectively. Note that $\mathcal S^\perp$ is the linear subspace parallel to $\textup{LC}$, in particular, $\Pi(X)-\pi(X)=\pi(0)$ for all $X$. The identity perturbation ansatz is $\pi(I)$.

% $M=(\left<v_i, v_j\right>^2)_{i,j\in [m]}$ and
% \begin{align*}
%     M&=A+B\\
%     A&=\left(1+\frac 1d\right)I+M_\alpha+M_\beta+M_D\\
%     B&=\frac{11^T+1\eta^T+\eta 1^T}{d}\\
%     M_\alpha[i,j]&=\sum_{a\neq b} v_i[a]v_i[b]v_j[a]v_j[b]\\
%     M_\beta[i,j]&=\sum_{a\in [d]}\left(v_i[a]^2-\frac 1d\right)\left(v_j[a]^2-\frac 1d\right)\\
%     M_D&=\textup{diag}\left(|v_i|^4-\frac 1d |v_i|^2-1\right)_{i\in [m]}
% \end{align*}. We can also write $M=\mathcal {LL}^*$.

% The iterative step is given by $Q_0=C_F(\pi(I)-I)$ and 
% $$Q_{i+1}=Q_i-\pi(F(Q_i))+F(Q_i).$$

% The orthogonal polynomials for the Marchenko-Pastur distribution are defined recusively
% \[
%     \mathfrak q_0^{(\gamma)}(x)=1,
%     \qquad
%     \mathfrak q_1^{(\gamma)}(x)=x-1,
% \]
% \[
%     \mathfrak q_{t+1}^{(\gamma)}(x)
%     =
%     \bigl(x-(1+\gamma)\bigr)\mathfrak q_t^{(\gamma)}(x)
%     -
%     \gamma\,\mathfrak q_{t-1}^{(\gamma)}(x).
% \]

% }


\end{document}
